\documentclass[11pt,a4paper]{article}
\pdfextension mapfile{+cm.map}
\pdfextension mapfile{+cmextra.map}
\pdfextension mapfile{+symbols.map}
\pdfextension mapfile{+euler.map}
\pdfextension mapfile{+latxfont.map}
\pdfextension mapfile{+rsfs.map}
\usepackage[margin=25mm,headheight=15pt]{geometry}
\usepackage{fontspec}
\setmainfont{Latin Modern Roman}
\usepackage{amsmath,amssymb,amsthm,amscd,mathtools,mathrsfs}
\usepackage{microtype,booktabs,longtable,array,xcolor,graphicx,enumitem,needspace,fancyhdr}
\usepackage[round,authoryear]{natbib}
\usepackage[unicode,colorlinks=true,linkcolor=blue!40!black,citecolor=blue!40!black,urlcolor=blue!40!black]{hyperref}
\usepackage{xurl}
\usepackage{listings}
\usepackage{tabularx}
\lstset{basicstyle=\ttfamily\small,breaklines=true,columns=fullflexible,keepspaces=true}
\usepackage{tikz}
\usepackage{tikz-cd}
\usetikzlibrary{arrows.meta,calc,patterns,positioning}
\usepackage{aliascnt}
\usepackage[nameinlink,noabbrev]{cleveref}
\allowdisplaybreaks[2]
\setcounter{MaxMatrixCols}{32}
\setlength{\emergencystretch}{3em}
\setlength{\parskip}{2pt}
\newtheorem{theorem}{Theorem}[section]
\newaliascnt{proposition}{theorem}
\newtheorem{proposition}[proposition]{Proposition}
\aliascntresetthe{proposition}
\crefname{proposition}{proposition}{propositions}
\newaliascnt{lemma}{theorem}
\newtheorem{lemma}[lemma]{Lemma}
\aliascntresetthe{lemma}
\crefname{lemma}{lemma}{lemmas}
\newaliascnt{corollary}{theorem}
\newtheorem{corollary}[corollary]{Corollary}
\aliascntresetthe{corollary}
\crefname{corollary}{corollary}{corollarys}
\newaliascnt{claim}{theorem}
\newtheorem{claim}[claim]{Claim}
\aliascntresetthe{claim}
\crefname{claim}{claim}{claims}
\newaliascnt{statement}{theorem}
\newtheorem{statement}[statement]{Statement}
\aliascntresetthe{statement}
\crefname{statement}{statement}{statements}
\theoremstyle{definition}
\newaliascnt{definition}{theorem}
\newtheorem{definition}[definition]{Definition}
\aliascntresetthe{definition}
\crefname{definition}{definition}{definitions}
\newaliascnt{remark}{theorem}
\newtheorem{remark}[remark]{Remark}
\aliascntresetthe{remark}
\crefname{remark}{remark}{remarks}
\newaliascnt{example}{theorem}
\newtheorem{example}[example]{Example}
\aliascntresetthe{example}
\crefname{example}{example}{examples}
\newaliascnt{problem}{theorem}
\newtheorem{problem}[problem]{Open problem}
\aliascntresetthe{problem}
\crefname{problem}{open problem}{open problems}
\newaliascnt{question}{theorem}
\newtheorem{question}[question]{Question}
\aliascntresetthe{question}
\crefname{question}{question}{questions}
\newaliascnt{auditrecord}{theorem}
\newtheorem{auditrecord}[auditrecord]{Verification record}
\aliascntresetthe{auditrecord}
\crefname{auditrecord}{verification record}{verification records}
\newaliascnt{test}{theorem}
\newtheorem{test}[test]{Verification}
\aliascntresetthe{test}
\crefname{test}{verification}{verifications}
\newaliascnt{conjecture}{theorem}
\newtheorem{conjecture}[conjecture]{Conjecture}
\aliascntresetthe{conjecture}
\crefname{conjecture}{conjecture}{conjectures}
\newcommand{\C}{\mathbb C}\newcommand{\R}{\mathbb R}\newcommand{\Z}{\mathbb Z}\newcommand{\Q}{\mathbb Q}\newcommand{\N}{\mathbb N}\newcommand{\F}{\mathbb F}
\newcommand{\CC}{\mathbb C}\newcommand{\RR}{\mathbb R}\newcommand{\ZZ}{\mathbb Z}\newcommand{\QQ}{\mathbb Q}\newcommand{\NN}{\mathbb N}
\newcommand{\PP}{\mathbb P}\newcommand{\Pp}{\mathbb P}\newcommand{\Pone}{\mathbb P^1}\newcommand{\HH}{\mathbb H}
\newcommand{\RP}{\mathbb{RP}}\newcommand{\CP}{\mathbb{CP}}
\newcommand{\OO}{\mathcal O}\newcommand{\cO}{\mathcal O}\newcommand{\cM}{\mathcal M}\newcommand{\cF}{\mathcal F}
\newcommand{\h}{\mathfrak h}\newcommand{\hz}{\mathfrak h_z}\newcommand{\Acon}{\mathfrak A}\newcommand{\wt}[1]{\widehat{#1}}
\newcommand{\red}{\mathrm{red}}\newcommand{\tor}{\mathrm{tor}}\newcommand{\dP}{\mathrm{dP}}
\newcommand{\gh}{\widehat\gamma}\newcommand{\uh}{\widehat u}\newcommand{\wh}{\widehat w}\newcommand{\dhat}{\widehat\delta}\renewcommand{\dh}{\widehat\delta}
\newcommand{\dd}{\mathrm d}\newcommand{\ii}{\mathrm i}\newcommand{\eps}{\varepsilon}\newcommand{\partt}{\mathrm{part}}\newcommand{\Bcirc}{B^\circ}\newcommand{\ddbar}{\partial\bar\partial}
\DeclareMathOperator{\Pic}{Pic}\DeclareMathOperator{\Alb}{Alb}\DeclareMathOperator{\Aut}{Aut}\DeclareMathOperator{\NS}{NS}
\DeclareMathOperator{\Hom}{Hom}\DeclareMathOperator{\End}{End}\DeclareMathOperator{\Ext}{Ext}\DeclareMathOperator{\Tor}{Tor}
\DeclareMathOperator{\Spec}{Spec}\DeclareMathOperator{\Proj}{Proj}\DeclareMathOperator{\Sym}{Sym}\DeclareMathOperator{\Diff}{Diff}
\DeclareMathOperator{\GL}{GL}\DeclareMathOperator{\SL}{SL}\DeclareMathOperator{\PSL}{PSL}\DeclareMathOperator{\Sp}{Sp}
\DeclareMathOperator{\im}{im}\DeclareMathOperator{\coker}{coker}\DeclareMathOperator{\rank}{rank}\DeclareMathOperator{\diag}{diag}
\DeclareMathOperator{\ord}{ord}\DeclareMathOperator{\vol}{vol}\DeclareMathOperator{\pr}{pr}\DeclareMathOperator{\id}{id}
\DeclareMathOperator{\divisor}{div}\DeclareMathOperator{\Tr}{Tr}\DeclareMathOperator{\Nm}{Nm}\DeclareMathOperator{\Sha}{Sha}
\DeclareMathOperator{\Imn}{Im}\DeclareMathOperator{\Ren}{Re}\DeclareMathOperator{\trdeg}{tr.deg}
\newcommand{\Uses}[1]{\par\smallskip\noindent\textit{Uses:} #1.\par\smallskip}
\newcommand{\Verified}{}\newcommand{\AuthorSelfCheck}{}\newcommand{\NeedsCheck}{}\newcommand{\Open}{\textbf{Open.}\ }

\newcommand{\StatementRecord}[4]{\par\smallskip\begingroup\upshape\raggedright\noindent\textbf{Label:} \nolinkurl{#1}.\par\noindent\textbf{Status:} #2.\par\noindent\textbf{Check record:} #3.\par\noindent\textbf{Uses:} #4.\par\endgroup\smallskip}
\newcommand{\RungZero}[1]{\par\noindent\textbf{Result and mechanism.} #1\par\smallskip}
\newcommand{\RungOne}{\par\Needspace{6\baselineskip}\noindent\textbf{Proof at the first rung.}\par\nobreak\smallskip}
\newcommand{\CurrentChecker}{GPT-6 Astra (OpenAI, ChatGPT Pro mode, effort Pro), 18 November 2026}
\newaliascnt{cited}{theorem}
\newtheorem{cited}[cited]{Cited result}
\aliascntresetthe{cited}
\crefname{cited}{cited result}{cited results}
\newaliascnt{refuted}{theorem}
\newtheorem{refuted}[refuted]{Correction}
\aliascntresetthe{refuted}
\crefname{refuted}{correction}{corrections}
\newaliascnt{computed}{theorem}
\newtheorem{computed}[computed]{Computation}
\aliascntresetthe{computed}
\crefname{computed}{computation}{computations}
\newaliascnt{openq}{theorem}
\newtheorem{openq}[openq]{Open question}
\aliascntresetthe{openq}
\crefname{openq}{open question}{open questions}

\providecommand{\rk}{\operatorname{rk}}
\providecommand{\Ann}{\operatorname{Ann}}
\providecommand{\Log}{\operatorname{Log}}
\providecommand{\vir}{\mathrm{vir}}


\newcommand{\ExternalPaperRef}[3]{\ifmmode\text{\upshape #1#2 (Paper~#3)}\else #1#2 (Paper~#3)\fi}
\newcommand{\CompendiumRef}[2]{\ifmmode\text{\upshape S6\_09, Paper~#1, label~\nolinkurl{#2}}\else S6\_09, Paper~#1, label~\nolinkurl{#2}\fi}

\begin{document}

\title{Complex fourfolds from a marked rank-six torus family}
\author{AI-produced mathematical text}
\date{19 November 2026}
\maketitle

\begin{abstract}
The marked rank-six monodromy with finite orders five and four and its
primitive maximal cusp produces compact complex fourfolds on
$3\mathbb{HP}^2\#(-\mathbb{HP}^2)$. Their explicitly normalized family
is parametrized by the half-plane $\operatorname{Im}\xi>c_*$, where
$c_*=1.0945402503114027801162682630549457037\ldots$ is defined by
a certified period minimum. Its members are pairwise nonisomorphic,
have algebraic dimension one and automorphism group
$\mathbb C\times\mathbb Z/2$, and have no antiholomorphic
isomorphisms between them. Wall's extended-form classification and
Kramer--Stolz's 2007 homotopy-sphere absorption theorem supply the
smooth classification input. The calculations here give the actual
integral cusp maps, the genus-two exterior-square period construction,
the family invariants and its endpoint obstruction, together with an
effective classification of the retained integral inputs. The hard
topological step is the coefficient-one torsion boundary map; the
analytic step is the exact real-transversality minimum. Historical
checks are recorded by statement. GPT-6 Astra Pro re-derived every
retained proof in this writing and reran the stated exact certificates;
no human proof reading is recorded in this writing.
\end{abstract}
\paragraph{AI production and check provenance.}
This paper is AI-produced. No human has checked its mathematics.
The recorded stages are GPT-5.6 Sol (OpenAI), in Codex at reasoning
effort ultra, 24 August--9 September 2026, with some contributions by
GPT-5.6 Sol (OpenAI) in ChatGPT Pro mode; GPT-6 Astra (OpenAI), in
ChatGPT Pro mode at effort Pro, for the research runs of 12--16
November, S6\_09 of 16--17 November, the continuing S6\_06 re-derivation
of 16--18 November, and this writing of 17--19 November 2026;
and Claude Fable 5.2 (Anthropic)
at effort high or max and Claude Opus 5.5 (Anthropic) at effort max
for the readings, checks and plans of 12--18 November 2026.
S6\_06\_BLIND is credited as a re-derivation with access to preceding
material. Each statement records the actual check scope and preserves
the names, effort settings and dates of its historical checks.

\paragraph{Effort names.}
The archive's literal settings are retained with their model, provider,
interface and date. OpenAI documents reasoning allocation controls
\cite[sections ``Reasoning effort'' and ``Reasoning mode'']{OpenAIReasoning2026}
and maps ChatGPT's Pro selection to GPT-6 Pro, powered by GPT-6 Astra
\cite[section ``Choose a thinking level'']{OpenAIChatGPTPro2026}.
Anthropic documents effort as control of response-token use, including
thinking \cite[sections ``How effort works'' and ``Effort levels'']{AnthropicEffort2026}.
The check event states what was read, re-derived or executed.

\paragraph{Public-domain dedication.}
This AI-produced text is dedicated to the public domain under
CC0 1.0 Universal \cite{CC01Universal}.



\paragraph{Known inputs and the precise additional calculations.}
The characteristic index polynomials are used in the explicit
almost-complex form explained in \S\ref{sec:V-known-inputs}.
The final smooth comparison uses the extended Wall form
\cite[Definition 3.6 and Theorem 3.7(1), pp.~12--14]{IVVWallForm}
and the absorption theorem
\cite[Theorem A, p.~178]{IVVKramerStolz}.
The open-base period construction uses the section form of the Oka
property \cite[Definition 7.7, p.~22; Proposition 7.13(d), p.~24]{IVVOkaOneV5}
after constructing the required local sprays. Each application retains
its marked lattice, source plane, finite quotient and cusp maps.

\paragraph{Reading the verification record.}
The current mathematical review covers all retained statement proofs,
including the complete cellular reduction in \cref{night:P1.19}.
Each proof's record distinguishes this review from earlier checks.
The current source openings identify the exact version and locator
used, and the computation records distinguish executed programs from
historical certificates. Typesetting and visual inspection have their
own records and confer no additional mathematical check.

\section{The fourfold and its family}
\label{sec:V-results}

We use the complex orientation and column homology vectors.  A marked
Kuhn cusp means the specified periodic unimodular triangulation, its
integral lattice marking, and its face identifications; these data are
part of the filling.  The two finite fibres have multiplicities five
and four.  The regular fibre is a complex three-torus.

\subsection{The integral eight-manifold}
\RungZero{The displayed rank-six monodromy and the specified affine
fillings produce a three-connected eight-manifold with integral middle
form $3\langle1\rangle\oplus\langle-1\rangle$.  The primitive cusp
collapses five fibre generators, while the actual order-four boundary
map kills each remaining order-two torsion class with coefficient one.}
\RungOne
\begin{enumerate}[label=\arabic*.]
\item The matrices $M_1=A^t$, $M_2=B^t$ have orders five and four, and
  $N=(M_1M_2)^{-1}-I$ has square zero and primitive image of rank three.
  The displayed unimodular marking takes $N$ to
  $\left(\begin{smallmatrix}0&I_3\\0&0\end{smallmatrix}\right)$;
  the invariant covector $\ell(x)=x_0$ proves freeness of every
  nonidentity affine power
  (\cref{maps:src:F.2.5}, \cref{night:P1.1}).
\item The cusp and finite quotient charts glue to a closed smooth
  manifold.  Killing the cusp angular lattice leaves the central fibre
  generator $h$ and the relations $5a=h$, $4b=-h$, $a+b=0$;
  their integral relation matrix is unimodular, so the fundamental
  group vanishes (\cref{maps:src:F.2.5}, \cref{night:P1.10}).
\item In the order-four cap the boundary base coordinate is
  $\theta=a+4k$.  Its primitive transverse loop is $(1,0)$ and its
  fibre loop is $(-4,1)$.  The four-prism pullback carries a torsion
  representative $w$ to $(w,-(I+T+T^2+T^3)w)=(w,0)$, so the actual
  transgression is $[w]$, with no extra factor four
  (\cref{mapsaudit:order-four-torsion-transgression}).
\item The three torsion transgressions are isomorphisms onto
  $\mathbb Z/2$.  Together with the primitive free transgression,
  the integral Leray complex gives
  $H_k=\mathbb Z$ for $k=0,8$, $H_4=\mathbb Z^4$, and zero otherwise;
  Hurewicz therefore gives three-connectivity
  (\cref{maps:src:F.2.6}, \cref{night:P1.12}, \cref{night:P1.19}).
\item The compatible almost complex structure has $c_1=c_3=0$.
  Todd integrality and the signature bound $|\sigma|\leq4$ force
  $c_2^2=2$, $p_1^2=8$, $p_2=14$ and $\sigma=2$.
  Its unimodular middle lattice is consequently the odd form
  $I_{3,1}$ (\cref{night:P1.13}).
\end{enumerate}

\subsection{The genus-two period construction}
\RungZero{The genus-two pencil
$y^2=x^5+\lambda^3x+\lambda^4$ gives the required rank-three period
planes.  The actual morphism is the projection of
$\ell_0\wedge\mathcal H_{\mathrm{dR}}$, followed by a holomorphic
splitting of the rank-six extension; a scalar shift removes its real
intersection with its conjugate.}
\RungOne
\begin{enumerate}[label=\arabic*.]
\item The pencil has one node at $\lambda_*=-256/3125$ and good
  reductions after root orders eight and five.  Its Hodge bundle on
  that root curve is $T^3\oplus T$, with
  $T=\mathcal O((2/5)[\infty]-(3/8)[0])$ and
  $\deg T=1/40$ (\cref{night:P1.4}).
\item In the exterior-square lattice, projection induces
  $j:\omega^\perp\to(\bigwedge^2\mathbb Z^4)/\mathbb Z\omega$
  with Smith factors $(1,1,1,1,2)$.  This exact index-two map identifies
  the rational systems, and the quotient is the dual of the primitive
  integral lattice.  The lower five-dimensional blocks of the given
  rank-six tuple are the quotient blocks under this map
  (\cref{night:P1.2}).
\item Put $P=p(\ell_0\wedge\mathcal H_{\mathrm{dR}})$, where
  $\ell_0=\langle dx/y\rangle$ and $p$ is projection modulo $\omega$.
  The wedge map has rank three and its image meets $\mathbb C\omega$
  trivially: a decomposable two-vector has square zero, whereas
  $\omega^2\ne0$.  This proves the required rank before a lift is
  chosen (\cref{night:P1.5}, \cref{maps:src:F.3.5}).
\item The orbifold extension over $P$ splits because
  $H^1(P^*)=0$, as computed from its two line-bundle extensions and
  their integral floor degrees.  A chosen splitting $\sigma$ and
  the canonical functional $\alpha:P\to\mathcal O$ give
  $F_\xi=(\sigma+\xi\alpha e)P$
  (\cref{maps:src:F.3.5}, \cref{night:P1.7}).
\item On the real line $P\cap\overline P$, choose $n$ with
  $\alpha(n)=1$ and write
  $\sigma(n)-\overline{\sigma(n)}=2ih e$.
  Then $F_\xi\cap\overline{F_\xi}=0$ exactly when
  $h+\operatorname{Im}\xi\ne0$.  At the cusp the period matrix is
  $Z=(\log q)/(2\pi i)I_3+C(q)$, with $C$ holomorphic and bounded;
  hence $h\to+\infty$ and a sufficiently large imaginary shift
  works on the whole completed base
  (\cref{maps:src:F.3.6}, \cref{night:P1.8}).
\end{enumerate}

\subsection{The complex structure on the connected sum}
\RungZero{The period planes extend over the two finite root charts and
the marked cusp, giving a compact complex fourfold diffeomorphic to
$3\mathbb{HP}^2\#(-\mathbb{HP}^2)$.  Its characteristic vector is
carried explicitly to $(1,1,1,1)$, after which the known smooth
classification and homotopy-sphere absorption identify the oriented
manifold.}
\RungOne
\begin{enumerate}[label=\arabic*.]
\item The lifted periods commute with the finite root actions and
  have the prescribed bounded cusp correction, so the same quotient
  and toroidal charts are holomorphic and proper.  Their collar
  identifications are the marked maps used in the smooth construction
  (\cref{maps:src:F.3.7}, \cref{night:P1.9}).
\item The preceding integral calculation gives the intersection
  lattice $I_{3,1}$.  The spin characteristic vector $p_1/2$ has
  square two, by the Chern-number calculation
  (\cref{night:P1.13}).
\item In an odd diagonal basis write its coordinates, after sign
  changes, as positive odd integers $(a,b,c,d)$.  A reflection in
  $(1,1,1,1)$ replaces $d$ by $2d-a-b-c$, whose absolute value is
  smaller when $d>1$.  Iteration terminates at $(1,1,1,1)$, proving
  an isometry of the lattice together with its characteristic vector
  (\cref{maps:src:F.4.1}).
\item Wall's oriented almost-diffeomorphism classification, in the
  precise extended-form statement cited there, identifies this pair
  with that of $3\mathbb{HP}^2\#(-\mathbb{HP}^2)$ up to a homotopy
  eight-sphere.  Kramer--Stolz's Theorem A absorbs that sphere in one
  $\mathbb{HP}^2$ summand, giving the stated oriented diffeomorphism
  (\cref{maps:src:F.4.1}).
\end{enumerate}

\subsection{The exact family, its invariants, and its ends}
For this paragraph $\sigma_b$ is the specified algebraic splitting,
and the marking and the Kuhn fan remain those of the construction.
Set
\[
 c_*=-h_{b,\min}=1.0945402503114027801162682630549457037\ldots,
 \qquad \mathcal U=\{\xi\in\mathbb C:\operatorname{Im}\xi>c_*\}.
\]
The exact definition of $c_*$ is the Frobenius period expression in
\cref{night:P1.18}; the decimal is a certified enclosure, not its
definition.  The comparison with the point-normalized parameter is
the full complex translation $\xi_b=\xi_{\rm original}-a_{\rm split}$.

\RungZero{The marked family is parametrized exactly by $\mathcal U$:
its members are pairwise nonisomorphic, have algebraic dimension one
and automorphism group $\mathbb C\times\mathbb Z/2$, and admit no
antiholomorphic isomorphism for this fan.  The real-intersection
function determines the parameter domain; the invariant line and
one half-period determine the automorphisms, while the actual
reflected cusp map fails to preserve a ray of the chosen fan.}
\RungOne
\begin{enumerate}[label=\arabic*.]
\item The ball-arithmetic proof covers the two finite root discs,
  the intervening collar and the cusp annulus, and proves that $h_b$
  has its unique minimum at the infinity root point.  The local
  positive Hessian proves strictness there; the remaining rectangles
  have strict certified lower bounds.  Thus the transversality test
  gives exactly $\mathcal U$ (\cref{night:P1.18}).
\item On the effective root curve of orders five and four put
  $L=\mathcal O((4/5)[\infty]-(3/4)[0])$.  The direct images are
  $R^qh_*\mathcal O=\bigwedge^q
  (\mathcal O\oplus L^{-1}\oplus L^{-2})$, and
  $K_M=h^*L^{-6}$.  The toric normalization complex proves the
  exterior-algebra assertion at the cusp, and the root-chart
  invariant cohomology proves it at the finite fibres
  (\cref{night:P1A.1}).
\item The logarithmic cotangent sequence and its explicit residue
  matrices give the Hodge table below.  The tangent sequence gives
  $h^q(T_M)=(1,1,1,2,0)$; differentiating the actual period lift
  gives a nonzero Kodaira--Spencer class, so the Kuranishi germ is
  smooth of dimension one (\cref{night:P1A.2}, \cref{night:P1A.3}).
\item An effective horizontal divisor contradicts $H^2(M,\mathbb R)=0$,
  so the meromorphic field is the base field $\mathbb C(\lambda)$.
  Every isomorphism therefore preserves the base and its three
  distinguished values.  The simultaneous linear centralizer is
  $\mathbb Q I$; the order-five affine norm excludes its negative
  integral unit.  The translation sheaf has identity component
  $\mathbb C$ and the displayed order-two half-period, proving the
  automorphism group and equality of parameters for isomorphic
  members (\cref{night:P1A.4}).
\item The two integral anti-intertwiners are $\pm R$.  One is excluded
  by the finite cap norm.  For the other, the actual based cusp
  correction acts on rays by $-P$, and takes $e_2$ to $(0,-1,1)$,
  outside the component's rays $\pm(\{0,1\}^3\setminus\{0\})$.
  The corresponding monomial map contracts a divisor, contradicting
  a biholomorphism (\cref{night:P1.17}).
\item A proposed proper Hausdorff filling across a boundary point of
  a full parameter disc recovers the intrinsic base from $K_M^{-10}$.
  Its marked periods would extend as $F_\xi$ across values where
  $h_b+\operatorname{Im}\xi=0$, contradicting smooth torus fibres
  (\cref{night:P1A.5}).  For a family on a punctured disc whose
  fibres belong to $\mathcal U$, the unique classifying map is
  holomorphic; its Cayley transform extends and the maximum principle
  keeps its limit inside $\mathcal U$.  After the explicit degree-two
  removal of the finite torsor, that family has a smooth filling
  (\cref{night:P1A.6}).
\end{enumerate}

With rows indexed by $p$ and columns by $q$, the Hodge table is
\[
\begin{array}{c|ccccc}
 h^{p,q}&0&1&2&3&4\\\hline
 0&1&1&0&0&0\\
 1&0&2&1&0&0\\
 2&0&1&6&1&0\\
 3&0&0&1&2&0\\
 4&0&0&0&1&1
\end{array}.
\]

\subsection{The effective input classification and the period decision}
\RungZero{Every supplied rank-six integral datum with its affine
orders, twists, clutch and marked primitive cusp has a terminating
conjugacy and integral-topology calculation.  A different, explicit
period-domain construction gives holomorphic transverse periods with
any prescribed finite cusp jet on the punctured base; the polarized
abelian-surface source is excluded by its actual two-line filtration
and the Euler-parity equation.}
\RungOne
\begin{enumerate}[label=\arabic*.]
\item Primitive square-zero cusp matrices have the explicit integral
  normal form and centralizer
  $\left(\begin{smallmatrix}P&Q\\0&P\end{smallmatrix}\right)$.
  Eigenvalue one is necessary for a free affine lift, leaving the
  finite cyclotomic list and its complex-compatibility filter
  (\cref{maps:src:F.2.1}, \cref{maps:src:F.2.2}).
\item The rational intertwiner algebra, its radical and its
  semisimple quotient reduce the simultaneous integral-conjugacy
  problem to the cited lattice-isomorphism algorithm in its verified
  hypothesis.  Searching by height and lexicographic order terminates
  by the input height, yielding a canonical representative and an
  explicit unimodular conjugator (\cref{maps:src:F.2.36}, \cref{maps:src:F.2.37}).
\item For that conjugator $R$, the coordinate map $x\mapsto Rx$
  transports $(M_i,v_i,\kappa,P)$ to
  $(RM_iR^{-1},Rv_i,R\kappa,RP)$.  Since
  $(RP)^{-1}(Rx)=P^{-1}x$, it intertwines the specified cusp charts
  as well as the affine caps and clutch.  The resulting finite
  cellular cone and group presentation compute all integral homology
  and the full fundamental group (\cref{maps:src:F.2.31}, \cref{maps:src:F.2.35}, \cref{maps:src:F.2.39}).
\item The exact two-point survivor tests and the three-point
  orbifold quotient give the complete constructive parametrization;
  the separate involution branch has 56 classes.  The displayed
  unbounded trace family proves why the terminating bound depends on
  the supplied input (\cref{maps:src:F.2.30}, \cref{maps:src:F.2.38}, \cref{maps:src:F.2.39}).
\item Nine explicitly given complex-line directions and a uniform
  singular-value bound put a spanning tube in the unpolarized period
  domain.  The verified local-spray approximation theorem, with the
  relative extension-point argument proved here, produces a
  holomorphic transverse period section and retains any prescribed
  bounded cusp jet (\cref{maps:src:F.3.8}, \cref{maps:src:F.3.10}, \cref{maps:src:F.3.11}).
\item For a polarized abelian-surface source, semisimplicity and the
  maximal square-zero cusp force the middle rank-four system to have
  no invariant or coinvariant vectors.  The actual filtration
  $0\subset\mathbb Q\subset E\subset T$, with middle quotient that
  system and top quotient $\mathbb Q$, then forces both relevant
  invariant ranks to be one; the integral Euler-parity relation
  excludes $b_1=b_2=0$ (\cref{maps:src:F.3.4}).
\end{enumerate}

% Additional notation transported from the night manuscript.
\providecommand{\LL}{\mathcal L}
\providecommand{\cT}{\mathcal T}
\providecommand{\HP}{\mathbb H\mathbb P}
\providecommand{\tr}{\operatorname{tr}}
\providecommand{\sgn}{\operatorname{sgn}}
\providecommand{\vir}{\mathrm{vir}}

\section{Geometric and characteristic conventions}\label{sec:V-known-inputs}
The index and signature polynomials used below are
\cite[equations (1.2)--(1.4), pp.~3--4]{ffEsoleKang}; the
rank-at-most-four integral classification is
\cite[Theorem 2.2.1, p.~61]{ffPetersSterk}.  In the almost
complex application, the Todd polynomial and the cotangent-twisted
polynomial are the corresponding elliptic-symbol indices, so their
integrality does not require an integrable complex structure.
For $c_1=c_3=0$ the substitution gives
\begin{equation}
\operatorname{Td}=\frac{3c_2^2-c_4}{720},\qquad
\chi^1=\frac{3c_2^2-31c_4}{180},\qquad
\sigma=\frac{3c_2^2+14c_4}{45}.
\label{maps:eq:F.1}
\end{equation}
In particular $\chi^1-4\operatorname{Td}=-c_4/6$, and
$\operatorname{Td}=(3\sigma-c_4)/48$.  For the rank-four middle
group and $c_4=6$ this forces $\sigma=2$, since
$\sigma\equiv2\pmod {16}$ and $|\sigma|\leq4$.
The cited lattice theorem then gives $I_{3,1}$ and the displayed
polynomials give the complete numerical record
\begin{equation}
 H^4(X;\mathbb Z)\cong I_{3,1},\qquad c_4=6,\qquad
 c_2^2=2,\qquad\sigma=2,\qquad\operatorname{Td}=0.
 \label{maps:eq:F.2}
\end{equation}
These substitutions are the form of the known inputs used in
the construction; the longer credited re-derivation remains
[S6\_09, \nolinkurl{maps:src:F.1.2}].

Mandal's connected-sum criterion is used only in its stated
orientation: $\ell\mathbb{HP}^2\#m(-\mathbb{HP}^2)$ has an
orientation-compatible almost complex structure exactly when
$\ell=2m+1$ and $m$ is odd
\cite[Corollary 1.4, p.~4, and its proof in \S5]{ffMandal}.
Its normalized tangent classes are those of
\cite[Lemma 2.1, p.~5, and Theorem 2.4, p.~7]{ffMandal}.
The credited application retained in the compendium is
[S6\_09, \nolinkurl{maps:src:F.1.3}].


\section{The marked cusp and preliminary lattice restrictions}



\begin{theorem}[Euler number of a Kuhn cusp]\label{maps:src:F.1.1}
\StatementRecord{maps:src:F.1.1}{proved}{Historical events: Claude Fable 5.2 (Anthropic), effort high, 16 November 2026: carried as background without fault found. GPT-6 Astra (OpenAI), ChatGPT Pro mode, effort Pro, S6\_09\_PAPERS, 17 November 2026: targeted proof check. Also checked in this writing. Current writing events: GPT-6 Astra (OpenAI), ChatGPT Pro mode, effort Pro, S6\_10\_PAPERS, 18 November 2026: targeted proof check. GPT-6 Astra (OpenAI), ChatGPT Pro mode, effort Pro, S6\_10\_PAPERS continuation, 19 November 2026: computation or formal certificate; independent re-proof; used and rechecked in this run. \emph{Recorded proof scope.} GPT-6 Astra (OpenAI, ChatGPT Pro mode), effort Pro, 2026-11-17; GPT-6 Astra (OpenAI, ChatGPT Pro mode), S6\_06\_BLIND, effort Pro, 2026-11-17; Claude Fable 5.2 (Anthropic), effort high, 16 November 2026; also checked in this writing by GPT-6 Astra (OpenAI), ChatGPT Pro mode, effort Pro, 18 November 2026}{the hypotheses, notation and definitions stated here; the explicit constructions and calculations in the proof}


Let $g:M^8\to S^2$ be a proper triangulable map, a $T^6$-bundle
off finitely many values.  Suppose its unique non-quotient special
fibre is the rank-three Kuhn toroidal fibre with period-lattice
index $d$, and all its other special reduced fibres are free
finite quotients of $T^6$.  Then $\chi(M)=6d$.

\end{theorem}


\begin{proof}
\begin{enumerate}[label=\arabic*.]
\item A fundamental cube has six unimodular tetrahedra.  Taking the
quotient by an index-$d$ sublattice multiplies every cell count
by $d$, so the quotient cell counts are $d(1,7,12,6)$.
A cell of dimension $r$ corresponds to a torus stratum
$(\mathbb C^*)^{3-r}$.  Only the $6d$ point strata contribute
to Euler characteristic.  Euler integration over the regular
base gives zero since $\chi(T^6)=0$, and a finite free torus
quotient has Euler characteristic $0/|G|=0$.  Additivity now
gives the formula.  This checks the rank-three instance of the
localization mechanism in \ExternalPaperRef{proposition~}{4.1}{VI} and \ExternalPaperRef{proposition~}{4.3}{VI}.  It requires the specified
toroidal filling, not just its unipotent monodromy matrix.

\end{enumerate}
\end{proof}




\begin{proposition}[the diagonal obstruction]\label{maps:src:F.1.4}
\StatementRecord{maps:src:F.1.4}{proved}{Historical events: Claude Fable 5.2 (Anthropic), effort high, 16 November 2026: carried as background without fault found. GPT-6 Astra (OpenAI), ChatGPT Pro mode, effort Pro, S6\_09\_PAPERS, 17 November 2026: targeted proof check. Also checked in this writing. Current writing events: GPT-6 Astra (OpenAI), ChatGPT Pro mode, effort Pro, S6\_10\_PAPERS, 18 November 2026: targeted proof check. GPT-6 Astra (OpenAI), ChatGPT Pro mode, effort Pro, S6\_10\_PAPERS continuation, 19 November 2026: independent re-proof; used and rechecked in this run. \emph{Recorded proof scope.} GPT-6 Astra (OpenAI, ChatGPT Pro mode), effort Pro, 2026-11-17; Claude Fable 5.2 (Anthropic), effort high, 16 November 2026; also checked in this writing by GPT-6 Astra (OpenAI), ChatGPT Pro mode, effort Pro, 18 November 2026}{the hypotheses, notation and definitions stated here; the explicit constructions and calculations in the proof}


If $A\in\operatorname{GL}_n(\mathbb Z)$ has no eigenvalue one,
every affine map $x\mapsto Ax+b$ of $\mathbb R^n/\mathbb Z^n$
has a fixed point.  Thus the diagonal elliptic finite monodromies
in Lead 7 admit no free affine lifts.

\end{proposition}


\begin{proof}
\begin{enumerate}[label=\arabic*.]
\item The real map $I-A$ is invertible.  A real solution of
$(I-A)x=b$ maps to a fixed point in the torus.  The diagonal
copies of the order-four and order-six elliptic rotations have
no eigenvalue one, so the conclusion applies to every translation
part, not only to integral or invariant choices.

\end{enumerate}
\end{proof}




\begin{proposition}[the cusp normal form]\label{maps:src:F.2.1}
\StatementRecord{maps:src:F.2.1}{proved}{Historical events: Claude Fable 5.2 (Anthropic), effort high, 16 November 2026: carried as background without fault found. GPT-6 Astra (OpenAI), ChatGPT Pro mode, effort Pro, S6\_09\_PAPERS, 17 November 2026: targeted proof check. Also checked in this writing. Current writing events: GPT-6 Astra (OpenAI), ChatGPT Pro mode, effort Pro, S6\_10\_PAPERS, 18 November 2026: targeted proof check. GPT-6 Astra (OpenAI), ChatGPT Pro mode, effort Pro, S6\_10\_PAPERS continuation, 19 November 2026: independent re-proof; used and rechecked in this run. \emph{Recorded proof scope.} GPT-6 Astra (OpenAI, ChatGPT Pro mode), effort Pro, 2026-11-17; GPT-6 Astra (OpenAI, ChatGPT Pro mode), S6\_06\_BLIND, effort Pro, 2026-11-17; Claude Fable 5.2 (Anthropic), effort high, 16 November 2026; also checked in this writing by GPT-6 Astra (OpenAI), ChatGPT Pro mode, effort Pro, 18 November 2026}{the hypotheses, notation and definitions stated here; the explicit constructions and calculations in the proof}


Under the integral index-one condition in F.2, the cusp matrix is
conjugate in $\operatorname{GL}_6(\mathbb Z)$ to
\[
 U=\begin{pmatrix}I_3&I_3\\0&I_3\end{pmatrix}.
\]
Its integral centralizer consists exactly of
$\left(\begin{smallmatrix}P&Q\\0&P\end{smallmatrix}\right)$
with $P\in\operatorname{GL}_3(\mathbb Z)$ and
$Q\in\operatorname{Mat}_{3\times3}(\mathbb Z)$.

\end{proposition}


\begin{proof}
\begin{enumerate}[label=\arabic*.]
\item Let $N=M_0-I$.  The index-one condition says that the image
is the saturated kernel lattice and that
$\Lambda/\ker N\to\operatorname{im}N$ is an isomorphism.
Choose a basis $a_1,a_2,a_3$ of this image and lifts $b_i$ with
$Nb_i=a_i$.  Saturation shows that the six vectors form an
integral basis.  It gives the displayed matrix.  Commuting a
block matrix with $N$ forces its lower-left block to vanish
and its two diagonal blocks to agree.  Invertibility over
$\mathbb Z$ is then equivalent to unimodularity of that block.

\end{enumerate}
\end{proof}




\begin{proposition}[a complete finite-order filter]\label{maps:src:F.2.2}
\StatementRecord{maps:src:F.2.2}{proved}{Historical events: Claude Fable 5.2 (Anthropic), effort high, 16 November 2026: carried as background without fault found. GPT-6 Astra (OpenAI), ChatGPT Pro mode, effort Pro, S6\_09\_PAPERS, 17 November 2026: targeted proof check. Also checked in this writing. Current writing events: GPT-6 Astra (OpenAI), ChatGPT Pro mode, effort Pro, S6\_10\_PAPERS, 18 November 2026: targeted proof check. GPT-6 Astra (OpenAI), ChatGPT Pro mode, effort Pro, S6\_10\_PAPERS continuation, 19 November 2026: computation or formal certificate; independent re-proof; used and rechecked in this run. \emph{Recorded proof scope.} GPT-6 Astra (OpenAI, ChatGPT Pro mode), S6\_06\_BLIND, effort Pro, 2026-11-17; Claude Fable 5.2 (Anthropic), effort high, 16 November 2026; also checked in this writing by GPT-6 Astra (OpenAI), ChatGPT Pro mode, effort Pro, 18 November 2026}{the hypotheses, notation and definitions stated here; \cref{maps:src:F.1.4}}


The linear order of a finite matrix in $\operatorname{GL}_6(\mathbb Z)$
which admits a free affine lift belongs to
\[
 \{1,2,3,4,5,6,8,10,12\}.
 \tag{F.3}\label{maps:eq:F.3}
\]
There are 38 possible degree-six cyclotomic characteristic
polynomials with an eigenvalue one, and 19 of them commute
with some real complex structure.  These are necessary types,
not 38 or 19 conjugacy classes of admissible tuples.

\end{proposition}


\begin{proof}
\begin{enumerate}[label=\arabic*.]
\item Proposition \ref{maps:src:F.1.4} requires a nonzero one-eigenspace.  The other
cyclotomic factors have total degree at most five.  The elementary
identity
\[
 \frac{\phi(n)^2}{n}=
 \prod_{p^a\parallel n}p^{a-2}(p-1)^2\geq\frac12
\]
implies $n\leq50$ when $\phi(n)\leq5$.
The complete enumeration through 50 gives
$n=1,2,3,4,5,6,8,10,12$.  Enumerating their multiplicities
with total degree six and at least one factor $T-1$ gives 38
polynomials.  A commuting complex structure exists on a real
finite-order module exactly when the multiplicities of the
real eigenvalues $1$ and $-1$ are even; nonreal pairs already
carry such a structure.  Imposing those two parity conditions
leaves 19.  Taking least common multiples of their root orders
gives exactly \eqref{maps:eq:F.3}, also for the 38-polynomial list.
The code prints every polynomial and every exterior invariant
rank.  This is a bounded exhaustive enumeration for the
stated filter.  It does not bound affine multiplicities when
the affine action has larger order than its linear part.

\end{enumerate}
\end{proof}




\begin{proposition}[Counterexample: determination by peripheral data alone]\label{maps:src:F.2.3}
\StatementRecord{maps:src:F.2.3}{false as stated}{Historical events: Claude Fable 5.2 (Anthropic), effort high, 16 November 2026: carried as background without fault found. GPT-6 Astra (OpenAI), ChatGPT Pro mode, effort Pro, S6\_09\_PAPERS, 17 November 2026: targeted proof check. Also checked in this writing. Current writing events: GPT-6 Astra (OpenAI), ChatGPT Pro mode, effort Pro, S6\_10\_PAPERS, 18 November 2026: targeted proof check. GPT-6 Astra (OpenAI), ChatGPT Pro mode, effort Pro, S6\_10\_PAPERS continuation, 19 November 2026: targeted proof check; used and rechecked in this run. \emph{Recorded proof scope.} Claude Fable 5.2 (Anthropic), effort high, 16 November 2026; also checked in this writing by GPT-6 Astra (OpenAI), ChatGPT Pro mode, effort Pro, 18 November 2026}{the hypotheses, notation and definitions stated here; the explicit constructions and calculations in the proof}


Monodromy matrices and affine twist vectors alone do not
determine the homology of a compact glued manifold.  A cusp
filling, or an equivalent relative topological specification,
is additional data.

\end{proposition}


\begin{proof}
\begin{enumerate}[label=\arabic*.]
\item Use the constructed fourfold $Z$ and blow up a smooth point
of the smooth locus of its cusp fibre.  The operation takes
place entirely inside that filling.  It leaves every regular
fibre, every finite quotient model, every peripheral monodromy
and every affine twist unchanged.  The total space remains
smooth and compact, and its map to the same base remains
holomorphic.  Blowing up a point in complex dimension four
replaces a ball by the disc bundle of $\mathcal O_{\mathbb P^3}(-1)$.
Mayer--Vietoris along their common $S^7$ boundary adds a copy
of $\mathbb Z$ in degrees two, four and six.  Thus the Euler
number increases by three and the homology changes.  The
peripheral index-one unipotent matrix is unaffected.
The stronger specification of a particular height-one Kuhn
filling excludes this example.  With that extra specification
the question remains a substantive classification problem;
we do not claim to refute it.

\end{enumerate}
\end{proof}




\begin{theorem}[the self-dual parity obstruction]\label{maps:src:F.3.1}
\StatementRecord{maps:src:F.3.1}{proved}{Historical events: Claude Fable 5.2 (Anthropic), effort high, 16 November 2026: carried as background without fault found. GPT-6 Astra (OpenAI), ChatGPT Pro mode, effort Pro, S6\_09\_PAPERS, 17 November 2026: targeted proof check. Also checked in this writing. Current writing events: GPT-6 Astra (OpenAI), ChatGPT Pro mode, effort Pro, S6\_10\_PAPERS, 18 November 2026: targeted proof check. GPT-6 Astra (OpenAI), ChatGPT Pro mode, effort Pro, S6\_10\_PAPERS continuation, 19 November 2026: independent re-proof; used and rechecked in this run. \emph{Recorded proof scope.} Claude Fable 5.2 (Anthropic), effort high, 16 November 2026; also checked in this writing by GPT-6 Astra (OpenAI), ChatGPT Pro mode, effort Pro, 18 November 2026}{the hypotheses, notation and definitions stated here; \cref{maps:src:F.1.1}}


Suppose a compact smooth eight-manifold has the torus fibration
and the specified Kuhn cusp of \ref{maps:src:F.1.1} with index one, and its
other reduced fibres are free quotients of complex three-tori
by finite holomorphic affine actions.  If its rational
$H^1$ monodromy representation is self-dual, then it cannot
satisfy $b_1=b_2=0$.  In particular no symplectic-monodromy
datum of F.2 can give the three-connected fourfold required
in F.4, for either $k=2$ or $k=3$ (or any number $k$ of finite
points).

\end{theorem}


\begin{proof}
\begin{enumerate}[label=\arabic*.]
\item Let $\mathcal V$ be the rank-six cohomology local system on
the punctured sphere and let $j$ be its inclusion.  On a finite
quotient fibre, averaging gives $R^1g_*\mathbb Q=(j_*\mathcal V)$
at that point.  At the Kuhn cusp, its normalization components
are simply connected and the degree-one branch complex is
the cochain complex of its dual $T^3$.  Specialization takes
its three deck classes to the covectors annihilating the
angular lattice, exactly the three-dimensional monodromy
invariant subspace.  Hence the same equality holds at the cusp.
It holds everywhere.

\item Put $a=\dim H^0(B,j_*\mathcal V)$ and
$b=\dim H^2(B,j_*\mathcal V)$.  The only differential that
can kill a class in $H^0(R^1)$ or in $H^2(R^0)$ is
$d_2:H^0(R^1)\to H^2(R^0)=\mathbb Q$.  If $b_1=b_2=0$,
it must be an isomorphism, so $a=1$.  Moreover
$H^1(B,j_*\mathcal V)$ must vanish because its position
$(1,1)$ has no possible incoming or outgoing differential.
Self-duality and compactly supported duality give $b=a=1$.

\item At the $i$th finite point, let $r_i$ be the dimension of the
fixed subspace in $H^1$.  The derivative is complex linear,
so its one-eigenspace has even real dimension.  Each $r_i$
is therefore even.  The cusp fixed dimension is three.
Euler additivity gives
\[
 \chi(B,j_*\mathcal V)=6(1-k)+3+\sum_{i=1}^k r_i,
\]
an odd integer.  Consequently
\[
 \dim H^1(B,j_*\mathcal V)=a+b-\chi(B,j_*\mathcal V)
\]
is odd and cannot vanish.  This is the required contradiction.
A preserved nondegenerate alternating form identifies a
representation with its dual, so symplectic monodromy is
included.  The explicit form
$(6\gamma\wedge\delta+u\wedge w)\oplus\omega_E$
checks this hypothesis for the example of Part E.

\end{enumerate}
\end{proof}




\begin{proposition}[the remaining rank patterns]\label{maps:src:F.3.2}
\StatementRecord{maps:src:F.3.2}{proved}{Historical events: Claude Fable 5.2 (Anthropic), effort high, 16 November 2026: carried as background without fault found. GPT-6 Astra (OpenAI), ChatGPT Pro mode, effort Pro, S6\_09\_PAPERS, 17 November 2026: targeted proof check. Also checked in this writing. Current writing events: GPT-6 Astra (OpenAI), ChatGPT Pro mode, effort Pro, S6\_10\_PAPERS, 18 November 2026: targeted proof check. GPT-6 Astra (OpenAI), ChatGPT Pro mode, effort Pro, S6\_10\_PAPERS continuation, 19 November 2026: computation or formal certificate; independent re-proof; used and rechecked in this run. \emph{Recorded proof scope.} Claude Fable 5.2 (Anthropic), effort high, 16 November 2026; also checked in this writing by GPT-6 Astra (OpenAI), ChatGPT Pro mode, effort Pro, 18 November 2026}{the hypotheses, notation and definitions stated here; the explicit constructions and calculations in the proof}


If a possibly non-self-dual holomorphic datum with the same
local models had $b_1=b_2=0$, put
$b=\dim H^0(B,\mathcal V^*)$.  For $k=2$, the only possible
unordered pairs $(r_1,r_2;b)$ are
\[
 (2,2;0),\qquad(2,4;2).
\]
For $k=3$, they are
\[
 (2,2,6;0),\ (2,4,4;0),\ (2,4,6;2),\ (4,4,4;2).
 \tag{F.4}\label{maps:eq:F.4}
\]

\end{proposition}


\begin{proof}
\begin{enumerate}[label=\arabic*.]
\item The preceding Leray argument still gives $a=1$ and vanishing
$H^1(B,j_*\mathcal V)$.  Euler additivity therefore gives
$b=\sum r_i+3-1-6(k-1)$.  Each $r_i$ belongs to
$\{2,4,6\}$ because a free affine lift requires an eigenvalue
one and complex linearity makes its multiplicity even.
Also $0\leq b\leq\min(3,r_1,\ldots,r_k)$, by inclusion in
each local invariant subspace.  Enumerating these three-value
lists gives exactly \eqref{maps:eq:F.4}, as independently printed by the
complete bounded filter.  No existence of a tuple or of a
period family with these ranks is asserted.

\end{enumerate}
\end{proof}




\begin{proposition}[resolution of the classification problem]\label{maps:src:F.2.4}
\StatementRecord{maps:src:F.2.4}{proved}{Historical events: Claude Fable 5.2 (Anthropic), effort high, 16 November 2026: carried as background without fault found. GPT-6 Astra (OpenAI), ChatGPT Pro mode, effort Pro, S6\_09\_PAPERS, 17 November 2026: targeted proof check. Also checked in this writing. Current writing events: GPT-6 Astra (OpenAI), ChatGPT Pro mode, effort Pro, S6\_10\_PAPERS, 18 November 2026: targeted proof check. GPT-6 Astra (OpenAI), ChatGPT Pro mode, effort Pro, S6\_10\_PAPERS continuation, 19 November 2026: used and rechecked in this run. \emph{Recorded proof scope.} Claude Fable 5.2 (Anthropic), effort high, 16 November 2026; also checked in this writing by GPT-6 Astra (OpenAI), ChatGPT Pro mode, effort Pro, 18 November 2026}{\cref{maps:src:F.2.16}, \cref{maps:src:F.2.30}, \cref{maps:src:F.2.31}, \cref{maps:src:F.2.35}, \cref{maps:src:F.2.36}, \cref{maps:src:F.2.37}, \cref{maps:src:F.2.38}, \cref{maps:src:F.2.39}}


The previously incomplete linear classification and all-twist
integral topology computation are completed by \ref{maps:src:F.2.31} and
\ref{maps:src:F.2.35}--39, with the fixed cusp filling and full marking
specified as required by \ref{maps:src:F.2.3}.  Every linear tuple has an
exact canonical representative and a returned unimodular
conjugator; the finite bound depends on the supplied input.
The unbounded trace family proves that no fixed global
height cutoff can replace this procedure.  All actual affine
orders, numerators, clutch vectors and cusp markings are
included in the general fundamental-group and integral
homology algorithms.  The complete constructive survivor
list, its integral tests and the common oriented smooth
type are given in \ref{maps:src:F.2.30}.

This conclusion does not claim an executed enumeration of
the infinite general class set.  The exhaustive finite
involution catalogue, the general comparison tests, the
three full topology integrations and the unified wrapper
checks have the exact executed scopes stated in \ref{maps:src:F.2.35}--39.
The earlier obstruction record is retained verbatim in
\path{records/fourfolds_F2_4_superseded.tex}; its finite-order
filter and homology-level obstructions were superseded by
the complete canonical and chain constructions.


\end{proposition}
\begin{proof}
\begin{enumerate}[label=\arabic*.]
\item The integral intertwiner module and its centralizer order in
\cref{maps:src:F.2.36} decide simultaneous integral conjugacy. Enumerating
tuples in increasing height and then lexicographic order, and applying
that decision to each candidate, gives the canonical tuple and the
returned unimodular intertwiner of \cref{maps:src:F.2.37}. The supplied
tuple itself occurs by its own height, proving termination for this input.
\item The affine power vectors, integral normal closure and marked
extension presentation in \cref{maps:src:F.2.31} give the full fundamental
group. The finite-cap, cusp and complement maps of
\cref{maps:src:F.2.35} form the actual gluing cone, whose finite integer
differentials determine every integral homology group by Smith reduction.
All orders, affine numerators, the clutch vector and the cusp basis enter
those maps explicitly.
\item The unbounded conjugacy invariant of \cref{maps:src:F.2.16}
rules out any uniform finite height box: a finite box contains only
finitely many tuples and hence only finitely many values of that trace.
The survivor criterion and its smooth classification are exactly
\cref{maps:src:F.2.30}; the involution catalogue and complete algorithmic
assembly are \cref{maps:src:F.2.38}, \cref{maps:src:F.2.39}.
\end{enumerate}
\end{proof}



\section{The integral eight-manifold}



\begin{theorem}[a non-self-dual topological candidate]\label{maps:src:F.2.5}
\StatementRecord{maps:src:F.2.5}{proved}{Historical events: Claude Fable 5.2 (Anthropic), effort high, 16 November 2026: carried as background without fault found; targeted proof check. GPT-6 Astra (OpenAI), ChatGPT Pro mode, effort Pro, S6\_09\_PAPERS, 17 November 2026: targeted proof check. Also checked in this writing. Current writing events: GPT-6 Astra (OpenAI), ChatGPT Pro mode, effort Pro, S6\_10\_PAPERS, 18 November 2026: computation or formal certificate; targeted proof check. GPT-6 Astra (OpenAI), ChatGPT Pro mode, effort Pro, S6\_10\_PAPERS continuation, 19 November 2026: computation or formal certificate; used and rechecked in this run. \emph{Recorded proof scope.} GPT-6 Astra (OpenAI, ChatGPT Pro mode), effort Pro, 2026-11-17; Claude Fable 5.2 (Anthropic), effort high, 2026-11-16; GPT-6 Astra (OpenAI, ChatGPT Pro mode), S6\_06\_BLIND, effort Pro, 2026-11-17; also checked in this writing by GPT-6 Astra (OpenAI), ChatGPT Pro mode, effort Pro, 18 November 2026}{the hypotheses, notation and definitions stated here; \cref{maps:src:F.2.1}; \ExternalPaperRef{theorem~}{6.3}{III}; \S\ref{sec:V-known-inputs}}


There is a closed smooth oriented eight-manifold with the prescribed
index-one Kuhn cusp, two free affine fillings of multiplicities $5,4$,
trivial fundamental group, and rational Betti numbers
\[
 (1,0,0,0,4,0,0,0,1). \tag{F.5}\label{maps:eq:F.5}
\]
It admits an almost complex structure.  Its signature is two and
its middle intersection lattice modulo torsion is
$3\langle1\rangle\oplus\langle-1\rangle$.
Integral $H_2,H_3$ and a holomorphic period realization are not
asserted here.

\end{theorem}


\begin{proof}
\begin{enumerate}[label=\arabic*.]
\item Define the two cohomological matrices
\[
A=\begin{pmatrix}
1&0&0&1&-1&1\\0&1&-1&-1&1&0\\0&1&0&-1&0&0\\
0&2&-1&-1&0&1\\0&2&0&-1&0&0\\0&1&1&-1&0&0
\end{pmatrix},\quad
B=\begin{pmatrix}
1&0&0&0&0&0\\0&0&-1&0&1&0\\0&0&-1&0&0&1\\
0&0&-1&0&0&0\\0&1&-1&0&0&0\\0&0&-1&1&0&0
\end{pmatrix}.
\]
Use homology monodromies $M_1=A^t$, $M_2=B^t$ and
$M_0=(M_1M_2)^{-1}$, with twists
\[
 v_1=(1,3,-2,-1,2,0)^t,\qquad v_2=(-1,0,0,0,0,0)^t .
\]
Exact multiplication gives orders five and four.  For $N=M_0-I$,
$N^2=0$, $\operatorname{rank}N=3$, and the nonzero Smith factors
are $1,1,1$.  The finite characteristic polynomials are
$(T-1)^2\Phi_5(T)$ and $(T-1)^2(T+1)^2(T^2+1)$.
Thus each finite matrix commutes with an orientation-compatible
complex structure.  The primitive invariant covector
$\ell(x)=x_0$ has $\ell(v_1)=1$ and $\ell(v_2)=-1$.
For $\tau_i(x)=M_ix+v_i/m_i$, every nonidentity power is free,
because its image under $\ell:T^6\to\mathbb R/\mathbb Z$ is
translation by a nonintegral multiple of $\pm1/m_i$.

\item The finite fillings
$(T^6\times\Delta)/\langle(\tau_i,\zeta_{m_i})\rangle$ are smooth.
Take the stipulated Kuhn filling at the third point, marked by
Proposition \ref{maps:src:F.2.1}.  The torus bundle over the intervening pair of
pants has monodromies $M_1,M_2,M_0$.  Translations are isotopic to
the identity, so its boundary bundles identify with those of the
three fillings.  Glue using the displayed twists and zero cusp
clutching.  This proves existence and compact smoothness.

\item For van Kampen, the normal closure of the cusp angular lattice is
already generated by the columns of
$[\,N\mid M_1N\mid M_2N\,]$.
Its Hermite basis is exactly $e_1,\ldots,e_5$.  The residual fibre
group is therefore the central group $\mathbb Z\langle h\rangle$,
and the remaining meridians satisfy
$5a=h$, $4b=-h$, $a+b=0$.
Their three-by-three relation matrix is unimodular.  Hence
$\pi_1=1$, with no uncomputed finite quotient.

\item The exact rational Leray input is
\[
\begin{array}{c|rrrrrrr}
q&0&1&2&3&4&5&6\\ \hline
h^0(j_*\bigwedge^q\mathcal V)&1&1&0&0&0&0&1\\
h^1(j_*\bigwedge^q\mathcal V)&0&0&0&3&0&0&0\\
h^2(j_*\bigwedge^q\mathcal V)&1&0&0&0&0&1&1
\end{array}
\]
where the cohomology monodromies are $A^{-1},B^{-1}$, whose
simultaneous invariant/coinvariant calculations agree with $A,B$.
The finite local invariant ranks are $(1,2,3,4,3,2,1)$ at both
points.  The cusp row is $(1,3,9,9,9,3,1)$, and Theorem \ExternalPaperRef{}{6.3}{III}
adds one degree-four cusp kernel.  The degree-one transgression
is an isomorphism since $\pi_1=1$.  Poincare duality forces the
degree-six transgression to be an isomorphism.  No other rational
differential is possible, giving \eqref{maps:eq:F.5}.

\item For the almost complex assertion, prescribe the existing complex
structures on the three fillings.  A constant linear complex
structure along each regular fibre is a section of the associated
bundle with fibre
$\operatorname{GL}_6^+(\mathbb R)/\operatorname{GL}_3(\mathbb C)$.
This fibre deformation retracts to
$\operatorname{SO}(6)/\operatorname{U}(3)\simeq\mathbb{CP}^3$.
It is connected and simply connected.  Relative obstruction
theory over the two-dimensional pair of pants therefore extends
the prescribed boundary sections.  A smooth section and compatible
horizontal distribution give an almost complex structure agreeing
with the fillings.  Rational $H^2,H^6$ vanish, so $c_1,c_3$
vanish rationally.  Theorem \S\ref{sec:V-known-inputs} and the unimodular lattice
classification give the signature and free form.

\item The exact program \texttt{fourfolds\_nonselfdual.py} verifies every
matrix, Smith form, order, and exterior rank above.  It verifies
this one tuple, not a bounded simultaneous-conjugacy classification.

\end{enumerate}
\end{proof}


\label{sec:fourfold}

We use column vectors.  The matrices $A,B$ in the supplied data act on
cohomology; $M_1=A^t$, $M_2=B^t$ act on
$\Lambda=\mathbb Z^6=H_1(T^6;\mathbb Z)$, and
$M_0=(M_1M_2)^{-1}$.  Positive based meridians satisfy
$g_1g_2g_0=1$.  The exterior-square basis of $\bigwedge^2\mathbb Z^4$
is $(01,02,03,12,13,23)$.  All assertions about an integral matrix below
are checked over $\mathbb Z$ or $\mathbb Q$, without floating-point
arithmetic.  The files in \texttt{code/part1/} contain the complete
matrices, including the changes of basis in every Smith computation.

\begin{theorem}\label{night:P1.1}
\StatementRecord{night:P1.1}{proved}{Historical events: GPT-6 Astra (OpenAI), ChatGPT Pro mode, effort Pro, S6\_06\_BLIND, 16--18 November 2026: targeted proof check. Also checked in this writing. Current writing events: GPT-6 Astra (OpenAI), ChatGPT Pro mode, effort Pro, S6\_10\_PAPERS, 18 November 2026: computation or formal certificate; targeted proof check. GPT-6 Astra (OpenAI), ChatGPT Pro mode, effort Pro, S6\_10\_PAPERS continuation, 19 November 2026: computation or formal certificate; independent re-proof; used and rechecked in this run. }{Definitions and the displayed constructions}
The finite homology monodromies have orders $5,4$, respectively, and
characteristic polynomials
$(T-1)^2\Phi _5(T)$ and $(T-1)^2(T+1)^2(T^2+1)$.
For $N=M_0-I$ we have $N^2=0$, $\operatorname{rank}N=3$, and
nonzero Smith factors $1,1,1$.  Moreover
\[
 \Lambda^{\langle M_1,M_2\rangle}=0,\qquad
 \Lambda_{\langle M_1,M_2\rangle}=\mathbb Z,
\]
where the quotient homomorphism is $\ell(x)=x_0$.
On the cohomology lattice the corresponding groups are $\mathbb Ze_0$
and $\mathbb Z/2$.  Every nonidentity power of each specified affine
action $\tau_i$ is free.
\end{theorem}
\begin{proof}
\begin{enumerate}[label=\arabic*.]
\item Exact multiplication gives
\[
N=\begin{pmatrix}
0&0&0&0&0&0\\0&0&0&1&0&-1\\0&0&0&0&1&1\\
-1&0&0&0&0&0\\1&0&0&0&0&0\\-1&0&0&0&0&0
\end{pmatrix},\qquad
S=\begin{pmatrix}
0&0&0&1&0&0\\0&1&0&0&0&0\\0&0&1&0&0&0\\
-1&0&0&0&1&0\\1&0&0&0&0&1\\-1&0&0&0&0&0
\end{pmatrix}.
\]
The matrix $S$ is unimodular and
$S^{-1}M_0S=\left(\begin{smallmatrix}I_3&I_3\\0&I_3\end{smallmatrix}\right)$.
The columns of $S$ are henceforth denoted
$(a_1,a_2,a_3,b_1,b_2,b_3)$, so that $Nb_j=a_j$.
This is also an integral certificate of the Smith assertion.
The Hermite basis of
$(M_1-I)\Lambda+(M_2-I)\Lambda$ is $(e_1,\ldots,e_5)$;
the stacked matrix of the two fixed-vector equations has rank six.
Transposing the matrices gives a fixed lattice $\mathbb Ze_0$ and
coinvariant Smith factors $(1,1,1,1,1,2)$.  The pairing
$\Lambda^\vee\times\Lambda\longrightarrow\mathbb Z$ identifies the
fixed cohomology vector $e_0$ with the covector $\ell$.

\item We have $M_iv_i=v_i$, $\ell(v_1)=1$, and $\ell(v_2)=-1$.
Thus $\tau_i^k(x)=M_i^kx+kv_i/m_i$, and $\tau_i^{m_i}=1$ on the
torus.  A fixed point of $\tau_i^k$ would give
$k\ell(v_i)/m_i=0$ in $\mathbb R/\mathbb Z$.  For
$0<k<m_i$ this is impossible.  The order and characteristic-polynomial
computations, and the full fixed-lattice calculations, are reproduced by
\texttt{tuple\_check.py}; its output records every power tested.

\end{enumerate}
\end{proof}




\begin{theorem}\label{night:P1.3}
\StatementRecord{night:P1.3}{proved}{Historical events: GPT-6 Astra (OpenAI), ChatGPT Pro mode, effort Pro, S6\_06\_BLIND, 16--18 November 2026: targeted proof check. Also checked in this writing. Current writing events: GPT-6 Astra (OpenAI), ChatGPT Pro mode, effort Pro, S6\_10\_PAPERS, 18 November 2026: computation or formal certificate; targeted proof check. GPT-6 Astra (OpenAI), ChatGPT Pro mode, effort Pro, S6\_10\_PAPERS continuation, 19 November 2026: computation or formal certificate; independent re-proof; used and rechecked in this run. }{Definitions and the displayed constructions}
The two finite filling pieces have first integral homology
$\mathbb Z^2$ and $\mathbb Z^2\oplus\mathbb Z/2$.
\end{theorem}
\begin{proof}
\begin{enumerate}[label=\arabic*.]
\item For $K_i=T^6/\langle\tau_i\rangle$, the universal covering deck group
has commuting translations $x^u$, $u\in\Lambda$, and a generator $a_i$,
with the complete relations
\[
a_ix^ua_i^{-1}=x^{M_iu},\qquad a_i^{m_i}=x^{v_i}.
\]
Indeed the affine lifts satisfy these relations, and they reduce every
word uniquely to $x^ua_i^r$, $0\leq r<m_i$.
The abelianization relation matrix has columns
$((M_i-I)e_j,0)$ and $(-v_i,m_i)$.
Smith reduction of these explicit seven-by-seven matrices gives the
groups above.  The disk quotient
$(T^6\times\Delta)/\langle(\tau_i,\zeta_{m_i})\rangle$
retracts onto $K_i$ by equivariant radial contraction.  Freeness supplies
ordinary smooth quotient charts by disjoint translates of small balls.

\end{enumerate}
\end{proof}


\begin{lemma}[The actual order-four torsion transgression]
\label{mapsaudit:order-four-torsion-transgression}
\StatementRecord{mapsaudit:order-four-torsion-transgression}{proved}{No earlier separately identified check is recorded. Checked only in this writing. Current writing events: GPT-6 Astra (OpenAI), ChatGPT Pro mode, effort Pro, S6\_10\_PAPERS, 18 November 2026: targeted proof check. GPT-6 Astra (OpenAI), ChatGPT Pro mode, effort Pro, S6\_10\_PAPERS continuation, 19 November 2026: used and rechecked in this run. \emph{Recorded proof scope.} GPT-6 Astra (OpenAI, ChatGPT Pro mode, effort Pro), S6\_09 writing, 17 November 2026; also checked in this writing by GPT-6 Astra (OpenAI), ChatGPT Pro mode, effort Pro, 18 November 2026}{\cref{maps:src:F.2.5}; the displayed cellular and boundary maps}

Use the order-four filling and the marked matrices of \cref{maps:src:F.2.5}.
Let $T=\bigwedge^{q-1}D^t$ on the integral vertical cochains of the
five-torus. If $w$ represents a torsion element of
$\operatorname{coker}(T-I)$, its point-supported degree-$q$ class has
Leray differential equal, up to the one boundary-orientation sign, to
the image of $w$ in the global degree-$(q-1)$ coinvariants. There is
no multiplicity-four factor. For each of $q=2,3,4$, that image is
the nonzero class in the global $\mathbb Z/2$ quotient.
\end{lemma}


\begin{proof}
\begin{enumerate}[label=\arabic*.]
\item In the original integral marking the required matrices are
\[
D=\begin{pmatrix}
0&0&0&1&0\\-1&-1&-1&-1&-1\\0&0&0&0&1\\
1&0&0&0&0\\0&1&0&0&0
\end{pmatrix},\qquad
P=\begin{pmatrix}
-1&0&-1&0&1\\-1&2&-1&1&3\\-1&1&-1&2&3\\
0&-1&0&-1&-1\\1&-1&2&-1&-3
\end{pmatrix}.
\]
Direct multiplication gives $\det P=\pm1$ and $DP=PC$, where
$C$ cyclically permutes the first four basis vectors and negates the
fifth. Thus $P$ changes $D$ to a four-cycle permutation of four circle
coordinates and inversion of the fifth.
Equip each circle with its one-vertex, one-edge cellular structure
and use the product structure. The resulting automorphism is cellular,
acts on degree-$j$ integral cochains by $\bigwedge^jD^t$, and its
fourth power is the identity as a map, not only up to homotopy.
The product cellular differential of the five-torus is zero.
Consequently the mapping-torus cellular cochain complex is the cone
of $D^*-I$: a suspension class in degree $q$ is represented by
the cochain $w$ on its base edge times the vertical
degree-$(q-1)$ cells, modulo $(T-I)$ of the vertical cochains.
This is an actual cellular representative of the local torsion
class; no splitting of a Wang exact sequence is being assumed.

\item The finite normal line is topologically trivial in the supplied
marked quotient, as follows directly from its affine action. Write
the invariant torus coordinate as $x_0\in\mathbb R/\mathbb Z$ and
use the generator
\[
 (x_0,y,z)\longmapsto(x_0-1/4,Dy,iz).
\]
The function $z_{\mathrm{tr}}=e^{2\pi i x_0}z$ is invariant, and so gives a
nowhere-zero smooth normal trivialization. On the boundary put
$a=-4x_0\bmod1$ and $k=\arg(z_{\mathrm{tr}})/(2\pi)\bmod1$.
The boundary is therefore $B_4\times S^1_k$, and the argument
of the original base coordinate $z^4$ is
\[
             \theta=a+4k.
\]
On the base two-torus choose the new integral generators
\[
 e_\theta=(1,0),\qquad e_h=(-4,1)
\]
in the old $(a,k)$ marking. Their determinant is one. The vertical
monodromy along $e_\theta$ is $T$, and along $e_h$ it is
$T^{-4}=I$. Thus the fibre over $\theta$ is the original
six-torus, with its five vertical directions and the primitive
circle $h$. In particular $e_\theta$ goes around the suspension
circle once; confusing it with a fourfold traverse would introduce
an incorrect factor.

\item The pullback of a cellular one-cochain with values in a local
system is computed by following its oriented edges and transporting
each value to the starting fibre. Applied to the word
$h=a^{-4}k$, this gives
\[
 (w,0)_{(a,k)}\longmapsto
 \left(w,-(T^{-1}+T^{-2}+T^{-3}+T^{-4})w\right)_{(\theta,h)}.
\tag{FT.1}\label{maps:eq:FT.1}
\]
One can also obtain this formula by writing the four successive
fundamental prisms: reversal of an $a$ edge contributes $-T^{-1}$,
and the next three reversals contribute $-T^{-2},-T^{-3},-T^{-4}$.
The final $k$ edge contributes zero. Since $T^4=I$, the sum in
\textup{\eqref{maps:eq:FT.1}} is $N_T=I+T+T^2+T^3$.

\item A torsion coinvariant has $w=(T-I)v$ for some rational vertical
cochain $v$. Hence
\[
 N_Tw=N_T(T-I)v=(T^4-I)v=0
\]
as an integral cochain. Formula \textup{\eqref{maps:eq:FT.1}} therefore says
that the restriction to the six-torus fibre is zero, while the
first boundary Wang component is exactly $[w]$. This calculation
is unaffected by the ambiguity $w\mapsto w+(T-I)b$: its change is
a coboundary on the suspension circle. The cellular construction
preserves vertical degree, so there is no omitted component in a
lower vertical degree. If the inverse monodromy convention is used,
the same argument replaces $w$ by $-T^{-1}w$, which is $-w$ in
coinvariants and hence the same element in the order-two target.

\item Use the two-open-cover \v{C}ech resolution for the base, with one
open set containing the finite centre and the other its complement,
and the cellular fibre resolution just specified. The local class
has zero fibre restriction on their annular intersection. Its next
component is the boundary one-cocycle $w$ computed above. The
connecting homomorphism identifies this with its class in degree-two
base cohomology. For the punctured sphere, that group is the global
coinvariant quotient: transporting a coefficient along the two
ordinary generators imposes exactly $(\bigwedge^{q-1}A-I)$ and
$(\bigwedge^{q-1}B-I)$. Thus the two-step differential of this
actual \v{C}ech--fibre complex is precisely
\[
 d_2(w)=\pm[w]\in
 \frac{\bigwedge^{q-1}\mathbb Z^6}
 {(\bigwedge^{q-1}A-I)\bigwedge^{q-1}\mathbb Z^6+
  (\bigwedge^{q-1}B-I)\bigwedge^{q-1}\mathbb Z^6}.
\]
The finite modifications of stalks do not alter this degree-two
base cohomology: the kernel and cokernel sheaves of the comparison
with the invariant-extension sheaf are supported at finitely many
points and have no positive cohomology.

\item For completeness the supplied local representatives, in increasing
lexicographic exterior coordinates of the vertical five-lattice, are
\[
\begin{array}{c|l}
q=2&(0,0,0,-1,-2),\\
q=3&(0,0,0,0,0,0,0,1,0,-1),\\
q=4&(0,0,0,0,0,0,1,0,0,0).
\end{array}
\]
For each row the norm is zero; twice the vector lies in the integral
image of $T-I$, and the vector itself does not. Its embedding in
the six-lattice fills the coordinate slots not containing the first
basis vector and is zero on the other slots. The global relation
matrix has Smith factors all one except for a single two; adjoining
that embedded vector changes every Smith factor to one. It follows
that its residue generates the global $\mathbb Z/2$. These are
independent integral lattice tests in the accompanying checker, in
addition to the original exact Smith computation. The geometric
identification of the differential was proved above and is not an
inference from the arithmetic tests. The three torsion differentials
of \ref{maps:src:F.2.6} are therefore isomorphisms.

\end{enumerate}
\end{proof}




\begin{theorem}[integral three-connectivity of the candidate]\label{maps:src:F.2.6}
\StatementRecord{maps:src:F.2.6}{proved}{Historical events: Claude Fable 5.2 (Anthropic), effort high, 16 November 2026: carried as background without fault found. Claude Fable 5.2 (Anthropic), effort high, 17 November 2026: carried as background without fault found. GPT-6 Astra (OpenAI), ChatGPT Pro mode, effort Pro, S6\_09\_PAPERS, 17 November 2026: targeted proof check. Also checked in this writing. Current writing events: GPT-6 Astra (OpenAI), ChatGPT Pro mode, effort Pro, S6\_10\_PAPERS, 18 November 2026: targeted proof check. GPT-6 Astra (OpenAI), ChatGPT Pro mode, effort Pro, S6\_10\_PAPERS continuation, 19 November 2026: used and rechecked in this run. \emph{Recorded proof scope.} GPT-6 Astra (OpenAI, ChatGPT Pro mode), effort Pro, 2026-11-17; Claude Fable 5.2 (Anthropic), effort high, 2026-11-17; Claude Fable 5.2 (Anthropic), effort high, 2026-11-16; GPT-6 Astra (OpenAI, ChatGPT Pro mode), S6\_06\_BLIND, effort Pro, 2026-11-17; also checked in this writing by GPT-6 Astra (OpenAI), ChatGPT Pro mode, effort Pro, 18 November 2026}{\cref{mapsaudit:order-four-torsion-transgression}; the hypotheses, notation and definitions stated here; \cref{maps:src:F.2.5}; \cref{maps:eq:F.5}; \ExternalPaperRef{theorem~}{6.3}{III}; \cref{maps:eq:F.6}}


The manifold of Theorem \ref{maps:src:F.2.5} has free integral homology with
Betti numbers \eqref{maps:eq:F.5}.  In particular $H_2=H_3=0$, and it is
three-connected.

\end{theorem}


\begin{proof}
\begin{enumerate}[label=\arabic*.]
\item We identify the integral torsion differentials explicitly.
For the finite fillings, Smith forms give
$H_1(B_5)=\mathbb Z^2$ and
$H_1(B_4)=\mathbb Z^2\oplus\mathbb Z/2$.
The covering images on integral degree-two and degree-three
cohomology have index one.  In degree two use
$B_m\to S^1$, given by $m\ell$, with fibre $T^5$.
A primitive invariant vertical loop $\lambda$ supplies the
remaining suspension two-cycle.  An invariant integral
two-form $\omega$ descends exactly when
$\omega(v,\lambda)/m$ is integral; vertical periods are
already integral.  These three values on integral invariant
bases are $(-1,0,-1)$ at five and $(0,0,1)$ at four.
In degree three the determinants of the invariant middle
skew pairings are $5^4$ and $4^4$.  Pullback multiplies the
unimodular downstairs pairing by $m$, so the covering indices
are one.  The degree-two/degree-four invariant pairing has
determinant of absolute value $5^3$, respectively $4^3$.
Since the degree-two covering index is one, the degree-four
index is one as well.  Theorem \ExternalPaperRef{}{6.3}{III} gives integral
surjectivity at the cusp.

\item For coefficient exterior degrees two, three and four, the
global coboundary matrix obtained by stacking
$\bigwedge^qA-I$ and $\bigwedge^qB-I$ has one Smith factor
two and all its other nonzero Smith factors one.
Its torsion cocycle restricts nontrivially at the order-four
point.  Thus the integral parabolic groups in these degrees
are torsion free, of ranks $0,3,0$.  In degree one the global
coboundary image is primitive; the finite covering cokernels
are $\mathbb Z/5,\mathbb Z/4$, and the invariant covector
$\ell$ maps to a generator of each.  The Chinese remainder
theorem therefore gives $H^1(B,R^1g_*\mathbb Z)=0$.
The global coinvariants in coefficient degrees one through
four are each $\mathbb Z/2$.

\item For the crucial torsion differentials, write
$M_2=\operatorname{diag}(1,D)$.  The unimodular matrix
\[
P=\begin{pmatrix}
-1&0&-1&0&1\\-1&2&-1&1&3\\-1&1&-1&2&3\\
0&-1&0&-1&-1\\1&-1&2&-1&-3
\end{pmatrix}
\]
conjugates $D$ to a four-cycle permutation and a negative
one-dimensional summand.  Hence $B_4$ is a mapping torus of
a cellular torus automorphism.  Its torsion class in degree
$q$ is represented by
$w\in\operatorname{Tor}\operatorname{coker}
(\bigwedge^{q-1}D^t-I)$.
The boundary is $B_4\times S^1$ with original base coordinate
$\theta=a+4k$.  Write $h=k$ on its finite-cover fibre; then
$a=-4h$.  The fibre has the vertical $T^5$ and the loop
$(-4,1)$ in the $(a,k)$ coordinates.
The restriction of the suspension class to that fibre is
the norm $(I+D^t+(D^t)^2+(D^t)^3)w$, which is zero because
a torsion coinvariant representative lies rationally in
$\operatorname{im}(D^t-I)$.
The remaining component in the boundary Wang sequence is
exactly the original vertical coinvariant $w$.
This description follows on cellular chains by taking the
four consecutive fundamental prisms; the monomial marking
makes their boundary the displayed norm.

\item The base Mayer--Vietoris connecting map therefore identifies
the point-supported torsion Leray differential with
\[
d_2(w)=[w]\in
\frac{\bigwedge^{q-1}\mathbb Z^6}
 {(\bigwedge^{q-1}A-I)\bigwedge^{q-1}\mathbb Z^6+
  (\bigwedge^{q-1}B-I)\bigwedge^{q-1}\mathbb Z^6}.
\tag{F.6}\label{maps:eq:F.6}
\]
Indeed fibre restriction vanishes, so its first nonzero
boundary restriction is the just-computed class in
$H^1(S^1,\bigwedge^{q-1}\mathcal V)$.
The connecting map quotients by the two global monodromy
relations, which is the two-step differential of the
Cech--fibre double complex.
Changing the sheaf at finitely many points by its finite
covering cokernels or point-supported kernels does not
change this degree-two base cohomology group.
Inverse monodromy conventions change the representatives
by units and do not change the nonzero classes below.

\item The exact Smith transformations in the code give the
following vertical torsion representatives, in lexicographic
exterior bases of $\mathbb Z^5$:
\[
\begin{split}
q=2:\ &(0,0,0,-1,-2),\\
q=3:\ &(0,0,0,0,0,0,0,1,0,-1),\\
q=4:\ &(0,0,0,0,0,0,1,0,0,0).
\end{split}
\]
Each has residue one in its global $\mathbb Z/2$ quotient.
Thus all three differentials in \eqref{maps:eq:F.6} are isomorphisms.
The relevant integral Leray groups are
\[
\begin{array}{c|ccccc}
q&0&1&2&3&4\\ \hline
H^0(R^q)&\mathbb Z&\mathbb Z&\mathbb Z/2&
 \mathbb Z/2&\mathbb Z\oplus\mathbb Z/2\\
H^1(R^q)&0&0&0&\mathbb Z^3&0\\
H^2(R^q)&\mathbb Z&\mathbb Z/2&\mathbb Z/2&
 \mathbb Z/2&\mathbb Z/2 .
\end{array}
\]
The free degree-one transgression is onto because $\pi_1=1$
and rational $H^2=0$: the universal coefficient theorem makes
$H^2$ torsion free, so the finite cokernel of this map
$\mathbb Z\to\mathbb Z$ must vanish.  Together with the three torsion
isomorphisms this gives $H^2=H^3=0$ and $H^4=\mathbb Z^4$;
the two free filtration pieces of $H^4$ have ranks three
and one.  The universal coefficient theorem now gives
$H_2=H_3=0$.  Hurewicz gives three-connectivity, and Poincare
duality determines all remaining integral groups.

\end{enumerate}
\end{proof}




\begin{theorem}\label{night:P1.10}
\StatementRecord{night:P1.10}{proved}{Historical events: GPT-6 Astra (OpenAI), ChatGPT Pro mode, effort Pro, S6\_06\_BLIND, 16--18 November 2026: targeted proof check. Also checked in this writing. Current writing events: GPT-6 Astra (OpenAI), ChatGPT Pro mode, effort Pro, S6\_10\_PAPERS, 18 November 2026: targeted proof check. GPT-6 Astra (OpenAI), ChatGPT Pro mode, effort Pro, S6\_10\_PAPERS continuation, 19 November 2026: used and rechecked in this run. }{\cref{night:P1.11}, \cref{night:P1.9}}
For every real parameter $c>-h_{\min}$ the corrected construction gives
a compact smooth complex fourfold $M_c$ and a proper holomorphic map
to $\mathbb P^1$ with the two specified multiple fibres and the
index-one Kuhn cusp.  With its complex orientation,
$\pi_1(M_c)=1$ and $\chi(M_c)=6$.
\end{theorem}
\begin{proof}
\begin{enumerate}[label=\arabic*.]
\item Evaluation of $\Lambda$ on $F_c$ maps $\Lambda_{\mathbb R}$
isomorphically onto the underlying real bundle of $F_c^\vee$:
the kernel is zero exactly by $F_c\cap\bar F_c=0$.
The image is consequently a full discrete lattice.  Quotient by its
holomorphic period sections gives a holomorphic torus bundle on the
nonspecial base.

\item On a finite root chart $u^{m_i}=t_i$, divide the pullback torus bundle by
\[
(x,u)\longmapsto(M_ix+v_i/m_i,\zeta_{m_i}u).
\]
Freeness proves smoothness, including the central fibre.  The special
divisor is $m_i$ times the reduced free quotient.  On the punctured
root disk translation by
$-(\log u/(2\pi i))v_i$, interpreted by its period section, identifies
this affine quotient with the original Jacobian bundle.  A change of
logarithm gives an integral lattice period; replacing $u$ by
$\zeta_{m_i}u$ cancels exactly the affine term $v_i/m_i$.
This proves the holomorphic gluing and its sign.
At the cusp use P1.9 and the exponential period map.  The zero section
$\zeta=0$ extends to $z=1$ in the dense torus of the component marked
by the vertex zero.  The resulting map is proper by a finite compact
cover of the base by the three local pieces and a compact remainder
with torus fibres.  All total-space charts and their overlaps have
already been specified, so this produces a smooth compact complex
fourfold, and in particular an almost complex structure.

\item For completeness we record every van Kampen relation.  The main
bundle has its zero section and fundamental group
$\Lambda\rtimes\langle g_1,g_2\rangle$.
The finite pieces impose
$g_i^{m_i}=v_i$ and $g_ix^ug_i^{-1}=x^{M_iu}$; the logarithmic gluing
above supplies these signs.  The cusp zero section kills
$g_0=(g_1g_2)^{-1}$ and the cusp collapse kills $\operatorname{im}N$.
The latter assertion follows directly from the toric moment model in
P1.11: angular cycles $a_j$ collapse and the base cycles $b_j$ survive.
Its fundamental group is $\mathbb Z^3$, since the simply connected
complete toric components and their connected toric intersections have
dual complex $\mathbb R^3/\mathbb Z^3$; alternatively, van Kampen on
the moment cells kills the angular loops generated by the primitive
facet normals and retains exactly the three base loops.
Thus the full presentation is
\[
\left\langle \Lambda,g_1,g_2\ \middle|\
 [\Lambda,\Lambda]=1,\quad
 g_ix^ug_i^{-1}=x^{M_iu},\quad
 g_1^5=v_1,\ g_2^4=v_2,\ g_1g_2=1,\ \operatorname{im}N=0
\right\rangle.
\]
The last relation is redundant after the preceding ones but records
the actual cusp map.  Since $g_2=g_1^{-1}$, the two induced actions are
inverse; their orders divide both five and four and are therefore one.
The remaining fibre group is $\mathbb Z$ through $\ell$.  With its
generator $z$, the additive relations are
$5g_1-z=0$, $4g_2+z=0$, $g_1+g_2=0$.
Their matrix
\[
\begin{pmatrix}5&0&-1\\0&4&1\\1&1&0\end{pmatrix}
\]
has determinant $-1$, proving triviality.
Euler additivity gives $\chi(M_c)=6$: smooth tori and finite free
torus quotients have Euler number zero, and P1.9 gives the remaining
contribution.

\end{enumerate}
\end{proof}




\begin{theorem}\label{night:P1.11}
\StatementRecord{night:P1.11}{proved}{Historical events: Claude Opus 5.5 (Anthropic), effort max, sweep 1, scope s6, 18 November 2026: computation or formal certificate. GPT-6 Astra (OpenAI), ChatGPT Pro mode, effort Pro, S6\_06\_BLIND, 16--18 November 2026: targeted proof check. Also checked in this writing. Current writing events: GPT-6 Astra (OpenAI), ChatGPT Pro mode, effort Pro, S6\_10\_PAPERS, 18 November 2026: targeted proof check. GPT-6 Astra (OpenAI), ChatGPT Pro mode, effort Pro, S6\_10\_PAPERS continuation, 19 November 2026: used and rechecked in this run. }{\cref{night:P1.9}}
The index-one central cusp fibre $C$ has torsion-free integral homology
with ranks $(1,3,9,9,10,3,1)$.
The collapse $\kappa:T^6\to C$ has primitive image in every degree;
it is onto except in degree four, where its cokernel is $\mathbb Z$.
Its image ranks are $(1,3,9,9,9,3,1)$.
In every degree its kernel is the saturation of
$(\bigwedge^kM_0-I)\bigwedge^k\Lambda$.
\end{theorem}
\begin{proof}
\begin{enumerate}[label=\arabic*.]
\item We first remove the unit matrix $C(q)$ topologically, by an explicit
deformation, without changing the integral marking.  Replace it by
$\theta C(q)$, $0\leq\theta\leq1$.  The action and properness estimates
of P1.9 are uniform in $\theta$.  Local quotient charts form a smooth
family of pairs, with the local equation $q=x_0x_1x_2x_3$ independent
of $\theta$.  Local lifts of $\partial/\partial\theta$, tangent to every
coordinate stratum and annihilating $q$, can be patched by a partition
of unity.  Over a smaller compact disk their flow exists for the whole
interval.  This gives a stratified diffeomorphism of the central
fibres and a diffeomorphism of the nearby fibres, preserving the
period marking.  The same construction on the real oriented blowup
preserves the collapse map.  We may therefore calculate at $C(q)=0$.

\item Here is a concrete topological model.  Put $G=4I_3-\mathbf1\mathbf1^t$.
Its positive eigenvalues are $1,4,4$.  The Voronoi polytope in dual
coordinates is cut out by
\[
|v^ty|\leq \tfrac12v^tGv
 \quad (0\ne v\in\{0,1\}^3).
\]
The value on a binary vector with $k$ nonzero entries is $k(4-k)$.
These fourteen primitive facet normals give the star of a vertex in
the Kuhn triangulation.  One checks its Delaunay tetrahedra by the
successive chains
$0,e_i,e_i+e_j,e_1+e_2+e_3$ for the six permutations $(i,j,k)$;
solving their three perpendicular-bisector equations and the fourteen
inequalities verifies precisely their common faces.  The translates
of this Voronoi polytope by $G\mathbb Z^3$ tile $\mathbb R^3$.

\item For a simplex $\sigma$ let $D_\sigma$ be its dual cell in this tiling,
and let
$L_\sigma=\operatorname{span}_{\mathbb Z}\{v-v_0:v\in\sigma\}$.
The compact toric component with this normal fan is the quotient of
its moment polytope times the angular torus by the subtori generated
by the facet normals that meet a given face.  This description follows
on each affine chart by writing its coordinates in polar form;
at a zero coordinate its circle argument is forgotten.  Gluing the
faces and then taking the period quotient gives a continuous
homeomorphism from the compact quotient
\[
\mathcal Q=
\bigl((\mathbb R^3/G\mathbb Z^3)\times
       (\mathbb R^3/\mathbb Z^3)\bigr)/\!\sim\ \longrightarrow C,
\]
where over the relative interior of $D_\sigma$ the equivalence
relation collapses the subtorus $L_\sigma\otimes\mathbb R/\mathbb Z$.
The map is bijective by the polar-coordinate description on each
toric stratum, so compactness and Hausdorffness make it a homeomorphism.
Its angular identity defines a continuous section of the projection to
the base three-torus.

\item The actual nearby-fibre collapse is the quotient map in this model.
To see the comparison, perform the real oriented blowup of the
normal-crossings divisor in each chart.  Retain the arguments of all
$x_j$, with their sum fixed by $\arg q$, and then forget the arguments
of coordinates whose radii are zero.  At a face the forgotten
subtorus is exactly $L_\sigma$.  The uncollapsed arguments and the
dual moment coordinates give the product of the angular and base
three-tori above.  In a chart the passage to a small smooth fibre is
obtained by increasing the nonnegative radii continuously so their
product is $|q|$, keeping the arguments and $\arg q$ fixed; this is
invertible after fixing the transverse radial parameter.  On common
faces the integral monomial transformations give exactly the same
angular quotient.  Patching the radial parameters by positive
partitions of unity gives the global identification, while forgetting
the arguments is unchanged.  Thus this comparison fixes the map on
the marked angular cycles $a_j$ and base cycles $b_j$, not merely
the Betti numbers of a putative central fibre.

\item Filter $\mathcal Q$ by the dimension $d$ of its dual base cells.
Over the interior of the dual cell of a $p$-simplex, $d=3-p$ and the
angular quotient has lattice
$Q_\sigma=\mathbb Z^3/L_\sigma$ of rank $d$.
Excision for the cell and its boundary gives the relative chain
groups of the resulting homology spectral sequence:
\[
E^1_{d,q}=\bigoplus_{\dim\sigma=3-d}\bigwedge^qQ_\sigma.
\]
Orient dual cells by the orientation of the base and the ordered
simplex.  The differential is the signed quotient map for each
codimension-one coface.  More explicitly, with integral projection
matrices $Q_\sigma$ and chosen right inverses $S_\sigma$,
\[
D_{p,q}[\tau,\sigma]=
 \sum_{\sigma=\partial_j\tau\ {\rm mod}\ \mathbb Z^3}
 (-1)^j\bigwedge^q(Q_\tau S_\sigma).
\tag{1.11.1}
\]
All $L_\sigma$ are primitive because the tetrahedra are unimodular.
Thus these projections and right inverses are integral.  Formula
(1.11.1) specifies every integer entry, including multiple incidences
under translation identifications.  The program
\texttt{cusp\_complex.py} enumerates the $26$ simplex classes, all
$74$ codimension-one incidences, and exports every matrix, basis, and
Smith certificate to \texttt{cusp\_complex.json}.

\item The nonzero $E^2_{d,q}$ ranks are
\[
\begin{array}{c|rrrr}
q\backslash d&0&1&2&3\\\hline
0&1&3&3&1\\
1&0&6&8&3\\
2&0&0&7&3\\
3&0&0&0&1
\end{array}
\]
and every group is free.  Support leaves only the possible higher
differential $d_2:E^2_{3,0}\to E^2_{1,1}$.
The angular identity section splits the base projection as a map of
filtered spaces, so the base fundamental class survives.  This kills
that differential.  All subsequent differentials have zero source or
target.  The filtration extensions are extensions of free abelian
groups and therefore split, proving the asserted integral homology.

\item The product-to-quotient map induces on $E^1$ the exterior powers of
the quotient projections.  The program \texttt{cusp\_collapse.py}
computes the induced maps on the above Smith bases; all cokernels are
free, zero except for one generator in total degree four.
Surjectivity and primitivity transfer from the associated graded
groups through the finite filtration: successively lift a class from
each filtration quotient, and note that a filtration with free
quotients has no torsion.  This proves the corresponding assertions
for the actual collapse, including its integral index.

\item For a fully marked certificate we represent a base cycle $b_J$ by
Poincar\'e duality from the Alexander--Whitney cup cocycle on the
complementary coordinates.  On an ordered $p$-simplex its coefficient
is
\[
\operatorname{sgn}(J,J^c)
 \prod_{r=1}^p(v_r-v_{r-1})_{(J^c)_r}.
\]
Tensoring with an angular exterior vector and applying (1.11.1)
gives the matrices exported by \texttt{cusp\_marked\_collapse.py}.
They annihilate $\bigwedge^kM_0-I$ exactly.  The actual collapse also
annihilates that operator, since turning the smoothing parameter
does not change its image in the central fibre.  The computed ranks
equal the ranks of the monodromy coinvariants; primitivity therefore
identifies both kernels with the saturation of its image.
Consequently the exported matrices are matrices of the actual
collapse in suitable integral target bases.  This last argument also
eliminates any ambiguity from choosing splittings of the filtration.

\end{enumerate}
\end{proof}




\begin{theorem}\label{night:P1.12}
\StatementRecord{night:P1.12}{proved}{Historical events: GPT-6 Astra (OpenAI), ChatGPT Pro mode, effort Pro, S6\_06\_BLIND, 16--18 November 2026: targeted proof check. Also checked in this writing. Current writing events: GPT-6 Astra (OpenAI), ChatGPT Pro mode, effort Pro, S6\_10\_PAPERS, 18 November 2026: targeted proof check. GPT-6 Astra (OpenAI), ChatGPT Pro mode, effort Pro, S6\_10\_PAPERS continuation, 19 November 2026: used and rechecked in this run. }{\cref{night:P1.11}}
The fourfold $M_c$ has torsion-free integral homology
\[
H_k(M_c;\mathbb Z)=
\begin{cases}
\mathbb Z,&k=0,8,\\
\mathbb Z^4,&k=4,\\
0,&k=1,2,3,5,6,7.
\end{cases}
\]
In particular it is three-connected.
\end{theorem}
\begin{proof}
\begin{enumerate}[label=\arabic*.]
\item We compute the complement $Y$ of the open cusp cap, then attach the
cusp using only the degrees needed to prove three-connectivity.
This avoids assigning a value to a boundary suspension that is not
determined by the fibre specialization alone.

\item Write $L_k=\bigwedge^k\Lambda$ and $W_{i,k}=\bigwedge^kM_i$.
The cellular model of the main torus bundle over a graph with two
oriented edges has
\[
G_k=L_k\oplus L_{k-1}\oplus L_{k-1},\qquad
d_G(f,x_1,x_2)=((W_{1,k-1}-I)x_1+(W_{2,k-1}-I)x_2,0,0).
\tag{1.12.1}
\]
The boundary mapping torus at $i$ has
$B_{i,k}=L_k\oplus L_{k-1}$ and
\[
d_{B_i}(f,x)=((W_{i,k-1}-I)x,0).
\tag{1.12.2}
\]
The finite boundary maps to $G$ by identity on $L_k$ and inclusion
of the suspension into its respective edge.  The cusp boundary,
corresponding to $(g_1g_2)^{-1}$, maps by
\[
j_{0,G}(f,x)=(f,-W_{2,k-1}W_{0,k-1}x,-W_{0,k-1}x).
\tag{1.12.3}
\]
This follows by traversing the second edge backwards and then the
first backwards, transporting the fibre chain to each initial end.
Indeed its boundary is $W_{0,k-1}-I$ because
\[
-(W_1-I)W_2W_0-(W_2-I)W_0=W_0-I.
\]

\item Here are full models for the finite pieces and their attaching maps.
Put $\varepsilon_i=\ell(v_i)$ and
$P_i=(\varepsilon_iv_i,e_1,\ldots,e_5)$.
This is unimodular.  It splits
$\Lambda=\mathbb Z(\varepsilon_iv_i)\oplus K$, $K=\ker\ell$,
and $P_i^{-1}M_iP_i=1\oplus R_i$.
The cap retracts onto the mapping torus of $R_i$ on $T_K$.
Its chains are
\[
E_{i,k}=\bigwedge^kK\oplus\bigwedge^{k-1}K,
\qquad d_{E_i}(f,x)=((\bigwedge^{k-1}R_i-I)x,0).
\tag{1.12.4}
\]
In these split coordinates, with $v=\varepsilon_iv_i$, the finite
boundary-to-cap chain map is specified on all summands by
\[
\begin{array}{c|c}
\text{boundary chain}&\text{cap chain}\\\hline
x\in\bigwedge^kK&(x,0)\\
v\wedge x&(0,\varepsilon_i\sum_{j=0}^{m_i-1}
                  (\bigwedge^{k-1}R_i)^jx)\\
g_i\wedge x&(0,x)\\
g_i\wedge v\wedge x&0.
\end{array}
\tag{1.12.5}
\]
Changing back uses
$\operatorname{diag}(\bigwedge^kP_i^{-1},
\bigwedge^{k-1}P_i^{-1})$.
Geometrically this map is
$(t,y,\theta)\mapsto[y,\theta+\varepsilon_im_it]$ into the mapping
torus.  The $t$-circle sweeps $m_i$ strips, giving the norm sum; the
$g_i$-circle sweeps one strip; their product maps to a single
base direction, giving zero on the last row.  This proves the map
including its sign and the possible finite-cap torsion.

\item We justify that these are integral cellular models, rather than a
formal use of invariant rational forms.  In the bases with columns
\[
K_1=\begin{pmatrix}
-1&-1&0&-1&0\\0&1&0&1&1\\0&1&0&0&0\\
0&-1&-1&0&-1\\0&0&1&0&0
\end{pmatrix},\qquad
K_2=\begin{pmatrix}
-1&0&-1&0&-1\\-1&1&0&1&-1\\0&1&-1&1&-1\\
0&-1&0&-1&1\\1&-1&1&0&1
\end{pmatrix},
\]
both unimodular, $R_1$ is the permutation of five coordinates in a
cycle, and $R_2$ is a four-cycle together with multiplication by $-1$
on the fifth coordinate.  The minimal product cell structure of
circles is preserved by these signed permutations; its cellular
boundary is zero and its action on cells is exactly the exterior
power.  Formula (1.12.5) is obtained by subdividing its swept
base interval into the indicated strips.  The common torus can be
marked into these two structures by based cellular approximations.
Their cellular maps are the exterior-power matrices because the
cellular differentials vanish.  Based maps between these aspherical
tori with a fixed map on $\pi_1$ have a contractible space of
homotopies: lifting to the universal covers gives equivariant maps
to a contractible vector space and linear interpolation relative to
the base point.  Thus the choices introduce no extra attaching class
on the graph.  Composing the two edge cylinders gives exactly
(1.12.3).  The explicit conjugators were found in the bounded search
$[-1,1]^5$ and then verified by exact multiplication in
\texttt{finite\_permutation\_models.py}.

\item Let $F_k:B_{1,k}\oplus B_{2,k}\to G_k\oplus E_{1,k}\oplus E_{2,k}$
have the main boundary maps as its first row, and the negatives of
(1.12.5) in the corresponding two cap rows.  The chain complex of
$Y$ is the mapping cone
\[
C_k(Y)=G_k\oplus E_{1,k}\oplus E_{2,k}
       \oplus B_{1,k-1}\oplus B_{2,k-1},\qquad
d_Y=\begin{pmatrix}d_G\oplus d_{E_1}\oplus d_{E_2}&F_{k-1}\\
0&-d_{B_1}\oplus-d_{B_2}\end{pmatrix}.
\tag{1.12.6}
\]
Equations (1.12.1)--(1.12.6), minors for exterior powers, and the input
matrices specify every entry of every differential.
The full arrays are exported to \texttt{complement\_export.json}; its
program verifies $d^2=0$ in all degrees and computes Smith
decompositions.  Chain ranks in degrees zero through eight are
$(3,22,71,132,155,118,57,16,2)$.
The resulting homology of $Y$ is free, with ranks
$(1,1,3,9,12,9,3,0,0)$.

\item The boundary $B_0$ has
\[
H_1(B_0)=\mathbb Z^4,\qquad H_2(B_0)=\mathbb Z^{12},\qquad
H_3(B_0)=\mathbb Z^{18}\oplus\mathbb Z/2.
\]
The homomorphism $j_Y:H_k(B_0)\to H_k(Y)$ from (1.12.3) is
surjective for $k=1,2,3$.  In degrees two and three, the program
\texttt{complement\_kernel\_check.py} supplies integral cycle bases
of its kernel whose entire suspension block in (1.12.2) is zero.
This is an exact zero matrix, not a vanishing after tensoring with
$\mathbb Q$.  Therefore these kernels are represented by cycles of
the nearby torus itself.  The restriction of the actual map to the
cusp cap on these kernels is consequently the already proved
collapse of P1.11, regardless of its value on other boundary
suspension cycles.

\item For transparency, in the Smith bases exported with these cycles the
combined map in degree one is
\[
\begin{pmatrix}20&1&0&0\\1&0&-1&-1\\1&0&0&0\\0&0&2&1\end{pmatrix},
\qquad\det=1.
\tag{1.12.7}
\]
The cusp meridian maps to zero here because the zero section extends.
In degree two $\ker j_Y=\mathbb Z^9$ and its map to $H_2(C)$ is
\[
\begin{pmatrix}
0&0&1&-1&0&-1&0&0&0\\
0&0&1&0&0&-1&0&1&0\\
0&0&1&-1&0&-2&0&1&0\\
0&0&0&0&-1&-1&0&0&-1\\
1&1&-1&0&0&1&1&0&0\\
1&2&0&0&1&0&0&0&0\\
-1&-1&0&0&0&0&0&0&0\\
0&0&0&0&1&0&0&0&0\\
0&0&0&0&0&-1&0&0&0
\end{pmatrix},\qquad\det=1.
\tag{1.12.8}
\]
In degree three $\ker j_Y=\mathbb Z/2\oplus\mathbb Z^9$.
Its collapse matrix, with the torsion generator first, is
\[
\begin{pmatrix}
0&1&0&1&0&0&0&0&0&1\\
0&1&-1&0&0&0&-1&0&-1&0\\
0&0&0&0&0&0&1&0&1&0\\
0&0&0&-1&1&-1&0&-1&0&0\\
0&0&0&0&0&0&0&1&0&0\\
0&0&0&0&0&0&1&0&0&0\\
0&0&0&0&0&1&0&1&0&0\\
0&0&1&0&1&-2&1&-1&2&0\\
0&1&0&-1&1&-2&0&-1&0&0
\end{pmatrix}.
\tag{1.12.9}
\]
The determinant of its last nine columns is one.  In particular the
torsion of the boundary and the second finite cap has been included
throughout the calculation.

\item We now apply Mayer--Vietoris to $M_c=Y\cup C_{\rm cap}$.
Because $j_Y$ is onto, subtracting a preimage of the $Y$-coordinate
identifies the cokernel of $(j_Y,j_C)$ with
$H_k(C)/j_C(\ker j_Y)$; its kernel is
$\ker j_Y\cap\ker j_C$.
Equations (1.12.7)--(1.12.9) therefore say that the combined map is an
isomorphism in degrees one and two and is onto in degree three.
The exact sequence yields $H_2(M_c)=H_3(M_c)=0$.
We already proved $\pi_1(M_c)=0$, and successive applications of
Hurewicz now give $\pi_2(M_c)=\pi_3(M_c)=0$.

\item Integral Poincar\'e duality and the universal coefficient theorem
give $H_5=H_6=H_7=0$.  They also show that $H_4$ is free: its possible
torsion would be $\operatorname{Ext}(H_4,\mathbb Z)$ in $H^5$, which
vanishes by duality with $H_3=0$.
Finally $\chi=6$ gives $\operatorname{rank}H_4=4$.
This proves all claimed integral homology groups without making an
assumption about the higher-degree cusp suspension chain term.

\end{enumerate}
\end{proof}




\begin{theorem}\label{night:P1.19}
\StatementRecord{night:P1.19}{proved}{Historical events: Claude Opus 5.5 (Anthropic), effort max, sweep 1, scope s6, 18 November 2026: computation or formal certificate. GPT-6 Astra (OpenAI), ChatGPT Pro mode, effort Pro, S6\_06\_BLIND, 16--18 November 2026: targeted proof check. Also checked in this writing. Current writing events: GPT-6 Astra (OpenAI), ChatGPT Pro mode, effort Pro, S6\_10\_PAPERS, 18 November 2026: computation or formal certificate; targeted proof check. GPT-6 Astra (OpenAI), ChatGPT Pro mode, effort Pro, S6\_10\_PAPERS continuation, 19 November 2026: used and rechecked in this run. }{\cref{night:P1.11}, \cref{night:P1.12}, \cref{night:P1.9}}
There is an independent computation from the full global cellular
complex.  Its chain ranks in degrees zero through eight are
\[
 (4,26,87,162,200,156,79,23,3),
\]
and its differentials, after the explicit integral chain equivalences
below, are the matrices of case one in
\path{code/part1/global_ext_cases.json}.
Their integral Smith calculation gives
\[
 H_k(M;\mathbb Z)=
 \begin{cases}
 \mathbb Z,&k=0,8,\\
 \mathbb Z^4,&k=4,\\
 0,&\text{otherwise}.
 \end{cases}
\]
This calculation uses local Poincar\'e--Lefschetz duality of the
smooth cusp cap to determine its secondary attaching class.  It does
not use global Poincar\'e duality or the homology conclusion of P1.12.
In particular, the map obtained by setting the whole cusp suspension
chain term to zero is not the actual relative cusp model; its precise
missing term is accounted for here.
\end{theorem}
\begin{proof}
\begin{enumerate}[label=\arabic*.]
\item We retain the actual complement complex and finite-cap maps of P1.12.
The issue left by that argument was the full cusp boundary map.
We first determine its homology maps, then its secondary chain term,
and finally give the explicit integral reduction used by the program.

\item Write the cusp lattice as $A\oplus B$, with ordered generators
$a_1,a_2,a_3,b_1,b_2,b_3$ and $Nb_i=a_i$.
The boundary is the mapping torus of
\[
 (b,a)\longmapsto(b,a+b)
\]
on $T_B^3\times T_A^3$.
For the pure toric model, its retraction to the central moment model
has the following explicit description.  On the dual cell of a Kuhn
simplex $\sigma$, choose any vertex $v$ of $\sigma$ and put
\[
 (b,a,\theta)\longmapsto
   [b,a+\theta(b-v)],\qquad0\leq\theta\leq1.
\tag{1.19.1}
\]
The brackets denote the angular quotient by $L_\sigma$.
Changing $v$ changes the displayed angular coordinate by a real
multiple of a vector in $L_\sigma$, so the map is unchanged.
Replacing $b$ by $b+n$ also replaces $v$ by $v+n$; hence the formula
is periodic.  At $\theta=1$ the integer translation $v$ disappears
on the angular torus, and the map agrees with the cusp monodromy
identification.  These observations prove that (1.19.1) is well-defined
on the boundary mapping torus.  In a monomial chart of P1.9 it is
exactly the operation of forgetting the arguments of the vanishing
coordinates, with the argument of $q$ recorded by $\theta$.
Thus it is the actual boundary retraction.  The deformation removing
the bounded period units in P1.11 transports this map by a homotopy
and preserves its homology class.

\item We prove that the suspension classes have zero image in cusp homology.
Let $D$ be the derivation of $\bigwedge^\bullet(A\oplus B)$ extending
$Nb_i=a_i$, $Na_i=0$.  Then $\bigwedge M_0=\exp D$ and
$\ker(\bigwedge M_0-I)=\ker D$ over $\mathbb Q$.
An integral vector in this kernel has the same property.
The summand with $p$ base generators and $q$ angular generators is
sent by $D$ to the summand with $(p-1,q+1)$; consequently each
homogeneous summand of an invariant vector is separately killed.
Let $E$ replace $a_i$ by $b_i$. In the orthonormal exterior basis
$E=D^t$ and $[D,E]=(q-p)I$ on that summand.  For $Dv=0$,
\[
 (q-p)\|v\|^2
 =\langle[D,E]v,v\rangle=\|Ev\|^2\geq0.
\]
Thus an invariant $k$-cycle has base degree at most $\lfloor k/2\rfloor$.

\item This assertion about cycles can be realized geometrically, rather
than only in a rational associated graded group.  Give $T_B^3$ its
minimal product CW structure and put
$F_p=T_A^3\times(T_B^3)^{(p)}$.
Its cellular differential is zero, so
$H_*(F_p;\mathbb Z)\to H_*(T^6;\mathbb Z)$ is injective.
The shear preserves $F_p$ and admits cellular approximation through
maps preserving this filtration.  The corresponding mapping-cylinder
cells give the integral complex
$\operatorname{Cone}(\bigwedge M_0-I)$.
An invariant cellular vector of base degree at most $p$ therefore
has a suspension lift whose entire projection to $T_B^3$ has dimension
at most $p$.  Equivalently, its closing boundary already bounds in
$F_p$, by the just-proved injectivity.  This construction fixes the
suspension representatives used in the Wang complex and is compatible
with the based product-cell comparison in P1.12.

\item Formula (1.19.1) preserves $b$, so the image of such a lift also has
base projection of dimension at most $p$.  It represents filtration
at most $p$ in the dual-cell filtration of P1.11.  Here is a direct
cellular justification of this last step.  In every dual cell of
dimension greater than $p$, choose an interior point outside that
projection.  Radial retraction from this point pushes the projected
cycle to the cell boundary.  On the central moment model it lifts
by keeping the angular coordinate fixed; on the boundary the angular
quotient is the canonical further quotient.  The lift is continuous
and fixes the boundary.  Performing these retractions in descending
cell dimension gives a homologous cycle on the dual $p$-skeleton.
There are finitely many cells modulo periods, so this operation
terminates.

\item The integral $E_2=E_\infty$ table of P1.11 has nonzero entries only
at angular degree $q\leq d$, where $d$ is the base-cell degree.
For a suspension of a $k$-cycle, the total degree is $k+1$, whereas
$d\leq p\leq\lfloor k/2\rfloor$. Its angular degree would therefore
be $k+1-d>d$, where that table is zero.  The suspension image
vanishes.  On the fibre summand the map is precisely the marked
collapse $\kappa_k$ already exported in
\texttt{cusp\_marked\_collapse.json}.  Thus the homology maps agree
with those of the chain map $f_0=(\kappa,0)$, in every degree.

\item Agreement on homology does not identify these two chain maps.
We now determine the remaining difference.  Integral Smith reduction
splits a finite free complex into free homology generators, contractible
unit pairs, and two-term pairs with differential equal to a torsion
invariant factor.  The cusp homology is free.  The boundary homology
has just one torsion group, $\mathbb Z/2$ in degree three.
Consequently, after maps on homology have been fixed, the difference
of two maps from the boundary complex to the cusp complex is
represented by one vector
\[
 \varepsilon\in H_4(C;\mathbb Z)/2H_4(C;\mathbb Z)
                 \cong(\mathbb Z/2)^{10}.
\tag{1.19.2}
\]
This follows directly by applying a map to the upper generator of
the two-term pair $\mathbb Z\xrightarrow{2}\mathbb Z$; changing a
chain homotopy changes its image by twice an integral vector.
There are no other secondary terms.  Thus no rational formality or
unproved equivariant formality is used at this point.

\item In the marked cusp degree-four basis, the primitive functional
\[
 \psi=(0,0,0,1,-1,1,1,-1,-1,1)
\tag{1.19.3}
\]
annihilates $\kappa_4$ and identifies its cokernel with $\mathbb Z$.
The exact multiplication $\psi\kappa_4=0$ and the primitive Smith
factors are exported.  Setting $\varepsilon=0$ gives relative
cusp homology $H_4(C,\partial)=\mathbb Z^{10}\oplus\mathbb Z/2$,
as computed in \texttt{cusp\_relative\_candidate.json}.
More intrinsically, the relative exact sequence in this degree is
an extension of $\mathbb Z^9\oplus\mathbb Z/2$ by the cokernel
$\mathbb Z$ detected by $\psi$.  Its torsion extension class is
exactly $\psi\varepsilon$ modulo two.

\item The actual cusp neighborhood is a compact oriented smooth
8-manifold with boundary: take the inverse image of a sufficiently
small closed disc, whose boundary circle avoids the critical value.
It retracts onto $C$.  Poincar\'e--Lefschetz duality and the already
computed free cusp homology give
\[
 H_4(\text{cusp cap},\partial;\mathbb Z)
       \cong H^4(C;\mathbb Z)=\mathbb Z^{10}.
\]
The preceding extension must therefore be nonsplit on its torsion
summand.  This determines the required secondary class for the global
calculation:
\[
                 \psi\varepsilon=1\pmod2.
\tag{1.19.4}
\]
This also proves the stated failure of the wholly zero suspension
chain map as a relative cusp model.

\item We next give the integral equivalence which reduces every class
satisfying (1.19.4) to one displayed matrix.  Write the degree-four
homology gluing map in the plus-sign convention as
\[
 j_4:H_4(\partial)\longrightarrow H_4(Y)\oplus H_4(C).
\]
Multiplication by $-1$ on the cusp chain summand is the chain
isomorphism from the minus-sign convention in the exported whole
diagram to this convention.  The matrix $j_4$ is $22$ by $18$,
has rank eighteen, and all its nonzero Smith factors are one.
Its cokernel is free of rank four. In its exported cokernel basis,
the map from $H_4(C)$ is
\[
 \begin{pmatrix}0\\0\\0\\\psi\end{pmatrix}.
\tag{1.19.5}
\]
This is a direct integer matrix calculation, independent of the
homology of $M$.

\item Let $e_4$ denote the fourth coordinate vector of $\mathbb Z^{10}$,
so $\psi e_4=1$.  The nine columns
$e_i-\psi_i e_4$, $i\ne4$, form an integral basis $K$ of $\ker\psi$;
$[e_4,K]$ is unimodular.  The file
\texttt{cusp\_ext\_reduction.json} gives integral lifts $W$ satisfying
\[
 j_4W=\begin{pmatrix}0\\K\end{pmatrix}.
\]
It also gives actual boundary cycles $X$ representing these lifts,
and integral complement chains $T$ with
\[
 d_YT=j_YX,\qquad f_0X=K.
\tag{1.19.6}
\]
Every entry is checked by multiplication in the program.

\item For clarity, the chain automorphisms implementing the reduction
are equally explicit.  Let $r$ be the exported integral degree-four
cocycle representing the torsion generator in $H^4(\partial;\mathbb Z)$.
In the ordered 35-dimensional boundary chain group it is
$r=-e_{17}^{t}-e_{18}^{t}$.  It satisfies
\[
 r d_{\partial,5}=0,\qquad rX=0,
 \qquad 2r=t\,d_{\partial,4},
\tag{1.19.7}
\]
where the integral row $t$ is exported as well.
For a column $v\in\mathbb Z^9$, the endomorphism
$\eta=(Xv)r$ of degree-four boundary chains, zero in all other
degrees, is a chain map and has $\eta^2=0$.
Thus $I+\eta$ is an integral chain automorphism with inverse
$I-\eta$. Precomposing the attaching map by it changes the
secondary vector by $Kv$, up to the explicitly given homotopy
$(Tv)r$ in the complement, by (1.19.6).
The last equation in (1.19.7) removes an even multiple of $e_4$
by the homotopy $e_4t$.  Decomposing $\varepsilon$ in the
unimodular basis $[e_4,K]$, equation (1.19.4) therefore reduces
the actual global attaching map to the canonical correction
\[
 f_4=f_{0,4}+e_4r,
 \qquad f_k=f_{0,k}\quad(k\ne4).
\tag{1.19.8}
\]
These formulas prove the chain equivalence; they do not merely
compare the ranks of the two complexes.

\item The program \texttt{global\_ext\_cases.py} inserts (1.19.8) into
the complete main-piece/three-cap mapping cone, retaining every
finite-cap differential and every boundary map of P1.12.
It checks $d^2=0$ over $\mathbb Z$ and exports every differential.
Its odd case is the actual global model by the preceding argument.
The even case is also computed as a diagnostic, and happens to give
the same global homology; it is not used to identify the local cap.
The independent program \texttt{global\_cell\_certificate.py}
exports the Smith left and right matrices, kernel bases, incoming
matrices in those bases, and their Smith left and right matrices
for the odd case.  The differential ranks in degrees one through
eight are
\[
                 (3,23,64,98,98,58,21,2).
\]
All nonzero incoming Smith factors in the cycle lattices are one.
Subtracting incoming and outgoing ranks from the displayed chain
ranks gives $(1,0,0,0,4,0,0,0,1)$, with no torsion.
This completes the full cellular comparison without using the
global duality conclusion of P1.12.

\end{enumerate}
\end{proof}


\section{The genus-two pencil and the lifted period planes}



\begin{theorem}\label{night:P1.4}
\StatementRecord{night:P1.4}{proved}{Historical events: GPT-6 Astra (OpenAI), ChatGPT Pro mode, effort Pro, S6\_06\_BLIND, 16--18 November 2026: targeted proof check. Also checked in this writing. Current writing events: GPT-6 Astra (OpenAI), ChatGPT Pro mode, effort Pro, S6\_10\_PAPERS, 18 November 2026: computation or formal certificate; targeted proof check. GPT-6 Astra (OpenAI), ChatGPT Pro mode, effort Pro, S6\_10\_PAPERS continuation, 19 November 2026: computation or formal certificate; targeted proof check; used and rechecked in this run. }{Definitions and the displayed constructions}
The genus-two pencil has one nodal fibre away from $0,\infty$, at
$\lambda_*=-256/3125$.  Good reduction at $0,\infty$ requires root
orders $8,5$.  On this root orbifold its Hodge bundle is
$\mathcal H=T^3\oplus T$, where
\[
T=\mathcal O\bigl((2/5)[\infty]-(3/8)[0]\bigr),\qquad
\deg T=1/40,
\]
and $T^3$ is the line generated meromorphically by $dx/y$.
\end{theorem}
\begin{proof}
\begin{enumerate}[label=\arabic*.]
\item The discriminant of $x^5+\lambda^3x+\lambda^4$ is
$\lambda^{15}(3125\lambda+256)$.
At $\lambda_*$ the double root is $64/625$;
the second $x$-derivative and first $\lambda$-derivative there are
nonzero.  The remaining roots are simple.  The odd-degree
hyperelliptic compactification is smooth at its branch point at
infinity.  This gives one node, a smooth total surface at that node,
and genus-one normalization.

\item At zero set $\lambda=u^8$, $x=u^6X$, $y=u^{15}Y$; at infinity set
$\lambda=u^{-5}$, $x=u^{-4}X$, $y=u^{-10}Y$.
The resulting smooth equations and differentials are
\[
\begin{array}{c|c|c|c}
 &\text{equation}&dx/y&x\,dx/y\\\hline
0&Y^2=X^5+X+u^2&u^{-9}dX/Y&u^{-3}X\,dX/Y\\
\infty&Y^2=X^5+uX+1&u^6dX/Y&u^2X\,dX/Y.
\end{array}
\]
Their central polynomials are square-free.  The deck maps at zero and
infinity are $(X,Y)\mapsto(iX,\zeta_8Y)$ and
$(X,Y)\mapsto(\zeta_5^{-1}X,Y)$, respectively.
The eigenvalues on $H^{1,0}$ are $(\zeta_8,\zeta_8^3)$ and
$(\zeta_5^{-1},\zeta_5^{-2})$.  Including conjugates gives characteristic
polynomials $T^4+1$ and $\Phi_5$, proving minimality of the root orders.
The fourth power at zero is the hyperelliptic involution, acting as
$-I$ on $H^1$ and as $I$ on its exterior square.

\item The two meromorphic differential generators have orbifold divisors
$(6/5)[\infty]-(9/8)[0]$ and
$(2/5)[\infty]-(3/8)[0]$.
At the node they are independent sections of the dualizing Hodge
extension: one has opposite nonzero residues at the two points of the
normalization, and the linear combination vanishing at the node is a
regular nonzero elliptic differential.  Thus there is no additional
divisor at $\lambda_*$.  These computations prove the asserted direct
sum and its degrees.  The discriminant and all substitutions are
checked in \texttt{pencil\_charts.py/.out}.

\end{enumerate}
\end{proof}




\begin{theorem}\label{night:P1.5}
\StatementRecord{night:P1.5}{false as stated}{Historical events: GPT-6 Astra (OpenAI), ChatGPT Pro mode, effort Pro, S6\_06\_BLIND, 16--18 November 2026: targeted proof check. Also checked in this writing. Current writing events: GPT-6 Astra (OpenAI), ChatGPT Pro mode, effort Pro, S6\_10\_PAPERS, 18 November 2026: targeted proof check. GPT-6 Astra (OpenAI), ChatGPT Pro mode, effort Pro, S6\_10\_PAPERS continuation, 19 November 2026: independent re-proof; used and rechecked in this run. }{Definitions and the displayed constructions}
The formula in the prompt with $P$ the image of
$\ell_0\wedge\mathcal H$, and $\mathcal H$ the rank-two Hodge bundle,
cannot produce a rank-three subbundle.  The required rank-three image
is instead the image of $\ell_0\wedge\mathcal H_{\mathrm{dR}}$, where
$\mathcal H_{\mathrm{dR}}$ is the full rank-four de Rham bundle.
\end{theorem}
\begin{proof}
\begin{enumerate}[label=\arabic*.]
\item The kernel of $\ell_0\otimes\mathcal H\to\bigwedge^2\mathcal H$
is $\ell_0\otimes\ell_0$, so its image has rank one.  Projection and
addition of a map into the invariant line cannot increase its domain
rank.  For the corrected domain there is an injective holomorphic map
\[
\ell_0\otimes(\mathcal H_{\mathrm{dR}}/\ell_0)
 \longrightarrow\bigwedge^2\mathcal H_{\mathrm{dR}},\quad
u\otimes[v]\longmapsto u\wedge v.
\]
Every nonzero image vector is decomposable and has wedge square zero;
$\omega^2\ne0$.  Projection modulo $\mathcal O\omega$ is therefore
injective on its rank-three image.  This specifies the comparison map
and the corrected bundle used below.

\end{enumerate}
\end{proof}




\begin{theorem}[the geometric rank-five quotient]\label{maps:src:F.3.5}
\StatementRecord{maps:src:F.3.5}{proved}{Historical events: Claude Fable 5.2 (Anthropic), effort high, 16 November 2026: targeted proof check. GPT-6 Astra (OpenAI), ChatGPT Pro mode, effort Pro, S6\_09\_PAPERS, 17 November 2026: targeted proof check. Also checked in this writing. Current writing events: GPT-6 Astra (OpenAI), ChatGPT Pro mode, effort Pro, S6\_10\_PAPERS, 18 November 2026: targeted proof check. GPT-6 Astra (OpenAI), ChatGPT Pro mode, effort Pro, S6\_10\_PAPERS continuation, 19 November 2026: used and rechecked in this run. \emph{Recorded proof scope.} GPT-6 Astra (OpenAI, ChatGPT Pro mode), effort Pro, 2026-11-17; GPT-6 Astra (OpenAI, ChatGPT Pro mode), S6\_06\_BLIND, effort Pro, 2026-11-17; Claude Fable 5.2 (Anthropic), effort high, 2026-11-16; also checked in this writing by GPT-6 Astra (OpenAI), ChatGPT Pro mode, effort Pro, 18 November 2026}{the hypotheses, notation and definitions stated here; \cref{maps:src:F.2.5}}


Put $\lambda_*=-256/3125$ and
$U=\mathbb P^1\setminus\{0,\infty,\lambda_*\}$.  The smooth projective
models of
\[
 C_\lambda:\quad y^2=x^5+\lambda^3x+\lambda^4
 \tag{F.7}\label{maps:eq:F.7}
\]
give a polarized rank-four rational cohomology system $\mathcal H$ on
$U$.  Its primitive exterior square is rationally isomorphic to
the rank-five quotient of the system with matrices $A,B$ in
Theorem \ref{maps:src:F.2.5} by its invariant first coordinate line.
On the orbifold $\mathcal B'$ obtained by taking roots of orders
five at infinity and eight at zero, the extended Hodge bundle is
\[
 \mathcal E=\mathcal T^3\oplus\mathcal T,\qquad
 \mathcal T=\mathcal O_{\mathcal B'}(2D_\infty-3D_0),\qquad
 \deg\mathcal T=\frac1{40}.
 \tag{F.8}\label{maps:eq:F.8}
\]
Here $D_\infty,D_0$ are the reduced root divisors; the extension at
$\lambda_*$ is the nilpotent canonical extension.

\end{theorem}


\begin{proof}
\begin{enumerate}[label=\arabic*.]
\item The discriminant of the degree-five polynomial in \eqref{maps:eq:F.7} is exactly
$\lambda^{15}(3125\lambda+256)$.  At $\lambda_*$ its unique double
root is $r=64/625$ and the coefficient of $(x-r)^2$ is
$10r^3\ne0$.  Its other roots are simple.  Thus the fibre has one
ordinary node and irreducible genus-one normalization.  The local
model is a smoothing of a nonseparating node; the vanishing circle
is a primitive nonseparating homology class.  Cutting and regluing
an annular neighborhood proves that its monodromy is a single
Dehn twist, and hence a primitive rank-one transvection on $H^1$.
In particular its logarithm $N_H$ has square zero and rank one.

\item At zero, the substitutions $\lambda=s^8$, $x=s^6X$, $y=s^{15}Y$
give
\[
 Y^2=X^5+X+s^2.
\]
This is smooth at $s=0$.  The generator $s\mapsto\zeta_8s$ at
fixed $(x,y)$ acts by $(X,Y)\mapsto(iX,\zeta_8Y)$, so its
eigenvalues on the two holomorphic differentials are
$\zeta_8,\zeta_8^3$.  The full rational characteristic polynomial
is $g(T)=T^4+1$.
At infinity put $\lambda=s^{-5}$, $x=s^{-4}X$, $y=s^{-10}Y$.
The resulting smooth family is
\[
 Y^2=X^5+sX+1.
\]
The deck generator acts by $(X,Y)\mapsto(\zeta_5^4X,Y)$,
with holomorphic eigenvalues $\zeta_5^4,\zeta_5^3$.  Its full
characteristic polynomial is $f(T)=T^4+T^3+T^2+T+1$.
Inverses have the same two rational characteristic polynomials.

\item We identify the representation, rather than merely its spectra.
Choose positive peripheral generators $T_1,T_2$ at infinity and
zero, in that order, with $T_1T_2$ inverse to nodal monodromy.
Set $R=T_1^{-1}$ and $S=T_2$.  Then $R-S$ has rank one and their
characteristic polynomials are $f,g$.  The following elementary
rank-one lemma fixes their simultaneous rational conjugacy class.
If $u$ spans $\operatorname{im}(R-S)$, it is cyclic for $R$:
the span of its iterates is stable under both $R$ and $S$, and
the irreducibility of $f$ makes that span all of $\mathbb Q^4$.
In the rational basis $u,Ru,R^2u,R^3u$, the matrix of $R$ is
the companion matrix of $f$, and $S=R-u\varphi$ for a covector
$\varphi$.  The determinant identity
\[
 g(T)=f(T)+\varphi\operatorname{adj}(TI-R)u
 \tag{F.9}\label{maps:eq:F.9}
\]
determines $\varphi$ uniquely: the four coordinates of
$\operatorname{adj}(TI-R)u$ are a basis of the polynomials of
degree at most three, as is seen from their successively distinct
leading degrees in the cyclic basis.  It follows that any two
such pairs are simultaneously rationally conjugate.  Apply the
lemma to the explicit companion pair
\[
 C_f=\begin{pmatrix}
 0&0&0&-1\\1&0&0&-1\\0&1&0&-1\\0&0&1&-1
 \end{pmatrix},\qquad
 C_g=\begin{pmatrix}
 0&0&0&-1\\1&0&0&0\\0&1&0&0\\0&0&1&0
 \end{pmatrix}.
 \tag{F.10}\label{maps:eq:F.10}
\]
Their difference has rank one and their characteristic polynomials
are $f,g$.  Thus $(T_1,T_2)$ is rationally conjugate to
$(C_f^{-1},C_g)$, with the chosen peripheral conventions.

\item Use the lexicographic basis in $\bigwedge^2\mathbb Q^4$.
Its common invariant vector is $(1,1,1,1,1,1)^t$.
Let $Q$ be the identity matrix of size six with first column
replaced by this vector.  It is unimodular.  Direct multiplication
gives the exact identities
\[
 \left(Q^{-1}\bigwedge^2(C_f^{-1})Q\right)_{2:6,2:6}
   =A_{2:6,2:6},\qquad
 \left(Q^{-1}\bigwedge^2(C_g)Q\right)_{2:6,2:6}
   =B_{2:6,2:6}.
 \tag{F.11}\label{maps:eq:F.11}
\]
They are checked over $\mathbb Z$ in
\texttt{fourfolds\_periods.py}, not inferred from eigenvalues.
The invariant line is the inverse of the polarization, up to a
nonzero rational scalar.  Projection along that line identifies
the quotient with the primitive exterior square $\mathcal W$.

\item The differentials $\omega_1=dx/y$ and $\omega_2=x\,dx/y$ span
the extended Hodge bundle everywhere.  At the node they are
independent dualizing differentials: $\omega_1$ has nonzero
opposite residues on the normalization and
$(x-r)\omega_1$ is a nonzero holomorphic differential there.
Here the extension assertion can be checked on the nodal
annulus, without invoking an additional extension theorem.
In coordinates $uv=t$ a relative dualizing differential has
a convergent Laurent expansion
$\sum_n a_n(t)u^n\,du/u$.  Integration on the vanishing circle
is $2\pi i a_0(t)$.  Integration across the annulus has the
single logarithmic term $a_0(t)\log t$; integration of every
nonconstant Laurent monomial and of the complementary fixed
arcs gives a holomorphic remainder.  Thus the period vector
has the form $\exp(sN_H)$ times a holomorphic vector,
$s=\log t/(2\pi i)$.  Untwisting by $\exp(-sN_H)$ is exactly
the nilpotent canonical extension used here.  The two independent
dualizing limits remain a basis of its Hodge subbundle.
Their isotropic pairing, zero away from zero, remains zero
at the limit by holomorphic continuation, so that limit is
Lagrangian.
At zero their expressions in the root coordinates are
$s^{-9}dX/Y$ and $s^{-3}X\,dX/Y$; at infinity they are
$s^6dX/Y$ and $s^2X\,dX/Y$.
Thus their divisor lines are
$\mathcal O(6D_\infty-9D_0)$ and
$\mathcal O(2D_\infty-3D_0)$, proving \eqref{maps:eq:F.8}.
The nonzero residue of $\omega_1$ also shows
\[
 N_H\ell_\infty\ne0,
 \qquad \ell=\mathcal T^3\subset\mathcal E,
 \tag{F.12}\label{maps:eq:F.12}
\]
where $\ell_\infty$ denotes its limit at the nodal cusp.
Indeed integration on the vanishing circle tends to
$2\pi i$ times that residue; the transvection sends the cohomology
class of a differential to that period times its vanishing
cohomology vector.

\end{enumerate}
\end{proof}




\begin{theorem}\label{night:P1.2}
\StatementRecord{night:P1.2}{proved}{Historical events: GPT-6 Astra (OpenAI), ChatGPT Pro mode, effort Pro, S6\_06\_BLIND, 16--18 November 2026: targeted proof check. Also checked in this writing. Current writing events: GPT-6 Astra (OpenAI), ChatGPT Pro mode, effort Pro, S6\_10\_PAPERS, 18 November 2026: targeted proof check. GPT-6 Astra (OpenAI), ChatGPT Pro mode, effort Pro, S6\_10\_PAPERS continuation, 19 November 2026: used and rechecked in this run. }{Definitions and the displayed constructions}
The lower blocks of $A,B$ are precisely the quotient representations of
$\bigwedge^2(C_f^{-1}),\bigwedge^2C_g$ in the supplied basis $Q$.
The integral quotient is the dual of the primitive exterior-square
lattice.  The natural map from the primitive lattice into that quotient
has cokernel $\mathbb Z/2$; these two integral representations are not
isomorphic.  Their rationalizations are isomorphic by that same map.
\end{theorem}
\begin{proof}
\begin{enumerate}[label=\arabic*.]
\item The two matrices obtained from their two-by-two minors are
\[
Q^{-1}\bigwedge^2(C_f^{-1})Q=
\begin{pmatrix}
1&-1&0&1&0&0\\0&1&-1&-1&1&0\\0&1&0&-1&0&0\\
0&2&-1&-1&0&1\\0&2&0&-1&0&0\\0&1&1&-1&0&0
\end{pmatrix},
\]
\[
Q^{-1}\bigwedge^2C_gQ=
\begin{pmatrix}
1&0&1&0&0&0\\0&0&-1&0&1&0\\0&0&-1&0&0&1\\
0&0&-1&0&0&0\\0&1&-1&0&0&0\\0&0&-1&1&0&0
\end{pmatrix}.
\]
Their lower five-by-five blocks are the asserted blocks, entry by entry.
Put $L=\bigwedge^2\mathbb Z^4$, let $(\ ,\ )$ be its unimodular wedge
pairing, and set $\omega=(1,1,1,1,1,1)^t$.
Then $(\omega,\omega)=2$ and the primitive sublattice
$P_{\mathbb Z}=\omega^\perp$ is the kernel of
$(1,-1,1,1,-1,1)$.  In the basis with columns
\[
\begin{pmatrix}
1&-1&-1&1&-1\\1&0&0&0&0\\0&1&0&0&0\\
0&0&1&0&0\\0&0&0&1&0\\0&0&0&0&1
\end{pmatrix},
\]
the inclusion followed by projection
$j:P_{\mathbb Z}\longrightarrow L/\mathbb Z\omega$ has matrix
\[
j=\begin{pmatrix}
0&1&1&-1&1\\-1&2&1&-1&1\\-1&1&2&-1&1\\
-1&1&1&0&1\\-1&1&1&-1&2
\end{pmatrix},\qquad \det j=2.
\]
Its Smith factors are $(1,1,1,1,2)$.  The unimodular pairing on $L$
identifies $L/\mathbb Z\omega$ with $P_{\mathbb Z}^\vee$;
under this identification $j$ is the pairing map.

\item To rule out a different integral intertwiner, we solve the two
simultaneous bilinear-form invariance equations.  Each solution space
over $\mathbb Q$ has dimension one.  The primitive integral forms have
determinants $-2$ on $P_{\mathbb Z}$ and $-16$ on the quotient.
A unimodular intertwiner would pull one primitive form to a rational
multiple of the other.  Having coefficient content one on both sides
forces that multiple to be $1$ or $-1$, which contradicts the different
determinants.  All these linear equations and their solutions occur in
\texttt{primitive\_lattice.py/.out}.  Thus the rational wording in the
data is correct, whereas the stronger integral identification requested
in item 1.4 has precisely the defect just exhibited.

\end{enumerate}
\end{proof}




\begin{theorem}\label{night:P1.7}
\StatementRecord{night:P1.7}{proved}{Historical events: GPT-6 Astra (OpenAI), ChatGPT Pro mode, effort Pro, S6\_06\_BLIND, 16--18 November 2026: targeted proof check. Also checked in this writing. Current writing events: GPT-6 Astra (OpenAI), ChatGPT Pro mode, effort Pro, S6\_10\_PAPERS, 18 November 2026: targeted proof check. GPT-6 Astra (OpenAI), ChatGPT Pro mode, effort Pro, S6\_10\_PAPERS continuation, 19 November 2026: used and rechecked in this run. }{\cref{night:P1.2}}
The integral $H^1$ monodromy of the pencil has a companion marking with
finite pair $(C_f^{-1},C_g)$, and cusp product $C_fC_g^{-1}$.
Its exterior-square quotient is consequently the integral quotient
representation appearing in P1.2.
\end{theorem}
\begin{proof}
\begin{enumerate}[label=\arabic*.]
\item Put $z=\lambda/\lambda_*$ and use the based-meridian convention
$g_\infty g_1g_0=1$ at $z=1/2$ of Beukers--Heckman,
Definition 3.1.  The two finite spectra are disjoint.  The nodal
monodromy is a rank-one Picard--Lefschetz transformation with primitive
vanishing cycle $\delta$ and coefficient $\pm1$.
Here primitivity follows by taking a transverse integral curve to the
nonseparating vanishing circle; coefficient one follows by turning the
annulus $xy=z-1$ once positively.  The sign depends only on the order
of the pairing and on homology versus cohomology.

\item Beukers--Heckman, Invent. Math. 95 (1989), Theorem 3.5 (Levelt), gives
the companion triple $(C_f,C_f^{-1}C_g,C_g^{-1})$.
The proof is defined over $\mathbb Q$ in this case: the intersection of
the three appropriate translates of the rational rank-three kernel
contains a rational vector whose four successive translates form a
basis.  The disjoint spectra imply irreducibility, since a common
proper invariant subspace or its quotient on which the rank-one
transformation is trivial would give a shared eigenvalue.  This
verifies the field of the resulting intertwiner as well as the
hypotheses of the theorem.

\item The integral upgrade uses the exact matrices
\[
d=(1,1,1,0)^t,\qquad
J=\begin{pmatrix}
0&-1&1&-1\\1&0&-1&1\\-1&1&0&-1\\1&-1&1&0
\end{pmatrix}.
\]
They satisfy $C_f^tJC_f=C_g^tJC_g=J$, $\det J=1$, and
$C_f^{-1}C_g=I-d(d^tJ)$.
The orbit $(d,C_fd,C_f^2d,C_f^3d)$ has alternating pairing
\[
\begin{pmatrix}0&-1&0&0\\1&0&-1&0\\0&1&0&-1\\0&0&1&0\end{pmatrix},
\]
of determinant one.  Under a rational intertwiner let $d$ map to
$a\delta$, and let the pulled-back intersection form be $cJ$.
The space of simultaneous invariant alternating forms has dimension
one by the exact invariance equations.  Comparing the rank-one
transvection formulas yields $c/a^2=\pm1$.
It follows that the four integral cycles
$\delta,H_\infty\delta,H_\infty^2\delta,H_\infty^3\delta$ have the
displayed unimodular pairing, with an overall sign.  Their sublattice
has index one in the unimodular integral $H^1$, since its determinant
is the square of that index.  Rescaling the intertwiner by $a$ thus
gives an integral marking.

\item The antidiagonal reversal matrix $R$ satisfies
$RC_fR^{-1}=C_f^{-1}$ and $RC_g^{-1}R^{-1}=C_g$.
Consequently the cusp product is $(C_gC_f^{-1})^{-1}$, matching
$(BA)^{-1}$ on the supplied cohomology system.  Its transpose is the
specified $M_0$.  The matrices and pairing calculation are recorded in
\texttt{companion\_markings.py/.out}.  The verified primary source is
\url{https://math.uchicago.edu/~drinfeld/Hypergeometric_motives/Beukers-Heckman.pdf},
pp.~331--332.

\end{enumerate}
\end{proof}




\begin{theorem}[a global nonpolarized period family]\label{maps:src:F.3.6}
\StatementRecord{maps:src:F.3.6}{proved}{Historical events: Claude Fable 5.2 (Anthropic), effort high, 16 November 2026: targeted proof check. GPT-6 Astra (OpenAI), ChatGPT Pro mode, effort Pro, S6\_09\_PAPERS, 17 November 2026: targeted proof check. Also checked in this writing. Current writing events: GPT-6 Astra (OpenAI), ChatGPT Pro mode, effort Pro, S6\_10\_PAPERS, 18 November 2026: targeted proof check. GPT-6 Astra (OpenAI), ChatGPT Pro mode, effort Pro, S6\_10\_PAPERS continuation, 19 November 2026: used and rechecked in this run. \emph{Recorded proof scope.} GPT-6 Astra (OpenAI, ChatGPT Pro mode), effort Pro, 2026-11-17; GPT-6 Astra (OpenAI, ChatGPT Pro mode), S6\_06\_BLIND, effort Pro, 2026-11-17; Claude Fable 5.2 (Anthropic), effort high, 2026-11-16; also checked in this writing by GPT-6 Astra (OpenAI), ChatGPT Pro mode, effort Pro, 18 November 2026}{the hypotheses, notation and definitions stated here; \cref{maps:eq:F.11}; \cref{maps:eq:F.14}; \cref{maps:eq:F.16}; \cref{maps:eq:F.19}; \cref{maps:eq:F.12}; \cref{maps:src:F.2.1}; \cref{maps:eq:F.22}; \cref{maps:eq:F.20}; \cref{maps:eq:F.21}}


The integral rank-six system with cohomology monodromies $A,B$
admits a holomorphic subbundle $\mathcal F$ of rank three over
$U$ satisfying $\mathcal F\cap\overline{\mathcal F}=0$.
It extends holomorphically on the finite root charts.  At the
cusp an integral marking puts its periods in the form
\[
 Z(t)=\frac{\log t}{2\pi i}I_3+C(t),
 \tag{F.13}\label{maps:eq:F.13}
\]
where $C$ is holomorphic at zero and bounded there.

\end{theorem}


\begin{proof}
\begin{enumerate}[label=\arabic*.]
\item Write the exact sequence supplied by \eqref{maps:eq:F.11} as
\[
 0\longrightarrow\mathbb Q e\longrightarrow\mathcal V
   \longrightarrow\mathcal W\longrightarrow0.
 \tag{F.14}\label{maps:eq:F.14}
\]
Here $e$ is the first coordinate vector.  The nilpotent canonical
extension at the cusp and the ordinary extensions on finite root
charts preserve this sequence.  This follows directly by using
the frames $\exp(-sN)$: each residue has only eigenvalue zero,
and the same exponential construction respects the inclusion and
the quotient.  At finite roots the statement is ordinary exactness
of bundles with trivialized local monodromy.  Denote these extended
bundles by $\mathcal H_D,\mathcal V_D,\mathcal W_D$.

\item Let $\omega\in\bigwedge^2\mathcal H_D$ be the real, flat
nondegenerate bivector inverse to the polarization, and let
$p:\bigwedge^2\mathcal H_D\to\mathcal W_D$ be primitive projection.
Define the holomorphic rank-three bundle
\[
 \mathcal P=p(\ell\wedge\mathcal H_D).
 \tag{F.15}\label{maps:eq:F.15}
\]
Projection is injective on $\ell\wedge\mathcal H_D$: every
nonzero bivector there is decomposable whereas $\omega$ is not.
This argument applies also at the cusp.  The polarization gives
$\mathcal H_D/\mathcal E\simeq\mathcal E^*$, including there,
since $\mathcal E$ is Lagrangian.  Hence
\[
 0\longrightarrow\mathcal T^4\longrightarrow\mathcal P
 \longrightarrow\mathcal O\oplus\mathcal T^2\longrightarrow0.
 \tag{F.16}\label{maps:eq:F.16}
\]
After dualizing, the end terms are
$\mathcal O\oplus\mathcal T^{-2}$ and $\mathcal T^{-4}$.

\item We compute the obstruction to lifting $\mathcal P$ in \eqref{maps:eq:F.14}.
For the coarse projection $\rho:\mathcal B'\to\mathbb P^1$,
the elementary root-chart calculation of invariant sections gives
\[
 \rho_*\mathcal T^k=
 \mathcal O_{\mathbb P^1}
 \left(\left\lfloor\frac{2k}{5}\right\rfloor+
       \left\lfloor\frac{-3k}{8}\right\rfloor\right).
 \tag{F.17}\label{maps:eq:F.17}
\]
Higher direct images vanish, because taking invariants under a
finite group in characteristic zero is exact.  For $k=0,-2,-4$
the displayed degrees are respectively $0,-1,-1$.  Their first
cohomology vanishes.  The dual of \eqref{maps:eq:F.16} therefore gives
$H^1(\mathcal B',\mathcal P^*)=0$.
The pullback of \eqref{maps:eq:F.14} to $\mathcal P$ consequently splits as
holomorphic bundles: choose a holomorphic lift
$\sigma:\mathcal P\to\mathcal V_D$ of its inclusion in
$\mathcal W_D$.  This is an actual extension-class calculation
on a compact orbifold, not an approximation or an Oka principle.
To specify the choice concretely, take a finite root-coordinate
cover with local bundle lifts $\sigma_i$ and write their
differences as $\xi_{ij}e$, where
$\xi_{ij}$ is a Cech cocycle in $\mathcal P^*$.
Choose any solution $b_i-b_j=\xi_{ij}$ of this explicitly
vanishing cohomology class and glue the maps
$\sigma_i-b_ie$.  Every such solution is allowed in the
construction that follows.

\item There is also a canonical holomorphic functional
$\alpha:\mathcal P\to\mathcal O$.
For $v\in\mathcal P$, its unique preimage in
$\ell\wedge\mathcal H_D$ is
\[
 j(v)=v+\alpha(v)\omega
 \tag{F.18}\label{maps:eq:F.18}
\]
under the primitive decomposition.
At a smooth point, the positive Hodge metric implies that
$\ell$ and $\bar\ell$ form a nondegenerate symplectic pair.
Their common wedge line is $\ell\wedge\bar\ell$.
The bivector $\omega$ does not belong to
$\ell\wedge\mathcal H+\bar\ell\wedge\mathcal H$, since its
component on the two-dimensional symplectic orthogonal complement
is nonzero.  Projection $p$ therefore preserves the intersection,
and
\[
 \mathcal P\cap\bar{\mathcal P}
      =p(\ell\wedge\bar\ell).
 \tag{F.19}\label{maps:eq:F.19}
\]
This is the complexification of a real line.  If $u$ spans
$\ell$, then $i u\wedge\bar u$ is real and its coefficient on
$\omega$ is nonzero, by positivity.  Thus $\alpha$ is real and
nonzero on that real intersection line.  There is a unique smooth
real section $n$ of it such that $\alpha(n)=1$.

\item Since the quotient of $\sigma(n)$ is real, there is a unique
real smooth function $h$ on $U$ such that
\[
 \sigma(n)-\overline{\sigma(n)}=2ih e.
 \tag{F.20}\label{maps:eq:F.20}
\]
For $c\in\mathbb R$ set
\[
 \mathcal F_c=(\sigma+ic\alpha e)(\mathcal P).
 \tag{F.21}\label{maps:eq:F.21}
\]
This is a holomorphic rank-three subbundle projecting
isomorphically to $\mathcal P$.  An intersection with its complex
conjugate must project to \eqref{maps:eq:F.19}.  On its real generator the two
lifts differ by $2i(h+c)e$.  It follows, including arbitrary
complex multiples of that generator, that
\[
 \mathcal F_c\cap\overline{\mathcal F_c}=0
       \quad\Longleftrightarrow\quad h+c\ne0.
 \tag{F.22}\label{maps:eq:F.22}
\]

\item We verify the required cusp property before choosing $c$.
In a complex symplectic basis $a,b,d_1,d_2$ of the limiting
$\mathcal H_D$, relation \eqref{maps:eq:F.12} permits us to arrange
$N_Hb=a$, $N_Ha=N_Hd_1=N_Hd_2=0$ and $\ell_\infty=\mathbb Cb$.
Then
\[
 \mathcal P_\infty=
 \langle k=p(b\wedge a),\ b\wedge d_1,\ b\wedge d_2\rangle,
 \quad
 N_W(b\wedge d_i)=a\wedge d_i,
 \quad N_Wk=0.
 \tag{F.23}\label{maps:eq:F.23}
\]
The line generated by $k$ is
$\ker N_W/\operatorname{im}N_W$; all other kernel directions
come from the displayed image.  The exact matrices give
$N_V^2=0$, $\operatorname{rank}N_V=3$ and
$\operatorname{rank}N_W=2$.
Every lift of an image vector $a\wedge d_i$ is killed by $N_V$,
because that vector has a lift in $\operatorname{im}N_V$ and
$N_Ve=0$.  Rank three now forces $N_V$ to take any lift of $k$
to a nonzero multiple of $e$.  Therefore
$N_V:\mathcal F_{c,\infty}\to\operatorname{im}N_V$ is an
isomorphism for every $c$.  In particular its domain is a
complement to $\ker N_V$.

\item The primitive index-one cusp has a real integral basis
$a_i,b_i$ with $N_Vb_i=a_i$, by Proposition \ref{maps:src:F.2.1}, also for the
dual or inverse convention.  In this basis the untwisted
subbundle near zero is the graph $b+C_c(t)a$, with $C_c$
holomorphic and bounded.  Restoring the flat frame by
$\exp(sN_V)$, $s=\log t/(2\pi i)$, gives \eqref{maps:eq:F.13}.
Thus $\operatorname{Im}Z=\operatorname{Im}s\,I_3+
\operatorname{Im}C_c$ is nonsingular on a sufficiently small
punctured disc.  Apply this first with $c=0$.
Equation \eqref{maps:eq:F.22} shows that $h$ is nonzero and has constant sign
on that connected punctured disc.

\item At the two finite root points, all ingredients in \eqref{maps:eq:F.19}--\eqref{maps:eq:F.20}
are smooth and $\alpha(n)=1$ remains valid, so $h$ extends and
is bounded.  On the compact complement of the cusp disc it has
a finite bound $K$.  Choose $c$ with the same sign as $h$ near
the cusp by the exact rule
$c=\operatorname{sign}(h(t_0))(K+2)$, where $t_0$ is any point
of that punctured disc and $K$ is the supremum of $|h|$ on
the compact complement.  Then $h+c$ is nonzero both on that
disc and on its complement.  Equations \eqref{maps:eq:F.21}--\eqref{maps:eq:F.22} prove
global real transversality, and the preceding cusp calculation
proves the bounded correction for this same $c$.

\item Finally, the fourth power of the order-eight action on
$\mathcal H$ is $-I$.  It acts trivially on
$\mathcal W,\mathcal P,\alpha$ and $\mathcal V$.
Hence these bundles, the lift and the constructed periods descend
from $\mathcal B'$ to the orbifold $\mathcal B$ with orders five
and four.  Indeed any equivariant bundle map between bundles with
trivial action of the extra order-two stabilizer descends on its
root chart.  This gives exactly the finite orders required for
the rank-six datum.

\end{enumerate}
\end{proof}




\begin{theorem}\label{night:P1.8}
\StatementRecord{night:P1.8}{proved}{Historical events: GPT-6 Astra (OpenAI), ChatGPT Pro mode, effort Pro, S6\_06\_BLIND, 16--18 November 2026: targeted proof check. Also checked in this writing. Current writing events: GPT-6 Astra (OpenAI), ChatGPT Pro mode, effort Pro, S6\_10\_PAPERS, 18 November 2026: computation or formal certificate; targeted proof check. GPT-6 Astra (OpenAI), ChatGPT Pro mode, effort Pro, S6\_10\_PAPERS continuation, 19 November 2026: used and rechecked in this run. }{\cref{night:P1.5}}
With the corrected rank-three bundle $P$ of P1.5, the supplied rank-six
system admits a holomorphic lifted family
$F_c=(\sigma+ic\alpha e)(P)$ on the punctured base.
For a specified splitting $\sigma$ there is a real analytic function
$h$ with finite attained minimum $h_{\min}$, and the exact domain of
real parameters giving $F_c\cap\overline{F_c}=0$ everywhere is
$c>-h_{\min}$.
At the cusp its periods have the form
$Z(q)=(\log q/(2\pi i))I_3+C(q)$ with $C$ holomorphic and bounded.
The construction extends to the two finite root charts and descends to
orders $5,4$ on the rank-six system.
\end{theorem}
\begin{proof}
\begin{enumerate}[label=\arabic*.]
\item Work first on the root orbifold with orders $5,8$, leaving the nodal
point as a logarithmic point.  The symplectic pairing and the Hodge
subbundle give an exact sequence
$0\to\mathcal H\to\mathcal H_{\mathrm{dR}}\to\mathcal H^\vee\to0$.
This remains exact at the node by the residue pairing in its plumbing
chart.  With $\ell_0=T^3$ the corrected bundle fits into
\[
0\longrightarrow T^4\longrightarrow P
 \longrightarrow\mathcal O\oplus T^2\longrightarrow0.
\]
Contraction in $\ell_0\wedge\mathcal H_{\mathrm{dR}}$ defines
$\alpha:P\to\mathcal O$, using the unique lift supplied by P1.5;
it is surjective.  Thus it is not an asserted contraction on the whole
quotient by $\omega$, where that formula would not be well defined.

\item For $a=2,4,-2,-4$ the invariant coarse direct image of $T^a$ has degree
$-1$, by taking the floors of the two coefficients of its rational
divisor.  Finite-stabilizer invariants are exact in characteristic zero
by averaging.  Consequently $H^0(T^a)=H^1(T^a)=0$ for these exponents.
The two-chart Laurent-series computation for $\mathcal O(-1)$ proves
this vanishing directly.  The extension therefore splits, and
\[
P\cong T^4\oplus T^2\oplus\mathcal O,\qquad
H^1(P^\vee)=0,\quad H^0(P^\vee)=\mathbb C\alpha,
\quad H^0(P)=\mathbb Cg,
\]
where $\alpha(g)=1$.

\item Let $V$ be the logarithmic extension of the supplied rank-six flat
bundle.  At the cusp use its nilpotent exponential frame, and on each
finite root chart use the trivial flat extension with the deck action.
These charts give the exact bundle sequence
$0\to\mathcal Oe\to V\to W\to0$.
Its pullback to $P$ splits because $H^1(P^\vee)=0$.
All holomorphic splittings are $\sigma+z\alpha e$, $z\in\mathbb C$.
To remove this scalar choice we require the $e$-coordinate of
$\sigma(g)$ to be zero in the fixed flat frame at $z=1/2$.

\item Choose a local generator $u$ of $\ell_0$ and put
$r=i\pi(u\wedge\bar u)\in W_{\mathbb R}$, $A=\alpha(r)>0$;
the sign is that of the polarization.  The symplectic orthogonal
decomposition into $\langle u,\bar u\rangle$ and its two-dimensional
complement proves $P\cap\bar P=\mathbb Cr$.
Indeed $u\wedge v-\bar u\wedge w=t\omega$ has zero component in the
exterior square of that complement on the left and a nonzero such
component on the right unless $t=0$.  With $t=0$ its intersection is
the line generated by $u\wedge\bar u$.
Define $h$ by the exact equality
\[
\sigma(r)-\overline{\sigma(r)}=2iAh e.
\]
Both $r$ and $A$ scale by $|a|^2$ when $u$ scales by $a$, so this is
well defined.  Comparing the unique lifts of a multiple of $r$ proves
\[
F_c\cap\overline{F_c}=0\quad\Longleftrightarrow\quad h+c\ne0.
\]

\item The sign and growth at the cusp follow from the actual extension
matrices.  Put $D=C_fC_g^{-1}-I$ and $N_V=(BA)^{-1}-I$.
Only the zeroth column of $D$ is nonzero, and it is $(0,1,1,1)^t$;
the top row of $N_V$ is $(0,0,0,-1,1,-1)$.
For arbitrary $a,b\in\mathbb C^4$, two-by-two minor expansion gives
\[
(N_V)_{0,*}(0,\pi(a\wedge Db))
 =\operatorname{contr}(a\wedge Db)=a_0b_0,
 \qquad N_W\pi(a\wedge Db)=0.
\]
These polynomial identities, with arbitrary symbols for the entries of
$a,b$, are checked in \texttt{cusp\_lift\_asymptotic.py}.
Write $u=\exp(sD)a(q)$, $s=\log q/(2\pi i)$ and $y=\operatorname{Im}s$.
The vector $a$ is holomorphic, and $a_0(0)\ne0$ because $dx/y$ has
nonzero nodal residue.  In the logarithmic frame we obtain
\[
\exp(-sN_W)r=i\pi(a\wedge\bar a)+2y\pi(a\wedge D\bar a),
\quad
A=i\operatorname{contr}(a\wedge\bar a)+2y|a_0|^2.
\]
The splitting is bounded in that frame.  On returning by
$\exp(sN_V)$, the only quadratic contribution to the imaginary
$e$-coordinate is $2|a_0|^2y^2$.  All other terms are $O(y)$,
uniformly for $\operatorname{Re}s$ in a fixed unit interval.
Dividing by $A=2|a_0|^2y+O(1)$ proves
\[
h(q)=y+O(1)=-\frac{\log|q|}{2\pi}+O(1).
\]
The function extends continuously over both finite root charts and
tends to $+\infty$ at the cusp.  Compactness away from that cusp gives
an attained minimum.  Connectedness and the intermediate value
theorem give $h(U)=[h_{\min},\infty)$, with finite-root points included
when testing the completed family.  Thus the domain is precisely the
displayed half-line.  Replacing $\sigma$ by $\sigma+ib\alpha e$ replaces
$h$ by $h+b$ and sends $F_c$ to $F_{c+b}$, proving the stated
comparison between choices.

\item At $q=0$ the map $N_V:F_{\log}\to\ker N_V$ is an isomorphism.
On the quotient its rank is two and its kernel is generated by
$\pi(a\wedge Da)$; the displayed scalar identity sends a lift of this
generator to $a_0^2e\ne0$.  In the integral Jordan marking, the
logarithmic three-plane is therefore a graph of a holomorphic matrix
$C(q)$; applying $\exp(sN_V)$ adds $sI_3$.  Evaluation on the homology
lattice gives the period matrix $(I_3,Z(q))$.
The fourth power of the order-eight deck transformation acts trivially
on $V$ and $P$, so the graph descends to the order-four chart.
The other chart has order five.  The affine translations are
holomorphic translations on the resulting tori, and their linear
parts preserve this graph by construction.

\end{enumerate}
\end{proof}


\paragraph{Scope of P1.8.}
The definition of $h$ is exact and specifies its normalization, its
range, and its cusp growth.  We have not evaluated $h_{\min}$ or
supplied the two independent thirty-digit period calculations requested
in item 1.4.  These numerical requirements remain separately recorded;
they are not consequences of the existence proof.

\begin{theorem}\label{night:P1.9}
\StatementRecord{night:P1.9}{proved}{Historical events: GPT-6 Astra (OpenAI), ChatGPT Pro mode, effort Pro, S6\_06\_BLIND, 16--18 November 2026: targeted proof check. Also checked in this writing. Current writing events: GPT-6 Astra (OpenAI), ChatGPT Pro mode, effort Pro, S6\_10\_PAPERS, 18 November 2026: targeted proof check. GPT-6 Astra (OpenAI), ChatGPT Pro mode, effort Pro, S6\_10\_PAPERS continuation, 19 November 2026: used and rechecked in this run. }{Definitions and the displayed constructions}
The periodic Kuhn fan, in the marking $S$, gives a smooth proper
holomorphic filling of the periods $Z(q)=sI_3+C(q)$ for a sufficiently
small disk.  Its central fibre is reduced with normal crossings.
Modulo the index-one period lattice it has respectively $1,7,12,6$
strata of dimensions $3,2,1,0$, each isomorphic to the corresponding
complex torus $(\mathbb C^*)^d$, and Euler number $6$.
An index-$d$ period sublattice gives $d$ times these stratum counts and
Euler number $6d$.
\end{theorem}
\begin{proof}
\begin{enumerate}[label=\arabic*.]
\item Let $\Sigma$ consist of the cones over all simplices of the periodic
Kuhn triangulation at height one, together with their faces.  For an
ordered maximal simplex $v_0,\ldots,v_3$, the columns $(v_j,1)$ form a
unimodular four-by-four matrix.  If $u_j$ are its dual integral basis,
the corresponding chart has coordinates $x_j=\chi^{u_j}$ and is
$\mathbb C^4$.  The base character is the sum of the four $u_j$, so
\[
q=x_0x_1x_2x_3.
\]
On a common face we invert the coordinates of omitted rays and change
dual bases; these are mutually inverse integral monomial maps.  These
explicit charts and overlaps construct a smooth toric space over the
disk with a reduced normal-crossings zero fibre.

\item The period group acts by
\[
(z,q)\longmapsto
\bigl(q^\lambda\exp(2\pi iC(q)\lambda)z,q\bigr),
\qquad \lambda\in\mathbb Z^3.
\]
The monomial part takes the cone over $\sigma$ to the cone over
$\sigma+\lambda$; the unit part is torus multiplication holomorphic in
$q$, which extends over each chart.  These maps compose according to
addition in $\mathbb Z^3$.  On the zero fibre a fixed point would
require a nonzero translation to stabilize a bounded simplex, which
is impossible.  For small nonzero $q$, the imaginary part of the
period matrix is $\operatorname{Im}(s)I_3+\operatorname{Im}C(q)$ and
is invertible; hence the action is free there too.

\item We give a compactness check for the infinite-chart construction.  Put
$u=\log|z|/\log|q|$.  A period changes $u$ by
$(I+E(q))\lambda$, where $E(q)=O(1/|\log|q||)$ uniformly.
On a small disk $(I+E)^{-1}$ has norm at most two.  Reduction of
$(I+E)^{-1}u$ modulo $\mathbb Z^3$ therefore gives a representative
with $u$ in a fixed bounded region.  Only finitely many simplices
meet that region.  In a simplex containing $u$, the numbers
$\log|x_j|/\log|q|$ are its barycentric coordinates; thus
$|x_j|\leq1$.  Finitely many closed polydisks in these charts give
compact representatives above a smaller closed base disk.  Every
central point is included in such representatives by translating its
finite star into the same bounded region and then using its complete
toric component.

\item The action is proper.  For compact coordinate neighborhoods approaching
the central fibre, only finitely many cones and their faces occur;
their translates can meet only when the translation is bounded.
Equivalently, the valuation estimate just given and the uniform bound
on $(I+E)^{-1}$ bound the translations for a sequence with $q\to0$.
On a closed annulus away from zero, properness is the usual
properness of a full lattice action.  These two checks give finitely
many intersecting translates of any pair of compact neighborhoods.
The quotient is therefore Hausdorff, is smooth by freeness, and is
proper over the disk by the compact representatives.  The map is flat
because $q$ is a nonzero divisor in each smooth chart.  Exponentiation
$z_j=\exp(2\pi i\zeta_j)$ identifies the nonzero fibres with
$\mathbb C^3/(\mathbb Z^3+Z(q)\mathbb Z^3)$.

\item A $p$-simplex supplies a stratum $(\mathbb C^*)^{3-p}$.
Enumeration of all faces of the six Kuhn tetrahedra, retaining their
translation identifications, gives $(1,7,12,6)$.
Additivity of compactly supported Euler characteristic and
$\chi_c(\mathbb C^*)=0$ leave six point contributions.
Inclusion of an index-$d$ period sublattice induces a finite covering
of these toric quotients.  Each simplex orbit has exactly $d$ lifts,
which proves the stated index dependence by an actual comparison map.

\end{enumerate}
\end{proof}




\begin{theorem}[compact holomorphic filling]\label{maps:src:F.3.7}
\StatementRecord{maps:src:F.3.7}{proved}{Historical events: Claude Fable 5.2 (Anthropic), effort high, 16 November 2026: carried as background without fault found. GPT-6 Astra (OpenAI), ChatGPT Pro mode, effort Pro, S6\_09\_PAPERS, 17 November 2026: targeted proof check. Also checked in this writing. Current writing events: GPT-6 Astra (OpenAI), ChatGPT Pro mode, effort Pro, S6\_10\_PAPERS, 18 November 2026: targeted proof check. GPT-6 Astra (OpenAI), ChatGPT Pro mode, effort Pro, S6\_10\_PAPERS continuation, 19 November 2026: used and rechecked in this run. \emph{Recorded proof scope.} GPT-6 Astra (OpenAI, ChatGPT Pro mode), effort Pro, 2026-11-17; GPT-6 Astra (OpenAI, ChatGPT Pro mode), S6\_06\_BLIND, effort Pro, 2026-11-17; Claude Fable 5.2 (Anthropic), effort high, 16 November 2026; also checked in this writing by GPT-6 Astra (OpenAI), ChatGPT Pro mode, effort Pro, 18 November 2026}{the hypotheses, notation and definitions stated here; \cref{maps:src:F.2.6}; \cref{maps:src:F.2.5}; \cref{maps:eq:F.24}; \cref{maps:eq:F.13}; \ExternalPaperRef{equation~}{{E.4}}{III}; \ExternalPaperRef{equation~}{{E.5}}{III}; \ExternalPaperRef{theorem~}{6.2}{III}; \ExternalPaperRef{theorem~}{6.3}{III}}


The smooth eight-manifold of Theorem \ref{maps:src:F.2.6} carries an integrable
complex structure.  There is a proper holomorphic map to
$\mathbb P^1$ whose regular fibres are complex three-tori,
whose two finite reduced fibres are the free affine quotients
of orders five and four specified in Theorem \ref{maps:src:F.2.5}, and whose
remaining fibre is the index-one Kuhn degeneration.

\end{theorem}


\begin{proof}
\begin{enumerate}[label=\arabic*.]
\item We first fix the dual convention.  The preceding construction
has cohomology monodromies $A,B$; Theorem \ref{maps:src:F.2.5} uses homology
monodromies $A^t,B^t$.  All three special values are real.
Choose a positive real basepoint and paths along the positive
real circle to infinity and to zero.  Complex conjugation fixes
these paths and reverses the respective peripheral meridians.
Pull back the preceding period family by conjugation of the
base and conjugate its complex fibres.  Its Hodge subbundle is
holomorphic, with local expression
$\overline{\mathcal F(\bar\lambda)}$, and its cohomology
monodromies are $A^{-1},B^{-1}$.  Its homology monodromies are
therefore exactly $A^t,B^t$.  The third monodromy is fixed by
the product relation, so this also fixes its cusp marking.
The primitive index-one condition and the bounded holomorphic
cusp correction are preserved, after replacing the real
angular basis by its negative if necessary.

\item For a holomorphic Hodge subbundle $\mathcal F$ with
$\mathcal F\cap\bar{\mathcal F}=0$, restriction of integral
homology functionals embeds the homology lattice discretely in
$\mathcal F^*$.  Real transversality makes its real span the
whole real vector space.  The quotient
$\mathcal F^*/\mathcal V_{\mathbb Z}^*$ is consequently a
holomorphic family of compact complex three-tori over $U$,
with its zero section.  This also proves local properness,
by a smoothly varying fundamental parallelepiped.

\item On the root disc at the finite point of order $m_i$, the family
extends smoothly and its deck map has linear part $M_i$.
The invariant vector $v_i$ of Theorem \ref{maps:src:F.2.5} gives a holomorphic
torsion section $v_i/m_i$.  Use the diagonal affine deck action
\[
 (z,s)\longmapsto(M_i z+v_i/m_i,\zeta_{m_i}s).
 \tag{F.24}\label{maps:eq:F.24}
\]
Here $M_i z$ denotes the induced complex-linear map
$\mathcal F_s^*\to\mathcal F_{\zeta_{m_i}s}^*$ from the marked
lattice action, not multiplication of a complex three-vector
by a six-by-six matrix.  Likewise $v_i/m_i$ means its evaluation
as a point of the target torus.
It has order $m_i$, since the $m_i$th translation is integral.
Every nonidentity power is fixed-point-free on the central
torus: the invariant integral covector $\ell$ takes $v_i$ to
$1$ or $-1$, so a fixed point would require the noninteger
$\pm j/m_i$ to vanish in $\mathbb R/\mathbb Z$.
Away from the central fibre the base rotation is already free.
The quotient is therefore smooth and proper over the coarse
disc, with precisely the required multiple fibre.

\item Its marked punctured family is holomorphically identified with
the original torus family by logarithmic translation.  On the
root cover use $r(s)=(\log s/(2\pi i))v_i$.  Then
$r(\zeta_{m_i}s)-M_i r(s)=v_i/m_i$, modulo the integral lattice;
the change of branch of $\log s$ is itself an integral vector.
Thus translation by $r$ conjugates the linear deck action to
\eqref{maps:eq:F.24}.  This fixes the affine clutching, including its sign.

\item At the cusp use the integral marking of \eqref{maps:eq:F.13}, the fan \ExternalPaperRef{}{({E.4})}{III}
and the action \ExternalPaperRef{}{({E.5})}{III} with $C_3=C$.
Its maximal charts are $\mathbb C^4$ with map
$t=z_0z_1z_2z_3$, and its transition maps are the toric changes
of coordinates multiplied by the holomorphic units
$\exp(2\pi i C(t)\lambda)$.
The proof of Theorem \ExternalPaperRef{}{6.2}{III} used only boundedness of $C$,
integrality of the standard cusp, and unimodularity of that
fan.  It therefore proves smoothness, proper discontinuity
and properness for this model without a symmetry condition
on $C$.  Glue by the displayed integral marking and use no
additional cusp translation.

\item The three punctured overlap maps are holomorphic isomorphisms
with the same marked torus family, so their cocycle condition
is literal.  Proper local models and a compact complementary
base give a Hausdorff compact total space.  Every chart is
smooth.  The affine vectors, cusp fan and clutching are exactly
those in Theorems \ref{maps:src:F.2.5}--\ref{maps:src:F.2.6}.  Varying their bounded corrections
and varying the fibre complex structures along a smooth path
does not change these marked smooth fillings: on the finite
root models the underlying real torus family is the fixed
local-system bundle; on the cusp use the stratum-preserving
parameter vector fields and uniform properness estimate of
Theorem \ExternalPaperRef{}{6.3}{III}.  Thus the resulting underlying smooth manifold
is the one whose integral topology was computed there.

\end{enumerate}
\end{proof}


\section{The oriented diffeomorphism type}

\paragraph{The retained diagnostic complex.}
The earlier file \texttt{global\_candidate\_complex.py} sets the entire
cusp suspension map to zero and has the same free Betti numbers.
That file and its output are retained as a diagnostic record.  They
are not used in P1.12 to identify the actual higher-degree gluing.
P1.19 supplies the requested independent full global cell computation.
It proves that the local secondary cusp term is nonzero, gives the
explicit integral reduction to the corrected global matrix, and
recovers every homology group without global Poincar\'e duality.

\begin{theorem}\label{night:P1.13}
\StatementRecord{night:P1.13}{proved}{Historical events: GPT-6 Astra (OpenAI), ChatGPT Pro mode, effort Pro, S6\_06\_BLIND, 16--18 November 2026: targeted proof check. Also checked in this writing. Current writing events: GPT-6 Astra (OpenAI), ChatGPT Pro mode, effort Pro, S6\_10\_PAPERS, 18 November 2026: targeted proof check. GPT-6 Astra (OpenAI), ChatGPT Pro mode, effort Pro, S6\_10\_PAPERS continuation, 19 November 2026: targeted proof check; used and rechecked in this run. }{\cref{night:P1.12}; \S\ref{sec:V-known-inputs}; \cite[equations (1.2)--(1.4), pp.~3--4]{ffEsoleKang}; \cite[Theorem 2.2.1, p.~61]{ffPetersSterk}}
With the complex orientation the fourfold has
\[
\sigma(M_c)=2,\qquad H^4(M_c;\mathbb Z)\cong I_{3,1},\qquad
w_2(M_c)=0,
\]
and there is a diagonal integral basis in which
$p_1/2=(1,1,1,1)$.
Its Chern numbers are
$c_4=6$, $c_2^2=2$, and
$c_1^4=c_1^2c_2=c_1c_3=0$; its Pontryagin numbers are
$p_1^2=8$, $p_2=14$.
\end{theorem}
\begin{proof}
\begin{enumerate}[label=\arabic*.]
\item P1.12 implies $c_1=c_3=0$ as integral cohomology classes;
$c_4[M_c]=\chi(M_c)=6$.  The Chern-root identities give
$p_1=-2c_2$, $p_2=c_2^2+2c_4$.
The symbol-index interpretation in \S\ref{sec:V-known-inputs},
with the opened polynomial \cite[equation (1.2), p.~3]{ffEsoleKang},
gives integrality of the Todd genus. Expanding
$\prod x_j/(1-e^{-x_j})$ to degree four gives
\[
\frac{3c_2^2-6}{720}\in\mathbb Z,
\quad\text{hence}\quad c_2^2=2+240k\quad(k\in\mathbb Z).
\]
The signature polynomial \cite[equation (1.4), p.~4]{ffEsoleKang}, applied to
this closed oriented smooth eight-manifold gives
\[
\sigma=\frac{7p_2-p_1^2}{45}
      =\frac{3c_2^2+84}{45}=2+16k.
\]
The rank-four middle form has $|\sigma|\leq4$, so $k=0$.
This proves the displayed characteristic numbers.

\item Duality makes the middle form unimodular.  Its signature is two and
its positive and negative indices are three and one.  The rank-at-most-four integral classification
\cite[Theorem 2.2.1, p.~61]{ffPetersSterk} applied to these
positive and negative indices gives the odd form \(I_{3,1}\).
Since $H^2(M_c;\mathbb Z/2)=0$, its $w_2$ is zero.
For the complex tangent bundle,
$w_4=c_2\bmod2=(p_1/2)\bmod2$.
The Wu identity in degree four, with $w_1=w_2=0$, says that $w_4$ is
characteristic for the middle pairing.  Thus the vector $p_1/2$ is
characteristic and has square $2$.

\item We give the required integral isometry of this pair explicitly as
a terminating procedure.  In a diagonal basis a characteristic vector
has four odd coordinates.  Change their signs to write it as
$(a,b,c,d)$ with all positive and
$a^2+b^2+c^2=d^2+2$.
For $d>1$,
\[
d<a+b+c<3d.
\]
The first inequality follows by squaring and using the positive cross
terms; the second follows from Cauchy--Schwarz and $d>1$.
Reflection in the integral norm-two vector $(1,1,1,1)$ sends this
vector to itself minus $(a+b+c-d)(1,1,1,1)$.
Its fourth coordinate is $2d-a-b-c$, whose absolute value is strictly
smaller than $d$.  The reflection preserves the integral lattice and
characteristic parity.  Repeating, with sign changes at each step,
terminates at $d=1$, when the norm equation forces $a=b=c=1$.
This supplies an integral isometry taking $p_1/2$ to $(1,1,1,1)$,
and proves the claimed pair, not just its square.

\end{enumerate}
\end{proof}




\begin{theorem}[a complex structure on the stated connected sum]\label{maps:src:F.4.1}
\StatementRecord{maps:src:F.4.1}{proved}{Historical events: Claude Fable 5.2 (Anthropic), effort high, 15 November 2026: carried as background without fault found. Claude Fable 5.2 (Anthropic), effort high, 16 November 2026: targeted proof check. Claude Fable 5.2 (Anthropic), effort high, 17 November 2026: carried as background without fault found. GPT-6 Astra (OpenAI), ChatGPT Pro mode, effort Pro, S6\_09\_PAPERS, 17 November 2026: targeted proof check. Also checked in this writing. Current writing events: GPT-6 Astra (OpenAI), ChatGPT Pro mode, effort Pro, S6\_10\_PAPERS, 18 November 2026: targeted proof check. GPT-6 Astra (OpenAI), ChatGPT Pro mode, effort Pro, S6\_10\_PAPERS continuation, 19 November 2026: targeted proof check; used and rechecked in this run. \emph{Recorded proof scope.} GPT-6 Astra (OpenAI, ChatGPT Pro mode), effort Pro, 2026-11-17; Claude Fable 5.2 (Anthropic), effort high, 2026-11-17; Claude Fable 5.2 (Anthropic), effort high, 2026-11-15; GPT-6 Astra (OpenAI, ChatGPT Pro mode), S6\_06\_BLIND, effort Pro, 2026-11-17; Claude Fable 5.2 (Anthropic), effort high, 2026-11-16; also checked in this writing by GPT-6 Astra (OpenAI), ChatGPT Pro mode, effort Pro, 18 November 2026}{the hypotheses, notation and definitions stated here; \cref{maps:src:F.3.7}; \cref{maps:src:F.2.6}; \S\ref{sec:V-known-inputs}; the explicit integral reflections below; \cite[Corollary 1.4, p.~4]{ffMandal}; \cite[Definition 3.6 and Theorem 3.7(1), pp.~12--14]{IVVWallForm}; \cite[Theorem A, p.~178]{IVVKramerStolz}}


The compact complex fourfold of Theorem \ref{maps:src:F.3.7} is three-connected,
has Euler number six, and is orientation-preservingly diffeomorphic to
\[
 3\mathbb H\mathbb P^2\ \#\ (-\mathbb H\mathbb P^2).
 \tag{F.25}\label{maps:eq:F.25}
\]
Its integral middle lattice is $I_{3,1}$ and, in a diagonal integral
basis, $p_1/2=(1,1,1,1)$.  In particular the conclusion is an exact
smooth diffeomorphism, with no remaining homotopy-sphere ambiguity.

\end{theorem}


\begin{proof}
\begin{enumerate}[label=\arabic*.]
\item \textbf{The middle form and its characteristic vector are fixed.}
The integral calculation in \cref{maps:src:F.2.6} gives a free
middle group of rank four.  The characteristic-number calculation
gives $I_{3,1}$, $c_2^2=2$, and $p_1=-2c_2$, since $H^2=0$.
The reduction of $c_2$ is $w_4$.  The Wu formula on this spin
eight-manifold says that $w_4$ is characteristic for the middle
form.  Hence $v=p_1/2=-c_2$ is a characteristic vector with
$v^2=2$.  It is primitive, since the square of any integer
dividing $v$ would divide two.

\item \textbf{Explicit lattice isometries carry $v$ to
$(1,1,1,1)$.}
Testing the characteristic condition on the four standard basis
vectors of $I_{3,1}$ makes every coordinate of $v$ odd.
Independent sign changes are integral isometries.  After those
changes write $v=(a,b,c,d)$ with positive odd coordinates and
\[
a^2+b^2+c^2=d^2+2.
\]
For $d>1$, put $s=a+b+c$.  We have $s>d$, since
$s^2\geq a^2+b^2+c^2=d^2+2$, and
$s\leq\sqrt{3(d^2+2)}<3d$.  The vector
$u=(1,1,1,1)$ has square two.  Its reflection is the explicit
integral isometry
\[
R_u(z)=z-(z,u)u.
\]
It replaces the fourth coordinate by $2d-s$, whose absolute
value is strictly smaller than $d$.  All four new coordinates
remain odd because $s-d$ is even; in particular the new fourth
coordinate is nonzero.  Taking absolute values by sign
isometries and iterating therefore terminates at $d=1$.
The equation then forces $a=b=c=1$.  The product of these
specified reflections and sign changes is the required
isometry of the lattice together with $p_1/2$.

\item \textbf{The connected sum has the same extended normal form.}
Let $M_0=3\mathbb{HP}^2\#(-\mathbb{HP}^2)$ with its indicated
orientation.  Its standard middle generators have form $I_{3,1}$
and $p_1(TM_0)=2\sum_i x_i$; the normalized tautological-bundle
calculation is \cite[Lemma 2.1, p.~5, and Theorem 2.4, p.~7]{ffMandal}.
Thus the preceding isometry matches both the intersection form
and $p_1$.

For completeness these two pieces of information determine the
unstable normal datum used in the classification.  The normal
rank-four bundle of an embedded middle sphere has Euler number
equal to its self-intersection and Pontryagin number equal to
the evaluation of the ambient $p_1$, since $p_1(TS^4)=0$.
In the clutching group
$\pi_3\operatorname{Spin}(4)=\mathbb Z\oplus\mathbb Z$ the two
numbers, in compatible generators, are
\[
e=a-b,\qquad p_1=-2(a+b).
\]
This homomorphism is injective: its zero equations give
$a=b$ and $a=-b$, hence $a=b=0$.  No further normal-bundle
choice remains after the form and $p_1$ are matched.

\item \textbf{The known smooth classification gives an almost
diffeomorphism.}
Wall's classification, in the explicit extended-form formulation
\cite[Definition 3.6, pp.~12--13, and Theorem 3.7(1),
pp.~13--14]{IVVWallForm}, identifies these oriented
three-connected eight-manifolds up to connected sum with an
oriented homotopy eight-sphere.  Its input is exactly the normal
form identified in the preceding step.  Thus the constructed
manifold is oriented diffeomorphic to $M_0\#\Sigma$ for some
$\Sigma\in\Theta_8$.

\item \textbf{The homotopy sphere is absorbed in a displayed
summand.}
Kramer--Stolz's Theorem A
\cite[Theorem A, p.~178]{IVVKramerStolz} applies to
$\mathbb{HP}^2$: it is simply connected, has dimension eight,
and has integral homology $\mathbb Z$ in degrees zero, four
and eight and zero otherwise.  It gives
$\mathbb{HP}^2\#\Sigma\cong\mathbb{HP}^2$.
This diffeomorphism preserves orientation because both sides
have signature one.  Associativity of the oriented connected
sum and this map give the explicit comparison
\[
M_0\#\Sigma
\cong(\mathbb{HP}^2\#\Sigma)\#\mathbb{HP}^2
       \#\mathbb{HP}^2\#(-\mathbb{HP}^2)
\cong M_0.
\]
Composing it with the almost-diffeomorphism proves the claimed
oriented diffeomorphism.
\end{enumerate}
\end{proof}


\section{The normalized connection and the exact parameter domain}



\begin{theorem}\label{night:P1.15}
\StatementRecord{night:P1.15}{proved}{Historical events: GPT-6 Astra (OpenAI), ChatGPT Pro mode, effort Pro, S6\_06\_BLIND, 16--18 November 2026: targeted proof check. Also checked in this writing. Current writing events: GPT-6 Astra (OpenAI), ChatGPT Pro mode, effort Pro, S6\_10\_PAPERS, 18 November 2026: computation or formal certificate; targeted proof check. GPT-6 Astra (OpenAI), ChatGPT Pro mode, effort Pro, S6\_10\_PAPERS continuation, 19 November 2026: used and rechecked in this run. }{\cref{night:P1.1}, \cref{night:P1.8}}
The rank-six flat extension, as a complex flat bundle with its invariant
line, has the following explicit algebraic connection model.  Set
$\omega_k=x^kdx/y$, $0\leq k\leq3$, and use the quotient frame
\[
(w_2,w_1,w_0,w_{-1},w_{-2})
 =(\omega_{01},\omega_{02},\omega_{03},\omega_{13},\omega_{23}),
\qquad \omega_{12}=-3\omega_{03}.
\]
These are meromorphic frames of
$(L^2,L,\mathcal O,L^{-1},L^{-2})$, where $L=T^2$ on the root
orbifold of orders $5,4$.  If a frame is written as a row of vectors,
its derivative is right multiplication by
\[
\Gamma_6=\begin{pmatrix}0&b\\0&\Gamma_W\end{pmatrix},\qquad
b=\frac1{\lambda-\lambda_*}
 \left(-\frac{25}{6\lambda^2},-\frac{125}{48\lambda},0,1,0\right),
\tag{1.15.1}
\]
where
\[
\resizebox{\textwidth}{!}{$\displaystyle
\Gamma_W=\frac1{3125\lambda+256}
\begin{pmatrix}
-8(625\lambda+48)/\lambda&-8&10\lambda&75\lambda^2/2&0\\
-125/(2\lambda)&-4(625\lambda+48)/\lambda&8&0&75\lambda^2/2\\
300/\lambda^2&-375/\lambda&0&-48&60\lambda\\
120/\lambda^3&0&125/(2\lambda)&4(625\lambda+48)/\lambda&-8\\
0&120/\lambda^3&50/\lambda^2&-125/(2\lambda)&8(625\lambda+48)/\lambda
\end{pmatrix}.$}
\tag{1.15.2}
\]
This model specifies a canonical algebraic splitting, denoted
$\sigma_b$.  The compatible contraction used for the real parameter is
$\alpha=(625\pi i/128)\alpha_{\rm alg}$, where
$\alpha_{\rm alg}(w_0)=4/3$ and it vanishes on $w_1,w_2$.
\end{theorem}
\begin{proof}
\begin{enumerate}[label=\arabic*.]
\item Reduction modulo exact differentials is completely explicit.  To
differentiate $\omega_k$, solve the polynomial identity
\[
-\tfrac12x^k\partial_\lambda f
 =f\sum_{j=0}^3G_{kj}x^j+R_k'f-\tfrac12R_kf',
 \qquad\deg R_k\leq4,
\]
where $f=x^5+\lambda^3x+\lambda^4$.  The last two terms are the
numerator of $d(R_k/y)$.  Matching the nine coefficients gives the
four-by-four Gauss--Manin matrix exactly.
At the branch point at infinity, a local coordinate $z=x^{-1/2}$
gives $\omega_k=-2z^{2-2k}(1+O(z^8))\,dz$.
Residues of a primitive times a differential give the alternating
pairings $J_{03}=4/3$, $J_{12}=4$, with their negatives in reverse
order and all other entries zero.  This form is nondegenerate.
The invariant bivector is therefore $3\omega_{03}+\omega_{12}$.
Taking the exterior-square connection and this quotient gives
(1.15.2), as checked symbolically in \texttt{de\_rham\_connection.py}.

\item The root-chart powers computed from $x^kdx/y$ are
$u^{6k-9}$ at zero and $u^{6-4k}$ at infinity.
Their nondegenerate pairing shows that the four resulting regular
frames form the full de Rham bundle on those charts.  At the node
the computed connection has a simple pole and square-zero residue;
hence these are the nilpotent logarithmic extension frames there.
Thus the four line bundles are
$T^3,T,T^{-1},T^{-3}$, and their quotient exterior-square lines are
the stated powers of $L$.

\item All $H^1(L^j)$ for $j=0,\pm1,\pm2$ vanish by the coarse-degree
calculation already used in P1.8.  Consequently every holomorphic
extension of this $W$ by $\mathcal O$ splits as a bundle.
In such a splitting its flat connection is an upper block extension
by a row $b$.  Regularity at the finite root charts and a logarithmic
pole at the cusp require, entry by entry,
\[
b=\frac1{\lambda-\lambda_*}
 (a\lambda^{-2},b'\lambda^{-1},0,c,d\lambda).
\]
For example the $L^j$ entry, after $\lambda=u^4$, has exponent
$3j+3+4p$ for a term $\lambda^p/ (\lambda-\lambda_*)$;
the infinity chart requires $-4j-1-5p\geq0$.
These two inequalities leave exactly the four displayed terms.
The equation that the full cusp residue have square zero is
$b_{\rm res}(\Gamma_W)_{\rm res}=0$.
Its solution space has basis
\[
(-25/6,-125/48,1,0),\qquad(5/3,-25/3,0,1).
\]
The sole nonconstant holomorphic splitting change is the row selecting
$w_0$, and its connection coboundary is the third row of $\Gamma_W$.
In these two basis vectors that row is $(-48b_1+60b_2)/3125$.
Thus there is one nontrivial flat extension direction modulo splitting
changes and rescaling of the invariant line.  The first vector has
rank-three square-zero cusp residue, so it represents that direction.
The supplied extension is nontrivial, since its dual has no common
invariant vector by P1.1.  It is therefore isomorphic to (1.15.1).
This proves the connection identification without inferring it from
the subsequent numerical matching.

\item For the last normalization, the kernel of the quotient cusp residue
on $(w_2,w_1,w_0)$ is generated by
$v=(-4096/1171875,64/1875,1)^t$.
It satisfies $b_{\rm res}(v)=625/192$ and
$\alpha_{\rm alg}(v)=4/3$.
The logarithmic monodromy is $2\pi i$ times this residue.
The exact cusp identity of P1.8 characterizes the contraction on this
line by $N_V(\widetilde v)=\alpha(v)e$.
Since the space of global contraction forms is one-dimensional, the
ratio $(625/192)/(4/3)=625/256$ fixes its scalar everywhere and gives
the displayed $625\pi i/128$.

\end{enumerate}
\end{proof}


\paragraph{An explicit formula for $h$ and the normalization comparison.}
Let $Y'=Y\Gamma_6$.  Normalize $Y(1)=X$ by requiring that continuation
around the specified meridians intertwine exactly with $A,B$, and that
$Xe_0=e_0$.  Existence follows from P1.15; uniqueness follows from the
scalar centralizer, whose exact linear equations are in
\texttt{centralizer\_torsion.py}.
At each nonspecial point there is a unique pair $v_2,v_1\in\mathbb C$
such that the lower five coordinates of
$Y(0,v_2,v_1,-i,0,0)^t$ are real.
The intersection-line proof in P1.8 and the normalization of $\alpha$
prove existence and uniqueness.  Therefore the following is a fully
specified period formula for the real function associated with
$\sigma_b$:
\[
h_b(\lambda)=\frac{96}{625\pi}
 \operatorname{Im}\bigl(Y_{01}v_2+Y_{02}v_1-iY_{03}\bigr).
\tag{1.15.3}
\]
Its periods for the real parameter $c$ are the three columns
$Y_{*,1},Y_{*,2},Y_{*,3}-(625\pi/96)c\,e_0$.

This is a different splitting normalization from the one chosen in
P1.8.  To compare them exactly, parallel-transport the fixed marking
from $\lambda=1$ to $\lambda=\lambda_*/2$ along a specified path in the
upper half-plane, and put
$a=(\sigma_b(g))_e$ at that point, where
$g=(96/(625\pi i))w_0$ and $\alpha(g)=1$.
Then the splitting of P1.8 is
$\sigma=\sigma_b-a\alpha e$, so $h=h_b-\operatorname{Im}a$.
For the full complex splitting parameter $\xi$, the comparison is
$\xi_b=\xi-a$.  This gives the actual translation between the two
families, including its possible real part; it does not silently
identify their real parameter slices.



\begin{theorem}\label{night:P1.16}
\StatementRecord{night:P1.16}{proved}{Historical events: GPT-6 Astra (OpenAI), ChatGPT Pro mode, effort Pro, S6\_06\_BLIND, 16--18 November 2026: targeted proof check. Also checked in this writing. Current writing events: GPT-6 Astra (OpenAI), ChatGPT Pro mode, effort Pro, S6\_10\_PAPERS, 18 November 2026: computation or formal certificate; targeted proof check. GPT-6 Astra (OpenAI), ChatGPT Pro mode, effort Pro, S6\_10\_PAPERS continuation, 19 November 2026: used and rechecked in this run. }{\cref{night:P1.15}, \cref{night:P1.18}, \cref{night:P1.8}}
In the canonical splitting $\sigma_b$ of P1.15,
\[
\begin{aligned}
-1.0342766224163217994942488453783461&<h_b(1),\\
h_b(1)&<-1.0342766224163217994942488453783460.
\end{aligned}
\]
In particular $h_b(1)<-1$.  The parameter $c=5/(3\pi)$ is outside the
global transversality domain.  The lifted period matrices at
$\lambda=1,2$ have been evaluated by two different continuation methods
agreeing to more than thirty decimal digits.
\end{theorem}
\begin{proof}
\begin{enumerate}[label=\arabic*.]
\item The file \texttt{certified\_lift\_bound.py} propagates complex balls
without dropping their radii.  We describe the explicit truncation
bound so the numerical inequality has a mathematical certificate.
On a Taylor step from $z$ to $z+h$, choose a disk of radius $R$ avoiding
the two finite poles.  Use a fixed positive diagonal matrix $S$ on
this disk.  Each entry of the connection is a polynomial divided by
$\lambda^p(\lambda-\lambda_*)$; the triangle inequality gives a
rigorous bound $M$ for the maximum row-sum norm of
$S^{-1}\Gamma_6S$ on the disk.  Along a radius, the integral equation
and Gronwall give a bound $\exp(MR)$ for the fundamental solution
normalized to the identity at its center.  Cauchy's coefficient
estimate then bounds the omitted Taylor tail after degree $N$ by
\[
\exp(MR)\frac{(|h|/R)^{N+1}}{1-|h|/R}.
\tag{1.16.1}
\]
The program adds this error to each entry with factor $S_i/S_j$.
All coefficient evaluations, matrix products, and linear solves retain
their outward-rounded complex-ball enclosures.

\item For both finite meridians we use rational polygonal paths, and choose
the order adaptively until (1.16.1) is below $10^{-65}$.
There are $152$ steps in total and the largest Taylor order is $239$.
The thirty independent equations determining the normalized
intertwiner are solved by interval matrix inversion; exact existence
and uniqueness are P1.15 and the scalar-centralizer calculation.
The real intersection vector in (1.15.3) is likewise obtained by an
invertible four-equation subsystem, whose exact consistency follows
from P1.8.  Substitution into (1.15.3) gives the certified enclosure
\[
\begin{aligned}
h_b(1)\in\bigl[&
-1.03427662241632179949424884537834605259496477790218821132\\
&\qquad\mathbin{\pm}8.11\cdot10^{-57}\bigr].
\end{aligned}
\]
which implies the stated slightly wider decimal bounds.
Every step radius, norm bound, order and tail bound is exported in
\texttt{certified\_lift\_bound.json}.  Since $5/(3\pi)<1$, the value
$h_b(1)+5/(3\pi)$ is negative, whereas $h_b+c$ tends to positive
infinity at the cusp.  A path joining these points must contain a
zero.  This proves inadmissibility, not merely failure of a numerical
test at one point.

\item For the independent period evaluation,
\texttt{frobenius\_period\_check.py} uses regular finite-root frames
and their deck eigencharacters, without integrating a closed loop.
At zero the frame powers are $(0,6,3,0,-3,-6)$ in $u$ with
$\lambda=u^4$, and at infinity they are
$(0,-8,-4,0,4,8)$ with $\lambda=u^{-5}$.
The transformed connections are regular at $u=0$ and have denominators
$3125u^4+256$ and $3125+256u^5$, respectively.
Their Taylor recursions therefore give the connection matrices from
the center to $u=1/4$ and $u=1/2$.  Conjugating the exact diagonal
deck characters by these matrices gives the finite monodromies.
Their normalized intertwiner agrees with the closed-loop calculation
to $4.5\cdot10^{-94}$ in maximum entry norm.
The resulting period planes at $\lambda=1,2$ and their real
transversality determinants are all recorded in the accompanying
JSON files; for example at $c=10$ the determinants are purely imaginary
with positive imaginary parts
\[
\begin{gathered}
88279007523.2167451846066607088299709479845\ldots,\\
3178499608.3637608517583236010752861132441\ldots.
\end{gathered}
\]
These last two local checks do not by themselves prove that $c=10$
is globally admissible.  Global admissibility is exactly P1.8's
inequality involving the minimum.  The independent root-chart
calculation also gives
\[
\begin{split}
h_b(\infty)&=-1.094540250311402780116268263054945703725533\ldots,\\
h_b(0)&=-0.931793376356523505430795105796011114990828\ldots.
\end{split}
\]
P1.18 proves, by a complete interval covering argument, that the
first value is the unique global minimum.

\end{enumerate}
\end{proof}




\begin{theorem}\label{night:P1.18}
\StatementRecord{night:P1.18}{proved}{Historical events: GPT-6 Astra (OpenAI), ChatGPT Pro mode, effort Pro, S6\_06\_BLIND, 16--18 November 2026: targeted proof check. Also checked in this writing. Current writing events: GPT-6 Astra (OpenAI), ChatGPT Pro mode, effort Pro, S6\_10\_PAPERS, 18 November 2026: computation or formal certificate; targeted proof check. GPT-6 Astra (OpenAI), ChatGPT Pro mode, effort Pro, S6\_10\_PAPERS continuation, 19 November 2026: used and rechecked in this run. }{\cref{night:P1.15}, \cref{night:P1.16}, \cref{night:P1.8}}
In the algebraic splitting $\sigma_b$, the function $h_b$ has its
unique global minimum at the finite-root point $\lambda=\infty$.
The minimum is the exact Frobenius period constant
\[
 h_{b,\min}=-\frac{96}{625\pi}\operatorname{Re}(XH_\infty)_{03},
 \qquad
 h_{b,\min}=-1.0945402503114027801162682630549457037\ldots,
\tag{1.18.1}
\]
where $X$ is the normalized integral intertwiner of P1.15 and
$H_\infty$ is the regular infinity-root limit of its fundamental
solution, normalized to be the identity at $\lambda=1$ before
multiplication by $X$.  Formula (1.18.1), the rational connection,
and those two normalizations specify the constant exactly.
A certified enclosure has radius $4.04\cdot10^{-56}$ about the
56-digit midpoint retained in
\path{code/part1/minimum_certificate_summary.json}.
Consequently the exact domain of completed transverse periods is
\[
 c>-h_{b,\min},\qquad
 \operatorname{Im}\xi>-h_{b,\min}
\]
for the real and complex parameters, respectively.
The member $c=10$ used in the numerical period checks is therefore
admissible on the whole completed base.

For the point-normalized splitting of P1.8, continued from
$1$ to $\lambda_*/2$ by the polygon
$1,1+i,\lambda_*/2+i,\lambda_*/2$, the exact comparison constant
$a_{\rm split}$ of (1.15.4) has certified value
\[
\begin{split}
 \operatorname{Re}a_{\rm split}
   &=-0.7099885679126427771723473455012308149\ldots,\\
 \operatorname{Im}a_{\rm split}
   &=-0.9134017059364153362143229223828373269\ldots.
\end{split}
\]
Thus the minimum in that splitting is
$-0.1811385443749874439019453406721083768\ldots$,
attained at the same base point.  Its real parameter interval is
\[
 c>0.1811385443749874439019453406721083768\ldots.
\]
The full parameter comparison remains
$\xi_b=\xi_{\rm original}-a_{\rm split}$; both its real and imaginary
parts are necessary to identify the two normalized families.
\end{theorem}
\begin{proof}
\begin{enumerate}[label=\arabic*.]
\item We give the error bounds and the complete covering argument behind
the interval computation.  This is a finite proof of the minimum,
not an inference from the preliminary grid search.  All coefficients,
accepted rectangles and local continuation bounds are retained in the
four \texttt{certified\_minimum\_*} programs and their JSON outputs.
The file \texttt{minimum\_certificate\_summary.json} records their
SHA-256 digests as well as sharp interval endpoints.

\item Put $a=|\lambda_*|$.  The normalized root coordinates are
\[
 z_\infty=(a/\lambda)^{1/5},\qquad
 z_0=(\lambda/a)^{1/4}.
\]
In the regular root frames of P1.16 their connections have the form
\[
 G(z)=\frac{\sum_{j=0}^{m-1}B_jz^j}{1+z^m},\qquad m=5,4.
\tag{1.18.2}
\]
The entries of $B_j$ are the exact algebraic numbers obtained by
substituting the positive real number $a^{\pm1/m}$ in the rational
connection.  The generator of the finite deck group acts on these
frames by their stated eigencharacters.  Since its flat matrix is
real and fixes $e$, it leaves $h_b$ invariant.  It therefore suffices
to cover the sector $0\leq\arg z\leq2\pi/m$ in each root disc.

\item Use the positive diagonal weights
\[
 (1,1/128,1/128,1/32,1/64,1/32)
\]
and the corresponding weighted maximum row norm.  Exact ball
arithmetic gives
\[
 \left\|\sum_j |B_j|\right\|<1.447\quad(m=5),\qquad
 \left\|\sum_j |B_j|\right\|<1.135\quad(m=4).
\]
Here the absolute values are entrywise, before taking the row norm.
Let $B$ denote the actual enclosed bound.  The scalar inequality
\[
 \frac1{1-t^m}\leq\frac1{m(1-t)}+\frac{m-1}{m},\qquad0\leq t<1,
\]
follows on multiplying by $m(1-t)(1+\cdots+t^{m-1})$;
the remaining assertion is
$\sum_{j=1}^{m-1}t^j\geq(m-1)t^m$.
Gronwall's inequality consequently bounds the fundamental solution
$U(0)=I$ on $|z|\leq R<1$ by
\[
 K_R=\exp\!\left[B\left(-\frac{\log(1-R)}m
                    +\frac{m-1}{m}R\right)\right].
\]
Cauchy's estimate bounds the tail after degree $n$ on $|z|\leq r<R$
by
\[
 K_R\frac{(r/R)^{n+1}}{1-r/R}.
\tag{1.18.3}
\]
The code uses $R=0.999$.  The coefficients are computed directly
from (1.18.2):
\[
 nU_n=-(n-m)U_{n-m}
       +\sum_{j=0}^{\min(m-1,n-1)}U_{n-1-j}B_j,
\]
with the first term omitted when $n<m$.
There is no derivative approximation in this recursion.
The regular limit frame is determined by certified continuation
from $\lambda=1$ to a positive root-chart point and inversion of
this convergent series.  At the infinity center it yields
\[
 h_b(\infty)\in
 [-1.094540250311402780116268263054945704
       \ \mathbin{\pm}\ 3\cdot10^{-37}].
\tag{1.18.4}
\]
The corresponding enclosure at zero has center
\[
-0.9317933763565235054307951057960111149908282005165194125
\]
and radius $7.91\cdot10^{-56}$.

\item We explain two numerical details which are needed to make these
bounds rigorous and efficient.  At a root center, let $H$ denote
the actual limiting period frame, and form the real matrix whose
five columns are
\[
 \operatorname{Im}H_{\mathrm{low},1},\quad
 \operatorname{Re}H_{\mathrm{low},1},\quad
 \operatorname{Im}H_{\mathrm{low},2},\quad
 \operatorname{Re}H_{\mathrm{low},2},\quad
 \operatorname{Im}H_{\mathrm{low},3}.
\]
Its certified inverse gives a constant real change of the lower five
flat coordinates.  Subtracting the corresponding real linear
combination of those coordinates from the top coordinate also
makes the first two top entries zero at the center.  This is an
explicit real invertible change fixing $e$.  It does not change
$h_b$: on the real intersection vector the extra top term is real,
so its imaginary part is zero.  This change greatly improves the
four real equations used to evaluate (1.15.3).

\item For a polynomial $p(z)=\sum c_jz^j$ with ball coefficients, evaluation
on a rectangle centered at $z_c$, of enclosing radius $\rho$, uses
its midpoint-coefficient polynomial at the point $z_c$.  Coefficient
uncertainty is bounded by
$\sum_j\operatorname{rad}(c_j)|z_c|^j$, and variation over the
rectangle by
$\rho\sum_j j|c_j|r^{j-1}$, where $r$ bounds the modulus on it.
The tail (1.18.3) is then added.  The point evaluation uses 800-digit
ball arithmetic; its rounding radius is retained.  These explicit
bounds avoid the excessive interval enlargement of long rectangular
Horner recursions.  No error radius is discarded.

\item On the square $|\operatorname{Re}z_\infty|,
|\operatorname{Im}z_\infty|\leq0.1$, differentiating the exact linear
system for the real intersection vector gives the following
certified bounds for the real Hessian of $h_b$:
\[
\begin{split}
 0.2898644&<h_{xx},h_{yy}<0.3946911,\\
 |h_{xy}|&<0.0524256,\qquad
 \det\operatorname{Hess}(h_b)>0.0757786.
\end{split}
\tag{1.18.5}
\]
The derivatives of the period frame are evaluated by
$Y'=YG$ and $Y''=Y(G^2+G')$; the derivatives of the intersection
vector follow by differentiating its invertible four-by-four real
system.  Finite rotation invariance forces $dh_b(0)=0$ in this
root chart.  Positive definiteness on the convex square therefore
proves a strict minimum at its center.

\item Outside that square, a recursive dyadic rectangle cover proves
$h_b>h_b(\infty)$ throughout $|z_\infty|\leq0.98$ in its sector.
It has 900 accepted rectangles, maximum subdivision depth eight.
The zero-root disc $|z_0|\leq0.98$ is certified by 463 rectangles,
also of maximum depth eight.  A rectangle is discarded only when
its entire intersection with the required sector and disc is empty,
or when it lies in the already proved local square.  Every remaining
rectangle is either certified by the strict interval inequality or
subdivided.  The complete records show that the subdivision terminates;
there is no unresolved rectangle.

\item It remains to cover a cusp neighborhood and the collar between these
two root discs.  Write
\[
 q=(\lambda-\lambda_*)/\lambda_*.
\]
In this coordinate the connection is $R/q+A(q)$, with $R^2=0$
and $A$ analytic on $|q|<1$.  A regular-singular factorization is
\[
 Y(q)=H\exp((\log q)R)P(q),\qquad P(0)=I.
\]
The recursion for its analytic factor is
\[
 (n-\operatorname{ad}_R)P_n
   =\sum_{j=0}^{n-1}P_{n-1-j}A_j,
 \qquad \operatorname{ad}_R(Q)=QR-RQ.
\]
Because $R^2=0$, the cube of this adjoint operator vanishes and its
inverse is the finite expression
$1/n+\operatorname{ad}_R/n^2+\operatorname{ad}_R^2/n^3$.
In an exact dyadic diagonal norm the code proves
$\|R\|<0.841$ and uses the sharper recorded bound
\[
 C=1+2\|R\|+4\|R\|^2<5.508.
\]
On the circle of radius $0.4$, the absolute-coefficient majorant of
$A$ is less than $4.895$.  The coefficient recursion is dominated
by the scalar equation $p'=Ca_{\rm abs}p$; hence
$\|P\|\leq\exp(C\cdot4.895\cdot0.4)$ there.  Cauchy's estimate
now gives the same explicit geometric tail bound as (1.18.3), with
this exponential in place of $K_R$ and $R=0.4$.

\item The constant $H$ is obtained by certified continuation to $q=0.01$
along a specified polygon in the upper half-plane.  The matrix
$2\pi iHRH^{-1}$ is integral, since it is the logarithm of the
corresponding marked local monodromy.  Its ball enclosure contains
exactly one integer matrix; that matrix and an integral Jordan
marking $J$ are exported.  They satisfy
\[
 J^{-1}(2\pi iHRH^{-1})J=\begin{pmatrix}0&I_3\\0&0\end{pmatrix}.
\]
At the fixed parameter $c_*=-h_b(\infty)$, the last three columns of
$J^{-1}HP(q)$ specifying $F_{c_*}$ have blocks $T(q),D(q)$.
The lower block is certified invertible on the stated cusp disc.
The period matrix is therefore
\[
 Z(q)=\frac{\log q}{2\pi i}I_3+C(q),\qquad C(q)=T(q)D(q)^{-1}.
\]
Put $B(q)=\operatorname{Im}C(q)$ and $y=-\log|q|/(2\pi)$.
Then
\[
 \det\operatorname{Im}Z
 =y^3+\operatorname{tr}(B)y^2
   +\Bigl(\sum_{i<j}(B_{ii}B_{jj}-B_{ij}B_{ji})\Bigr)y+\det B.
\tag{1.18.6}
\]
On the whole disc $|q|\leq0.01$, interval substitution in (1.18.6),
followed by $y=-\log(0.01)/(2\pi)+v$, gives four strictly positive
coefficients in $v$; their certified lower bounds, in increasing
order, are $0.8111$, $3.6097$, $3.5346$, and $1$.
The sharp endpoints are in \texttt{cusp\_cubic\_bounds.json}.
Since $v\geq0$, the determinant is positive
throughout the punctured disc.  Near the cusp $h_b+c_*>0$ by P1.8;
continuity and transversality then give this same strict inequality
throughout that disc.  A separate finite cover of
$0.01\leq|q|\leq0.25$, using the same analytic factor and (1.18.6),
contains 7033 certified rectangles, of maximum depth seven.  It
extends the inequality to the whole punctured disc $|q|\leq0.25$.

\item The only remaining region, expressed in the infinity-root chart, is
\[
 0.98\leq |z_\infty|\leq0.98^{-4/5},\qquad
 0\leq\arg z_\infty\leq2\pi/5,
 \qquad |q|\geq0.25,
\]
where $q=-z_\infty^{-5}-1$.
Each rectangle in this collar is reached from radius $0.95$ by
Taylor continuation of (1.18.2).  The local analytic radius is
$0.449$ times the distance to the nearest root of $1+z^5$.
The numerator is bounded in the same diagonal norm and the
denominator by the product of the five distances to those roots.
Gronwall and Cauchy give an explicit tail for every step, exactly
as in P1.16.  There are 427 steps and 117 accepted rectangles,
of maximum depth eight; each proves $h_b>h_b(\infty)$.

\item These regions cover every point of the completed base except the
cusp itself, where $h_b$ tends to positive infinity.  Every inequality
away from the infinity center is strict.  This proves both the
global minimum and its uniqueness, and P1.8 gives the exact parameter
domain.

\item Finally, \texttt{normalization\_shift.py} evaluates
$a_{\rm split}=96Y_{03}(\lambda_*/2)/(625\pi i)$ by the same
certified continuation, with the polygon stated above.
The real and imaginary radii are respectively less than
$4.31\cdot10^{-56}$ and $3.15\cdot10^{-56}$.
Subtracting its imaginary part from (1.18.4) gives the stated
point-normalized minimum and interval.  This is the explicit
splitting morphism already established in P1.15, now numerically
certified.

\end{enumerate}
\end{proof}


\section{Cohomology, deformations and isomorphisms}

We retain the complex orientation and the integral marking of P1.1.
Let $\widehat B$ be the root orbifold of orders $5,4$, let
$h:M\to\widehat B$ be the map supplied by the two multiple divisors,
and let $\pi:\widehat B\to\mathbb P^1$ be its coarse map.
Write $L=T^2$; in the pencil coordinate it is
\[
 L=\mathcal O_{\widehat B}\bigl(\tfrac45[\infty]
                  -\tfrac34[0]\bigr),\qquad \deg L=\tfrac1{20}.
\]
The cusp point has line bundle $L^{20}$, and
$T_{\widehat B}=L^9$. These identities include the stabilizer
characters: their quotients have integral divisors of degree zero
on $\mathbb P^1$. Denote the cusp by $W$, its normalization by $Y$,
and its reduced singular locus by $D$.

\begin{theorem}\label{night:P1A.1}
\StatementRecord{night:P1A.1}{proved}{Historical events: GPT-6 Astra (OpenAI), ChatGPT Pro mode, effort Pro, S6\_06\_BLIND, 16--18 November 2026: targeted proof check. Also checked in this writing. Current writing events: GPT-6 Astra (OpenAI), ChatGPT Pro mode, effort Pro, S6\_10\_PAPERS, 18 November 2026: computation or formal certificate; targeted proof check. GPT-6 Astra (OpenAI), ChatGPT Pro mode, effort Pro, S6\_10\_PAPERS continuation, 19 November 2026: used and rechecked in this run. }{\cref{night:P1.10}, \cref{night:P1.8}, \cref{night:P1.9}}
For every admissible member, its vertical translation bundle on
$\widehat B$ and its first coherent direct image are
\[
 E_0=F^\vee=\mathcal O\oplus L^{-1}\oplus L^{-2},\qquad
 G=R^1h_*\mathcal O_M=\mathcal O\oplus L^{-1}\oplus L^{-2}.
\]
Moreover $R^qh_*\mathcal O_M=\bigwedge^qG$ for $0\le q\le3$,
$K_M=h^*L^{-6}$, and
\[
 (h^0(\mathcal O_M),\ldots,h^4(\mathcal O_M))=(1,1,0,0,0).
\]
\end{theorem}
\begin{proof}
\begin{enumerate}[label=\arabic*.]
\item The evaluation of the marked lattice on the lifted plane $F$
identifies the Lie bundle of the smooth tori with $F^\vee$.
P1.8 gives $F=\mathcal O\oplus L\oplus L^2$; its lift as a bundle
has the same splitting. The logarithmic toric vector fields extend
this Lie bundle over the cusp, so the identification holds on the
whole orbifold, including both finite root charts.

\item On the smooth locus the constant-cocycle description of torus
cohomology gives $G=V/F$. In the rank-five orthogonal quotient,
$K=\ker(\alpha:P\to\mathcal O)=L\oplus L^2$ is the radical of
the restricted pairing on $P$. Indeed every wedge containing the
fixed line pairs to zero with another such wedge except through
their contraction components; removing that scalar component
leaves exactly $K$. The nondegenerate pairing consequently
identifies $W/P$ with $K^\vee$. Thus
\[
 0\longrightarrow\mathcal O e\longrightarrow V/F
   \longrightarrow L^{-1}\oplus L^{-2}\longrightarrow0.
\]
For every integer $j$, descent by invariant local Laurent series gives
\begin{equation}\label{eq:P1Adegree}
 \pi_*L^j=\mathcal O_{\mathbb P^1}(d(j)),\qquad
 d(j)=\lfloor4j/5\rfloor+\lfloor-3j/4\rfloor.
\end{equation}
The extension group $H^1(L\oplus L^2)$ therefore vanishes, so this
sequence splits as a holomorphic vector-bundle sequence.

\item We verify that this quotient is the coherent extension at the cusp.
The normalization resolution of $\mathcal O_W$ has one toric
threefold, seven toric surfaces, twelve toric curves and six points.
All positive coherent cohomology of those smooth complete toric
strata vanishes. Their constants form precisely the cellular
cochain complex of the dual three-torus, with the incidences in
\path{code/part1/cusp_complex.json}. Its cohomology is
$\bigwedge^q\mathbb C^3$, and its cup product is the exterior product.
The invariant flat cohomology $\ker N$ maps to these classes by
constant Cech cocycles on the toroidal covering. A component
translation in that cover evaluates the corresponding lattice
coordinate. In the nilpotent logarithmic frame, $\ker N$ maps
isomorphically to $V/F$, since $N:F\to\ker N$ is an isomorphism.
Thus these are the actual extending quotient classes, rather than
an unrelated basis of the same dimension.

\item Flatness over each smooth root disc follows because its base is a
DVR and the smooth total space is torsion free over it. Properness
was proved in P1.9--P1.10. The coherent fibre dimensions are the
constant binomial dimensions on smooth tori, finite root fibres,
and the cusp. Grauert's \S7 Satz~5, with these checked hypotheses,
gives local freeness and base change. The cup-product maps are
isomorphisms on every fibre, including the preceding central
cellular model, so they identify the remaining direct images with
$\bigwedge^qG$.

\item There is an exact sequence of locally free logarithmic bundles
\begin{equation}\label{eq:P1Alogcotangent}
0\longrightarrow h^*L^{11}\longrightarrow\Omega_M^1(\log W)
 \longrightarrow h^*F\longrightarrow0.
\end{equation}
In a chart $q=x_0\cdots x_r$, its first arrow sends
$dq/q$ to $\sum dx_i/x_i$ and its quotient is the height-zero
logarithmic cotangent bundle of the torus. These formulas also
prove exactness on every stratum. Taking determinants and
cancelling $\mathcal O_M(W)=h^*L^{20}$ yields
$K_M=h^*L^{11+3-20}=h^*L^{-6}$.

\item For any bundle which is a sum of powers of $L$, projection and
the exact finite averaging projector compute its direct images
by tensoring with $\bigwedge^qG$ and applying
\eqref{eq:P1Adegree}. All exponents $-1,\ldots,-5$ have coarse
degree $-1$, while exponent zero has degree zero. The Leray
sequence on the curve then gives the asserted structure-sheaf
cohomology. This includes the nonreduced coarse fibres, since
the computation was made before exact finite descent.

\end{enumerate}
\end{proof}




\begin{theorem}\label{night:P1A.2}
\StatementRecord{night:P1A.2}{proved}{Historical events: GPT-6 Astra (OpenAI), ChatGPT Pro mode, effort Pro, S6\_06\_BLIND, 16--18 November 2026: targeted proof check. Also checked in this writing. Current writing events: GPT-6 Astra (OpenAI), ChatGPT Pro mode, effort Pro, S6\_10\_PAPERS, 18 November 2026: computation or formal certificate; targeted proof check. GPT-6 Astra (OpenAI), ChatGPT Pro mode, effort Pro, S6\_10\_PAPERS continuation, 19 November 2026: used and rechecked in this run. }{\cref{night:P1.15}, \cref{night:P1A.1}}
The complete Hodge table of every admissible fourfold is the
following, with rows indexed by $p$ and columns by $q$:
\[
\begin{array}{c|ccccc}
h^{p,q}&0&1&2&3&4\\\hline
0&1&1&0&0&0\\
1&0&2&1&0&0\\
2&0&1&6&1&0\\
3&0&0&1&2&0\\
4&0&0&0&1&1
\end{array}
\]
\end{theorem}
\begin{proof}
\begin{enumerate}[label=\arabic*.]
\item All maps needed for the two nontrivial rows can be evaluated at
the cusp. We give their matrices and the reason evaluation
determines their global ranks. In the algebraic splitting of
P1.15 write the three generators of $F$ as $(f_0,f_1,f_2)$ of
weights $(0,1,2)$, and those of $G$ as $(g_0,g_1,g_2)$ of
weights $(0,-1,-2)$. The residue of the quotient connection is
\[
R:F_0\longrightarrow G_0,\qquad
R=\begin{pmatrix}0&A&B\\C&0&D_0\\E&D_0&0\end{pmatrix},
\]
where
\begin{align*}
A&=-125(192z-3125)/12288,&
B&=-78125(96z+3125)/393216,\\
C&=-125/512,&E&=78125/32768,\\
D_0&=-146484375/2097152.
\end{align*}
Here the lifted first vector is $w_0+(4/3)ze$.
These entries result from the exact Gauss--Manin reduction of
P1.15, with the quotient projection taken after lifting $F$.
They are independently recomputed in
\path{code/part1/analytic_cohomology.py}.

\item The connecting map of \eqref{eq:P1Alogcotangent}, on relative
cohomology, sends $f\otimes\eta$ to $R(f)\wedge\eta$ at the
cusp. To see the actual extension class, trivialize logarithmic
forms on the toric covering. A deck translation changes a
relative logarithmic form by its period derivative times
$dq/q$. The constant Cech class of this derivative is exactly
the residue of the quotient connection. This proves the
formula with the indicated frames. It also proves its exterior
power version, with the usual alternating signs.

\item The floor formula gives the following absolute cohomology:
\[
\begin{array}{c|ccccc}
&h^0&h^1&h^2&h^3&h^4\\\hline
h^*F&1&3&2&0&0\\
h^*L^{11}&0&2&3&1&0\\
h^*(\bigwedge^2F)&0&2&3&1&0\\
h^*(L^{11}\otimes F)&2&6&7&3&0.
\end{array}
\]
Every group in this display comes from global sections of the
indicated direct image; none comes from $H^1$ of a line on the
coarse base. Every contributing summand has coarse degree zero,
so evaluation at the cusp is injective on its space of sections.
For the equal-dimensional source and target spaces used below
it is an isomorphism onto the displayed weight subspaces.
Thus fibrewise evaluation computes the actual ranks of the
absolute connecting maps.

\item For one-forms, the three middle source classes are $f_i\otimes g_i$.
In the ordered target $(g_0g_1,g_0g_2,g_1g_2)$, the matrix is
\[
C_1=\begin{pmatrix}-C&A&0\\-E&0&B\\0&-D_0&D_0\end{pmatrix},
\qquad
\det C_1=
\frac{476837158203125}{281474976710656}(384z+3125).
\]
The degree-zero connecting map is nonzero because $R(f_0)$
has nonzero $g_1,g_2$ coefficients. The degree-two connecting
map, from the classes $f_1g_0g_1,f_2g_0g_2$ to $g_0g_1g_2$,
has rank one because $D_0\ne0$.

\item We must exclude the sole possible zero of $\det C_1$ on the
whole admissible parameter domain. P1.15 identifies
$z=-(625\pi/128)c$ on the real slice in its algebraic splitting.
The exact exceptional parameter is therefore $c=5/(3\pi)<1$.
The certified continuation in
\path{code/part1/certified_lift_bound.py} proves
\[
\begin{aligned}
h_b(1)\in\bigl[&
-1.03427662241632179949424884537834605259496477790218821132\\
&\qquad\mathbin{\pm}8.11\cdot10^{-57}\bigr].
\end{aligned}
\]
so $h_b(1)+5/(3\pi)<0$. Since $h_b\to+\infty$ at the cusp,
the intermediate value theorem makes this parameter
nontransverse somewhere on the base. For the full complex
parameter the determinant has the same single zero, and the
same argument excludes it. The certification propagates complex
balls and adds an explicit Cauchy bound for every omitted
Taylor tail; no midpoint replacement discards an error radius.
Its exact connection and all step bounds are retained with the
program. Consequently the three connecting ranks are $1,3,1$.
The long exact sequence of \eqref{eq:P1Alogcotangent} gives
\[
h^q(\Omega_M^1(\log W))=(0,1,1,0,0).
\]

\item The residue sequence
\[
0\to\Omega_M^1\to\Omega_M^1(\log W)\to\nu_*\mathcal O_Y\to0
\]
is exact even at self-intersections: in a local branch chart it
records independently the coefficient of each $dx_i/x_i$ on
$x_i=0$. The normalization $Y$ is smooth complete toric, so its
only coherent structure cohomology is $H^0=\mathbb C$.
The one-form row follows immediately.

\item For two-forms take the second exterior power of
\eqref{eq:P1Alogcotangent}. Its short exact sequence has left
term $h^*(L^{11}\otimes F)$ and right term
$h^*(\bigwedge^2F)$. In degree two the three source classes
are $f_if_j\otimes g_ig_j$, $i<j$. The target is
$F\otimes g_0g_1g_2$. Its connecting matrix is
\[
C_2=\begin{pmatrix}-D_0&D_0&0\\E&0&-B\\0&-C&A\end{pmatrix},
\qquad \det C_2=-\det C_1.
\]
Thus it is an isomorphism, and
$h^3(\Omega_M^2(\log W))=1$.

\item For completeness the ordinary/logarithmic comparison in this
degree is the exact residue resolution
\[
0\to\Omega_M^2\to\Omega_M^2(\log W)
 \to\nu_*\Omega_Y^1(\log\partial Y)
 \to\nu^{[2]}_*\mathcal O_{W^{[2]}}\to0.
\]
The last space is the disjoint normalization of the seven
double surfaces. Local exactness is verified by writing a
logarithmic two-form as a sum of
$a_{ij}\,dx_i/x_i\wedge dx_j/x_j$ and terms with ordinary
differentials. The two residues of an $ij$ term have opposite
sign; their cancellation is the sole compatibility. Every local
double-surface coefficient lifts to one branch, which proves
surjectivity at the last term.
Both final sheaves have zero positive cohomology: the first
is the pushforward of the trivial rank-three toric logarithmic
bundle, and the second has complete toric surface components.
Writing $Q=\Omega_M^2(\log W)/\Omega_M^2$, we obtain
$H^q(Q)=0$ for $q\ge2$. Hence
$h^3(\Omega_M^2)=1$.
In degree zero the last residue map is zero: the fourteen toric
boundary divisors of $Y$ occur in seven opposite pairs; their
primitive normal vectors have opposite signs and their gluing
is a torus translation. This gives the zero $7$-by-$3$ residue
matrix on global logarithmic one-forms. Thus $H^0(Q)=\mathbb C^3$.
The two global logarithmic two-forms come from the
$L^{12},L^{13}$ summands of $L^{11}\otimes F$; their residues
are the independent $f_1,f_2$ directions at the cusp.
Therefore $h^0(\Omega_M^2)=0$.

\item Serre duality on this compact smooth complex fourfold gives
$h^{2,1}=h^{2,3}=1$. Finally the analytic Riemann--Roch theorem
in Hirzebruch, Theorem~25.3.1 and its explicit derivation
in \S25.4(a), pp.~188--190, applies to the Dolbeault elliptic
operator of $\Omega_M^p$. No Kähler hypothesis occurs there.
Using $c_1=c_3=0$, the Chern-root product gives
\[
\chi_y(M)=
\frac{(1+y)^4}{240}\,c_2^2+
\left(\frac{(-1+10y-y^2)^2}{144}
       -\frac{(1+y)^4}{120}\right)c_4
=-y+4y^2-y^3.
\]
The program verifies this polynomial exactly. It gives
$\chi(\Omega_M^2)=4$, hence $h^{2,2}=6$.
P1A.1 and Serre duality supply all other entries of the table.

\end{enumerate}
\end{proof}




\begin{theorem}\label{night:P1A.3}
\StatementRecord{night:P1A.3}{proved}{Historical events: GPT-6 Astra (OpenAI), ChatGPT Pro mode, effort Pro, S6\_06\_BLIND, 16--18 November 2026: targeted proof check. Also checked in this writing. Current writing events: GPT-6 Astra (OpenAI), ChatGPT Pro mode, effort Pro, S6\_10\_PAPERS, 18 November 2026: computation or formal certificate; targeted proof check. GPT-6 Astra (OpenAI), ChatGPT Pro mode, effort Pro, S6\_10\_PAPERS continuation, 19 November 2026: used and rechecked in this run. }{Definitions and the displayed constructions}
For every admissible fourfold,
\[
 (h^0(T_M),h^1(T_M),h^2(T_M),h^3(T_M),h^4(T_M))
 =(1,1,1,2,0).
\]
The complex parameter family has a nonzero Kodaira--Spencer map
at every member and supplies a smooth one-dimensional miniversal
Kuranishi germ.
\end{theorem}
\begin{proof}
\begin{enumerate}[label=\arabic*.]
\item The critical ideal of $h$ is $\mathcal I_D$. Locally
$q=x_0\cdots x_r$, its generators are the products omitting
one factor. A vector field in the kernel of $dq$ has coefficients
$a_i=x_i b_i$ in the crossing directions and
$\sum b_i=0$; the remaining ordinary directions are free.
Thus the kernel is the translation bundle $h^*E_0$, and
\[
0\to h^*E_0\to T_M\to J\to0,
\qquad J=\mathcal I_D\otimes h^*L^9.
\]
The conductor square gives
\[
0\to\mathcal O_W\to\nu_*\mathcal O_Y\oplus\mathcal O_D
 \to\nu_*\mathcal O_{\partial Y}\to0.
\]
The boundary is connected and anticanonical in the toric
threefold $Y$. Its coherent cohomology is $(1,0,1)$, by
$0\to K_Y\to\mathcal O_Y\to\mathcal O_{\partial Y}\to0$
and Serre duality. The singular locus $D$ is connected and
has dimension two, as is immediate from the connected
double-face incidence graph. The map on constants is
$(a,b)\mapsto a-b$ and is onto. In higher degree the conductor
sequence makes $H^1(\mathcal O_W)\to H^1(\mathcal O_D)$
an isomorphism. The map
$H^2(\mathcal O_{\partial Y})\to H^3(\mathcal O_W)$ is onto,
since $H^3(\mathcal O_D)=H^3(\mathcal O_Y)=0$; both sides
have dimension one, so it is an isomorphism. Consequently
$H^2(\mathcal O_W)\to H^2(\mathcal O_D)$ is also an
isomorphism and $h^q(\mathcal O_D)=(1,3,3)$.

\item Projection and \eqref{eq:P1Adegree} give
\[
h^q(h^*E_0)=(1,1,0,0,0),\qquad
h^q(h^*L^9)=(1,2,1,0,0).
\]
The maps from the latter groups to $H^q(\mathcal O_D)$
are evaluation at the cusp followed by the just-proved
restriction isomorphisms. Each contributing coarse summand
has degree zero, so these evaluations are injective, with
ranks $1,2,1$. The critical-ideal sequence gives
$h^q(J)=(0,0,1,2,0)$. The tangent dimensions and the
isomorphism $H^1(h^*E_0)\to H^1(T_M)$ follow.

\item Differentiating the actual periods gives, on the universal
torus covering, the Cech cocycle of constant translations
obtained by differentiating each lattice period. Under the
constant-cocycle description of fibre cohomology this is
the map $F\to V/F$ sending $f$ to $\alpha(f)e$.
It is nonzero and occupies precisely the trivial summand
of $E_0\otimes G$. The Leray edge map identifies this summand
with $H^1(h^*E_0)$, so the global Kodaira--Spencer class is
nonzero. The finite and toroidal charts depend holomorphically
on the parameter and supply the actual compact family,
so this cocycle is the derivative of that family on every chart.

\item Douady's analytic versality theorem, VII.8, applies to the
compact complex manifold. The Kuranishi germ embeds in its
one-dimensional tangent space. A classifying map from the
actual parameter disc has nonzero derivative, hence its
composite with that embedding is an invertible one-variable
analytic germ. Its pullback on $\mathbb C\{t\}$ is injective,
so the defining ideal of the Kuranishi germ is zero, including
nilpotents. This proves smoothness and miniversality.

\end{enumerate}
\end{proof}




\begin{theorem}\label{night:P1A.4}
\StatementRecord{night:P1A.4}{proved}{Historical events: GPT-6 Astra (OpenAI), ChatGPT Pro mode, effort Pro, S6\_06\_BLIND, 16--18 November 2026: targeted proof check. Also checked in this writing. Current writing events: GPT-6 Astra (OpenAI), ChatGPT Pro mode, effort Pro, S6\_10\_PAPERS, 18 November 2026: computation or formal certificate; targeted proof check. GPT-6 Astra (OpenAI), ChatGPT Pro mode, effort Pro, S6\_10\_PAPERS continuation, 19 November 2026: computation or formal certificate; targeted proof check; used and rechecked in this run. }{\cref{night:P1.1}}
The fourfold has algebraic dimension one and
\[
\Aut^0(M)=\mathbb C,\qquad
\Aut(M)=\mathbb C\times\mathbb Z/2.
\]
In a fixed holomorphic splitting, two members of the full
admissible complex parameter half-plane are biholomorphic
exactly when their parameters are equal.
\end{theorem}
\begin{proof}
\begin{enumerate}[label=\arabic*.]
\item An effective divisor dominating the base would restrict to
a nonempty effective divisor on a general complex torus.
Pairing its Chern class with the square of a Kähler form
on that torus is strictly positive. This contradicts
$H^2(M;\mathbb Z)=0$. Thus every divisor is vertical.
The pole divisor of a meromorphic function is vertical;
its restriction to a general compact fibre is holomorphic
and therefore constant. The graph and properness of $f$
descend it to a meromorphic function on $\mathbb P^1$.
Conversely pullbacks from the base occur. Hence the
meromorphic function field is $\mathbb C(t)$ and $a(M)=1$.
Every automorphism consequently induces a Möbius map on
the base. The two distinct multiplicities and the cusp
identify all three values separately, so this map is the
identity.

\item On smooth fibres an automorphism is affine. Its integral
linear part is locally constant and commutes with $M_1,M_2$.
Solving the two exact matrix equations in
\path{code/part1/analytic_automorphisms.py} gives the rational
centralizer $\mathbb Q I$; its integral invertible elements
are $\pm I$. At the order-five filling an affine lift with
linear part $\epsilon I$ would have to conjugate the same
finite affine generator. For its translation $b$ this
requires
\[
(I-M_1)b=(1-\epsilon)v_1/5\pmod\Lambda.
\]
Applying the invariant integral covector $\ell$ makes the
left side zero. For $\epsilon=-1$ the right side is
$2/5\pmod\mathbb Z$, which is nonzero. Thus every
automorphism is a translation.

\item Let $\mathcal A$ be the sheaf of translations extending to
the specified cusp model and finite root charts. There is
an exact analytic sheaf sequence
\[
0\to j_*\Lambda\to E_0\to\mathcal A\to0
\]
on the root orbifold. Here $j$ omits the cusp and the stalk
of $j_*\Lambda$ there is $\ker N$.
To check the last surjectivity at the cusp, lift a translation
to the toroidal cover. It preserves the fan precisely when
its valuation is an integral translation of the periodic
fan. Subtract the corresponding period; the remaining
characters are nonvanishing holomorphic functions on a
disc and possess logarithms. These give a local section
of $E_0$. The kernel consists exactly of the invariant
integral periods. This also proves that the sequence
accounts for all extendable translations, including
possible component permutations.

\item The cohomology of $j_*\Lambda$ in degree one is calculated
from three positive meridians. For the finite points use
the integer kernels of $I+M_i+\cdots+M_i^{m_i-1}$;
for the cusp use the primitive image of $N$.
Writing these column bases as $L_1,L_2,L_0$, the cocycle
and coboundary matrices are
\[
D_1=(L_1\ M_1L_2\ M_1M_2L_0),\qquad
D_0=\begin{pmatrix}L_1^{-1}(M_1-I)\\
L_2^{-1}(M_2-I)\\L_0^{-1}(M_0-I)\end{pmatrix}.
\]
Each inverse here denotes the unique coordinates in the
displayed integral column basis, not a square matrix
inverse. The exact $6$-by-$11$ and $11$-by-$6$ integer
matrices, their bases and Smith operations are exported
in \path{code/part1/analytic_automorphisms.json} and the
underlying parabolic program. They satisfy $D_1D_0=0$ and
give $\ker D_1/\im D_0=\mathbb Z/2$. The cocycle relations
are the finite-order relations and the product of the
three meridians; the cusp condition is precisely the
vanishing of its restriction modulo $(M_0-I)\Lambda$.
Thus this complex computes the required orbifold sheaf
cohomology, with its finite stabilizer torsion retained.

\item Now $H^0(j_*\Lambda)=0$ by P1.1,
$H^0(E_0)=\mathbb C$, and $H^1(E_0)=0$ by the floor
calculation. The exact sequence gives
\[
0\to\mathbb C\to H^0(\mathcal A)\to\mathbb Z/2\to0.
\]
The half-period
$\tau=(0,1,1,1,1,1)^t/2$ is fixed modulo $\Lambda$
by every monodromy matrix. It extends over the finite
quotients because the translations commute there.
At the cusp $N\tau$ is integral; primitiveness of $N$
provides $\lambda\in\Lambda$ with $N\lambda=N\tau$.
The representative $\tau-\lambda\in\ker N$ is a bounded
angular half-period and extends on every toric chart.
It is a nonzero order-two global section. Since the
additive group $\mathbb C$ has no nonzero torsion, it
splits the displayed exact sequence. Translations
commute, giving the asserted full group.

\item For an isomorphism between two members, the same intrinsic
base argument applies. Its integral linear part still
lies in the common centralizer. Either scalar $\pm I$
preserves a complex period plane exactly by equality
of the two planes. The graphs
$(\sigma+\xi\alpha e)(P)$ are equal precisely when their
parameters $\xi$ agree, since $e\notin F$ and $\alpha$
is nonzero. Translation does not change this condition.
This proves the parameter assertion.

\end{enumerate}
\end{proof}




\begin{theorem}\label{night:P1.17}
\StatementRecord{night:P1.17}{proved}{Historical events: GPT-6 Astra (OpenAI), ChatGPT Pro mode, effort Pro, S6\_06\_BLIND, 16--18 November 2026: targeted proof check. Also checked in this writing. Current writing events: GPT-6 Astra (OpenAI), ChatGPT Pro mode, effort Pro, S6\_10\_PAPERS, 18 November 2026: targeted proof check. GPT-6 Astra (OpenAI), ChatGPT Pro mode, effort Pro, S6\_10\_PAPERS continuation, 19 November 2026: used and rechecked in this run. }{\cref{night:P1.1}}
For the cusp marking $S$ fixed in P1.1 and the corresponding Kuhn fan,
none of the completed fourfolds in this family admits an
anti-holomorphic automorphism.  In particular none admits a real
structure.  More strongly, two members with this same cusp marking
are never anti-biholomorphic.
\end{theorem}
\begin{proof}
\begin{enumerate}[label=\arabic*.]
\item We first recover the base intrinsically.  An effective divisor on
$M_c$ cannot be horizontal: its first Chern class vanishes by
$H^2(M_c;\mathbb Z)=0$, whereas its nonzero effective restriction to a
general torus fibre would have positive pairing with the square of a
flat K\"ahler class.  The pole divisor of any meromorphic function is
therefore vertical.  On every regular fibre outside that divisor the
function is holomorphic and hence constant.  It factors through a
meromorphic function on the base, which is rational on
$\mathbb P^1$.  Thus the function field is exactly $\mathbb C(\lambda)$.
An anti-biholomorphism consequently induces an anti-M\"obius map of
the base.  The two distinct multiplicities $5,4$ and the singular cusp
support distinguish the three special values, so this base map fixes
$\infty,0,\lambda_*$.  As those points are real, the map is precisely
$\lambda\mapsto\bar\lambda$.

\item A holomorphic map between compact complex tori is affine: after
lifting to their vector spaces, its derivative is periodic and
therefore constant, since it descends to holomorphic functions on a
compact torus.  Applying this to an anti-holomorphic map after
conjugating its source gives the same assertion for its anti-linear
lift.  Its induced integral map on $H_1$ is constant along the
regular base.  Complex conjugation reverses the two finite meridians,
so this matrix must simultaneously conjugate $M_i$ to $M_i^{-1}$.
Solving these exact integral linear equations gives only $R$ and
$-R$ among unimodular matrices, where
\[
R=\begin{pmatrix}
1&0&0&0&0&0\\
2&0&0&0&-1&0\\
-1&0&-1&0&0&0\\
-1&0&0&-1&0&0\\
2&-1&0&0&0&0\\
-1&1&1&1&1&1
\end{pmatrix},\qquad R^2=I.
\]
Indeed the simultaneous intertwiner space over $\mathbb Q$ has
dimension one and this primitive matrix is unimodular; determinant
one in absolute value permits only the scalar factors $\pm1$.
The equations and this certificate are in
\texttt{centralizer\_torsion.py} and
\texttt{real\_structure\_check.py}.

\item The finite affine cap excludes $G=R$.  Normalize the pullback by
$t=u^{m_i}$; any base-preserving isomorphism lifts to this canonical
root cover and conjugates its deck generator to the inverse.
At its central torus the translation part $t_i$ would satisfy
\[
(I-M_i^{-1})t_i=-(v_i+Gv_i)/m_i\pmod\Lambda.
\]
Applying $\ell$, and using $\ell R=\ell$, gives
$0=-2\ell(v_i)/m_i$ in $\mathbb R/\mathbb Z$, impossible already
for $m_i=5$.  Hence $G=-R$ is necessary.

\item We now use the cusp marking, retaining the conjugation of based
meridians.  Conjugation sends
$g_0=(g_1g_2)^{-1}$ to $g_2g_1=g_2g_0^{-1}g_2^{-1}$.
Thus the matrix in a fixed local cusp frame is
$H=M_2^{-1}G$, and it anti-conjugates $M_0$.
Exact multiplication in the marking $S$ gives
\[
S^{-1}HS=
\begin{pmatrix}
1&0&0&-1&-1&0\\
0&1&0&-2&0&0\\
1&-1&-1&1&-1&-1\\
0&0&0&-1&0&0\\
0&0&0&0&-1&0\\
0&0&0&-1&1&1
\end{pmatrix}
=\begin{pmatrix}P&Q\\0&-P\end{pmatrix},\quad
P=\begin{pmatrix}1&0&0\\0&1&0\\1&-1&-1\end{pmatrix}.
\tag{1.17.1}
\]
Changing the based cusp path by a cusp loop multiplies by a power of
$M_0$ and does not change $P$.

\item We explain why an arbitrary translation cannot remove the obstruction
in (1.17.1).  In the additive torus coordinates whose first periods
are the $a_j$, the anti-linear map must have the form
$\zeta' =P\bar\zeta+t(q)$, because its values on these three real
periods already determine its linear part.  In multiplicative
coordinates it is
$z'=u(q)\bar z^{-P}$.
The source zero section limits to a smooth point of the central fibre.
An isomorphism takes it to another smooth point.  The central fibre
is irreducible, and its smooth locus is exactly the dense torus of
its toric normalization: every toric boundary divisor is a
normal-crossings self-identification.  In a chart about this image,
the translation therefore has the form $q^v$ times nonvanishing
holomorphic units, with $v\in\mathbb Z^3$.  Translating by the period
$v$ changes the selected component chart; the units act by torus
translation.  Neither changes the linear action $-P$ on the fan.

\item The primitive rays of the normal fan of that component are
\[
\{\pm v:0\ne v\in\{0,1\}^3\}.
\]
An extending torus monomial isomorphism must permute these rays.
This can be seen directly on a one-parameter arc $z=t^v$ meeting the
generic point of its boundary divisor: if the image vector lies in
the interior of a cone of dimension at least two, that divisor is
contracted, which is impossible for an isomorphism.  Complex
conjugation itself preserves the toric charts, so the same criterion
applies to the anti-monomial map above.
But $-Pe_2=(0,-1,1)^t$ is not one of the displayed rays.
This is a contradiction.  We have exhausted both possible integral
linear parts, so no anti-biholomorphism exists.
The assertion is explicitly about the chosen $S$ and its fan; it does
not identify different toroidal compactifications without a map
between them.

\end{enumerate}
\end{proof}


\section{The boundary of the family}

Fix one splitting and write the complex parameter as $\xi$, so that
$F_\xi=(\sigma+\xi\alpha e)(P)$. Put
$c_*=-\min h$ and $\mathcal H=\{\operatorname{Im}\xi>c_*\}$.
The real family in the prompt is $\xi=ic$.
By an analytic extension across a boundary point we mean a proper
Hausdorff complex-analytic family on a full complex disc whose
restriction to the nonempty open part lying in $\mathcal H$ is
the prescribed family. At infinity we use the coordinate $u=1/\xi$.
This definition permits singular central spaces, multiplicities,
and additional central components.

\begin{theorem}\label{night:P1A.5}
\StatementRecord{night:P1A.5}{proved}{Historical events: GPT-6 Astra (OpenAI), ChatGPT Pro mode, effort Pro, S6\_06\_BLIND, 16--18 November 2026: targeted proof check. Also checked in this writing. Current writing events: GPT-6 Astra (OpenAI), ChatGPT Pro mode, effort Pro, S6\_10\_PAPERS, 18 November 2026: targeted proof check. GPT-6 Astra (OpenAI), ChatGPT Pro mode, effort Pro, S6\_10\_PAPERS continuation, 19 November 2026: used and rechecked in this run. }{\ExternalPaperRef{theorem~}{17.4}{I}}
There is no such analytic extension across any finite boundary
point of $\mathcal H$, or across infinity. The same conclusion
holds after any finite ramified change of the local boundary
coordinate. In particular neither end of the real interval
$c_*<c<+\infty$ has a proper Hausdorff analytic extension in
this sense.
\end{theorem}
\begin{proof}
\begin{enumerate}[label=\arabic*.]
\item We show first that a purported extension would carry the marked
torus fibration on its whole punctured parameter disc, rather
than only on its original one-sided open set.
Take the reduced component containing that open family,
normalize it and resolve it by the analytic resolution theorems
used in P2.4. The comparison map is proper and is an isomorphism
on the given smooth open set. Proper images of the singular,
exceptional and critical loci in the parameter disc are analytic.
They omit a nonempty open set, hence are discrete. Shrinking
the disc retains only its centre. Thus its punctured restriction
is a proper holomorphic submersion with smooth compact
fourfold fibres.

\item On the resolved total space use the tenth inverse power of the
relative canonical line. Its proper direct image is coherent.
Remove torsion; on a smooth one-dimensional base the resulting
sheaf is locally free. Its generic rank is four, since on the
original open family $K^{-10}=f^*\mathcal O_{\mathbb P^1}(3)$
and the fibres of $f$ are connected. The analytic cohomology
and base-change complex of Grauert has constant matrix ranks
off a discrete set. Shrinking again makes every punctured
fibre comparison an isomorphism. No assertion of base change
at the central fibre is needed.

\item Evaluation of these four sections has no base points on the
original open set. The proper image of its base locus is
again discrete, so it has none over the remaining punctured
disc. It defines a proper map to $\mathbb P^3$ over that
disc. The image is a family of rational normal cubics there.
Here is the comparison in detail. Resolve the graph of the
meromorphic evaluation map over the full disc and take its
proper analytic image in $\mathbb P^3\times\Delta$.
This modification is an isomorphism wherever the evaluation
has no base points, hence on every punctured fibre now under
consideration. The image is projective over the disc. After
removing the torsion supported at the centre, its coherent
structure sheaf is flat away from a discrete subset. Its
Hilbert polynomial there is $3n+1$, since that is the
polynomial on the original nonempty open parameter set.
The singular and critical images are proper analytic subsets
of the parameter disc and omit that open set. Shrinking
removes them away from the centre. We therefore obtain a
smooth family of degree-three genus-zero curves on the
punctured disc; the nondegenerate embedding supplied by its
four independent sections is its complete cubic linear system.
This argument allows the basis of those sections to vary.

\item The recovered map to this smooth relative curve has connected
fibres. Indeed its Stein factorization has a finite factor
from a normal space to the normal relative curve. That finite
map has degree one on the original open parameter set, where
the given fibration has connected fibres. Its generic degree
is consequently one. It is finite and birational onto a normal
space, hence an isomorphism. This also proves connectedness
for the regular threefold fibres used below.

\item The critical image of the resulting map to this curve family
is proper analytic. Its components which dominate the
parameter disc are exactly the three continuations of the
original special values: any additional dominating component
would meet the original open parameter set. Their degrees
are one, because that is their degree on the original open
set. The multiplicities five and four distinguish the first
two, and the third is the cusp. They therefore give three
labelled sections on the punctured disc. Their cross-ratio
trivializes the genus-zero family as
$\mathbb P^1\times\Delta^*$, with the same coordinate as
the original fibration. Any vertical critical-image components
lie over a discrete set of parameters and have already been
removed. Hence the regular part is a proper holomorphic
submersion over
$(\mathbb P^1\setminus\{0,1,\infty\})\times\Delta^*$.

\item Every fibre of this last submersion is a complex three-torus.
Indeed any two points in its connected base are joined by
a path covered by finitely many coordinate discs. Along
each such disc the restricted family is a proper holomorphic
submersion. Catanese's Theorem~2.1, \emph{Deformation in the
large of some complex manifolds, I}, p.~264, says that a
deformation in the large of a complex torus of dimension
$n$ remains a complex torus of dimension $n$. The definition
on p.~261 is exactly a chain of these proper holomorphic
submersions; it imposes no Kähler assumption on the total
space. Starting on the original open set checks all its
hypotheses with $n=3$.

\item Pass to the universal cover of this regular base and fix the
integral $H^1$ marking by continuation from the given open
family. The holomorphic period planes of these tori give
a map to $\operatorname{Gr}(3,6)$. On the original open
set it is exactly $F_\xi$. Both this period map and the
explicit graph $F_\xi$ are holomorphic on the connected
cover, so the identity theorem makes them equal throughout.
The explicit graph remains a rank-three holomorphic plane
for every finite complex $\xi$, regardless of transversality;
at $\xi=\infty$ it has the Grassmannian limit
$\mathbb C e+\ker\alpha$. No choice of a continuation
marking changes the transversality of a real integral
period plane.

\item A full disc about a finite boundary point contains parameters
with $\operatorname{Im}\xi<c_*$. For any such parameter,
$-\operatorname{Im}\xi>\min h$ is attained by $h$ at a
regular base point: join a minimum point to a punctured
cusp neighbourhood, where $h$ is arbitrarily large, and
avoid the two isolated finite special points in this path.
The intermediate value theorem gives a regular point with
$h+\operatorname{Im}\xi=0$. At that point
$F_\xi\cap\overline{F_\xi}\ne0$, contradicting the period
decomposition of the just-proved complex torus fibre.
The same contradiction occurs in every punctured
neighbourhood of $u=0$: the function $\xi=1/u$ takes
values below the horizontal boundary line there. A finite
ramified reparameterization still maps an open neighbourhood
onto a neighbourhood of the boundary point, or onto a
punctured neighbourhood of infinity. It therefore retains
parameters on both sides and gives the same contradiction.
This proves both nonexistence assertions with all the
allowed central singularities retained.

\end{enumerate}
\end{proof}




\begin{theorem}\label{night:P1A.6}
\StatementRecord{night:P1A.6}{proved}{Historical events: GPT-6 Astra (OpenAI), ChatGPT Pro mode, effort Pro, S6\_06\_BLIND, 16--18 November 2026: targeted proof check. Also checked in this writing. Current writing events: GPT-6 Astra (OpenAI), ChatGPT Pro mode, effort Pro, S6\_10\_PAPERS, 18 November 2026: targeted proof check. GPT-6 Astra (OpenAI), ChatGPT Pro mode, effort Pro, S6\_10\_PAPERS continuation, 19 November 2026: independent re-proof; used and rechecked in this run. }{\cref{night:P1A.3}, \cref{night:P1A.4}}
A holomorphic punctured-disc family whose fibres are all
members of this fourfold family has an interior parameter
limit. After a base change of degree at most two, it is
isomorphic to the pullback of the constructed family by
a holomorphic map of the full disc to $\mathcal H$, and
therefore has a proper smooth filling. Such a family
cannot approach either parameter end.
\end{theorem}
\begin{proof}
\begin{enumerate}[label=\arabic*.]
\item At every punctured parameter, P1A.3 supplies local
miniversality and a local holomorphic classifying map to
$\mathcal H$. P1A.4 makes this map unique on each geometric
point. Thus these local maps agree on overlaps and define
one holomorphic function $\xi:\Delta^*\to\mathcal H$.
There is no unmentioned quotient by parameter translations:
P1A.4 proved their absence by the integral centralizer.

\item The Cayley coordinate
\[
 z=\frac{\xi-i(c_*+1)}{\xi-i(c_*-1)}
\]
has $|z|<1$. It extends holomorphically across the puncture
by the removable-singularity theorem. Its value there has
modulus strictly less than one, since a boundary value
would force $z$ to be constant by the maximum-modulus
principle. The inverse Cayley map therefore extends
$\xi$ with value in $\mathcal H$.

\item Local identifications with the pulled-back constructed
family differ by its automorphisms. These form the constant
holomorphic group $\mathbb C\times\mathbb Z/2$ over the
parameter half-plane: the translation field and the
half-period in P1A.4 extend over the entire family.
The finite component torsor on $\Delta^*$ is killed by
the double cover. The remaining additive holomorphic
torsor is zero because $H^1(\Delta^*,\mathcal O)=0$;
the punctured disc is Stein and its additive Cousin
problem has a solution. Thus after that cover the local
identifications glue to a global isomorphism. Pulling
back the already constructed proper smooth family along
the extended classifying map supplies the stated filling.

\end{enumerate}
\end{proof}


\section{The effective integral input classification}



\begin{theorem}[the orbifold quotient]\label{maps:src:F.2.7}
\StatementRecord{maps:src:F.2.7}{proved}{Historical events: Claude Fable 5.2 (Anthropic), effort high, 16 November 2026: carried as background without fault found. Claude Fable 5.2 (Anthropic), effort high, 17 November 2026: carried as background without fault found. GPT-6 Astra (OpenAI), ChatGPT Pro mode, effort Pro, S6\_09\_PAPERS, 17 November 2026: targeted proof check. Also checked in this writing. Current writing events: GPT-6 Astra (OpenAI), ChatGPT Pro mode, effort Pro, S6\_10\_PAPERS, 18 November 2026: targeted proof check. GPT-6 Astra (OpenAI), ChatGPT Pro mode, effort Pro, S6\_10\_PAPERS continuation, 19 November 2026: used and rechecked in this run. \emph{Recorded proof scope.} Claude Fable 5.2 (Anthropic), effort high, 2026-11-17; Claude Fable 5.2 (Anthropic), effort high, 16 November 2026; also checked in this writing by GPT-6 Astra (OpenAI), ChatGPT Pro mode, effort Pro, 18 November 2026}{the hypotheses, notation and definitions stated here; the explicit constructions and calculations in the proof}


For the specified torus fillings in F.2, with affine
multiplicities $m_1,\ldots,m_k$, the fundamental group of the
total space has a canonical surjection onto
\[
 \langle x_1,\ldots,x_k\mid
       x_i^{m_i}=1,\ x_1\cdots x_k=1\rangle.
 \tag{F.27}\label{maps:eq:F.27}
\]
Consequently a simply connected example with $k=2$ requires
$\gcd(m_1,m_2)=1$.  There is no simply connected example with
$k=3$ and $m_1,m_2,m_3\geq2$.  An affine multiplicity-one
point is an ordinary smooth point and reduces to the problem
with one fewer multiple fibre.

\end{theorem}


\begin{proof}
\begin{enumerate}[label=\arabic*.]
\item Over the complement of small discs at the special values,
the torus bundle projects surjectively on fundamental groups
to the punctured base.  Quotient its fundamental group by the
normal closure of all regular-fibre loops.  At a finite
filling its root cover is a smooth torus bundle over a disc,
and its free deck group is cyclic of order $m_i$.
The covering exact sequence shows that the quotient of the
filling group by its fibre subgroup is precisely this cyclic
group.  The boundary base meridian maps to its generator.
Thus adjoining that filling imposes exactly
$x_i^{m_i}=1$ after the fibre loops have been killed.
At the index-one Kuhn cusp the fundamental group is generated
by the three deck loops of its dual torus; these are regular
fibre loops.  Its base meridian contracts.  Hence its filling
imposes the ordinary meridian relation and leaves no extra
generator in this quotient.  The base relation gives the
product relation in \eqref{maps:eq:F.27}.  Van Kampen proves that \eqref{maps:eq:F.27} is
the resulting quotient, in particular the asserted surjection.

\item For two points, substitute $x_2=x_1^{-1}$; the quotient is
$\mathbb Z/\gcd(m_1,m_2)$.
For three points of orders at least two it is a nontrivial
triangle group.  One can see nontriviality without assuming
a classification of the total space.  Take a geodesic triangle
with angles $\pi/m_1,\pi/m_2,\pi/m_3$ in the sphere, Euclidean
plane or hyperbolic plane according as their sum is greater
than, equal to or less than $\pi$.
In the spherical case the possible orders are, after sorting,
$(2,2,n),(2,3,3),(2,3,4),(2,3,5)$, and the corresponding
spherical triangles exist.  The Euclidean and hyperbolic
angle conditions give the other triangles directly.
If $r_1,r_2,r_3$ are reflections in the opposite sides,
the rotations $r_2r_3,r_3r_1,r_1r_2$ have the stated
orders and product one.  They therefore give a nontrivial
representation of \eqref{maps:eq:F.27} into the orientation-preserving
isometry group.  Its quotient cannot be trivial.
Finally a free cyclic action of order one is the identity
construction on a root disc of order one, hence an ordinary
torus bundle filling and not a multiple fibre.

\end{enumerate}
\end{proof}




\begin{proposition}[exclusion of an involutory finite part]\label{maps:src:F.2.8}
\StatementRecord{maps:src:F.2.8}{proved}{Historical events: Claude Fable 5.2 (Anthropic), effort high, 16 November 2026: carried as background without fault found. GPT-6 Astra (OpenAI), ChatGPT Pro mode, effort Pro, S6\_09\_PAPERS, 17 November 2026: targeted proof check. Also checked in this writing. Current writing events: GPT-6 Astra (OpenAI), ChatGPT Pro mode, effort Pro, S6\_10\_PAPERS, 18 November 2026: targeted proof check. GPT-6 Astra (OpenAI), ChatGPT Pro mode, effort Pro, S6\_10\_PAPERS continuation, 19 November 2026: used and rechecked in this run. \emph{Recorded proof scope.} Claude Fable 5.2 (Anthropic), effort high, 16 November 2026; also checked in this writing by GPT-6 Astra (OpenAI), ChatGPT Pro mode, effort Pro, 18 November 2026}{the hypotheses, notation and definitions stated here; \cref{maps:src:F.2.7}; \cref{maps:src:F.2.15}; \cref{maps:src:F.3.1}; \cref{maps:eq:F.3}}


A simply connected datum with $b_2=0$, $k=2$ and the stated
finite affine fillings has neither finite linear
monodromy of order one or two.  Its unordered pair of finite
linear orders must therefore belong to
\[
 \{(3,4),(3,5),(3,8),(3,10),(4,5),
       (5,6),(5,8),(5,12)\}.
 \tag{F.28}\label{maps:eq:F.28}
\]
The affine multiplicities may still be larger than these
linear orders.

\end{proposition}


\begin{proof}
\begin{enumerate}[label=\arabic*.]
\item An identity finite generator would make the product of the
two finite matrices finite, contrary to its nontrivial
unipotent cusp.  Suppose the cohomology generators are
$A,B$ with $A^2=I$, where the other finite linear order is
odd.  It must be odd because the linear orders divide the
coprime affine multiplicities of Theorem \ref{maps:src:F.2.7}.
Write $U=AB$; the cusp condition gives $(U-I)^2=0$.
Multiplication of this identity on the left by $A$ and on
the right by $B^{-1}$ gives
\[
 BA+AB^{-1}=2I.
 \tag{F.29}\label{maps:eq:F.29}
\]
For $v\in\ker(B-I)$ we obtain $(B+I)Av=2v$.
As $B$ has odd order, $B+I$ is invertible, so $Av=v$.
Every $B$-fixed vector is therefore globally fixed.
The finite free affine action has a nonzero fixed
linear subspace, and \ref{maps:src:F.2.15} makes its real dimension even; it is
at least two.  Consequently the global cohomology invariant
rank is at least two.  The Leray argument of Theorem \ref{maps:src:F.3.1}
shows that $b_1=b_2=0$ requires that rank to be exactly one,
a contradiction.  Inverse or transpose monodromy conventions
give the same identity and fixed-space conclusion.

\item Finally take the complete finite-order list \eqref{maps:eq:F.3}, discard
one and two, and keep its coprime unordered pairs.
The eight pairs in \eqref{maps:eq:F.28} are exactly the result of that
finite enumeration.  This is only a necessary filter, not
a classification of relative positions or affine twists.

\end{enumerate}
\end{proof}




\begin{proposition}[the only first-cohomology pattern for two points]\label{maps:src:F.2.9}
\StatementRecord{maps:src:F.2.9}{proved}{Historical events: Claude Fable 5.2 (Anthropic), effort high, 16 November 2026: carried as background without fault found. GPT-6 Astra (OpenAI), ChatGPT Pro mode, effort Pro, S6\_09\_PAPERS, 17 November 2026: targeted proof check. Also checked in this writing. Current writing events: GPT-6 Astra (OpenAI), ChatGPT Pro mode, effort Pro, S6\_10\_PAPERS, 18 November 2026: targeted proof check. GPT-6 Astra (OpenAI), ChatGPT Pro mode, effort Pro, S6\_10\_PAPERS continuation, 19 November 2026: used and rechecked in this run. \emph{Recorded proof scope.} Claude Fable 5.2 (Anthropic), effort high, 16 November 2026; also checked in this writing by GPT-6 Astra (OpenAI), ChatGPT Pro mode, effort Pro, 18 November 2026}{the hypotheses, notation and definitions stated here; \cref{maps:src:F.3.1}; \cref{maps:src:F.2.15}}


For a simply connected two-multiple-fibre datum with $b_2=0$,
the rational
cohomology system $V$ has one invariant, no coinvariants,
and fixed rank two at each finite generator.  Its quotient
$W=V/\mathbb Qe$ by the invariant line has neither invariants
nor coinvariants.  If one finite characteristic polynomial
on $W$ is $(T-1)\Phi_d(T)$ with $\deg\Phi_d=4$, then $W$
is irreducible over $\mathbb Q$.

\end{proposition}


\begin{proof}
\begin{enumerate}[label=\arabic*.]
\item Write $a=\dim V^\Gamma$, $b=\dim(V^*)^\Gamma$ and $f_i$
for the finite fixed ranks.  Theorem \ref{maps:src:F.3.1} gives $a=1$,
vanishing first parabolic cohomology, and
$f_1+f_2=4+b$.  By \ref{maps:src:F.2.15} each $f_i$ is positive and even, so $b$ is
even.  The maximal square-zero cusp has
$\ker N=\operatorname{im}N$.  Therefore every globally
fixed vector belongs to $\operatorname{im}N$, whereas every
globally fixed covector kills $\operatorname{im}N$.
The global invariant and coinvariant spaces pair to zero.

\item Restriction of a globally invariant covector to the fixed
space of a finite generator is injective: averaging that
generator is the projection onto its fixed space, and the
covector is unchanged by averaging.  Its restriction kills
the $a$ global invariants in that space.  Hence
$a+b\leq f_i$ for each $i$.
If $b\geq2$ is even, the even ranks satisfy
$f_i\geq b+2$, so $4+b=f_1+f_2\geq2b+4$, a contradiction.
Thus $b=0$ and $f_1=f_2=2$.

\item The quotient $W$ has no invariant covectors, since its dual
embeds into $V^*$.  If it had a fixed line, its preimage in
$V$ would be an extension of the trivial line by the trivial
line.  Every finite generator on this two-dimensional
space would be unipotent of finite order, hence identity.
This would give two invariants in $V$, a contradiction.
Thus $W$ also has no invariant vectors.

\item If one finite characteristic polynomial is
$(T-1)\Phi_d(T)$ with the second factor irreducible of degree
four, every proper nonzero rational submodule has dimension
one or four.  In dimension one that generator acts as the
identity.  Since the product of the two finite generators
is unipotent, their induced one-dimensional characters are
inverse; hence the other generator also acts trivially.
This would be an invariant line.  In dimension four the
one-dimensional quotient is trivial by the same argument,
giving a coinvariant.  Both have just been excluded.

\end{enumerate}
\end{proof}


\paragraph{The primitive exterior-square bridge used in the proof.}
Here is an elementary construction of the classical five-dimensional
orthogonal representation, with the precise unipotent comparison.
Let $E=\mathbb C^4$ have a symplectic basis $e_1,e_2,f_1,f_2$.
The invariant bivector $\omega=e_1\wedge f_1+e_2\wedge f_2$ has
nonzero square. The wedge pairing on $\bigwedge^2E$ is symmetric
and nondegenerate. Its orthogonal complement
$E_0=\omega^\perp$ therefore has dimension five and is nondegenerate.
Preservation of $\omega^2$ gives determinant one on $E$;
$\det(\bigwedge^2g)=(\det g)^3$ and the fixed line
$\mathbb C\omega$ therefore give determinant one on $E_0$.
The action gives
\[
 \rho:\operatorname{Sp}(E)\longrightarrow\operatorname{SO}(E_0).
\]
Its differential is injective. Indeed, a symplectic infinitesimal
transformation acting trivially on $E_0$ also acts trivially on
$\mathbb C\omega$, hence on all of $\bigwedge^2E$. In a basis,
the equations $Xe_i\wedge e_j+e_i\wedge Xe_j=0$ for every $i<j$
first make every off-diagonal coefficient of $X$ zero, and then give
$X_{ii}+X_{jj}=0$ for every distinct pair. Since there are four
indices, these equations make every diagonal coefficient zero.
Both Lie algebras have dimension ten, so the image contains a
neighbourhood of the identity in $\operatorname{SO}(E_0)$.
That complex special orthogonal group is connected. For completeness,
orthogonal elimination expresses an orthogonal matrix as a product
of reflections in vectors of nonzero norm: to fix a nonisotropic
vector, reflection in the difference of its image and itself works
when that difference is nonisotropic; otherwise reflection in their
sum followed by reflection in the vector works. At least one of
sum and difference is nonisotropic. Continue on its nondegenerate
orthogonal complement. Determinant one makes the number of
reflections even. The nonisotropic vectors of a complex vector
space of dimension at least two form a path-connected set, since
two generic complex line segments can avoid the zero set of the
norm polynomial except at finitely many parameter values, around
which the paths can be detoured. Varying one vector in a pair of
reflections to the other connects that pair to the identity.
The one-dimensional final complement contributes only the already
accounted-for sign. This proves connectedness. Thus the open image
of $\rho$ is the whole special orthogonal group.

The kernel is $\{I,-I\}$. If an element acts trivially on
$E_0$, it also fixes $\omega$ and hence every bivector.
It preserves every two-plane, as follows by applying it to a
nonzero decomposable bivector of that plane. Intersections of
two-planes show that it preserves every line. A linear map
preserving every line is scalar: apply it to the coordinate lines
and their pairwise sums. Its scalar has square one. In particular
every finite orthogonal element of order $m$ has lifts whose
$m$th powers belong to $\{I,-I\}$; the lifts are finite and
semisimple.

The symplectic transvection fixing $e_1,e_2,f_2$ and sending
$f_1$ to $f_1+e_1$ acts on the following primitive basis:
\[
 e_1\wedge e_2,\quad e_1\wedge f_2,\quad
 e_2\wedge f_1,\quad f_1\wedge f_2,\quad
 e_1\wedge f_1-e_2\wedge f_2.
\]
Its difference from the identity sends the third vector to
the negative of the first and the fourth to the second, and
annihilates the other three. Its orthogonal Jordan type is
therefore $(2,2,1)$. Conversely, for an orthogonal unipotent
$I+N$ of that type, $N^2=0$ implies $N^*=-N$, because
$(I+N)^*=(I+N)^{-1}=I-N$. Its image is an isotropic plane.
After choosing a basis of that plane and a dual isotropic
plane, use the isometry on their residual nondegenerate line
to make the identification of two such decompositions have
determinant one. Changing its sign changes that determinant
and commutes with $N$, which kills the residual line.
In these decompositions, $N$ is determined by a nonzero alternating form on
the dual two-dimensional space. All such forms are equivalent
under its general linear group. Extending that change by the
inverse transpose on the first plane gives an orthogonal
change of basis; its determinant is one. Thus these unipotents
form one special-orthogonal conjugacy class. Surjectivity of
$\rho$ lifts the conjugacy. The two lifts of $I+N$ are a
symplectic transvection and its negative. Hence, after changing
one lift's sign, the product of the two finite lifts in the
argument is precisely a rank-one symplectic transvection.

Finally a symplectic diagonal element with eigenvalues
$x,x^{-1},y,y^{-1}$ acts on $E_0$ with eigenvalues
\[
 1,\quad xy,\quad x/y,\quad y/x,\quad (xy)^{-1}.
\]
This gives the exact eigenvalue comparison used above and
completes the required bridge without changing any assertion
about rational or integral conjugacy.

\begin{theorem}[exclusion of three further order pairs]\label{maps:src:F.2.10}
\StatementRecord{maps:src:F.2.10}{proved}{Historical events: Claude Fable 5.2 (Anthropic), effort high, 16 November 2026: carried as background without fault found. GPT-6 Astra (OpenAI), ChatGPT Pro mode, effort Pro, S6\_09\_PAPERS, 17 November 2026: targeted proof check. Also checked in this writing. Current writing events: GPT-6 Astra (OpenAI), ChatGPT Pro mode, effort Pro, S6\_10\_PAPERS, 18 November 2026: targeted proof check. GPT-6 Astra (OpenAI), ChatGPT Pro mode, effort Pro, S6\_10\_PAPERS continuation, 19 November 2026: used and rechecked in this run. \emph{Recorded proof scope.} Claude Fable 5.2 (Anthropic), effort high, 16 November 2026; also checked in this writing by GPT-6 Astra (OpenAI), ChatGPT Pro mode, effort Pro, 18 November 2026}{the hypotheses, notation and definitions stated here; \cref{maps:src:F.2.9}; \cref{maps:eq:F.31}}


No simply connected datum with $k=2$, $b_2=0$, and finite linear orders
$(3,5)$, $(3,8)$ or $(3,10)$ has the required maximal
square-zero cusp.  Among such data, the
remaining necessary unordered linear order pairs are
\[
 \{(3,4),(4,5),(5,6),(5,8),(5,12)\}.
 \tag{F.30}\label{maps:eq:F.30}
\]
When one order is five, the other finite characteristic
polynomial cannot be $(T-1)^2\Phi_4(T)^2$ or
$(T-1)^2\Phi_6(T)^2$.

\end{theorem}


\begin{proof}
\begin{enumerate}[label=\arabic*.]
\item Let the second order be $d\in\{5,8,10\}$.
Proposition \ref{maps:src:F.2.9} gives an irreducible rational rank-five
quotient $W$, with finite characteristic polynomials
\[
 (T-1)(T^2+T+1)^2,
 \qquad (T-1)\Phi_d(T).
\]
Since the invariant line of $V$ lies in $\operatorname{im}N_V$,
the quotient cusp has Jordan type $(2,2,1)$.
All three local conjugacy classes are self-dual.  Their
centralizer dimensions are respectively $9,5,13$.
For the last value, the conjugate partition is $(3,2)$,
so the dimension is $3^2+2^2$; the finite values are the
sums of squares of their eigenvalue multiplicities.
Consequently the invariant-sheaf Euler formula applied to
$\operatorname{Hom}(W,W^*)$ gives
\[
 \chi\bigl(\mathbb P^1,j_*\operatorname{Hom}(W,W^*)\bigr)
      =-25+9+5+13=2.
 \tag{F.31}\label{maps:eq:F.31}
\]
If $W$ were not self-dual, irreducibility would make both
$\operatorname{Hom}_\Gamma(W,W^*)$ and
$\operatorname{Hom}_\Gamma(W^*,W)$ zero.  These are the
degree-zero group and the dual of the degree-two group in
\eqref{maps:eq:F.31}.  The Euler characteristic would then be nonpositive,
a contradiction.  Thus $W$ has an invariant nondegenerate
rational bilinear form.  Its symmetric part is nonzero,
since an alternating form cannot be nondegenerate in odd
dimension; irreducibility then makes that symmetric part
nondegenerate.  Both generators have determinant one.

\item Over $\mathbb C$ use
$\operatorname{Spin}_5\simeq\operatorname{Sp}_4$ and its
explicit degree-two representation on the primitive exterior
square.  Choose finite semisimple lifts $R,S$ of the two
orthogonal generators, choosing one of their signs so that
$RS$ is the unipotent lift of their unipotent product.
The orthogonal Jordan type $(2,2,1)$ lifts to a rank-one
transvection.  This can be verified directly, as follows.
For a rank-two square-zero skew-adjoint operator on a
five-dimensional quadratic space, its image is a maximal
isotropic plane, and its nonzero skew form on the dual plane
determines one orbit under the complex orthogonal group.
A rank-one symplectic transvection in dimension four maps,
under the primitive exterior square, to precisely that orbit:
its logarithm has two length-two blocks and one trivial block.
The exact exterior matrices checking these ranks and square
zero are included in \texttt{fourfolds\_order\_pairs.py}.

\item The first lift has, up to a common sign, eigenvalues
$1,1,\zeta_3,\zeta_3^{-1}$.  Indeed a symplectic torus with
eigenvalues $x,x^{-1},y,y^{-1}$ has primitive exterior
eigenvalues $1,xy,x/y,y/x,(xy)^{-1}$; substitution of the
two equal primitive-third pairs gives the asserted list.
In particular $R$ has a repeated eigenvalue.
The lift $S$ has order dividing $2d$, so it has no eigenvalue
of order three or six.  Nor can it have eigenvalue $1$ or
$-1$: a symplectic pair $x=x^{-1}=\pm1$ would make two of
the nontrivial primitive exterior eigenvalues repeat, whereas
$\Phi_d$ has four distinct roots.
Thus $R^{-1}$ and $S$ have coprime characteristic polynomials.

\item We finish with an elementary consequence of the rank-one
companion argument.  Two matrices with coprime characteristic
polynomials whose difference has rank one are both cyclic.
To prove this, let $u$ span the difference image.
A common invariant subspace either meets the image nontrivially,
in which case the two induced quotient matrices agree, or is
killed by the difference, in which case their restrictions
agree.  Coprimality forces that subspace to be respectively
the whole space or zero.  The span of the iterates of $u$
under either matrix is stable under both and is nonzero;
it is therefore the whole space, proving cyclicity.
Apply this to $R^{-1},S$, since
$R^{-1}-S=R^{-1}(I-RS)$ has rank one.
But a finite semisimple matrix with a repeated eigenvalue
is not cyclic.  This contradicts the displayed spectrum of
$R$, and excludes all three order pairs.
The same proof gives the final assertion: replace the two
primitive-third pairs on $W$ by two primitive-fourth or
primitive-sixth pairs.  The local centralizer dimension
is still nine, so the Hom Euler characteristic is again two.
The spin lift has, up to sign, eigenvalues
$1,1,\zeta_j,\zeta_j^{-1}$ for $j=4$ or $6$.
The other lift has order dividing ten and has no eigenvalue
$1$ or $-1$ because its primitive exterior spectrum is
simple.  Their spectra are disjoint, and cyclicity gives
the same contradiction.

\end{enumerate}
\end{proof}




\begin{proposition}[rational uniqueness of the extension]\label{maps:src:F.2.11}
\StatementRecord{maps:src:F.2.11}{proved}{Historical events: Claude Fable 5.2 (Anthropic), effort high, 16 November 2026: carried as background without fault found. GPT-6 Astra (OpenAI), ChatGPT Pro mode, effort Pro, S6\_09\_PAPERS, 17 November 2026: targeted proof check. Also checked in this writing. Current writing events: GPT-6 Astra (OpenAI), ChatGPT Pro mode, effort Pro, S6\_10\_PAPERS, 18 November 2026: targeted proof check. GPT-6 Astra (OpenAI), ChatGPT Pro mode, effort Pro, S6\_10\_PAPERS continuation, 19 November 2026: used and rechecked in this run. \emph{Recorded proof scope.} Claude Fable 5.2 (Anthropic), effort high, 16 November 2026; also checked in this writing by GPT-6 Astra (OpenAI), ChatGPT Pro mode, effort Pro, 18 November 2026}{the hypotheses, notation and definitions stated here; \cref{maps:src:F.2.9}; \cref{maps:eq:F.32}}


Let $W$ be a rational rank-five representation of a
twice-punctured affine line, with two finite generators
having fixed rank one and product with unipotent Jordan type
$(2,2,1)$.  Suppose $W$ has neither invariants nor
coinvariants.  Up to rescaling the invariant line, there is
exactly one rational extension
$0\to\mathbb Qe\to V\to W\to0$ in which the two finite
generators retain their finite orders and the cusp has
square-zero logarithm of rank three.  Thus, for the
two-point candidates of \ref{maps:src:F.2.9}, the remaining rational
representation problem is exactly the problem of the
rank-five quotients.  This statement does not identify
integral lattices in those representations.

\end{proposition}


\begin{proof}
\begin{enumerate}[label=\arabic*.]
\item Write the two quotient generators as $T_1,T_2$ and the
extension generators in block form
$\left(\begin{smallmatrix}1&c_i\\0&T_i\end{smallmatrix}\right)$.
If $T_i$ has order $m_i$, finite order of the extension is
equivalent to $c_i(I+T_i+\cdots+T_i^{m_i-1})=0$.
Over $\mathbb Q$ this says $c_i\in\operatorname{row}(T_i-I)$,
a space of dimension four.  The two rows give eight
parameters.  Changes of splitting by a row $z\in W^*$
give $c_i\mapsto c_i+z(T_i-I)$; their space has dimension
five because $W^*$ has no invariants.  Hence extensions
split at the two finite cyclic groups form a vector space
of dimension three.

\item Restriction to the cusp takes this space to
$W^*/\operatorname{row}(T_0-I)$, also of dimension three.
Its kernel is the first parabolic cohomology of $W^*$.
That kernel vanishes: the invariant-sheaf Euler characteristic
is $-5+1+1+3=0$, and its degree-zero and degree-two groups
vanish by the assumed absence of invariants and coinvariants.
Thus cusp restriction is an isomorphism over $\mathbb Q$.

\item Put $N=T_0-I$ on $W$, so $N^2=0$ and $\operatorname{rank}N=2$.
The cusp extension has nilpotent part
$\left(\begin{smallmatrix}0&c\\0&N\end{smallmatrix}\right)$.
Its square is zero exactly when $cN=0$.
The corresponding subspace of cusp extension classes is
\[
 \frac{\operatorname{ann}(\operatorname{im}N)}
      {\operatorname{row}N},
 \tag{F.32}\label{maps:eq:F.32}
\]
which has dimension $3-2=1$.
A class in \eqref{maps:eq:F.32} is zero exactly when $c$ vanishes on
$\ker N$; those are precisely the extensions whose cusp
nilpotent part still has rank two.
Every nonzero class raises that rank to three and has square
zero.  Rescaling $e$ acts transitively on the nonzero rational
classes in this one-dimensional space.  This proves both
existence and uniqueness in the stated rational category.

\end{enumerate}
\end{proof}




\begin{theorem}[an integral obstruction for a rational $(5,12)$ family]\label{maps:src:F.2.12}
\StatementRecord{maps:src:F.2.12}{proved}{Historical events: Claude Fable 5.2 (Anthropic), effort high, 16 November 2026: carried as background without fault found. GPT-6 Astra (OpenAI), ChatGPT Pro mode, effort Pro, S6\_09\_PAPERS, 17 November 2026: targeted proof check. Also checked in this writing. Current writing events: GPT-6 Astra (OpenAI), ChatGPT Pro mode, effort Pro, S6\_10\_PAPERS, 18 November 2026: targeted proof check. GPT-6 Astra (OpenAI), ChatGPT Pro mode, effort Pro, S6\_10\_PAPERS continuation, 19 November 2026: used and rechecked in this run. \emph{Recorded proof scope.} Claude Fable 5.2 (Anthropic), effort high, 16 November 2026; also checked in this writing by GPT-6 Astra (OpenAI), ChatGPT Pro mode, effort Pro, 18 November 2026}{the hypotheses, notation and definitions stated here; \cref{maps:src:F.2.5}; \cref{maps:eq:F.35}}


Let a rational rank-six cohomology representation have generators
\[
 A=\begin{pmatrix}1&0\\0&A_5\end{pmatrix},\qquad
 B=\begin{pmatrix}1&c\\0&B_5\end{pmatrix},\qquad
 c=(0,0,0,1,0),
\]
where
\[
 A_5=\begin{pmatrix}
 1&0&1&0&0\\1&0&0&0&0\\-1&0&0&1&1\\
 1&-1&0&-1&0\\0&1&0&0&0
 \end{pmatrix},\qquad
 B_5=\begin{pmatrix}
 0&1&0&0&0\\0&0&0&0&1\\1&-2&0&0&0\\
 0&1&0&-1&-1\\0&0&1&1&2
 \end{pmatrix}.
 \tag{F.33}\label{maps:eq:F.33}
\]
For every invariant integral lattice with an index-one maximal
square-zero cusp, and every free affine filling of multiplicities
$m_1,m_2$ divisible respectively by five and twelve, a simply
connected total space has
\[
 \operatorname{Tor}_3 H_2
       \simeq\mathbb Z/3^{v_3(m_2)}.
 \tag{F.34}\label{maps:eq:F.34}
\]
In particular none of these data is three-connected.
The standard lattice nevertheless has a primitive cusp,
free affine fillings of orders five and twelve, and a simply
connected total space with rational Betti vector
$(1,0,0,0,4,0,0,0,1)$.

\end{theorem}


\begin{proof}
\begin{enumerate}[label=\arabic*.]
\item All matrix identities and integral lattices used below are
printed and checked exactly in
\texttt{fourfolds\_512\_probe.py}.
The generators have orders five and twelve, characteristic
polynomials $(T-1)^2\Phi_5$ and $(T-1)^2\Phi_{12}$,
and $(AB-I)^2=0$ with three nonzero Smith factors equal to one.
Their common invariant line is the first coordinate line.
For homology matrices $A^t,B^t$ put
\[
 u_1=(1,0,0,0,0,0)^t,\qquad
 u_2=(1,-1,1,-1,0,-1)^t,
 \qquad \ell=(1,0,0,0,0,0).
\]
They are invariant vectors with $\ell(u_i)=1$.
The cusp vanishing lattice and its two finite translates
generate the saturated lattice $\ker\ell$.
Choose twists $v_1=-2u_1$, $v_2=5u_2$.
Their $\ell$-values are units modulo five and twelve, proving
every nonidentity affine power free.
After killing the vanishing lattice the group relations are
$5a=-2h$, $12b=5h$, $a+b=0$; their determinant has absolute
value $|-24+25|=1$.  Thus the total space is simply connected.
The exact rational exterior calculations give
\[
 h^0=(1,1,0,0,0,0,1),\qquad
 h^1=(0,0,0,3,0,0,0),\qquad
 h^2=(1,0,0,0,0,1,1)
\]
for the invariant sheaves, and the same rational Leray
argument as in \ref{maps:src:F.2.5} gives the asserted Betti vector.

\item We first prove that changing the lattice cannot remove the
obstruction at the prime three.  Write $R=\mathbb Z_3$ and
$K=\mathbb Q_3$.  The standard quotient lattice $W_R$ defined
by $A_5,B_5$ has irreducible reduction modulo three.
Indeed $\Phi_5$ is irreducible over $\mathbb F_3$, since
three has order four modulo five.  Every proper nonzero
$A_5$-stable subspace therefore has dimension one or four.
On a one-dimensional subrepresentation or quotient its
$A_5$ character is one; the unipotent product forces the
$B_5$ character to be one too.  The stacked and horizontal
matrices for $A_5-I,B_5-I$ both have rank five modulo three,
as certified in the code, ruling out those two possibilities.

\item This irreducible reduction makes $W_R$ the unique stable
$R$-lattice up to homothety.  To see this, rescale another
lattice into $W_R$ but not into $3W_R$.  Its reduction is
a nonzero invariant subspace, hence all of $W_R/3W_R$.
Nakayama's lemma, or the elementary Smith form of the
inclusion, then makes it equal to $W_R$.

\item Normalize the quotient of any stable lattice in $V_K$ to
be $W_R$.  Its intersection with the invariant line is
$3^rRe$ for some integer $r$.
The representation of the free group on the two generators
has extension class
$\alpha\in H^1(\Gamma,W_R^*)$.
This cohomology group is torsion free: the coboundary map
from $W_R^*$ into two copies has full column rank modulo
three, since $W_R^*$ has no invariants modulo three.
The class $\alpha$ is primitive in this free $R$-module.
Otherwise the extension would split modulo three and $V$
would have a trivial quotient there.  The horizontal
matrix $[A-I\ B-I]$ instead has Smith factors
$(1,1,1,1,1,2)$, so it has full rank modulo three.

\item In a basis starting with $3^re$, the extension class is
$3^{-r}\alpha$.  Integrality is therefore possible exactly
when $r\leq0$.  Changes of splitting introduce no further
classes: the coboundary image is primitive, so an integral
coboundary over $K$ is already an integral coboundary over
$R$.  Thus these lattices are represented by multiplying
the extension rows by $3^{-r}$.
On the quotient cusp $N_W$ has primitive rank-two image.
The extra cusp map is the nonzero map
$\ker N_W/\operatorname{im}N_W\to Re$.
It is a unit for the displayed lattice, because its full
cusp has Smith factors $(1,1,1)$; for the lattice indexed
by $r$ it is multiplied by $3^{-r}$.
Hence the full cusp is primitive exactly when $r=0$.
We have proved uniqueness of the primitive-cusp lattice
over $R$, without relying on a bounded lattice search.

\item It remains to compute the obstruction on that lattice.
We first justify the arbitrary-twist comparison integrally.
For any global invariant lattice take the primitive invariant
homology covector $\ell$ and quotient the van Kampen group
by every fibre loop in $\ker\ell$.
The finite presentations become
$[a,h]=[b,h]=1$, $a^{m_1}=h^{r_1}$,
$b^{m_2}=h^{r_2}$, where $r_i=\ell(v_i)$.
The cusp meridian gives $ab=h^{-\kappa}$ for an integer
$\kappa$ determined by its translation clutching
($\kappa=0$ for the stated standard clutching).
The cusp angular lattice lies in $\ker\ell$, so it adds
no further relation to this quotient.  In particular this
is an abelian quotient with determinant
$m_1m_2\kappa+m_2r_1+m_1r_2$.
Simple connectivity makes that determinant $\pm1$.
Reduction modulo each $m_i$ proves both
$\gcd(m_1,m_2)=1$ and $\gcd(r_i,m_i)=1$.

\item The two finite fillings have free $R$-cohomology in the
degrees needed here.  In the invariant marking whose first
homology vector over $R$ has $\ell$-value one, the map
$m_i\ell$ from the reduced finite fibre to $\mathbb R/\mathbb Z$
has connected fibre $T^5$: the affine generator permutes
its $m_i$ preimage components by addition of $r_i$ modulo
$m_i$, transitively by the proved coprimality.
Its monodromy on vertical homology is $D_i^{s_i}$, where
$s_ir_i\equiv1\pmod {m_i}$, possibly followed by a translation.
Thus it is a mapping torus of a power prime to the linear
order.  The translation is isotopic to the identity and
does not change its Wang maps.  At the second point, over
$R$ we may take the invariant first vector to be $v_2/r_2$;
its denominator is prime to three and does not change the
$R$-homology lattice.  At the first point the covering
degree $m_1$ is prime to three, so averaging and the transfer
identify its $R$-cohomology with the invariant fibre
cohomology directly.
In degrees one and two the matrices
$D_i-I$ and $\bigwedge^2D_i-I$ have only unit Smith factors.
This proves the stated freeness by the Wang sequence.

\item The finite degree-two covering cokernel is read from
integral periods.  Vertical cycles have integral periods.
The remaining primitive two-cycle is the commuting torus
formed by a base suspension and a primitive fixed vertical
loop $\lambda_i$; its covering period is
$\omega(u_i,\lambda_i)/m_i$ for an invariant integral
two-form $\omega$.
More explicitly the homological Wang sequence is
\[
 0\longrightarrow\operatorname{coker}(D_i^{s_i}-I\mid H_2(T^5;R))
 \longrightarrow H_2(B_i;R)
 \longrightarrow\ker(D_i^{s_i}-I\mid H_1(T^5;R))
 \longrightarrow0.
\]
The last group has rank one, generated by $\lambda_i$.
The commuting suspension torus maps to that generator;
at the second point the invariant vector $v_2/r_2$ may be used after taking
$R$-linear combinations of integral cycles.
Therefore vertical cycles and this one torus generate
$H_2(B_i;R)$, proving sufficiency of the period test.
For integral bases of the invariant two-forms, the numerator
rows at the two points are
\[
 (0,0,5),\qquad (3,1,1).
 \tag{F.35}\label{maps:eq:F.35}
\]
The vertical Wang quotient has no torsion, so those cycles
generate the full free second homology and the period test
is sufficient as well as necessary.

\item Since $3\mid m_2$, the already-proved coprimality gives
$3\nmid m_1$ and makes $r_2$ a unit in $R$.
Changing that unit twist or adding a fixed vertical vector
does not change the three-primary cokernel in \eqref{maps:eq:F.35}: the
former multiplies its generating period by a unit, and the
latter has zero wedge with the fixed vertical loop.
The first finite point has zero covering cokernel over $R$;
the second has $R/m_2R$.

\item The global invariant degree-two group is zero, and the
degree-two global coboundary image is primitive, with
vanishing rational parabolic first cohomology.  Therefore
$H^1(\mathbb P^1,R^2g_*R)=R/m_2R$.
The global degree-one coinvariants vanish over $R$, and
there are no degree-two or degree-three point-supported
kernels: the finite cap groups just proved free, and the
Kuhn cusp has its sole kernel in degree four.
The degree-three global invariant group also vanishes.
Consequently the only nonzero total-degree-three Leray
term over $R$ is $E_2^{1,2}=R/m_2R$; it has no possible
differential.  Thus $H^3(g^{-1}(\mathbb P^1);R)=R/m_2R$.
The universal coefficient theorem gives precisely \eqref{maps:eq:F.34}.

\end{enumerate}
\end{proof}




\begin{theorem}[two further rational families excluded integrally]\label{maps:src:F.2.13}
\StatementRecord{maps:src:F.2.13}{proved}{Historical events: Claude Fable 5.2 (Anthropic), effort high, 16 November 2026: carried as background without fault found. GPT-6 Astra (OpenAI), ChatGPT Pro mode, effort Pro, S6\_09\_PAPERS, 17 November 2026: targeted proof check. Also checked in this writing. Current writing events: GPT-6 Astra (OpenAI), ChatGPT Pro mode, effort Pro, S6\_10\_PAPERS, 18 November 2026: targeted proof check. GPT-6 Astra (OpenAI), ChatGPT Pro mode, effort Pro, S6\_10\_PAPERS continuation, 19 November 2026: used and rechecked in this run. \emph{Recorded proof scope.} Claude Fable 5.2 (Anthropic), effort high, 16 November 2026; also checked in this writing by GPT-6 Astra (OpenAI), ChatGPT Pro mode, effort Pro, 18 November 2026}{the hypotheses, notation and definitions stated here; \cref{maps:src:F.2.11}; \cref{maps:src:F.2.5}; \cref{maps:src:F.2.12}}


Let $C_f,C_g$ denote the degree-four companion matrices.
Take $f=\Phi_5$ and either
\[
 \text{case A: }g=\Phi_{12},\qquad
 \text{case B: }g=\Phi_3\Phi_6.
 \tag{F.36}\label{maps:eq:F.36}
\]
Let $W$ be the quotient of
$\bigwedge^2(C_f^{-1},C_g)$ by its invariant line, and
let $V$ be its unique rational extension of \ref{maps:src:F.2.11}.
For every invariant integral lattice in $V$ with primitive
maximal square-zero cusp, every simply connected free affine
filling, with multiplicities $m_1,m_2$ divisible by five and
six respectively, fails three-connectivity.
More precisely, put $r=v_3(m_2)\geq1$.
In case A, $\operatorname{Tor}_3 H_3$ has order
$3^{2r-1}$.  In case B,
$\operatorname{Tor}_3 H_2\simeq\mathbb Z/3^r$.
For $m_2=6$ the respective obstructing group is $\mathbb Z/3$.

\end{theorem}


\begin{proof}
\begin{enumerate}[label=\arabic*.]
\item The complete integer representatives are supplied by the
companion and Smith-basis construction in
\texttt{fourfolds\_56\_probe.py} and
\texttt{fourfolds\_56b\_probe.py}.
In each case the source cusp is a rank-one transvection;
the primitive exterior quotient has cusp type $(2,2,1)$,
finite fixed ranks one, and no invariants or coinvariants.
The second finite characteristic polynomial on $V$ is
$(T-1)^2(T+1)^2\Phi_3$ in case A and
$(T-1)^2(T+1)^2\Phi_6$ in case B.
The exact extension equations give an integral primitive
rank-three cusp in both cases.

\item For completeness these are not merely rational candidates
with impossible twist lattices.  In the displayed integral
bases of the probe, change lattice by the rational basis
matrix $Q$ whose diagonal entries are two and whose only
off-diagonal entries are $Q_{1,4}=Q_{1,6}=1$.
The resulting generators $Q^{-1}AQ,Q^{-1}BQ$ remain
integral and have primitive cusps.
The invariant homology vectors of $\ell$-value one can be
chosen as
\[
\begin{array}{c|cc}
 &u_1&u_2\\\hline
 \mathrm A&(1,3,-3,0,1,0)&(1,-1,1,2,-1,1)\\
 \mathrm B&(1,0,0,1,0,1)&(1,0,0,1,0,0).
\end{array}
\]
Use twists $u_1,-u_2$ and multiplicities five and six.
All nonidentity affine powers are free, and the cusp normal
closure is the saturated $\ker\ell$.
The residual determinant $6-5=1$ proves simple connectivity.
The rational Leray ranks are the same as in \ref{maps:src:F.2.5}.

\item The complete three-adic lattice argument in \ref{maps:src:F.2.12} applies
to each case, not only to these representatives.
The reason is checked by the same exact integer matrices:
the rank-five quotient has no invariants or coinvariants
modulo three, its order-five polynomial is
$(T-1)\Phi_5$ with irreducible $\Phi_5$ modulo three, the
global extension class is primitive, and the rank-six
cusp has unit Smith factors.
Thus each representation has exactly one primitive-cusp
$\mathbb Z_3$-lattice up to homothety.
Both the quotient argument forcing $r_i$ prime to $m_i$
and the finite mapping-torus comparison of \ref{maps:src:F.2.12} apply.

\item Set $R=\mathbb Z_3$.  All finite cap cohomology groups
through degree four are free over $R$: the nonunit Smith
factors of their vertical exterior differences, when
present, are two.  The global coboundary and coinvariant
Smith factors in degrees one through four likewise have
no factor three.  The rational invariant and parabolic
ranks are those printed in \ref{maps:src:F.2.12}.  In particular,
the global invariant groups in degrees two, three and
four vanish, their relevant coinvariants vanish over $R$,
and the degree-two parabolic first cohomology is zero.

\item For the transformed integral lattices just specified, the
finite degree-two suspension periods at order six are
\[
 \mathrm A:\ (-2,-1,-1),\qquad
 \mathrm B:\ (1/3,0,1/3).
 \tag{F.37}\label{maps:eq:F.37}
\]
The order-five periods are integral and its covering degree
is a unit in $R$ for any simply connected multiplicities.
In case B the numerator row at six is $(2,0,2)$; replacing
six by $m_2$ gives the cyclic cokernel $R/3^rR$.
The total-degree-three Leray argument of \ref{maps:src:F.2.12} gives
$H^3(P;R)=R/3^rR$, proving the asserted $H_2$ obstruction.

\item In case A the obstruction instead occurs in degree four.
The invariant rank-four middle cohomology lattice of the
sixfold torus cover has skew-pairing determinant $6^2$.
This determinant is computed exactly, with the integral
invariant basis, in \texttt{fourfolds\_56a\_newlattice.py}.
The free middle pairing downstairs is unimodular over $R$.
If $I$ is the index of the pullback lattice in the invariant
lattice, pullback degree $m_2$ gives
\[
 I^2\,6^2=m_2^4
 \quad\text{up to a unit of }R.
 \tag{F.38}\label{maps:eq:F.38}
\]
Thus the degree-three finite covering cokernel has
three-primary order $3^{2r-1}$.
There are no global invariant degree-three forms to
quotient this cokernel out.  The integral parabolic first
cohomology in degree three is free over $R$, so the torsion
of $H^1(B,R^3g_*R)$ is precisely that finite cokernel.
There are no differentials to or from this middle base
column.  In total degree four the other surviving term
is the free rank-one point-supported cusp class;
the degree-two base term is zero over $R$.
The resulting extension by a free module splits.
Therefore $\operatorname{Tor}_3H^4(P;R)$ has order
$3^{2r-1}$, and the universal coefficient theorem identifies
it with $\operatorname{Tor}_3H_3(P;\mathbb Z)$.
This proves the claimed obstruction in case A.

\end{enumerate}
\end{proof}




\begin{theorem}[four order pairs and eight local spectral pairs]\label{maps:src:F.2.14}
\StatementRecord{maps:src:F.2.14}{proved}{Historical events: Claude Fable 5.2 (Anthropic), effort high, 16 November 2026: carried as background without fault found. GPT-6 Astra (OpenAI), ChatGPT Pro mode, effort Pro, S6\_09\_PAPERS, 17 November 2026: targeted proof check. Also checked in this writing. Current writing events: GPT-6 Astra (OpenAI), ChatGPT Pro mode, effort Pro, S6\_10\_PAPERS, 18 November 2026: targeted proof check. GPT-6 Astra (OpenAI), ChatGPT Pro mode, effort Pro, S6\_10\_PAPERS continuation, 19 November 2026: used and rechecked in this run. \emph{Recorded proof scope.} Claude Fable 5.2 (Anthropic), effort high, 16 November 2026; also checked in this writing by GPT-6 Astra (OpenAI), ChatGPT Pro mode, effort Pro, 18 November 2026}{the hypotheses, notation and definitions stated here; \cref{maps:src:F.2.15}; \cref{maps:src:F.2.9}; \cref{maps:src:F.2.7}; \cref{maps:eq:F.40}; \cref{maps:src:F.2.2}}


Every simply connected two-finite-point datum with $b_2=0$ has
one finite linear order equal to five and the other equal
to four, six, eight or twelve.
After suppressing the common factor $(T-1)^2$, its unordered
pair of finite characteristic polynomials is exactly one
of the following necessary possibilities:
\[
\begin{array}{c|c|l}
 \text{orders}&\text{order-five factor}&\text{other factor}\\\hline
 (5,4)&\Phi_5&\Phi_2^2\Phi_4\\
 (5,6)&\Phi_5&\Phi_2^2\Phi_3,\ \Phi_2^2\Phi_6,\ \Phi_3\Phi_6\\
 (5,8)&\Phi_5&\Phi_8\\
 (5,12)&\Phi_5&\Phi_{12},\ \Phi_3\Phi_4,\ \Phi_4\Phi_6.
\end{array}
 \tag{F.39}\label{maps:eq:F.39}
\]
Its rank-five quotient is rationally irreducible.
In particular all order-three pairs, including $(3,4)$,
are impossible.  The table is a necessary spectral filter;
it does not identify the relative rational or integral
positions of the matrices having those polynomials.

\end{theorem}


\begin{proof}
\begin{enumerate}[label=\arabic*.]
\item By \ref{maps:src:F.2.15} the finite actions are locally complex-compatible.
By \ref{maps:src:F.2.9}, the rank-six cohomology system $V$ has one
invariant, no coinvariants, and finite fixed ranks two.
Its quotient $W$ has neither invariants nor coinvariants.
Every rational one-dimensional subrepresentation or
quotient character of $W$ is trivial: its two finite
orders divide the coprime affine multiplicities of \ref{maps:src:F.2.7},
so at least one is odd and that character value is one;
the unipotent product forces the other value to be one.
Thus $W$ and $W^*$ have no invariant rational line of any
character.

\item The space of invariant alternating forms on $W$ has
dimension at most one.  Indeed a nonzero such form cannot
have rank four, because its radical would be a stable
rational line in $W$.  It therefore has rank two and is
a decomposable element of $\bigwedge^2W^*$.
If two independent invariant rank-two forms existed,
their sum would also have rank at most two; hence their
wedge product would be zero.  Their two supporting
planes in $W^*$ would be distinct stable planes meeting
in a line.  That stable rational line contradicts the
preceding paragraph.  This proves the bound without an
irreducibility assumption.

\item The dual exterior sequence
\[
 0\longrightarrow\bigwedge^2W^*
 \longrightarrow\bigwedge^2V^*
 \longrightarrow W^*\longrightarrow0
\]
and $(W^*)^\Gamma=0$ show that the coinvariant dimension
$b_2'$ of $\bigwedge^2V$ is at most one.
On the other hand its invariant dimension $a_2'$ is zero.
To see this, the corresponding $H^0(R^2g_*\mathbb Q)$
term has no point-supported rational kernel at a finite
fibre or in cusp degree two, and its only possible outgoing
differential has target $H^2(R^1g_*\mathbb Q)=0$.
It therefore survives in total degree two, which vanishes.

\item Let $d_i$ be the dimension fixed by the $i$th finite
generator on $\bigwedge^2V$.  The cusp fixed dimension is
nine, so the invariant-sheaf Euler identity gives
\[
 -15+d_1+d_2+9
       =a_2'-h^1(j_*\bigwedge^2V)+b_2'\leq1.
 \tag{F.40}\label{maps:eq:F.40}
\]
Each finite generator has fixed rank two on $V$.
Its complementary four-dimensional real representation
commutes with a complex structure and has no fixed vectors.
It either has two distinct nonreal inverse pairs, one
nonreal pair with multiplicity two, two negative real
eigenvalues and a nonreal pair, or four negative eigenvalues.
The respective invariant dimensions in degree two are
$3,5,3,7$.  In particular $d_i$ is odd and at least three.
The left side of \eqref{maps:eq:F.40} is therefore nonnegative and even,
so it must be zero.  Thus $d_1=d_2=3$.

\item The complete cyclotomic filter \ref{maps:src:F.2.2} now leaves precisely
the finite polynomials with two fixed eigenvalues and
the indicated degree-two fixed rank three.
Keeping coprime unordered pairs gives the eight rows
expanded in \eqref{maps:eq:F.39}.  This finite enumeration is reproduced
from the complete saved 38-type list by
\texttt{fourfolds\_order\_pairs.py}.
An order-three generator would instead have characteristic
polynomial $(T-1)^2\Phi_3^2$, hence $d_i=5$, explaining
its exclusion directly.
Every remaining pair contains $(T-1)^2\Phi_5$, so \ref{maps:src:F.2.9}
proves rational irreducibility of $W$.

\end{enumerate}
\end{proof}




\begin{proposition}[automatic local complex compatibility]\label{maps:src:F.2.15}
\StatementRecord{maps:src:F.2.15}{proved}{Historical events: Claude Fable 5.2 (Anthropic), effort high, 16 November 2026: carried as background without fault found. GPT-6 Astra (OpenAI), ChatGPT Pro mode, effort Pro, S6\_09\_PAPERS, 17 November 2026: targeted proof check. Also checked in this writing. Current writing events: GPT-6 Astra (OpenAI), ChatGPT Pro mode, effort Pro, S6\_10\_PAPERS, 18 November 2026: targeted proof check. GPT-6 Astra (OpenAI), ChatGPT Pro mode, effort Pro, S6\_10\_PAPERS continuation, 19 November 2026: used and rechecked in this run. \emph{Recorded proof scope.} Claude Fable 5.2 (Anthropic), effort high, 16 November 2026; also checked in this writing by GPT-6 Astra (OpenAI), ChatGPT Pro mode, effort Pro, 18 November 2026}{the hypotheses, notation and definitions stated here; \cref{maps:src:F.2.7}; \cref{maps:src:F.1.4}}


Every simply connected datum of F.2 with two finite affine
fillings automatically has finite linear actions commuting
with real complex structures.  This assertion concerns the
individual finite actions and does not produce simultaneous
global periods.

\end{proposition}


\begin{proof}
\begin{enumerate}[label=\arabic*.]
\item Let $m_i$ be the order of the free cyclic affine action
defining the $i$th filling.  Its linear order divides $m_i$.
By \ref{maps:src:F.2.7} the integers $m_1,m_2$ are coprime, so at least one
of them is odd.  The corresponding linear determinant is
one, because a finite real matrix of odd order cannot have
determinant minus one.  The relation
$M_1M_2M_0=I$ and the unipotence of $M_0$ imply
$\det M_1\det M_2=1$, so both determinants are one.

\item A finite-order real matrix is semisimple.  Its nonreal
eigenvalues occur in conjugate pairs, and its only possible
real eigenvalues are $1$ and $-1$.  Determinant one says
that the $-1$ eigenspace has even real dimension.  Since the
total dimension is six, the $1$ eigenspace also has even
real dimension.  Choose a complex structure on each of
these real eigenspaces by pairing a real basis.  On each
nonreal isotypical real summand, multiplication by $i$ in
its complex eigenspace supplies a commuting complex
structure.  Their direct sum commutes with $M_i$.
The fixed eigenspace is nonzero because the affine action
is free, as in \ref{maps:src:F.1.4}, so its even dimension is at least two.

\item Each chosen structure makes $\mathbb R^6/\mathbb Z^6$ a
complex torus and the given affine map holomorphic; thus it
supplies precisely the local complex-compatibility
hypothesis used in the subsequent reductions.  For three genuine
finite fillings \ref{maps:src:F.2.7} has already excluded simple
connectivity without this argument.

\end{enumerate}
\end{proof}


Theorem \ref{maps:src:F.2.10} supplies an independent spin-representation
proof for three of the excluded order-three pairs and for
the repeated fourth- and sixth-root cases.
Within the third row of the order-six spectral possibilities,
$\Phi_3\Phi_6$, no exclusion follows from F.2.13: that theorem
treats particular rational representations with the first
two order-six spectra.  Likewise \ref{maps:src:F.2.12} excludes one
rational representation of order pair $(5,12)$, not its
entire spectral locus.

\begin{theorem}[an unbounded integral family]\label{maps:src:F.2.16}
\StatementRecord{maps:src:F.2.16}{proved}{Historical events: Claude Fable 5.2 (Anthropic), effort high, 16 November 2026: carried as background without fault found. GPT-6 Astra (OpenAI), ChatGPT Pro mode, effort Pro, S6\_09\_PAPERS, 17 November 2026: targeted proof check. Also checked in this writing. Current writing events: GPT-6 Astra (OpenAI), ChatGPT Pro mode, effort Pro, S6\_10\_PAPERS, 18 November 2026: computation or formal certificate; targeted proof check. GPT-6 Astra (OpenAI), ChatGPT Pro mode, effort Pro, S6\_10\_PAPERS continuation, 19 November 2026: used and rechecked in this run. \emph{Recorded proof scope.} Claude Fable 5.2 (Anthropic), effort high, 16 November 2026; also checked in this writing by GPT-6 Astra (OpenAI), ChatGPT Pro mode, effort Pro, 18 November 2026}{the hypotheses, notation and definitions stated here; \cref{maps:eq:F.U1}; \cref{maps:eq:F.U6}; \cref{maps:src:F.2.5}; \cref{maps:eq:F.U8}}


There are infinitely many simultaneous rational conjugacy classes
of integral rank-six tuples of the linear type $(5,4)$ in F.2,
with primitive maximal square-zero cusp.  An explicit infinite
subfamily has free affine lifts of orders five and four.
Apart from the parameter zero, the rank-five irreducible quotient
in the family is not self-dual.  Thus no finite bound on rational
conjugacy classes is valid for the unrestricted input tuples.
This does not assert infinitely many three-connected survivors.

\end{theorem}


\begin{proof}
\begin{enumerate}[label=\arabic*.]
\item All the following expressions are explicit.  Let $v\in\mathbb Z$,
$q=v^2-v+1$, and let $A$ be the permutation matrix
$Ae_j=e_{j+1}$ on $\mathbb Z^5$, with subscripts taken modulo five.
Put
\[
 X=(e_0,\ e_2+ve_4),\qquad
 Y=\begin{pmatrix}
 0&-q&-v^2&-1&v\\
 0&1-vq&-vq&v^2-v&q
 \end{pmatrix},\qquad N=XY,
 \tag{F.U1}\label{maps:eq:F.U1}
\]
and set $C=(I+N)A$, $B=C^{-1}=A^{-1}(I-N)$.
Direct multiplication gives
\[
 YX=0,\quad Y(A+I)^{-1}X=I_2,\quad
 A^5=C^4=I_5,\quad
 \chi_C(x)=(x-1)(x+1)^2(x^2+1).
 \tag{F.U2}\label{maps:eq:F.U2}
\]
The image of $X$ is primitive, and $Y$ is onto $\mathbb Z^2$.
For the latter assertion, $2(A+I)^{-1}X$ is integral, so the
cokernel of $Y$ is killed by two; reduction of \eqref{maps:eq:F.U1} in the
two cases $v\equiv0,1\pmod2$ has rank two.  Thus $N$ has
rank two, $N^2=0$, and its two nonzero Smith factors are one.

\item The common invariant space of $A,B$ is zero: the only
$A$-fixed line is generated by $\mathbf1=(1,1,1,1,1)^t$,
and the first entry of $Y\mathbf1$ is $-2q\ne0$.
There are no common invariant covectors either, since
$\mathbf1^tX=(1,v+1)$ and $Y$ has rank two.
The characteristic polynomial of $A$ is $(x-1)\Phi_5(x)$.
A proper rational invariant subspace must therefore have
dimension one or four.  A one-dimensional subrepresentation
or quotient has $A$ acting as one; the unipotent product then
forces $B$ also to act as one.  The two vanishing assertions
exclude both possibilities.  The rank-five pair is irreducible.

\item Here is an explicit primitive rank-six extension, including
the integral lattice.  Let $P$ have columns
$\mathbf1/q,e_1,e_2,e_3,e_4$, and put
$A'=P^{-1}AP$, $B'=P^{-1}BP$.  These matrices are integral:
$A\mathbf1=\mathbf1$ and
\[
 Y\mathbf1=2q\binom{-1}{1-v}
\]
show that the index-$q$ overlattice is stable.  Define
\[
 s=v^2+v-1,\qquad r=-\frac{(v+1)(v^2-2)}2,
 \qquad
 \beta=(s,\ s-vr,\ v^2s+vr,\ r,\ -vs-r).
 \tag{F.U3}\label{maps:eq:F.U3}
\]
The number $r$ is integral for every integer $v$.
Set
\[
 \widehat A=\begin{pmatrix}1&0\\0&A'\end{pmatrix},
 \qquad
 \widehat B=\begin{pmatrix}1&\beta\\0&B'\end{pmatrix}.
 \tag{F.U4}\label{maps:eq:F.U4}
\]
They have orders five and four, and
$(\widehat A\widehat B-I)^2=0$.
To check primitivity without a parameter bound, write
$X'=P^{-1}X$, $Y'=YP$.  Two rows of $X'$ have determinant
$-1$.  The quotient of $\mathbb Z^5$ by its image has
coordinates indexed by $0,3,4$, and the induced matrix $Y'$
there is
\[
 \begin{pmatrix}-2&-1&v\\2(1-v)&v^2-v&q\end{pmatrix}.
\]
Its kernel has primitive generator
$(-v^3+v-1,\,2,\,2(1-v^2))^t$: its first coordinate is
odd and the second is two.  Lift this to
\[
 z=(-v^3+v-1,0,0,2,2(1-v^2))^t.
\]
The same three entries are its signed two-by-two minors,
whose gcd is one; hence $Y'$ is onto $\mathbb Z^2$.
Then $Y'z=0$ and
\[
 \beta z=s(v^3-v-1)+2v^2r=1.
 \tag{F.U5}\label{maps:eq:F.U5}
\]
Thus the rank-two primitive cusp in the quotient acquires a
primitive third image direction in the extension.  All three
nonzero Smith factors of $\widehat A\widehat B-I$ are one.

\item For later homology comparisons use the following constant
change of lattice:
\[
 L=\frac14\begin{pmatrix}
 -2&3&-5&-1&5&-3\\
 0&2&-2&2&-2&2\\
 0&-4&0&0&4&-4\\
 0&0&0&-4&4&-4\\
 0&0&4&-4&4&-4\\
 0&0&0&-4&4&0
 \end{pmatrix},\quad
 A_v=L^{-1}\widehat A L,\quad B_v=L^{-1}\widehat B L.
 \tag{F.U6}\label{maps:eq:F.U6}
\]
Here $\det L=-1/4$, $L^{-1}$ is integral, and all entries of
$A_v,B_v$ are polynomials in $v$ with integer coefficients.
The primitive-cusp assertion is unchanged at odd primes;
at two, direct reduction for $v=0,1$ gives rank three.
At $v=0$, \eqref{maps:eq:F.U6} is exactly the pair in Theorem \ref{maps:src:F.2.5}.

\item For the infinite subfamily $v\in20\mathbb Z$, take homology
monodromies $M_1=A_v^t$, $M_2=B_v^t$ and
$M_0=(M_1M_2)^{-1}$.  The reversed product has the same
primitive square-zero cusp type.  Write
$a=(1,3,-2,-1,2,0)^t$ and $e=(1,0,0,0,0,0)^t$, and set
\[
 w_1=\frac15\sum_{j=0}^4(A_v^t)^j a,
 \qquad
 w_2=-\frac14\sum_{j=0}^3(B_v^t)^j e.
 \tag{F.U7}\label{maps:eq:F.U7}
\]
These vectors are integral.  Indeed, the polynomial matrices
reduce to their parameter-zero values modulo five and four,
respectively, and $a,e$ are fixed by the corresponding
parameter-zero matrices.  They are fixed by $M_1,M_2$, with
first coordinates $1,-1$.  Both monodromies preserve the
first-coordinate covector $\ell$.  Thus the affine actions
$x\mapsto M_i x+w_i/m_i$, with $(m_1,m_2)=(5,4)$, are free:
every nontrivial power translates the $\ell$-circle nontrivially.

\item Finally, let $H=NA|_{\operatorname{im}N}$ and
$H'=NA^{-1}|_{\operatorname{im}N}$.  The construction gives
\[
 \det H=1-v^2q,\quad \det H'=1+v^2q,
 \qquad\operatorname{tr}(A^2C^2)=v^5-v^4+3.
 \tag{F.U8}\label{maps:eq:F.U8}
\]
These are simultaneous-conjugacy invariants; the last is
unbounded on the positive multiples of twenty.  If the
irreducible rank-five representation were self-dual, its
invariant bilinear form would have a nonzero symmetric part,
which is nondegenerate by irreducibility; an alternating
nondegenerate form is impossible in odd dimension.  Relative
to that symmetric form, $A,C$ are orthogonal and
$N^*=-N$, so $I+tN$ is orthogonal for every $t$.
The odd-dimensional characteristic polynomial of $(I+tN)A$
is then anti-reciprocal.  Its quadratic coefficient in $t$ is
\[
 x^2\bigl((\det H)x-\det H'\bigr),
\]
by the determinant lemma, so anti-reciprocity forces
$\det H=\det H'$.  Formula \eqref{maps:eq:F.U8} contradicts this whenever
$v\ne0$.

\item All the matrix identities, coefficient integrality, both
mod-two checks, and the affine formula are checked symbolically
by \texttt{code/fourfolds\_unbounded\_family.py}.  Its separate
Smith diagnostics over $-6\le v\le10$ are explicitly bounded
and are not the basis for the universal assertions above.

\end{enumerate}
\end{proof}




\begin{proposition}[an odd torsion obstruction in this lattice]\label{maps:src:F.2.17}
\StatementRecord{maps:src:F.2.17}{proved}{Historical events: Claude Fable 5.2 (Anthropic), effort high, 16 November 2026: carried as background without fault found. GPT-6 Astra (OpenAI), ChatGPT Pro mode, effort Pro, S6\_09\_PAPERS, 17 November 2026: targeted proof check. Also checked in this writing. Current writing events: GPT-6 Astra (OpenAI), ChatGPT Pro mode, effort Pro, S6\_10\_PAPERS, 18 November 2026: computation or formal certificate; targeted proof check. GPT-6 Astra (OpenAI), ChatGPT Pro mode, effort Pro, S6\_10\_PAPERS continuation, 19 November 2026: used and rechecked in this run. \emph{Recorded proof scope.} Claude Fable 5.2 (Anthropic), effort high, 16 November 2026; also checked in this writing by GPT-6 Astra (OpenAI), ChatGPT Pro mode, effort Pro, 18 November 2026}{the hypotheses, notation and definitions stated here; \cref{maps:eq:F.U6}; \cref{maps:src:F.2.19}; \cref{maps:eq:F.U5}; \cref{maps:src:F.2.12}}


For the integral tuple $(A_v,B_v)$ in \eqref{maps:eq:F.U6}, any simply
connected filling with free finite affine caps and the fixed
primitive Kuhn cusp has a subgroup
\[
 \mathbb Z/(v^2-v+1)\ \subset\operatorname{Tor}H_2(M;\mathbb Z).
 \tag{F.U9}\label{maps:eq:F.U9}
\]
In particular, in this specified family of lattices only
$v=0$ and $v=1$ can be three-connected, regardless of the
affine multiplicities and twists.  The stronger exclusion for every stable lattice in each
rational class is proved in Theorem \ref{maps:src:F.2.19} below.

\end{proposition}


\begin{proof}
\begin{enumerate}[label=\arabic*.]
\item First compute the global coinvariants in the hatted lattice.
In the rank-five quotient, the $A'$ relations identify
$f_1=f_2=f_3=f_4=t$ and $qf_0=5t$.
The number $q=v^2-v+1$ is relatively prime to five, as its
five residue values show.  Thus these coinvariants are
infinite cyclic, with $f_0=5h$ and $t=qh$.
The $B'$ relations become $q\kappa_j h$, where
\[
 (\kappa_0,\ldots,\kappa_4)=
 (2v^2,\ v(v^3+v-1),\ v(v^3+v+1),
       -v^3+v+1,\ -v^3-v-1).
 \tag{F.U10}\label{maps:eq:F.U10}
\]
Their gcd is one: it divides $2v$ by subtracting the second
and third entries, and it divides $v^3-v-1$ by the fourth;
the latter is odd and is $-1$ modulo $v$.
Hence the rank-five global coinvariants are $\mathbb Z/q$.
The extension line is killed in rank-six coinvariants by
\eqref{maps:eq:F.U5}, since
$(\widehat A\widehat B-I)(0,z)^t=e$.
Consequently the hatted rank-six coinvariants are also
$\mathbb Z/q$.  The lattice change \eqref{maps:eq:F.U6} is an isomorphism
over every $\mathbb Z_{(p)}$ with $p$ odd, so the odd part of
the actual rank-six global coinvariants is precisely
$\mathbb Z/q$.

\item Let $V$ be the degree-one fibre cohomology local system.
The only difference between $R^1g_*\mathbb Z$ and $j_*V$
is a finite cokernel at finite caps; the cusp degree-one
specialization is onto.  Thus
\[
 H^2(\mathbb P^1,R^1g_*\mathbb Z)
 =H^2(\mathbb P^1,j_*V)=V_\Gamma,
\]
including its $q$-primary part just computed.  The space of
global degree-two invariants is zero: the irreducible odd
rank-five quotient has no invariant alternating tensor,
and the extension filtration
$0\to W\to\bigwedge^2V\to\bigwedge^2W\to0$
gives $(\bigwedge^2V)^\Gamma=0$.

\item It remains to check every possible source of the differential
$d_2:H^0(\mathbb P^1,R^2g_*\mathbb Z)
\to H^2(\mathbb P^1,R^1g_*\mathbb Z)$.
The cusp has no kernel in degree two.  At a finite cap the
kernel is torsion in $H^2$ of that cap, namely torsion in its
first homology.  Simple connectivity of the total filling
forces the distinguished affine numerator to be prime to the
actual affine multiplicity, by the determinant argument of
\ref{maps:src:F.2.12}.  The cap is therefore a mapping torus of a torus
automorphism whose order divides the linear order, five or
four.  Its first-homology torsion is supported only at the
primes five or two, respectively: averaging over the finite
order kills the torsion of the vertical coinvariants.
Hence the source of this $d_2$ has no torsion at any prime
dividing $q$, since $q$ is odd and prime to five.

\item The subgroup $\mathbb Z/q$ in bidegree $(2,1)$ survives.
There are no longer differentials on a two-dimensional base,
and this terminal filtration term is a subgroup of
$H^3(M;\mathbb Z)$.  Universal coefficients identify its
torsion with a subgroup of $\operatorname{Tor}H_2(M;\mathbb Z)$.
This proves \eqref{maps:eq:F.U9}.  Finally, $v^2-v+1=1$ for an integer $v$
exactly when $v=0$ or $v=1$.

\end{enumerate}
\end{proof}




\begin{theorem}[a survivor with a non-self-dual quotient]\label{maps:src:F.2.18}
\StatementRecord{maps:src:F.2.18}{proved}{Historical events: Claude Fable 5.2 (Anthropic), effort high, 16 November 2026: carried as background without fault found. GPT-6 Astra (OpenAI), ChatGPT Pro mode, effort Pro, S6\_09\_PAPERS, 17 November 2026: targeted proof check. Also checked in this writing. Current writing events: GPT-6 Astra (OpenAI), ChatGPT Pro mode, effort Pro, S6\_10\_PAPERS, 18 November 2026: computation or formal certificate; targeted proof check. GPT-6 Astra (OpenAI), ChatGPT Pro mode, effort Pro, S6\_10\_PAPERS continuation, 19 November 2026: used and rechecked in this run. \emph{Recorded proof scope.} Claude Fable 5.2 (Anthropic), effort high, 16 November 2026; also checked in this writing by GPT-6 Astra (OpenAI), ChatGPT Pro mode, effort Pro, 18 November 2026}{the hypotheses, notation and definitions stated here; \cref{maps:eq:F.U6}; \cref{maps:src:F.3.11}; \cref{maps:src:F.2.5}; \cref{maps:src:F.2.16}; \cref{maps:src:F.2.6}; \cref{maps:eq:F.6}; \cref{maps:src:F.4.1}; \cref{maps:eq:F.U8}}


The parameter $v=1$ in \eqref{maps:eq:F.U6}, with affine orders five and four
and twists
\[
 w_1=(1,3,-2,-1,2,0)^t,\qquad
 w_2=(-1,0,0,0,0,0)^t,
\]
gives a three-connected smooth eight-manifold with
$H_4=\mathbb Z^4$ and all other reduced homology zero except
$H_8=\mathbb Z$.  It is almost complex and is oriented
diffeomorphic to $3\mathbb{HP}^2\#(-\mathbb{HP}^2)$.
Its irreducible rank-five quotient has no invariant bilinear
form, so integral three-connectivity does not force that
quotient to be orthogonal.  It has bounded-cusp holomorphic
periods by Theorem \ref{maps:src:F.3.11}; simultaneous finite holomorphic
extensions are not asserted here.

\end{theorem}


\begin{proof}
\begin{enumerate}[label=\arabic*.]
\item The first cohomological matrix is the same $A$ as in
Theorem \ref{maps:src:F.2.5}.  The second is
\[
 B_1=\begin{pmatrix}
 1&0&0&0&0&0\\
 0&0&-1&0&1&-1\\
 0&0&-1&0&0&0\\
 0&0&-2&0&1&-2\\
 0&1&-2&0&1&-3\\
 0&0&-1&1&0&-1
 \end{pmatrix}.
 \tag{F.U11}\label{maps:eq:F.U11}
\]
Take $M_1=A^t$, $M_2=B_1^t$ and
$M_0=(M_1M_2)^{-1}$.  The displayed twists are fixed integral
vectors with $\ell(w_1)=1$, $\ell(w_2)=-1$.
They therefore give free affine actions.  The cusp is
primitive by \ref{maps:src:F.2.16}.  Exact Hermite reduction gives
\[
 \operatorname{col}_{\mathbb Z}
 [\,M_0-I\mid M_1(M_0-I)\mid M_2(M_0-I)\,]
       =\langle e_1,\ldots,e_5\rangle.
\]
Van Kampen leaves only $5a=h$, $4b=-h$, $a+b=0$,
whose relation matrix is unimodular.  This proves
$\pi_1=1$ for the specified gluing, including zero cusp
clutching.

\item For completeness we give the full new integral certificates
needed to reuse the already proved differential calculation
in \ref{maps:src:F.2.6}.  The rational exterior invariant and parabolic ranks
are exactly the table there.  The finite cap first homology
groups are $\mathbb Z^2$ and $\mathbb Z^2\oplus\mathbb Z/2$.
On integral invariant two-form bases their suspension periods
divided by the covering orders are, respectively,
$(-1,0,-1)$ and $(0,0,1)$.  The determinants of the invariant
degree-three skew pairings are $5^4$ and $4^4$, while the
degree-two/degree-four pairings have absolute determinants
$5^3$ and $4^3$.  Thus all the finite covering indices used
in degrees two, three and four are one.

\item For exterior coefficient degrees two, three and four, the
global coboundary Smith form has one factor two and all other
nonzero factors one.  Its torsion class has peripheral
triviality flags $(\mathrm{true},\mathrm{false},\mathrm{true})$,
so the parabolic groups are torsion free, of ranks $0,3,0$.
The degree-one coboundary is primitive and the degree-one
finite covering cokernels are killed by the common invariant
covector and the Chinese remainder theorem, exactly as in
\ref{maps:src:F.2.6}.  The global coinvariants in degrees one through four
are each $\mathbb Z/2$.

\item Write $M_2=\operatorname{diag}(1,D)$.  This time use the
unimodular marking
\[
 P_1=\begin{pmatrix}
 -1&1&-1&0&-1\\
 -1&2&-1&0&-1\\
 -1&0&0&2&-1\\
 1&-1&0&-1&1\\
 0&0&2&-1&1
 \end{pmatrix},\qquad \det P_1=-1.
\]
It satisfies $DP_1=P_1\operatorname{diag}(\text{four-cycle},-1)$.
The cap therefore has the same cellular mapping-torus model
used to prove \eqref{maps:eq:F.6}.  With the unchanged twist $w_2=-e_0$,
the boundary coordinates remain $\theta=a+4k$, $h=k$;
the suspension representative restricts to $w$ in the base
direction and to $-\operatorname{Norm}_4(w)=0$ in the
auxiliary fibre direction.  Consequently the coefficient in
the torsion differential is again one.

\item The new vertical representatives in lexicographic exterior
bases are
\[
\begin{split}
 q=2:\ &(0,-1,-1,-1,0),\\
 q=3:\ &(0,0,1,0,0,0,0,0,0,1),\\
 q=4:\ &(0,0,0,0,0,0,0,-1,-1,0).
\end{split}
\]
Each has residue one in its global coinvariant
$\mathbb Z/2$.  Thus all three torsion differentials in \eqref{maps:eq:F.6}
are isomorphisms.  The integral Leray table and its free
degree-one differential are therefore exactly those of
\ref{maps:src:F.2.6}.  They give $H^2=H^3=0$ and $H^4=\mathbb Z^4$.
Universal coefficients, Hurewicz and Poincare duality give
the stated integral homology and three-connectivity.

\item The relative almost-complex extension over the pair of pants
in \ref{maps:src:F.2.5} applies: the finite matrices have the same even
real eigenspace dimensions and the cusp is the same marked
Kuhn type.  The purely smooth characteristic-vector and
normal-invariant argument in \ref{maps:src:F.4.1} then gives the stated
oriented diffeomorphism type.  No integrability at the finite
punctures is used for this smooth classification.

\item Finally, \eqref{maps:eq:F.U8} has $(\det H,\det H')=(0,2)$ at $v=1$,
so its irreducible rank-five quotient is not self-dual.
The exact program
\texttt{code/fourfolds\_v1\_homology.py} runs every matrix,
Smith, Hermite, suspension-period, pairing and differential
certificate just displayed.  Its monomial marking was found
inside the stated finite cube and then verified directly;
the proof uses that explicit marking, not search completeness.

\end{enumerate}
\end{proof}




\begin{theorem}[a residual obstruction for every lattice]\label{maps:src:F.2.19}
\StatementRecord{maps:src:F.2.19}{proved}{Historical events: Claude Fable 5.2 (Anthropic), effort high, 16 November 2026: carried as background without fault found. GPT-6 Astra (OpenAI), ChatGPT Pro mode, effort Pro, S6\_09\_PAPERS, 17 November 2026: targeted proof check. Also checked in this writing. Current writing events: GPT-6 Astra (OpenAI), ChatGPT Pro mode, effort Pro, S6\_10\_PAPERS, 18 November 2026: computation or formal certificate; targeted proof check. GPT-6 Astra (OpenAI), ChatGPT Pro mode, effort Pro, S6\_10\_PAPERS continuation, 19 November 2026: used and rechecked in this run. \emph{Recorded proof scope.} Claude Fable 5.2 (Anthropic), effort high, 16 November 2026; also checked in this writing by GPT-6 Astra (OpenAI), ChatGPT Pro mode, effort Pro, 18 November 2026}{the hypotheses, notation and definitions stated here; \cref{maps:src:F.2.16}; \cref{maps:eq:F.U6}; \cref{maps:src:F.2.6}; \cref{maps:src:F.2.18}; \cref{maps:eq:F.U12}; \cref{maps:src:F.2.17}; \cref{maps:eq:F.U1}}


Let a rational rank-six datum have finite linear orders five
and four, a primitive maximal cusp on a stable integral lattice,
and invariant line $L$, with irreducible rank-five quotient $W$
and no rational global coinvariants.  Suppose that at a prime
$p\ne2,5$ the semisimplified reduction of a stable lattice in
$W$ is the sum of the trivial module and an absolutely
irreducible module of dimension four.  Then no stable integral
lattice in that rational datum, with primitive maximal cusp,
and no free finite affine twists can give
$\pi_1=1$ and $H_2=0$ simultaneously.

For every $v\notin\{0,1\}$, the rational family in \ref{maps:src:F.2.16}
satisfies this hypothesis at each prime dividing
$v^2-v+1$.  Consequently the exclusion of those parameters
holds for \emph{all} stable lattices and all twists, not just
the lattice in \eqref{maps:eq:F.U6}.  The two remaining parameters are the
verified survivors of \ref{maps:src:F.2.6} and \ref{maps:src:F.2.18}.

\end{theorem}


\begin{proof}
\begin{enumerate}[label=\arabic*.]
\item We first give the lattice-independent mechanism.  For any
stable lattice $V_\mathbb Z$, use its primitive invariant
line $L_\mathbb Z$ and quotient lattice $W_\mathbb Z$.
The semisimplified reduction of a rational representation is
independent of its stable lattice.  One elementary way to see
this is to compare two commensurable lattices by their finite
quotient $T$.  The exact sequence obtained by reduction modulo
$p$ has end terms $T[p]$ and $T/pT$, which have the same
composition factors by the filtration of $T$ by powers of $p$.
Hence the two lattice reductions have equal classes in the
Grothendieck group of finite-length modules.  Chaining such
comparisons proves the assertion for any two lattices.

\item Thus $W_\mathbb Z/pW_\mathbb Z$ has length two, with one
trivial constituent and one irreducible four-dimensional
constituent.  It has either a trivial submodule or a trivial
quotient (and possibly both).

\item Suppose first that it has a trivial submodule.  Its inverse
image in $V_\mathbb Z/pV_\mathbb Z$ is an extension of a
trivial module by a trivial module.  Each finite generator
acts on this two-dimensional space by a unipotent matrix.
Since its order is five or four, prime to $p$, it acts as
the identity.  There are therefore at least two global
invariant vectors in degree-one fibre cohomology modulo $p$.
Write their dimension as $a_p\ge2$; this is also the
dimension of the global coinvariants of fibre homology over
$\mathbb F_p$, by the perfect pairing between fibre homology
and cohomology.

\item The abelianized van Kampen presentation now shows
\[
 b_1(M;\mathbb F_p)\ge a_p-1.
 \tag{F.U12}\label{maps:eq:F.U12}
\]
Indeed, after the monodromy relations the fibre contributes
$a_p$ dimensions, and the two finite meridians contribute
two more.  The two finite power relations and the single
base relation impose at most three linear relations.
The cusp introduces no further fibre relation in this
abelianization.  On the dual homology lattice its nilpotent
operator is still primitive and maximal, so its angular lattice is
$\ker(M_0-I)=\operatorname{im}(M_0-I)$ by primitivity, and
the latter image already vanishes among the global
monodromy relations.  Its base meridian is the base relation
just counted.  Changing affine twists changes coefficients
of these three relations, not their number.  Thus \eqref{maps:eq:F.U12}
contradicts simple connectivity.

\item Suppose instead that the reduction of $W$ has a trivial
quotient.  Its integral global coinvariants have nonzero
$p$-primary torsion, since they have rational rank zero.
The natural map from the global coinvariants of $V$ to
those of $W$ is an isomorphism.  To check its kernel, a
primitive generator $e$ of $L_\mathbb Z$ lies in
$\ker(M_0-I)=\operatorname{im}(M_0-I)$, and hence is already
zero among global coinvariants.  Therefore $H^2(R^1)$ has
nonzero $p$-primary torsion.  If the total space were simply
connected, the finite-cap argument of \ref{maps:src:F.2.17} would make
every possible torsion source in $H^0(R^2)$ supported only
at two and five.  The invariant part of $\bigwedge^2V$ is
zero by irreducibility of the odd-dimensional quotient.
Consequently the nonzero $p$-primary subgroup survives to
$H^3(M;\mathbb Z)$ and then to
$\operatorname{Tor}H_2(M;\mathbb Z)$, again a contradiction.
This proves the general criterion.

\item For its application to \ref{maps:src:F.2.16}, let $p\mid q=v^2-v+1$.
We have $p\ne2,5$.  In the original cycle lattice of \eqref{maps:eq:F.U1},
$Y\mathbf1=2q(-1,1-v)^t$, so $\mathbf1$ gives a trivial
submodule after reduction.  On the four-dimensional quotient
the first finite generator has characteristic polynomial
$\Phi_5$, and $C=B^{-1}$ has polynomial
$(x+1)^2(x^2+1)$.  This quotient is absolutely irreducible.
Here is a characteristic-independent proof over
$\overline{\mathbb F}_p$.

\item A one-dimensional composition factor would make the
eigenvalues of $A$ and $C$ equal, because $CA^{-1}$ is
unipotent.  An eigenvalue of order dividing both five and
four must be one, which is absent from $\Phi_5$.
If the quotient were reducible, its only remaining possible
composition-factor dimensions would be two and two.
On each factor, the determinants of $A$ and $C$ must agree
and have orders dividing five and four, so both determinants
are one.  The two $C$ eigenvalue pairs must therefore be
$(-1,-1)$ and $(i,-i)$ in some order.  On the former
factor $C=-I$, since its order is prime to $p$.
Then the unipotence of $CA^{-1}=-A^{-1}$ would force the
eigenvalues of $A$ to be $-1$, impossible for $\Phi_5$.
This proves absolute irreducibility and establishes the
required residual composition factors at every $p\mid q$.
When $v\notin\{0,1\}$, the integer $q$ is greater than one
and has such a prime divisor.  The criterion applies.

\end{enumerate}
\end{proof}




\begin{theorem}[the complete rational $(5,4)$ normal form]\label{maps:src:F.2.20}
\StatementRecord{maps:src:F.2.20}{proved}{Historical events: Claude Fable 5.2 (Anthropic), effort high, 16 November 2026: carried as background without fault found. GPT-6 Astra (OpenAI), ChatGPT Pro mode, effort Pro, S6\_09\_PAPERS, 17 November 2026: targeted proof check. Also checked in this writing. Current writing events: GPT-6 Astra (OpenAI), ChatGPT Pro mode, effort Pro, S6\_10\_PAPERS, 18 November 2026: computation or formal certificate; targeted proof check. GPT-6 Astra (OpenAI), ChatGPT Pro mode, effort Pro, S6\_10\_PAPERS continuation, 19 November 2026: used and rechecked in this run. \emph{Recorded proof scope.} Claude Fable 5.2 (Anthropic), effort high, 16 November 2026; also checked in this writing by GPT-6 Astra (OpenAI), ChatGPT Pro mode, effort Pro, 18 November 2026}{the hypotheses, notation and definitions stated here; \cref{maps:src:F.2.11}; \cref{maps:eq:F.R4}; \cref{maps:eq:F.R6}; \cref{maps:eq:F.R7}; \cref{maps:eq:F.R8}}


Consider ordered rational rank-five pairs $(A,C)$ with
\[
 \chi_A=(x-1)\Phi_5,
 \qquad \chi_C=(x-1)(x+1)^2(x^2+1).
\]
Put
\[
 N=CA^{-1}-I,
 \qquad N^2=0,\quad\operatorname{rank}N=2,
\]
where $A^5=C^4=I$ and the pair is irreducible over $\mathbb Q$.
Their simultaneous rational conjugacy classes are parametrized
bijectively by
\[
 \bigl\{[r:s:t]\in\mathbb P^2(\mathbb Q):D(r,s,t)\ne0\bigr\}
 \ \sqcup\ \mathbb A^1(\mathbb Q),
 \tag{F.R1}\label{maps:eq:F.R1}
\]
where
\[
\begin{split}
 D={}&r^3-4r^2s-r^2t+7rs^2+rst-3rt^2\\
     &{}-7s^3+11s^2t-11st^2+7t^3.
\end{split}
 \tag{F.R2}\label{maps:eq:F.R2}
\]
The second summand lies over the single exceptional point
$[4:1:3]$ of $D=0$.  The construction below is explicit and
uses only a ten-variable linear system.  By \ref{maps:src:F.2.11} it also
parametrizes the rational rank-six maximal-cusp extensions
of all these pairs, with $B=C^{-1}$.
This is a rational classification; the stable integral
lattices, affine twists and surviving homology still require
the arithmetic conditions of F.2.

\end{theorem}


\begin{proof}
\begin{enumerate}[label=\arabic*.]
\item Write $K=\mathbb Q(\zeta)$, $\zeta^4+\zeta^3+\zeta^2+
\zeta+1=0$, and identify the rational space with
$\mathbb Q\oplus K$.  The matrix of $A$ is multiplication by
$(1,\zeta)$:
\[
 A=\operatorname{diag}(1,J),\qquad
 J=\begin{pmatrix}
 0&0&0&-1\\1&0&0&-1\\0&1&0&-1\\0&0&1&-1
 \end{pmatrix},\qquad R=(A+I)^{-1}.
 \tag{F.R3}\label{maps:eq:F.R3}
\]
Its rational centralizer is $\mathbb Q^*\times K^*$.
Let $I_N=\operatorname{im}N$.  If $X$ is any basis matrix
for this plane and $N=XY$, then $YX=0$.
The identity
\[
 C+I=(A+I)+XYA
\]
shows that its two-dimensional kernel has the same dimension
as the kernel of
$I_2+YARX=I_2-YRX$.  Hence
\[
 YRX=I_2,
 \qquad I_N\cap RI_N=0.
 \tag{F.R4}\label{maps:eq:F.R4}
\]
Conversely the two matrix equations in \eqref{maps:eq:F.R4}, together with
$YX=0$, give precisely that two-dimensional $-1$ eigenspace.

\item The plane $I_N$ projects nontrivially to $\mathbb Q$:
otherwise the first-coordinate covector is invariant under
both $A$ and $C$, contrary to irreducibility.
It contains no nonzero vector of $\mathbb Q\oplus0$,
because $R$ acts as $1/2$ on that line and \eqref{maps:eq:F.R4} would fail.
Write a basis of the plane as $(0,u),(1,v)$, where
$u\in K^*$.  Multiplication on $K$ by $u^{-1}$ puts it into
\[
 I_N=\langle(0,1),(1,\beta)\rangle_{\mathbb Q}.
\]
Changing the second basis vector by a rational multiple of
the first changes $\beta$ by a rational constant.
Two such planes are in the same centralizer orbit exactly
when $\beta'=a\beta+b$, with $a\in\mathbb Q^*$ and
$b\in\mathbb Q$.  Indeed a centralizer element carrying
one to the other must preserve the kernel of the projection
to $\mathbb Q$, so its $K^*$ component is itself rational;
the displayed affine change follows, and every such change
is realized.  We may therefore take
\[
 \beta=r\zeta+s\zeta^2+t\zeta^3,
 \qquad [r:s:t]\in\mathbb P^2(\mathbb Q),
 \qquad
 X=\begin{pmatrix}0&1\\1&0\\0&r\\0&s\\0&t\end{pmatrix}.
 \tag{F.R5}\label{maps:eq:F.R5}
\]
The triple is nonzero; a rational $\beta$ would make the
plane contain the forbidden first-coordinate line.
The stabilizer of such a plane in the centralizer consists
only of common scalar matrices.  Thus it introduces no
further identifications of the resulting nilpotent operators.

\item Write $Y$ as an unknown two-by-five matrix.  It is determined
by the following ten linear equations:
\[
 YX=0,\qquad YRX=I_2,\qquad
 \operatorname{tr}(YAX)=-1,\qquad
 \operatorname{tr}(YA^{-1}X)=1.
 \tag{F.R6}\label{maps:eq:F.R6}
\]
The first eight are \eqref{maps:eq:F.R4}.  The last two are exactly
$\operatorname{tr}C=\operatorname{tr}C^{-1}=-1$, since
$\operatorname{tr}A=\operatorname{tr}A^{-1}=0$ and
$C^{-1}=A^{-1}(I-XY)$.
They suffice to obtain the required characteristic
polynomial: the two $-1$ roots and determinant one leave a
monic cubic of constant term $-1$, and the two traces fix
that cubic to $(x-1)(x^2+1)$.
These remaining three roots are distinct, and the $-1$
eigenspace already has dimension two, so $C$ is semisimple
and $C^4=I$.  Also $YRX=I_2$ ensures that $N$ has rank two.
Thus every solution of \eqref{maps:eq:F.R6} gives all the asserted local
linear conditions.

\item In row-major order on the ten entries of $Y$, the determinant
of the coefficient matrix in \eqref{maps:eq:F.R6} is $D/4$.
For $D\ne0$ this proves uniqueness and gives the explicit
formula $Y=\operatorname{reshape}(E^{-1}b)$, where $E$ is
the coefficient matrix of the displayed linear equations
and $b$ their displayed right side.
The exact expanded matrix and all ten rational functions
are printed by \texttt{fourfolds\_relative\_normalform.py}.
There is no search or degree cutoff in this calculation.

\item We check the entire exceptional locus, rather than discard it.
Set
\[
 Q_1=r^2-rs-rt+s^2+t^2,
 \qquad Q_2=r^2-2rs-3rt+4s^2-4st+4t^2.
\]
The first column of the unique generic solution is
\[
 y_{10}=\frac{2Q_1Q_2}{D},\qquad
 y_{20}=\frac{2(r-s-t)Q_1}{D}.
 \tag{F.R7}\label{maps:eq:F.R7}
\]
Here $y_{10},y_{20}$ are the first column of $Y$.
The two equalities after multiplication by $D$ are polynomial
linear combinations of \eqref{maps:eq:F.R6}, by its adjugate matrix.
The program prints the corresponding two polynomial rows
and verifies those identities, so they remain necessary
also when $D=0$.
Moreover
\[
 2Q_1=(r-s)^2+(r-t)^2+s^2+t^2>0
\]
for every nonzero rational triple.  Consequently consistency
on $D=0$ forces $r=s+t$ and $Q_2=0$.
Substitution gives
\[
 D(s+t,s,t)=(t-3s)(2t-s)^2,
 \qquad Q_2(s+t,s,t)=(3s-t)(s-2t).
 \tag{F.R8}\label{maps:eq:F.R8}
\]
There are precisely two resulting projective points,
$[3:2:1]$ and $[4:1:3]$.
At the former, $[X\mid RX]$ has rank less than four, so
\eqref{maps:eq:F.R4} is impossible.  At the latter the coefficient and
augmented matrices of \eqref{maps:eq:F.R6} both have rank nine, and every
solution is
\[
 Y_z=\begin{pmatrix}
 4z+16/5&0&-2z-1&4/5-2z&2z\\
 2z-2/5&0&-z&2/5-z&z
 \end{pmatrix},\qquad z\in\mathbb Q.
 \tag{F.R9}\label{maps:eq:F.R9}
\]
Direct substitution proves both necessity and sufficiency.
This establishes the exceptional affine line in \eqref{maps:eq:F.R1}.

\item Every pair just constructed is irreducible over $\mathbb Q$.
An $A$-stable proper rational subspace has dimension one or
four.  A one-dimensional subrepresentation or quotient must
have trivial character, since its finite eigenvalues of
orders dividing five and four must agree under the
unipotent product.  The only $A$-invariant covector is the
first-coordinate covector.  Its product with $N$ is the
second row of $Y$, which is nonzero by $YRX=I_2$.
The only $A$-invariant vector is the first-coordinate vector;
its image under $N$ vanishes exactly if both entries in
\eqref{maps:eq:F.R7} vanish.  For $D\ne0$, positivity of $Q_1$ and
\eqref{maps:eq:F.R8} rule this out.  On the exceptional line the two
entries $4z+16/5$ and $2z-2/5$ never vanish simultaneously.
There are therefore neither invariant vectors nor covectors,
which excludes both possible dimensions.

\item Finally, rescaling $(r,s,t)$ is exactly the centralizer
change already described; it gives a conjugate pair.
Conversely any conjugacy of two normalized pairs centralizes
$A$ and hence identifies their projective parameters as proved
above.  At the exceptional point its stabilizer is scalar,
so distinct $z$ give distinct pairs.  This proves the stated
bijection, including completeness and uniqueness.
The saved program checks all determinant, adjugate,
characteristic-polynomial and finite-order identities over
the rational function field, and checks the exceptional
line symbolically for every $z$.

\end{enumerate}
\end{proof}




\begin{corollary}[all good-prime lattices are forced]\label{maps:src:F.2.21}
\StatementRecord{maps:src:F.2.21}{proved}{Historical events: Claude Fable 5.2 (Anthropic), effort high, 16 November 2026: carried as background without fault found. GPT-6 Astra (OpenAI), ChatGPT Pro mode, effort Pro, S6\_09\_PAPERS, 17 November 2026: targeted proof check. Also checked in this writing. Current writing events: GPT-6 Astra (OpenAI), ChatGPT Pro mode, effort Pro, S6\_10\_PAPERS, 18 November 2026: targeted proof check. GPT-6 Astra (OpenAI), ChatGPT Pro mode, effort Pro, S6\_10\_PAPERS continuation, 19 November 2026: used and rechecked in this run. \emph{Recorded proof scope.} Claude Fable 5.2 (Anthropic), effort high, 16 November 2026; also checked in this writing by GPT-6 Astra (OpenAI), ChatGPT Pro mode, effort Pro, 18 November 2026}{the hypotheses, notation and definitions stated here; \cref{maps:src:F.2.20}; \cref{maps:src:F.2.19}}


For a $(5,4)$ datum with $\pi_1=1$ and $H_2=0$, the
rank-five quotient has absolutely irreducible reduction
at every prime $p\ne2,5$.  At each such prime its stable
$\mathbb Z_p$-lattice is unique up to homothety.
The same is true of the rank-six extension among lattices
with primitive maximal cusp.  Thus, once a rational pair
in \ref{maps:src:F.2.20} admits one survivor, all remaining integral-lattice
ambiguity is supported at the primes two and five.
This assertion supplies no finite bound on those bad-prime
lattices or on the rational parameters.

\end{corollary}


\begin{proof}
\begin{enumerate}[label=\arabic*.]
\item Work first over an algebraic closure of $\mathbb F_p$.
Both finite generators are semisimple because $p\ne2,5$.
A one-dimensional composition factor must be trivial:
its $A$ and $C$ eigenvalues agree by the unipotent product,
and their orders divide five and four.
There can be at most one such factor, since the $1$
eigenvalue of $A$ occurs once.
If it occurs, the remaining four-dimensional module is
absolutely irreducible by the spectral argument in \ref{maps:src:F.2.19}.
The only other possible reducible composition pattern is
$2+3$.
On each factor, determinants of $A$ and $C$ agree and
have orders dividing five and four, hence are one.
The two-dimensional $C$ eigenvalue pair is therefore
$(-1,-1)$ or $(i,-i)$.
The first would give $C=-I$ on that factor, contradicting
the unipotence of $CA^{-1}$ and the absence of $-1$
among the eigenvalues of $A$.
In the second case the three-dimensional factor has
$C$ eigenvalues $1,-1,-1$, so $C^2=I$ there; it has
an $A$-fixed vector because the other factor has determinant
one and cannot contain the simple $1$ eigenvalue of $A$.

\item Here is the required elementary involutory identity in the
present characteristic.  Put $B=A^{-1}$ on this factor.
From $(CB-I)^2=0$ and $C^2=I$ one obtains
\[
 BC+CB^{-1}=2I.
\]
For $Bv=v$ this says $(B+I)Cv=2v$.
The operator $B+I$ is invertible, since $B^5=I$ and
$p\ne2$, so $Cv=v$.  The resulting common fixed line
contradicts simplicity of the three-dimensional factor.
We have proved that any reducible residual representation
has exactly the constituents $1$ and an absolutely
irreducible four-dimensional module.  Theorem \ref{maps:src:F.2.19}
excludes that possibility for a survivor, proving absolute
irreducibility at every indicated prime.

\item Let $R=\mathbb Z_p$ and $K_p=\mathbb Q_p$.
For two stable quotient lattices, multiply one by a power
of $p$ until it is contained in the other but not in its
$p$-multiple.  Its nonzero residual image must then be the
whole quotient by absolute irreducibility.  Nakayama's
lemma, or the Smith form of the inclusion, makes the
lattices equal.  This proves quotient-lattice uniqueness.

\item For completeness the extension-lattice comparison is also
integral.  Fix a primitive-cusp lattice $V_R$, with invariant
line $Re$ and quotient $W_R$.
The coboundary map
\[
 W_R^*\longrightarrow W_R^*\oplus W_R^*
\]
is primitive, because its reduction has trivial kernel.
Thus its cokernel, the degree-one cohomology of the free
group on the two generators, is torsion free, and an
integral coboundary over $K_p$ is already integral over $R$.
The extension class $\alpha$ is primitive in this cokernel.
Indeed, divisibility by $p$ would make the extra cusp map
\[
 \ker N_W/\operatorname{im}N_W\longrightarrow Re
\]
divisible by $p$, contrary to the three unit Smith factors
of the full cusp.  Here $N_W$ has primitive image:
it is the quotient of the primitive image of $N_V$ by
$Re\subset\operatorname{im}N_V$.

\item Normalize any other stable lattice by the already proved
quotient homothety so that its quotient is $W_R$.
Its invariant line is $p^rRe$ for some integer $r$.
The extension class in that basis is $p^{-r}\alpha$;
primitivity makes it integral only when $r\leq0$.
Primitivity of the coboundary map eliminates any additional
rational change of splitting.  The extra cusp map is then
multiplied by $p^{-r}$, so its Smith factor is a unit exactly
when $r=0$.  The lattices coincide after the common homothety.
This proves the full assertion and retains the bad-prime
and rational-parameter problems explicitly.

\end{enumerate}
\end{proof}




\begin{theorem}[a terminating lattice algorithm for each rational pair]\label{maps:src:F.2.22}
\StatementRecord{maps:src:F.2.22}{proved}{Historical events: Claude Fable 5.2 (Anthropic), effort high, 16 November 2026: carried as background without fault found. GPT-6 Astra (OpenAI), ChatGPT Pro mode, effort Pro, S6\_09\_PAPERS, 17 November 2026: targeted proof check. Also checked in this writing. Current writing events: GPT-6 Astra (OpenAI), ChatGPT Pro mode, effort Pro, S6\_10\_PAPERS, 18 November 2026: computation or formal certificate; targeted proof check. GPT-6 Astra (OpenAI), ChatGPT Pro mode, effort Pro, S6\_10\_PAPERS continuation, 19 November 2026: used and rechecked in this run. \emph{Recorded proof scope.} Claude Fable 5.2 (Anthropic), effort high, 16 November 2026; also checked in this writing by GPT-6 Astra (OpenAI), ChatGPT Pro mode, effort Pro, 18 November 2026}{the hypotheses, notation and definitions stated here; \cref{maps:src:F.2.14}; \cref{maps:eq:F.L3}; \cref{maps:eq:F.L5}; \cref{maps:src:F.2.11}; \cref{maps:eq:F.L6}; \cref{maps:eq:F.L7}}


Let $(A,C)$ be any irreducible rational rank-five pair in
one of the eight spectral cases of \ref{maps:src:F.2.14}, with
$A^5=I$, $B=C^{-1}$, and $(AB-I)^2=0$ of rank two.
There is an explicit terminating algorithm which decides
whether the pair preserves an integral lattice, and, if so,
lists all such lattices up to simultaneous integral conjugacy.
It then lists all rank-six extensions with the same finite
orders and primitive maximal square-zero cusp, up to
simultaneous integral conjugacy.  The bounds below depend
on the input rational pair and are proved to lose no lattice.
They do not bound the rational parameter space or classify
all affine twists and resulting homology.

\end{theorem}


\begin{proof}
\begin{enumerate}[label=\arabic*.]
\item We give the algorithm and its termination proof.
The fixed line of $A$ has dimension one.  Choose a nonzero
vector $e$ on it and a nonzero fixed covector $\varphi$.
Irreducibility implies that the vectors $ge$, with $g$ a
word in $A,C$, span the whole rational space; the analogous
statement holds for the covectors $\varphi h$.
In fact words of length at most four suffice.  The spans of
words of lengths at most $j$ increase from dimension one;
if one step does not increase the span, it is invariant
under both generators and hence already the whole space.
There are only four possible increases.

\item Choose words $g_1,\ldots,g_5$ and $h_1,\ldots,h_5$ giving
bases in this way, and put
\[
 P=I+A+\cdots+A^4,\qquad T_{ij}=g_iPh_j.
 \tag{F.L1}\label{maps:eq:F.L1}
\]
The rank-one operator $P$ is a nonzero scalar multiple of
$e\varphi$.  Hence the 25 matrices $T_{ij}$ are a rational
basis of $\operatorname{Mat}_5(\mathbb Q)$.
Each belongs to $\mathcal R=\mathbb Z[A,C]$.
Let $\mathcal R_0$ be their integral span and set
\[
 \mathcal R_0^\#=
 \{X:\operatorname{tr}(XT)\in\mathbb Z
                     \text{ for all }T\in\mathcal R_0\}.
 \tag{F.L2}\label{maps:eq:F.L2}
\]
The trace pairing on the full matrix algebra is
nondegenerate, as is seen from the matrix units, so this
is another full rational lattice, computed by inverting
the 25-by-25 Gram matrix.

\item If any stable lattice exists, every element of $\mathcal R$
acts integrally on it.  Consequently
$\operatorname{tr}(XY)\in\mathbb Z$ for all
$X,Y\in\mathcal R$, and necessarily
\[
 \mathcal R_0\subset\mathcal R\subset\mathcal R_0^\#.
 \tag{F.L3}\label{maps:eq:F.L3}
\]
A nonintegral Gram entry therefore proves nonexistence.
Otherwise let
$d=[\mathcal R_0^\#:\mathcal R_0]$, the absolute Gram
determinant.  Start with the lattice generated by
$\mathcal R_0$ and $I$, and repeatedly replace a current
lattice $L$ by $L+AL+CL$, using exact Hermite reduction.
If it leaves $\mathcal R_0^\#$, \eqref{maps:eq:F.L3} proves that there
is no stable lattice.
While it remains inside, every strict enlargement reduces
its index by a factor at least two.
There can be at most $\lfloor\log_2d\rfloor$ such
enlargements starting from $\mathcal R_0$, so a bound of
$1+\lfloor\log_2d\rfloor$ tests after adjoining $I$ is
sufficient.  Stabilization gives exactly $\mathcal R$:
the stable lattice contains $I$ and all words in $A,C$,
and every element introduced belongs to their ring.
It is then an order, and
$\Lambda_0=\mathcal R\mathbb Z^5$ is an explicitly
computed stable lattice.  This proves the decision step,
including termination in the nonexistence case.

\item Use a basis of $\Lambda_0$ from now on, so
$\mathcal R\subset\operatorname{Mat}_5(\mathbb Z)$.
Smith reduction gives an integer $c\ge1$ with
\[
 c\operatorname{Mat}_5(\mathbb Z)\subset\mathcal R;
 \tag{F.L4}\label{maps:eq:F.L4}
\]
one may take the largest Smith factor of the inclusion
of the order in the full matrix lattice.
Normalize any stable lattice by a rational homothety so
that it is contained in $\mathbb Z^5$ and the gcd of all
coordinates of a lattice basis is one.
The operators $cE_{ij}$ preserve it.  The coordinate values
of its vectors generate $\mathbb Z$, so Bezout's identity,
applied to $cE_{ij}$, shows that it contains every $ce_i$.
Therefore every normalized lattice lies in the finite interval
\[
 c\mathbb Z^5\subset\Lambda\subset\mathbb Z^5.
 \tag{F.L5}\label{maps:eq:F.L5}
\]
Conversely every lattice in this interval stable under
$A,C$ is retained by the algorithm.
The only rational operators commuting with the pair are
scalars: they act by a scalar on the one-dimensional fixed
line of $A$, and then by the same scalar on the spanning
vectors $g_ie$.  Thus rational equivalence of two stable
lattices is precisely homothety.  The coordinate-gcd
normalization is unique, so \eqref{maps:eq:F.L5} gives the integral
conjugacy classification without further identifications.

\item Here is a finite enumeration bound which is also implemented.
For each $p^a\Vert c$, enumerate all subspaces of
$\mathbb F_p^5$ in reduced echelon form.
Beginning with $\mathbb Z^5$, pass through $a$ levels of
stable lattices between $pH$ and $H$, retaining only those
containing $p^a\mathbb Z^5$ and identifying equal Hermite
forms.  Every desired lattice occurs: for any target
$\Lambda$, the chain
$\Lambda+p^j\mathbb Z^5$, $0\le j\le a$, is such a
chain and ends at $\Lambda$.
If $s_p$ denotes the number of all subspaces of
$\mathbb F_p^5$, there are at most $s_p^a$ branches before
deduplication.  Intersect the resulting lattices for the
different primes and discard those of coordinate gcd
greater than one.  The Chinese remainder theorem proves
that this enumerates \eqref{maps:eq:F.L5} completely.

\item For each quotient lattice $W_{\mathbb Z}$ so obtained,
put $N_W=AB-I$.
Its image must be primitive: for a primitive full cusp,
$\operatorname{im}N_V=\ker N_V$ contains the primitive
invariant line, and quotienting that image by the line
gives a primitive image in $W_{\mathbb Z}$.
Reject a quotient whose two nonzero Smith factors are
not one.  For the remaining quotients use extension rows
$c_1,c_2\in\mathbb Z^{1\times5}$ and impose
\[
\begin{gathered}
 c_1(I+A+\cdots+A^4)=0,
 \qquad c_2(I+B+\cdots+B^{m-1})=0,\\
 (c_2+c_1B)N_W=0,
 \qquad m=\operatorname{ord}(B).
\end{gathered}
 \tag{F.L6}\label{maps:eq:F.L6}
\]
These are integral linear equations.  Their integral kernel
modulo the coboundaries
\[
 z\longmapsto\bigl(z(A-I),z(B-I)\bigr),
 \qquad z\in\mathbb Z^{1\times5},
 \tag{F.L7}\label{maps:eq:F.L7}
\]
is a finitely generated group $G$ of rank one by \ref{maps:src:F.2.11}.
Compute its free and torsion factors by Smith reduction.
Choose an integral generator $w$ of the free rank-one
quotient $\ker N_W/\operatorname{im}N_W$.
The homomorphism
\[
 \kappa:G\longrightarrow\mathbb Z,
 \qquad [(c_1,c_2)]\longmapsto(c_2+c_1B)w
 \tag{F.L8}\label{maps:eq:F.L8}
\]
is well defined and nonzero: \eqref{maps:eq:F.L6} kills the image, and
\eqref{maps:eq:F.L7} kills $w$; rational nonvanishing is \ref{maps:src:F.2.11}.
The full cusp has unit Smith factors exactly when
$\kappa=\pm1$.
Indeed the quotient rank-two image is already primitive,
and the remaining invariant-line image is generated by
that one integer.
If $\operatorname{im}\kappa\ne\mathbb Z$, reject the
lattice.  Otherwise the inverse image of $1$ is a finite
coset of $\operatorname{Tor}G$, and lists all extensions.
Changing the sign of the invariant-line generator identifies
$\kappa=-1$ with $\kappa=1$.
There are no further identifications: an integral isomorphism
on the quotient is $\pm I$, on the invariant line is
$\pm1$, and, after $\kappa=1$ is fixed, their ratio is one.
The remaining row is exactly an integral change of splitting
\eqref{maps:eq:F.L7}.
This proves completeness and termination of the extension
step as well.

\item The order, trace-dual bound, conductor and finite-interval
steps are implemented in
\texttt{fourfolds\_lattice\_algorithm.py};
the extension-kernel and Smith steps are implemented in
\texttt{fourfolds\_two\_classes\_lattices.py} with the
finite order as an argument.  Their applications and exact
exhaustive bounds are recorded next.

\end{enumerate}
\end{proof}




\begin{theorem}[complete lattice classes for the two family survivors]\label{maps:src:F.2.23}
\StatementRecord{maps:src:F.2.23}{proved}{Historical events: Claude Fable 5.2 (Anthropic), effort high, 16 November 2026: carried as background without fault found. GPT-6 Astra (OpenAI), ChatGPT Pro mode, effort Pro, S6\_09\_PAPERS, 17 November 2026: targeted proof check. Also checked in this writing. Current writing events: GPT-6 Astra (OpenAI), ChatGPT Pro mode, effort Pro, S6\_10\_PAPERS, 18 November 2026: computation or formal certificate; targeted proof check. GPT-6 Astra (OpenAI), ChatGPT Pro mode, effort Pro, S6\_10\_PAPERS continuation, 19 November 2026: used and rechecked in this run. \emph{Recorded proof scope.} Claude Fable 5.2 (Anthropic), effort high, 16 November 2026; also checked in this writing by GPT-6 Astra (OpenAI), ChatGPT Pro mode, effort Pro, 18 November 2026}{the hypotheses, notation and definitions stated here; \cref{maps:src:F.2.16}; \cref{maps:src:F.2.6}; \cref{maps:src:F.2.18}; \cref{maps:eq:F.U1}; \cref{maps:eq:F.L4}; \cref{maps:eq:F.L5}; \cref{maps:eq:F.L6}; \cref{maps:eq:F.L7}; \cref{maps:src:F.2.12}; \cref{maps:eq:F.L10}; \cref{maps:eq:F.U6}; \cref{maps:src:F.2.19}}


For each of the two rational representations $v=0$ and $v=1$
of \ref{maps:src:F.2.16}, there are exactly two homothety classes of stable
rank-five quotient lattices, and exactly three integral
conjugacy classes of rank-six extensions with primitive
maximal cusp.  Exactly one of those three can have simply
connected finite affine fillings.  It is the integral class
already realized in \ref{maps:src:F.2.6}, respectively \ref{maps:src:F.2.18}.
Thus arbitrary stable-lattice changes produce no further
linear survivor classes anywhere in the family \ref{maps:src:F.2.16}.
This assertion retains the classification of all affine
twists and multiplicities as a separate problem.

\end{theorem}


\begin{proof}
\begin{enumerate}[label=\arabic*.]
\item Apply the preceding algorithm in the original five-cycle
lattice \eqref{maps:eq:F.U1}.  For both parameters the order is obtained
after four strict enlargements.  Its Smith factors in the
full matrix lattice are 21 factors one and four factors two;
thus the conductor $c=2$ in \eqref{maps:eq:F.L4} is exact.
All 374 subspaces of $\mathbb F_2^5$, including zero and
the whole space, are enumerated, with no height restriction
or neighbor-depth assumption.
After imposing stability and coordinate gcd one, the two
lattice representatives are
\[
 \mathbb Z^5,\qquad
 \Lambda_* =2\mathbb Z^5+\mathbb Z(1,1,1,1,1)^t.
 \tag{F.L9}\label{maps:eq:F.L9}
\]
The latter has index $16$.
The generic prime-by-prime implementation of \eqref{maps:eq:F.L5} and
an independent bit-vector enumeration give exactly the
same two Hermite forms.

\item For $\mathbb Z^5$, the extension group in \eqref{maps:eq:F.L6}--\eqref{maps:eq:F.L7}
is $\mathbb Z$ and $\kappa$ is a unit on its generator.
For $\Lambda_*$, that group is
$\mathbb Z\oplus\mathbb Z/2$, again with primitive
$\kappa$.  Hence there are exactly one and two primitive
extensions, respectively.
For each extension the program computes the image of the
primitive global invariant homology covector $\ell$ on
the two finite fixed lattices.  The resulting positive
divisibilities are
\[
\begin{array}{c|c|c}
 \text{quotient lattice}&\text{extension class}&
       (\operatorname{div}_1,\operatorname{div}_2)\\\hline
 \mathbb Z^5&\text{unique}&(1,2)\\
 \Lambda_*&\text{first torsion coset}&(1,1)\\
 \Lambda_*&\text{second torsion coset}&(1,2).
\end{array}
 \tag{F.L10}\label{maps:eq:F.L10}
\]
The table is identical for $v=0$ and $v=1$.
All six full matrices, the integral kernel bases, the
Smith factors and the finite divisibilities are computed
exactly by \texttt{fourfolds\_two\_classes\_lattices.py}
and saved in its output and JSON certificate.

\item For a simply connected filling the determinant quotient
of \ref{maps:src:F.2.12} requires $\gcd(\ell(v_i),m_i)=1$ at each
finite cap.  The second affine multiplicity is divisible
by four.  A divisibility of two in \eqref{maps:eq:F.L10} therefore
excludes that class for every choice of twists and larger
multiplicities.
The sole remaining class has an explicit integral
intertwiner of determinant $\pm1$ with the corresponding
pair in \eqref{maps:eq:F.U6}, as printed and checked by the program.
Those pairs already have the verified simply connected,
three-connected fillings.  \ref{maps:src:F.2.19} excludes every other
rational parameter for all lattices.
Together these facts prove all the claims, including the
absence of any additional linear survivor class in the
unbounded family.

\end{enumerate}
\end{proof}




\begin{theorem}[every surviving twist in the unbounded family]\label{maps:src:F.2.24}
\StatementRecord{maps:src:F.2.24}{proved}{Historical events: Claude Fable 5.2 (Anthropic), effort high, 16 November 2026: carried as background without fault found. GPT-6 Astra (OpenAI), ChatGPT Pro mode, effort Pro, S6\_09\_PAPERS, 17 November 2026: targeted proof check. Also checked in this writing. Current writing events: GPT-6 Astra (OpenAI), ChatGPT Pro mode, effort Pro, S6\_10\_PAPERS, 18 November 2026: computation or formal certificate; targeted proof check. GPT-6 Astra (OpenAI), ChatGPT Pro mode, effort Pro, S6\_10\_PAPERS continuation, 19 November 2026: used and rechecked in this run. \emph{Recorded proof scope.} Claude Fable 5.2 (Anthropic), effort high, 16 November 2026; also checked in this writing by GPT-6 Astra (OpenAI), ChatGPT Pro mode, effort Pro, 18 November 2026}{the hypotheses, notation and definitions stated here; \cref{maps:src:F.2.12}; \cref{maps:src:F.2.16}; \cref{maps:src:F.2.23}; \cref{maps:src:F.2.20}; \cref{maps:src:F.2.19}; \cref{maps:eq:F.U6}; \cref{maps:src:F.2.5}; \cref{maps:src:F.2.18}; \cref{maps:eq:F.T4}; \cref{maps:src:F.2.6}; \cref{maps:eq:F.T6}; \cref{maps:eq:F.T7}; \cref{maps:eq:F.T5}; \cref{maps:eq:F.T8}; \cref{maps:eq:F.T3}; \cref{maps:eq:F.T2}; \cref{maps:eq:F.T9}}


Use the fixed Kuhn cusp and the affine translation clutching
convention of \ref{maps:src:F.2.12}.  In the entire rational family \ref{maps:src:F.2.16},
allow every stable lattice, every free affine finite lift,
and every pair of actual finite multiplicities.
The data with $\pi_1=1$ and $H_2=H_3=0$ are precisely the
following, up to the integral linear classification already
proved and the translation identifications described below:
\[
\begin{gathered}
 v\in\{0,1\},\qquad
 \text{the unique eligible integral class of \ref{maps:src:F.2.23}},\\
 (m_1,m_2)=(5,4),\qquad
 20\kappa+4r_1+5r_2=\pm1.
\end{gathered}
 \tag{F.T1}\label{maps:eq:F.T1}
\]
Here $r_i=\ell(v_i)$ is the invariant affine numerator, and
$\kappa\in\mathbb Z$ is the invariant coordinate of the
cusp translation clutching, with signs fixed by the displayed
van Kampen relation below.  The remaining fixed-direction
components of the twists impose no condition.
Every manifold in \eqref{maps:eq:F.T1} is oriented diffeomorphic to its
corresponding known survivor and hence to
$3\mathbb{HP}^2\#(-\mathbb{HP}^2)$.
This completes the survivor classification for this entire
unbounded rational family; it does not classify the other
rational points of \ref{maps:src:F.2.20} or the other spectral loci.

\end{theorem}


\begin{proof}
\begin{enumerate}[label=\arabic*.]
\item \ref{maps:src:F.2.19} excludes all other rational parameters, and \ref{maps:src:F.2.23}
excludes all other integral linear classes.
We therefore work with either known integral pair, in the
marking of \eqref{maps:eq:F.U6}.  Its primitive invariant homology
covector is $\ell(x)=x_0$.
Each finite fixed lattice has an integral basis
$u_i,\lambda_i$, where $\ell(u_i)=1$,
$\ell(\lambda_i)=0$, and $\lambda_i$ is primitive in
$\ker\ell$.
We may write the affine numerator as
$v_i=r_i u_i+s_i\lambda_i$.
This covers arbitrary finite affine lifts: translating the
origin makes their translation vector invariant by averaging,
and its $m_i$th iterate is then the integral vector $v_i$.
The resulting boundary translation change is included in the
specified affine clutching.

\item For both linear pairs, the normal closure of the cusp angular
lattice in the fibre lattice is exactly $\ker\ell$, by
\ref{maps:src:F.2.5} and \ref{maps:src:F.2.18}.  Van Kampen therefore leaves the abelian
presentation
\[
 m_1a=r_1h,\qquad m_2b=r_2h,\qquad
 a+b+\kappa h=0.
 \tag{F.T2}\label{maps:eq:F.T2}
\]
This also defines the sign of $\kappa$.
Its determinant is, up to sign,
\[
 p=m_1m_2\kappa+m_2r_1+m_1r_2.
 \tag{F.T3}\label{maps:eq:F.T3}
\]
Thus simple connectivity is exactly $p=\pm1$.
It forces $\gcd(m_1,m_2)=1$ and
$\gcd(r_i,m_i)=1$ by reduction modulo each multiplicity.

\item We first prove that a surviving datum has the minimal
multiplicities.  Set $n_1=5$, $n_2=4$ and $m_i=n_i k_i$.
In the splitting $\Lambda=\mathbb Zu_i\oplus\ker\ell$,
write the finite linear map as $1\oplus D_i$.
The reduced finite fibre is a mapping torus of $T^5$ over
the circle $m_i\ell$; its fibre is connected because
$r_i$ is prime to $m_i$.
Its vertical monodromy is a power of $D_i$ prime to $n_i$,
possibly followed by a translation, which changes no
integral Wang group.

\item The primitive fixed vertical loop $\lambda_i$ gives the
only additional free two-cycle beyond vertical cycles.
Consequently an invariant integral two-form $\omega$ on
the covering torus descends integrally exactly when
\[
 \frac{\omega(v_i,\lambda_i)}{m_i}
       =\frac{r_i\omega(u_i,\lambda_i)}{m_i}
       \quad\text{is an integer}.
 \tag{F.T4}\label{maps:eq:F.T4}
\]
For precision, the torus made from the affine generator
and $\lambda_i$ projects with degree $r_i$ to the mapping-
torus base, whereas the ordinary $u_i$ suspension projects
with degree $m_i$.  A Bezout combination is a primitive
suspension.  Its integrality condition is equivalent to
\eqref{maps:eq:F.T4}, because the ordinary covering periods are already
integral and $\gcd(r_i,m_i)=1$.
The Wang sequence shows that these vertical cycles and one
primitive suspension generate the free second homology;
there is no omitted period condition.

\item The integral invariant two-form computations in \ref{maps:src:F.2.6} and
\ref{maps:src:F.2.18} give
\[
 \gcd_{\omega}\omega(u_i,\lambda_i)=n_i.
 \tag{F.T5}\label{maps:eq:F.T5}
\]
They are also independently printed for both pairs by
\texttt{fourfolds\_family\_twists.py}.
Hence the covering-image cokernel in degree two at the
$i$th finite point is exactly $\mathbb Z/k_i$.
The global invariant group of $\bigwedge^2\mathcal V$
and its integral first parabolic group both vanish.
The skyscraper cokernel sequence therefore gives
\[
 H^1(B,R^2g_*\mathbb Z)
       =\mathbb Z/k_1\oplus\mathbb Z/k_2.
 \tag{F.T6}\label{maps:eq:F.T6}
\]
Point-supported torsion kernels do not change this degree-one
base cohomology, and the cusp specialization is onto.
The term \eqref{maps:eq:F.T6}, in bidegree $(1,2)$, survives every Leray
differential on the two-dimensional base.
All rational terms of total degree three vanish, independently
of the multiplicities and twists, by the same exterior
local-system table as \ref{maps:src:F.2.6}.  Thus $H^3$ is finite.
If either $k_i>1$, its surviving nonzero filtration quotient
makes $H^3$ nonzero; universal coefficients then give
nonzero torsion in $H_2$.
This proves $m_1=5$, $m_2=4$ for every survivor.

\item We now identify all twists with these multiplicities.
The following integral fact permits removal of each $s_i$:
\[
 \operatorname{Norm}_{n_i}:
 \ker\ell\longrightarrow\mathbb Z\lambda_i
 \quad\text{is onto}.
 \tag{F.T7}\label{maps:eq:F.T7}
\]
Here the norm is $I+D_i+\cdots+D_i^{n_i-1}$.
One can read \eqref{maps:eq:F.T7} directly from \eqref{maps:eq:F.T5}.
Let $\varphi_i$ generate the primitive invariant covector
lattice on $\ker\ell$.
The splitting by $u_i$ shows that the gcd in \eqref{maps:eq:F.T5} is
$|\varphi_i(\lambda_i)|$: its forms are
$\ell\wedge\varphi_i$ plus invariant vertical forms.
Thus that value is $n_i$.
The norm has the form
$x\mapsto d_i\varphi_i(x)\lambda_i$ for an integer
$d_i$; evaluating $\varphi_i$ gives
$|d_i|n_i=n_i$, hence $|d_i|=1$.
This proves surjectivity.  The saved program additionally
gives integral preimages $q_i\in\ker\ell$ with
$\operatorname{Norm}_{n_i}q_i=\lambda_i$.

\item Put
\[
 w_i=-\frac{s_i}{n_i}\sum_{j=0}^{n_i-1}jM_i^j q_i.
\]
Then
\[
 (I-M_i)w_i=-s_i\lambda_i/n_i+s_iq_i.
 \tag{F.T8}\label{maps:eq:F.T8}
\]
Translation by $w_i$ conjugates the finite affine action
with numerator $r_i u_i+s_i\lambda_i$ to that with
numerator $r_i u_i$: the remaining $s_iq_i$ is integral.
It gives a diffeomorphism of the entire finite cap.
In the standard linear boundary marking, its translation
section is
$w_i+\theta s_i\lambda_i/n_i$.
Its endpoint mismatch with the linear deck map is exactly
$s_iq_i$ by \eqref{maps:eq:F.T8}, so its integral boundary residue has
$\ell$-coordinate zero.

\item We explain why these local changes extend simultaneously,
including the prescribed cusp marking.
Let $P$ be the complementary pair of pants and
$\mathcal T\to P$ its flat torus group bundle.
A section on the boundary extends over the relative
one-skeleton because the torus is connected.
The only remaining obstruction is the two-cell obstruction
with values in the fibre fundamental-group local system:
\[
 H^2(P,\partial P;\Lambda)
        \simeq H_0(P;\Lambda)=\Lambda_\Gamma.
 \tag{F.T9}\label{maps:eq:F.T9}
\]
The equality follows from the oriented relative cellular
complex, equivalently integral Poincare--Lefschetz duality
with local coefficients.  Its boundary residues map to
the indicated global coinvariants.
For both linear pairs the horizontal monodromy relation
matrix has five unit Smith factors, so
$\Lambda_\Gamma=\mathbb Z$, detected exactly by $\ell$.
This is independently checked in the saved program.
Thus the residues $s_iq_i$, with the identity section at the
cusp, have zero obstruction and extend over $P$.
A section may be smoothed relative to its boundary collars;
translation by it is a fibre-preserving diffeomorphism.
Together with the finite-cap maps it removes all $s_i$
without changing the cusp.
The same argument removes all components of translation
clutching whose $\ell$-coordinate is zero.

\item It remains to handle the integer numerators and the remaining
cusp coordinate.  If $p=-1$, the congruences from \eqref{maps:eq:F.T3}
give
\[
 r_1=1+5a,\qquad r_2=-1+4b,
 \qquad a+b+\kappa=0.
 \tag{F.T10}\label{maps:eq:F.T10}
\]
Adding $5a u_1$ or $4b u_2$ changes no finite affine action;
in its linear boundary marking it changes the translation
residue by $a u_1$ or $b u_2$.
The difference of the cusp marking has residue of
$\ell$-coordinate $\kappa$, by \eqref{maps:eq:F.T2}.
Their total obstruction in \eqref{maps:eq:F.T9} is therefore
$a+b+\kappa=0$.
The same extension argument identifies the glued manifold
with the original twists $(u_1,-u_2)$ and zero cusp
clutching.
If $p=1$, the identical calculation starts from
$(-u_1,u_2)$.

\item Finally simultaneous fibre inversion identifies these two
sign choices by an orientation-preserving diffeomorphism.
It conjugates every affine translation to its negative.
It extends over the Kuhn cusp: the lattice involution
$(x,h)\mapsto(-x,h)$ preserves its toric fan, since a
simplex given by a chain in a unit cube is taken to the
chain with reversed coordinate order in the opposite cube.
It normalizes the cusp lattice action and therefore descends
to its quotient, while fixing the base parameter.  The resulting toric map is holomorphic,
and in the regular real six-torus inversion has determinant
$(-1)^6=1$.  All translation maps also preserve orientation.
Thus every datum in \eqref{maps:eq:F.T1} gives the stated oriented smooth
type and has $H_2=H_3=0$, proving sufficiency as well as
necessity.

\end{enumerate}
\end{proof}




\begin{theorem}[the universal degree-two suspension obstruction]\label{maps:src:F.2.25}
\StatementRecord{maps:src:F.2.25}{proved}{Historical events: Claude Fable 5.2 (Anthropic), effort high, 16 November 2026: carried as background without fault found. GPT-6 Astra (OpenAI), ChatGPT Pro mode, effort Pro, S6\_09\_PAPERS, 17 November 2026: targeted proof check. Also checked in this writing. Current writing events: GPT-6 Astra (OpenAI), ChatGPT Pro mode, effort Pro, S6\_10\_PAPERS, 18 November 2026: computation or formal certificate; targeted proof check. GPT-6 Astra (OpenAI), ChatGPT Pro mode, effort Pro, S6\_10\_PAPERS continuation, 19 November 2026: used and rechecked in this run. \emph{Recorded proof scope.} Claude Fable 5.2 (Anthropic), effort high, 16 November 2026; also checked in this writing by GPT-6 Astra (OpenAI), ChatGPT Pro mode, effort Pro, 18 November 2026}{the hypotheses, notation and definitions stated here; \cref{maps:src:F.2.14}; \cref{maps:src:F.2.7}; \cref{maps:src:F.2.9}; \cref{maps:src:F.2.17}; \cref{maps:src:F.2.12}; \cref{maps:eq:F.T4}; \cref{maps:eq:F.S3}; \cref{maps:eq:F.S6}; \ExternalPaperRef{theorem~}{6.3}{III}; \cref{maps:eq:F.S2}; \cref{maps:eq:F.S7}; \cref{maps:eq:F.S5}; \cref{maps:eq:F.S8}}


Every datum of F.2 with $\pi_1=1$ and $H_2=H_3=0$ has
exactly two genuine finite fillings and has minimal affine
multiplicities.  After ordering the order-five point first,
its finite characteristic polynomials are one of precisely
the following two remaining necessary possibilities:
\[
\begin{array}{c|c|c}
 (m_1,m_2)&\chi_{M_1}&\chi_{M_2}\\\hline
 (5,4)&(x-1)^2\Phi_5&(x-1)^2\Phi_2^2\Phi_4\\
 (5,6)&(x-1)^2\Phi_5&(x-1)^2\Phi_2^2\Phi_3.
\end{array}
 \tag{F.S1}\label{maps:eq:F.S1}
\]
In particular all order-eight and order-twelve spectral
loci, and the other two order-six spectral loci of \ref{maps:src:F.2.14},
are excluded for every rational relative position, every
stable lattice and every free affine twist.
There is no multiplicity search bound in this statement:
the only possible actual affine multiplicities are the
two displayed pairs.

\end{theorem}


\begin{proof}
\begin{enumerate}[label=\arabic*.]
\item Theorem \ref{maps:src:F.2.7} excludes three genuine finite fillings.
For two fillings, \ref{maps:src:F.2.9} and \ref{maps:src:F.2.14} give a primitive global
invariant homology covector $\ell$, finite fixed rank two,
and an irreducible rank-five cohomology quotient.
In particular
\[
 H^0(B,j_*\bigwedge^2\mathcal V)=0.
 \tag{F.S2}\label{maps:eq:F.S2}
\]
This is also immediate from the exterior filtration and
irreducibility in odd dimension, as in \ref{maps:src:F.2.17}.
Simple connectivity gives the determinant quotient of
\ref{maps:src:F.2.12}, so each invariant affine numerator $r_i=\ell(v_i)$
is prime to the actual multiplicity $m_i$, and the two
multiplicities are coprime.
As before one may conjugate any finite affine lift by a
translation so that its numerator $v_i$ is integral and
fixed by its linear part.

\item Consider one finite point.  Let $M$ have linear order $n$,
so $n\mid m$, and let its complementary degree-four factor
be $f(x)$, with
$\chi_M=(x-1)^2f(x)$.
The image of $\ell$ on the finite fixed lattice is $g\mathbb Z$
for some $g\mid n$: averaging $I+M+\cdots+M^{n-1}$ on
any vector of $\ell$-value one proves this divisibility.
Since $g\mid r$ and $\gcd(r,m)=1$, we have $g=1$.
Choose a fixed integral vector $u$ with $\ell(u)=1$.
It gives an integral invariant splitting
\[
 \Lambda=\mathbb Zu\oplus H,
 \qquad H=\ker\ell,\qquad M=1\oplus D,
 \qquad\chi_D=(x-1)f(x).
 \tag{F.S3}\label{maps:eq:F.S3}
\]
The order of $D$ is $n$.
Let $\lambda$ generate its primitive fixed lattice and
let $\varphi$ generate its primitive invariant covector
lattice.  Choose their signs so that
\[
 h=\varphi(\lambda)>0.
 \tag{F.S4}\label{maps:eq:F.S4}
\]
Their rational pairing is nonzero, since the finite map
is semisimple and its averaging projection has rank one.

\item Both $\operatorname{Norm}_n(D)$ and $f(D)$ are integral
operators, have image on the line $\mathbb Q\lambda$,
and have invariant defining covector proportional to
$\varphi$.  Primitivity of $\lambda$ and $\varphi$
therefore implies
\[
 \operatorname{Norm}_n(D)=\frac n h\lambda\varphi,
 \qquad
 f(D)=\frac{f(1)}h\lambda\varphi,
 \qquad h\mid n,\quad h\mid f(1).
 \tag{F.S5}\label{maps:eq:F.S5}
\]
For example, the integer coefficient in the second formula
is seen by applying the integral operator to a vector on
which $\varphi$ is one; evaluating on $\lambda$ determines
that coefficient.  No lattice splitting over a localization
or semisimplification modulo a prime is used here.

\item Write $v=ru+s\lambda$.  The quotient finite fibre is a
connected $T^5$ mapping torus over the circle $m\ell$,
because $r$ is prime to $m$.
Its vertical cycles already have integral periods for an
invariant integral two-form on the covering torus.
The Wang sequence leaves one additional primitive free
suspension cycle, from $\lambda$.
Exactly as in \eqref{maps:eq:F.T4}, descent of an invariant form $\omega$
is equivalent to
$m\mid r\omega(u,\lambda)$, and hence to
$m\mid\omega(u,\lambda)$.
In the integral splitting \eqref{maps:eq:F.S3}, every invariant two-form
is an integral multiple of $\ell\wedge\varphi$ plus an
invariant vertical two-form.  Its possible values on
$(u,\lambda)$ are therefore exactly $h\mathbb Z$.
It follows that the degree-two specialization image at
this point has cokernel
\[
 C^2_p\simeq\mathbb Z/(m/h).
 \tag{F.S6}\label{maps:eq:F.S6}
\]
The integer $m/h$ is defined because $h\mid n\mid m$.
Torsion in the finite fibre cohomology is a specialization
kernel, not an additional restriction on these free periods.
The primitive suspension and vertical cycles generate the
free second homology, so \eqref{maps:eq:F.S6} is the full cokernel.

\item Let $I^2$ be the image sheaf of
$R^2g_*\mathbb Z\to j_*\bigwedge^2\mathcal V$.
The cusp specialization is onto by \ExternalPaperRef{}{6.3}{III}, and all kernels
and cokernels at finite points are skyscraper sheaves.
Thus $H^1(R^2g_*\mathbb Z)=H^1(I^2)$.
Equation \eqref{maps:eq:F.S2} and the cokernel exact sequence give an
injection
\[
 \bigoplus_{p\text{ finite}}C^2_p
           \lhook\joinrel\longrightarrow H^1(B,R^2g_*\mathbb Z).
 \tag{F.S7}\label{maps:eq:F.S7}
\]
The Leray term of bidegree $(1,2)$ has no incoming or
outgoing differential on a real two-dimensional base.
If $H_2=H_3=0$, the universal coefficient theorem gives
$H^3=0$, so that term must vanish.
Every cyclic cokernel in \eqref{maps:eq:F.S7} is therefore zero.
Combining \eqref{maps:eq:F.S5}--\eqref{maps:eq:F.S6} yields the universal necessary
condition
\[
 m=h=n,\qquad n\mid f(1).
 \tag{F.S8}\label{maps:eq:F.S8}
\]
This both fixes the multiplicity and tests the spectrum.

\item For completeness, apply \eqref{maps:eq:F.S8} to the entire eight-case
list of \ref{maps:src:F.2.14}.  The order-five factor gives $\Phi_5(1)=5$.
At the other point the exact values are
\[
\begin{array}{c|c|r|r}
 n&f&f(1)&n/\gcd(n,f(1))\\\hline
 4&\Phi_2^2\Phi_4&8&1\\
 6&\Phi_2^2\Phi_3&12&1\\
 6&\Phi_2^2\Phi_6&4&3\\
 6&\Phi_3\Phi_6&3&2\\
 8&\Phi_8&2&4\\
 12&\Phi_{12}&1&12\\
 12&\Phi_3\Phi_4&6&2\\
 12&\Phi_4\Phi_6&2&6.
\end{array}
 \tag{F.S9}\label{maps:eq:F.S9}
\]
The last column is a lower bound for the nonzero local
covering-cokernel order already at the smallest allowed
multiplicity.  Only the first two rows pass.
The complete table is checked by
\texttt{fourfolds\_suspension\_filter.py}; its input is
exactly these eight spectra, with no matrix-entry or
relative-position cutoff.  This proves \eqref{maps:eq:F.S1} and all
claimed universal exclusions.

\end{enumerate}
\end{proof}




\begin{theorem}[the complete rational remaining $(5,6)$ locus]\label{maps:src:F.2.26}
\StatementRecord{maps:src:F.2.26}{proved}{Historical events: Claude Fable 5.2 (Anthropic), effort high, 16 November 2026: carried as background without fault found. GPT-6 Astra (OpenAI), ChatGPT Pro mode, effort Pro, S6\_09\_PAPERS, 17 November 2026: targeted proof check. Also checked in this writing. Current writing events: GPT-6 Astra (OpenAI), ChatGPT Pro mode, effort Pro, S6\_10\_PAPERS, 18 November 2026: computation or formal certificate; targeted proof check. GPT-6 Astra (OpenAI), ChatGPT Pro mode, effort Pro, S6\_10\_PAPERS continuation, 19 November 2026: used and rechecked in this run. \emph{Recorded proof scope.} Claude Fable 5.2 (Anthropic), effort high, 16 November 2026; also checked in this writing by GPT-6 Astra (OpenAI), ChatGPT Pro mode, effort Pro, 18 November 2026}{the hypotheses, notation and definitions stated here; \cref{maps:eq:F.S1}; \cref{maps:eq:F.R2}; \cref{maps:src:F.2.25}; \cref{maps:eq:F.R3}; \cref{maps:eq:F.R5}; \cref{maps:src:F.2.20}; \cref{maps:eq:F.R6}; \cref{maps:eq:F.R7}; \cref{maps:eq:F.R8}; \cref{maps:eq:F.R12}; \cref{maps:src:F.2.22}}


The irreducible rational rank-five pairs in the second row
of \eqref{maps:eq:F.S1}, with square-zero rank-two cusp, are classified
up to simultaneous rational conjugacy by
\[
 \{[r:s:t]\in\mathbb P^2(\mathbb Q):
             D(r,s,t)\ne0,\ Q_-(r,s,t)\ne0\}
       \ \sqcup\ \mathbb A^1(\mathbb Q),
 \qquad Q_-=r^2-rs+st,
 \tag{F.R10}\label{maps:eq:F.R10}
\]
where $D$ is the same cubic as in \eqref{maps:eq:F.R2}.
The exceptional line again lies over $[4:1:3]$.
Thus the two rational spectral loci left by \ref{maps:src:F.2.25} both
have complete explicit normal forms, including their
entire rational exceptional loci.

\end{theorem}


\begin{proof}
\begin{enumerate}[label=\arabic*.]
\item Use precisely the space $\mathbb Q\oplus\mathbb Q(\zeta_5)$,
matrices $A,R,X$ and projective plane normalization in
\eqref{maps:eq:F.R3}--\eqref{maps:eq:F.R5}.  The centralizer and image-plane arguments
in \ref{maps:src:F.2.20} require only irreducibility and the two-dimensional
$-1$ eigenspace of $C$; both hold here.
Consequently every pair has a unique projective image-plane
parameter, with only scalar stabilizer.
The unknown $Y$ now satisfies
\[
 YX=0,\qquad YRX=I_2,\qquad
 \operatorname{tr}(YAX)=-2,\qquad
 \operatorname{tr}(YA^{-1}X)=2.
 \tag{F.R11}\label{maps:eq:F.R11}
\]
Its coefficient matrix is unchanged from \eqref{maps:eq:F.R6}, so has
determinant $D/4$.
The two traces and the $-1$ eigenspace imply
\[
 \chi_C=(x-1)(x+1)^2(x^2+x+1).
\]
Its other three roots are distinct, and the $-1$ eigenspace
has its full multiplicity, so $C^6=I$.
The cusp $XY$ still has square zero and rank two.
For $D\ne0$, the ten linear equations give the unique
explicit solution.

\item The two first-column adjugate identities become
\[
 D y_{10}=2Q_-Q_2,
 \qquad D y_{20}=2(r-s-t)Q_-,
 \tag{F.R12}\label{maps:eq:F.R12}
\]
with the same $Q_2$ as in \eqref{maps:eq:F.R7}.
For $D\ne0$, these entries vanish simultaneously exactly
when $Q_-=0$.
Indeed, if $Q_-\ne0$ they would instead force
$r=s+t$ and $Q_2=0$, and \eqref{maps:eq:F.R8} would give $D=0$.
The invariant-covector argument of \ref{maps:src:F.2.20} is unchanged.
Thus the pair is rationally irreducible exactly off the
conic $Q_-=0$ in the generic locus.

\item No nonzero rational triple satisfies $D=Q_-=0$.
If $s=0$, the conic forces $r=0$, and then $D=7t^3$.
If $s\ne0$, rescale to $s=1$ and put
$r=u$, $t=u-u^2$.  Direct substitution gives
\[
 D(u,1,u-u^2)=
 -(7u^2-11u+7)(u^4-u^3+u^2-u+1).
 \tag{F.R13}\label{maps:eq:F.R13}
\]
The first factor has discriminant $-75$ and positive
leading coefficient.  The second is strictly positive,
as the identity
\[
 u^4-u^3+u^2-u+1
 =\left(u^2-\frac u2\right)^2
  +\frac34\left(u-\frac23\right)^2+\frac23
\]
shows.  In particular neither has a rational zero.
On the exceptional cubic the necessary identities \eqref{maps:eq:F.R12}
therefore force the same two points $[3:2:1]$ and $[4:1:3]$
as before.  The first is again impossible because
$[X\mid RX]$ has rank less than four.
At the second, the complete solution is
\[
 Y_z=\begin{pmatrix}
 4z+26/5&0&-2z-2&-2z-1/5&2z+1\\
 2z-2/5&0&-z&2/5-z&z
 \end{pmatrix},\qquad z\in\mathbb Q.
 \tag{F.R14}\label{maps:eq:F.R14}
\]
Its first-column entries cannot vanish simultaneously,
so each pair on this line is irreducible.
There are no further identifications, by the scalar
stabilizer argument already proved in \ref{maps:src:F.2.20}.
This establishes the claimed bijection.

\item The exact program is
\texttt{fourfolds\_relative\_six\_probe.py}.
It checks the modified characteristic polynomials, finite-order
identities, first-column adjugate expressions, conic
substitution and entire exceptional line symbolically.
Its auxiliary trace-parameter computations for the other
repeated-negative spectrum are recorded as well, although
that spectrum is now excluded integrally by \ref{maps:src:F.2.25}.
No finite sample is used to infer a rational classification.
The finite lattice and extension algorithm of \ref{maps:src:F.2.22} applies
to every point of \eqref{maps:eq:F.R10}; it does not by itself decide
which rational parameters survive the remaining integral
homology tests.

\end{enumerate}
\end{proof}




\begin{theorem}[the entire remaining order-six locus is excluded]\label{maps:src:F.2.27}
\StatementRecord{maps:src:F.2.27}{proved}{Historical events: Claude Fable 5.2 (Anthropic), effort high, 16 November 2026: carried as background without fault found. GPT-6 Astra (OpenAI), ChatGPT Pro mode, effort Pro, S6\_09\_PAPERS, 17 November 2026: targeted proof check. Also checked in this writing. Current writing events: GPT-6 Astra (OpenAI), ChatGPT Pro mode, effort Pro, S6\_10\_PAPERS, 18 November 2026: computation or formal certificate; targeted proof check. GPT-6 Astra (OpenAI), ChatGPT Pro mode, effort Pro, S6\_10\_PAPERS continuation, 19 November 2026: used and rechecked in this run. \emph{Recorded proof scope.} Claude Fable 5.2 (Anthropic), effort high, 16 November 2026; also checked in this writing by GPT-6 Astra (OpenAI), ChatGPT Pro mode, effort Pro, 18 November 2026}{the hypotheses, notation and definitions stated here; \cref{maps:eq:F.S1}; \cref{maps:src:F.2.25}; \cref{maps:eq:F.S3}; \cref{maps:eq:F.S5}; \cref{maps:eq:F.X2}; \cref{maps:eq:F.X3}; \cref{maps:eq:F.S7}; \cref{maps:eq:F.X5}; \ExternalPaperRef{equation~}{{E.9}}{III}}


Every simply connected datum in the second row of \eqref{maps:eq:F.S1},
at every rational relative position and every stable integral
lattice with primitive maximal cusp, has
\[
                 \operatorname{Tor}_3 H_3(M;\mathbb Z)\ne0
\]
if its degree-two suspension cokernels vanish.
Consequently no such datum is three-connected.  Together with
\ref{maps:src:F.2.25} this leaves only the minimal order pair $(5,4)$ as a
possible three-connected datum in the full problem F.2.
\end{theorem}


\begin{proof}
\begin{enumerate}[label=\arabic*.]
\item We use cohomology with coefficients in the discrete valuation
ring $R=\mathbb Z_3$.  Suppose the degree-two suspension
cokernels vanish; otherwise \ref{maps:src:F.2.25} already gives nonzero
$H_2$ or $H_3$.  That theorem forces $m=n=h=6$ at the
order-six point, in the notation \eqref{maps:eq:F.S3}--\eqref{maps:eq:F.S5}.
We first determine its local integral representation over $R$.

\item The splitting by the fixed vector $u$ is integral:
$\Lambda_R=Ru\oplus H_R$.  On $H_R$, the finite map $D$
has characteristic polynomial
$(x-1)(x+1)^2\Phi_3(x)$.  Since two is a unit in $R$,
the idempotents $(I\pm D^3)/2$ give an integral direct sum
\[
 H_R=L_+\oplus L_-,\qquad
 \operatorname{rank}L_+=3,\quad
 \operatorname{rank}L_-=2.
\]
On $L_-$ the map is exactly $-I$, by its rational spectrum.
On $L_+$ it has order three.  The primitive fixed vector
$\lambda$ lies in $L_+$, and the primitive invariant
covector $\varphi$ vanishes on $L_-$, so their restrictions
remain primitive.  Equation \eqref{maps:eq:F.S5} with $h=6$ gives
\[
 2(I+D+D^2)|_{L_+}=\lambda\varphi|_{L_+}.
 \tag{F.X1}\label{maps:eq:F.X1}
\]
This operator is nonzero modulo three.  But modulo three
$I+D+D^2=(D-I)^2$ and $(D-I)^3=0$.  The reduction of $D$
on the rank-three space therefore has one Jordan block
of length three.  Choose $q\in L_+$ whose image has
$(D-I)^2q\ne0$ modulo three.  Then $q,Dq,D^2q$ are a
basis modulo three and hence an $R$-basis.  In this basis
$D$ is the three-cycle permutation matrix, since $D^3=I$.
We have proved, without a global lattice classification,
that the full finite representation over $R$ is
\[
                   1\oplus\operatorname{Perm}_3\oplus(-I_2).
 \tag{F.X2}\label{maps:eq:F.X2}
\]
Its dual has the same form, by reversing the three-cycle.

\item Let $u,a,b,c,d,e$ be a cohomology basis for \eqref{maps:eq:F.X2}, with
$a,b,c$ cycled and $d,e$ negated.  A full integral basis
of the invariant degree-three lattice is
\[
 u(ab+bc+ca),\qquad abc,\qquad ude,\qquad(a+b+c)de.
 \tag{F.X3}\label{maps:eq:F.X3}
\]
Here juxtaposition denotes exterior multiplication.
To check completeness, terms with an odd number of
$d,e$ are negated by the cube of the finite map and
have no invariants.  The terms with zero or two such
factors give exactly the four orbit sums in \eqref{maps:eq:F.X3},
each primitive.  With volume $uabcde$, their skew
pairing has one block with coefficient three and one
block with coefficient a unit.  Thus its determinant
has three-adic valuation two.  The exact certificate
\texttt{fourfolds\_universal\_56\_obstruction.py} checks
the full invariant kernel, its primitivity, and the
pairing independently by exterior multiplication and
six-column determinants; its Smith factors are
$(1,1,3,3)$.

\item Let $B_6$ be the actual reduced finite fibre.  Its torus
cover has degree six.  Its free middle cohomology has
rank four and unimodular skew pairing, by integral
Poincare duality.  If $I$ denotes the index of its
pullback lattice in the invariant lattice, multiplication
of pairings by the covering degree gives
\[
 2v_3(I)+2=4v_3(6)=4,
 \qquad v_3(I)=1.
 \tag{F.X4}\label{maps:eq:F.X4}
\]
Consequently the three-primary degree-three covering
cokernel is $R/3R$.  Translations of the affine action
do not enter this calculation: the cover degree,
linear invariant lattice and oriented pairings are
those of the actual free quotient for every allowed
affine numerator.

\item We explain why this local group must appear in total
homology even when the global integral lattice varies.
Let $\mathcal V$ be the rational rank-six cohomology
local system.  It is an extension of the irreducible
rank-five system $W$ by a trivial line.  The determinant
of $W$ is trivial.  Irreducibility in odd dimension
implies that neither $\bigwedge^2 W$ nor its dual
has a nonzero invariant: an invariant alternating
form would have an invariant radical and would have
to be nondegenerate, impossible in odd dimension.
Also $\bigwedge^3W\simeq\bigwedge^2W^*$ by its
trivial determinant.  The exterior filtration therefore
gives
\[
 (\bigwedge^3\mathcal V)^\Gamma=0,\qquad
 ((\bigwedge^2\mathcal V)^*)^\Gamma=0.
 \tag{F.X5}\label{maps:eq:F.X5}
\]
These rational vanishing statements remain true over
$\mathbb Q_3$, since they are kernels of rational
matrices.

\item As in \eqref{maps:eq:F.S7}, take the image of
$R^3g_*R\to j_*\bigwedge^3\mathcal V_R$.
Its kernel and cokernel are point-supported sheaves.
The first vanishing in \eqref{maps:eq:F.X5} makes the local group
$R/3R$ inject into $H^1(B,R^3g_*R)$.  Bidegree $(1,3)$
has no incoming or outgoing Leray differential on
the two-dimensional base, so this injection survives
to $E_\infty^{1,3}$.

\item There is one extension issue, and it cannot remove
this torsion.  The second vanishing in \eqref{maps:eq:F.X5} says
that $H^2(B,R^2g_*R)$ has rational rank zero.
Point-supported changes do not alter this degree-two
base cohomology; it is the coinvariant module, as
in \ExternalPaperRef{}{({E.9})}{III}.  It is a finitely generated torsion $R$-module
and hence a finite group.  Therefore
$E_\infty^{2,2}$ is finite as well.  In the filtration
of $H^4(M;R)$, take the inverse image of the injected
$R/3R\subset E_\infty^{1,3}$.  It is an extension
of $R/3R$ by the finite group $E_\infty^{2,2}$.
It is thus a nonzero finite three-primary subgroup
of $H^4(M;R)$; in particular it cannot be absorbed
as a finite quotient of a free filtration kernel.
Universal coefficients identify this torsion with
the three-primary torsion of $H_3(M;\mathbb Z)$.
This proves the claim at every rational position and
every integral lattice, and completes the exclusion.

\end{enumerate}
\end{proof}




\begin{theorem}[universal fundamental group and twist normalization]\label{maps:src:F.2.28}
\StatementRecord{maps:src:F.2.28}{proved}{Historical events: Claude Fable 5.2 (Anthropic), effort high, 16 November 2026: carried as background without fault found. GPT-6 Astra (OpenAI), ChatGPT Pro mode, effort Pro, S6\_09\_PAPERS, 17 November 2026: targeted proof check. Also checked in this writing. Current writing events: GPT-6 Astra (OpenAI), ChatGPT Pro mode, effort Pro, S6\_10\_PAPERS, 18 November 2026: targeted proof check. GPT-6 Astra (OpenAI), ChatGPT Pro mode, effort Pro, S6\_10\_PAPERS continuation, 19 November 2026: used and rechecked in this run. \emph{Recorded proof scope.} Claude Fable 5.2 (Anthropic), effort high, 16 November 2026; also checked in this writing by GPT-6 Astra (OpenAI), ChatGPT Pro mode, effort Pro, 18 November 2026}{the hypotheses, notation and definitions stated here; \cref{maps:src:F.2.25}; \cref{maps:eq:F.S5}; \cref{maps:eq:F.T8}; \cref{maps:src:F.2.24}}


In the remaining $(5,4)$ locus, fix the marked standard Kuhn
cap and allow affine translation clutching.
Let $\Lambda=\mathbb Z^6$ be fibre homology and let
$\Lambda_\Gamma$ be its global monodromy coinvariants.
Write $\Lambda_\Gamma=\mathbb Zh\oplus T$, with $T$ finite,
and let $\ell$ be the primitive invariant covector detecting
the free factor.  The following are complete reductions for
any datum that could satisfy $\pi_1=1$ and $H_2=H_3=0$.

The actual multiplicities must be $(5,4)$, and the local
vertical fixed pairings must be $h_1=5,h_2=4$, as in \ref{maps:src:F.2.25}.
For any affine twists and clutching the fundamental group
is abelian.  If its free presentation determinant is nonzero,
its order is
\[
 |\pi_1|=|20\kappa+4r_1+5r_2|\,|T|,
 \qquad r_i=\ell(v_i).
 \tag{F.T11}\label{maps:eq:F.T11}
\]
If that determinant is zero the group is infinite.
In particular simple connectivity is equivalent to
\[
 T=0,\qquad 20\kappa+4r_1+5r_2=\pm1.
 \tag{F.T12}\label{maps:eq:F.T12}
\]
Whenever these conditions and $h_1=5,h_2=4$ hold, every
such twist choice is oriented diffeomorphic to the normalized
filling with twists $(u_1,-u_2)$ and zero cusp clutching,
where $u_i$ is any integral fixed vector with $\ell(u_i)=1$.
Thus each eligible integral linear datum has exactly one
normalized smooth filling to test for the remaining
homology conditions; no affine-twist enumeration remains.

\end{theorem}


\begin{proof}
\begin{enumerate}[label=\arabic*.]
\item The image of the fibre fundamental group is a normal
abelian subgroup in the van Kampen presentation.
After quotienting by it, the two finite meridians satisfy
$a^5=b^4=ab=1$.
That quotient is trivial, because five and four are coprime.
The entire fundamental group is therefore equal to the
image of the fibre group and is abelian.
This argument does not assume that the normal closure of
cusp vanishing cycles already exhausts $\ker\ell$.

\item In abelianization the monodromy relations give
$\Lambda_\Gamma$.
The primitive cusp angular lattice is the image of its
nilpotent operator, and this image is already zero among
those monodromy relations.
The remaining presentation is
\[
 5a=v_1,\qquad4b=v_2,\qquad a+b+v_0=0
 \quad\text{in }\Lambda_\Gamma\oplus\mathbb Za\oplus\mathbb Zb,
 \tag{F.T13}\label{maps:eq:F.T13}
\]
where $\ell(v_0)=\kappa$ fixes the sign convention.
Use any Smith presentation of $T$ and split the free part
$\mathbb Zh$.
The full relation matrix is block triangular, with the
presentation of $T$ as one diagonal block and the
three-by-three free matrix of determinant
$\pm(20\kappa+4r_1+5r_2)$ as the other.
Its determinant gives \eqref{maps:eq:F.T11}; when the latter free
matrix has rank less than three the group has positive
rank and is infinite.
For determinant $\pm1$, elimination of the three free
variables leaves precisely $T$, even if the twists have
nonzero torsion components.
This proves the exact simple-connectivity criterion.

\item We give the global twist comparison to make its scope
precise.  Under \eqref{maps:eq:F.T12}, the finite fixed covector image
is one at each point, so choose integral fixed $u_i$ with
$\ell(u_i)=1$.  The primitive vertical fixed vector
$\lambda_i$ and invariant covector $\varphi_i$ satisfy
$\varphi_i(\lambda_i)=m_i$, by the assumed maximal local
pairings.  Formula \eqref{maps:eq:F.S5} makes the finite norm onto
$\mathbb Z\lambda_i$.
Consequently every vertical numerator component
$s_i\lambda_i$ is removed by the explicit cap translation
\eqref{maps:eq:F.T8}, whose boundary residue has $\ell$-coordinate zero.

\item For the complementary pair of pants the full relative
obstruction group for extending boundary translation
sections is
$H^2(P,\partial P;\Lambda)=\Lambda_\Gamma$.
Now $T=0$, so it is exactly $\mathbb Z$ detected by
$\ell$, not just a rational or free quotient.
Prescribing the identity translation at the cusp therefore
allows all the vertical finite changes to extend over $P$.
The same argument removes the zero-$\ell$ component of
any cusp translation clutching.
For determinant $-1$, write
$r_1=1+5a$, $r_2=-1+4b$; then \eqref{maps:eq:F.T12} says
$a+b+\kappa=0$.
The two integer changes of finite boundary marking and
the remaining cusp change thus also have zero total
relative obstruction and extend.
This gives the normalized twists $(u_1,-u_2)$.
For determinant $1$ the same argument gives their negatives.
The orientation-preserving fibre inversion extending over
the Kuhn fan, proved in \ref{maps:src:F.2.24}, identifies the two choices.
All maps used in the comparison are cap translations or
translations by smooth torus-bundle sections over $P$;
the prescribed cusp is retained throughout.
This proves the claimed complete twist normalization.

\end{enumerate}
\end{proof}




\begin{theorem}[a complete finite integral survivor test]\label{maps:src:F.2.29}
\StatementRecord{maps:src:F.2.29}{proved}{Historical events: Claude Fable 5.2 (Anthropic), effort high, 16 November 2026: carried as background without fault found. GPT-6 Astra (OpenAI), ChatGPT Pro mode, effort Pro, S6\_09\_PAPERS, 17 November 2026: targeted proof check. Also checked in this writing. Current writing events: GPT-6 Astra (OpenAI), ChatGPT Pro mode, effort Pro, S6\_10\_PAPERS, 18 November 2026: computation or formal certificate; targeted proof check. GPT-6 Astra (OpenAI), ChatGPT Pro mode, effort Pro, S6\_10\_PAPERS continuation, 19 November 2026: used and rechecked in this run. \emph{Recorded proof scope.} Claude Fable 5.2 (Anthropic), effort high, 16 November 2026; also checked in this writing by GPT-6 Astra (OpenAI), ChatGPT Pro mode, effort Pro, 18 November 2026}{the hypotheses, notation and definitions stated here; \cref{maps:src:F.2.20}; \cref{maps:src:F.2.22}; \cref{maps:src:F.2.25}; \cref{maps:src:F.2.28}; \cref{maps:src:F.2.6}; \cref{maps:eq:F.HC7}; \ExternalPaperRef{theorem~}{6.3}{III}; \cref{maps:src:F.2.16}}


Take any rational pair in \ref{maps:src:F.2.20} and any of its integral,
primitive-cusp rank-six extensions supplied by \ref{maps:src:F.2.22}.
Write $M_1,M_2$ for geometric transport on fibre homology
$\Lambda=\mathbb Z^6$, in the peripheral convention
$M_0M_1M_2=I$.
First require the preliminary conditions
\[
 \Lambda_\Gamma=\mathbb Z,\qquad
 \ell(\Lambda^{M_i})=\mathbb Z,\qquad h_1=5,\quad h_2=4 .
 \tag{F.HC1}\label{maps:eq:F.HC1}
\]
Here $\ell$ is the primitive global invariant covector, and
$h_i$ is the integral fixed-vector/fixed-covector pairing
on $\ker\ell$ from \ref{maps:src:F.2.25}.
Failure of any condition excludes the datum.
If they hold, use the unique normalized affine filling of
\ref{maps:src:F.2.28}, with twists $(u_1,-u_2)$ and zero cusp clutching.
The following six finite integer-matrix tests are necessary
and sufficient for this filling to have $H_2=H_3=0$.

Let
\[
 T_i=M_i^{-t},\quad T_0=(T_1T_2)^{-1},\qquad
 G_q=\operatorname{coker}
 [\,\bigwedge^qT_1-I\mid\bigwedge^qT_2-I\,]
 \quad(1\le q\le4).
 \tag{F.HC2}\label{maps:eq:F.HC2}
\]
All $G_q$ are finite.  Let $P_q$ be the integral parabolic
cohomology group of the exterior local system
$\bigwedge^q\mathcal V$, for $q=2,4$.
Let $t_q\in G_{q-1}$, $2\le q\le5$, be the natural global
coinvariant image of the unique nonzero torsion suspension
class at the order-four cap, as specified below.
Let $c_4\in G_3$ be the natural image of the unique nonzero
torsion class in
$\operatorname{coker}(\bigwedge^3T_0-I)$.
Then the six tests are
\[
\begin{array}{c|ccc}
 &\text{first test}&\text{second test}&\text{third test}\\ \hline
 H^3=0&G_1=\langle t_2\rangle&P_2=0&t_3\ne0\\
 H^5=0&G_3=\langle t_4,c_4\rangle&P_4=0&t_5\ne0 .
\end{array}
\tag{F.HC3}\label{maps:eq:F.HC3}
\]
Every group, class, and equality here is computed by
finitely many integer kernels, Smith forms, and exterior
minors.  This is a decision criterion for every rational
parameter and every integral lattice in the complete
normal form; it makes no bounded-height assertion about
the rational parameter set.

\end{theorem}


\begin{proof}
\begin{enumerate}[label=\arabic*.]
\item We first establish the uniform local calculations needed
for the criterion, so the test does not assume that a
general finite matrix is globally conjugate to the example
in \ref{maps:src:F.2.6}.
The fixed vectors $u_i$ split the homology lattice as
$\mathbb Zu_i\oplus H_i$, where $H_i=\ker\ell$.
Write $D_i=M_i|_{H_i}$.
The maximal pairing $h_i=m_i$ gives
\[
 \operatorname{Norm}_{m_i}(D_i)=\lambda_i\varphi_i ,
 \tag{F.HC4}\label{maps:eq:F.HC4}
\]
with both factors primitive.

\item At five, work over $\mathbb Z_5$.
Modulo five the norm is $(D_1-I)^4$ and is nonzero.
Choose $q$ with $\operatorname{Norm}_5q$ primitive.
The five vectors $q,D_1q,\ldots,D_1^4q$ are independent
modulo five: otherwise the nilpotent Jordan block generated
by $q$ would have length less than five.
They are therefore a $\mathbb Z_5$ basis.
Thus $D_1$ is the regular five-cycle permutation locally.

\item At four, work over $\mathbb Z_2$.
The norm is $(D_2-I)^3$ modulo two and is nonzero.
The vectors $q,D_2q,D_2^2q,D_2^3q$ similarly generate
a primitive regular four-cycle submodule $Q$ of rank four.
The quotient is rank one with $D_2=-1$, by the
characteristic polynomial.
For a lift $z$ of its generator, put $a=(D_2+I)z\in Q$.
Since $D_2^4=I$,
\[
 (D_2^3-D_2^2+D_2-I)a=0.
\]
On the regular four-cycle module the kernel of the
displayed alternating norm is exactly $(D_2+I)Q$,
integrally.  One can see this by writing its four
coefficients; equivalently $D_2+I$ has three nonzero
Smith factors, all one, and the indicated rational image.
Subtracting a suitable element of $Q$ from $z$ makes
$D_2z=-z$.
Consequently
\[
 M_1\sim 1\oplus\operatorname{perm}_5\quad\text{over }\mathbb Z_5,
 \qquad
 M_2\sim1\oplus\operatorname{perm}_4\oplus(-1)
 \quad\text{over }\mathbb Z_2 .
 \tag{F.HC5}\label{maps:eq:F.HC5}
\]
These local comparisons apply to cap cohomology, not just
to abstract representation lattices.
Indeed the integral lattice of intertwiners is the kernel
of an integer matrix and its tensor with $\mathbb Z_p$
is the lattice of local intertwiners.
Approximate a local invertible intertwiner by an integral
one whose determinant is a $p$-adic unit.
Keeping the $u_i$ summand fixed produces an equivariant
torus isogeny of degree prime to $p$, preserving the
normalized translation $\pm u_i/m_i$.
The resulting cap and boundary coverings induce
isomorphisms on $\mathbb Z_p$ cohomology: the vertical
torus maps are cohomology isomorphisms over $\mathbb Z_p$,
so the mapping-torus Wang sequences give the assertion.
At primes not dividing $m_i$, transfer already makes
finite-cover pullback an isomorphism onto invariants.

\item Here are the full required integral results for these
two monomial models.
The ranks of their invariant exterior lattices in
degrees zero through six are
$(1,2,3,4,3,2,1)$.
The order-five cap has no cohomology torsion in degrees
two through five.
The order-four cap has exactly $\mathbb Z/2$ in each
of those four degrees.
This follows from the Wang sequence for the mapping
torus of the vertical $D_i$:
its torsion in degree $q$ is
\[
 \operatorname{Tor}\operatorname{coker}
       (\bigwedge^{q-1}D_i^t-I).
 \tag{F.HC6}\label{maps:eq:F.HC6}
\]
In particular it introduces no additional torsion primes.

\item Integral pullback from either finite cap is onto the
invariant lattice in degrees two through five.
For degree two this is precisely the $m_i/h_i=1$
suspension calculation of \ref{maps:src:F.2.25}.
For completeness, the determinants of the pairings between
the invariant lattices in complementary degrees
$(1,5),(2,4),(3,3)$ are respectively
$m_i,m_i^3,m_i^4$ in absolute value.
The degree-one cover image has index $m_i$.
The free downstairs complementary pairings are unimodular,
and pullback multiplies each pairing by the covering
degree $m_i$.
The middle pairing first gives degree-three index one.
The degree-two index one then gives degree-four index one;
the degree-one index $m_i$ gives degree-five index one.
These determinant identities and all the Wang Smith
forms are verified by the fixed, exhaustive computation
in \texttt{fourfolds\_local\_caps.py}, for both orders and
all exterior degrees zero through six.
Thus no unexamined finite matrix is being inferred from
an example.

\item We next identify the boundary maps.
Choose any unimodular matrix $P_2=[u_2,H_2]$ realizing
the splitting at four, and put
$J=P_2^{-t}|_{\mathbb Z^5}$ for the inclusion of vertical
cohomology.
Choose the unique nonzero torsion representative
\[
 w_q\in\operatorname{coker}
            (\bigwedge^{q-1}D_2^t-I),\qquad
 t_q=[(\bigwedge^{q-1}J)w_q]\in G_{q-1}.
 \tag{F.HC7}\label{maps:eq:F.HC7}
\]
A Smith transformation supplies an integral $w_q$.
Its norm is exactly zero, since it lies rationally in
the image of $\bigwedge^{q-1}D_2^t-I$.
The finite boundary coordinates of \ref{maps:src:F.2.6} satisfy
$\theta=a+4k$, $h=k$, hence $a=\theta-4h$.
The pullback of the suspension cocycle has $\theta$
component $w_q$ and $h$ component
$-\operatorname{Norm}_4(w_q)=0$.
This holds in every exterior degree, including $q=5$.
It is an integral cellular assertion for the monomial
model, where the torus cellular differentials vanish;
the prime-to-two isogeny comparison transports it to
every eligible cap.
The base connecting map therefore gives exactly the
class $t_q$ in \eqref{maps:eq:F.HC7}.
Different representative choices change it by a global
monodromy relation.  Any sign change is immaterial for
an order-two class.

\item There is one further boundary map, from the free
degree-four cusp kernel.
The primitive cusp is integrally conjugate to
$\left(\begin{smallmatrix}I&I\\0&I\end{smallmatrix}\right)$.
Its degree-three exterior difference has ten Smith
factors one and one Smith factor two; its cokernel is
$\mathbb Z^9\oplus\mathbb Z/2$.
The other exterior differences in degrees one, two,
four and five have torsion-free cokernels.
The exact matrices and a torsion representative are
saved by \texttt{fourfolds\_cusp\_boundary.py}.
Let $C$ be the smooth cusp cap and $Y=\partial C$.
The boundary Wang sequence gives
$\operatorname{Tor}H^4(Y)=\mathbb Z/2$.
In the Poincare--Lefschetz exact sequence its image in
\[
 H^5(C,Y)\cong H_3(C)=\mathbb Z^9
\]
is zero, so this torsion class lifts to $H^4(C)$.
Its restriction to an ordinary fibre is zero.
By \ExternalPaperRef{}{6.3}{III} the kernel of integral specialization is
a primitive copy $K=\mathbb Z$ in degree four.
The lift is in $K$, and the image of this cyclic group
in $H^4(Y)$ contains the nonzero torsion class.
A cyclic image containing nonzero torsion cannot
also have a nonzero free part.
Its primitive generator therefore restricts to the
unique nonzero torsion class of $H^4(Y)$.
Consequently the cusp contribution to $d_2$ is
exactly $c_4$ as defined above.  In particular this
column must not be discarded merely because it
vanishes in the first examples.

\item It remains to prove that these are all the tests.
For the irreducible odd-dimensional quotient $W$
of $\mathcal V$, invariant alternating forms and
bivectors vanish: a nonzero one would give an
isomorphism $W\cong W^*$ with an odd-dimensional
nondegenerate alternating pairing.
The exact exterior sequences for
$0\to\mathbb Q\to\mathcal V_{\mathbb Q}\to W\to0$,
together with $\det W=1$, give
\[
 H^0(B,j_*\bigwedge^q\mathcal V)=0
       \quad(2\le q\le5),\qquad
 G_q\otimes\mathbb Q=0\quad(1\le q\le4).
 \tag{F.HC8}\label{maps:eq:F.HC8}
\]
For the endpoint assertions, the extension cannot split:
on $\mathbb Q\oplus W$ a square-zero nilpotent has rank at
most two, whereas the stipulated rank is three.
An invariant covector of $\mathcal V_{\mathbb Q}$ would
split this extension, since $W$ has no trivial quotient.
Thus $(\mathcal V_{\mathbb Q}^*)^\Gamma=0$, which gives
both $H^0(\bigwedge^5\mathcal V_{\mathbb Q})=0$ and
$G_1\otimes\mathbb Q=0$.
Here and below the first assertion is integral as well,
since the coefficient lattices are torsion free.
The local invariant ranks in exterior degrees two
and four are $3,3,9$ at the three punctures.
The Euler formula for the rank-fifteen local system
therefore gives $P_2\otimes\mathbb Q=P_4\otimes\mathbb Q=0$.
In degree one it similarly gives rational parabolic
cohomology zero.
The only global degree-one section is the invariant
line.  Simple connectivity and the universal
coefficient theorem make its free transgression
onto $H^2(B,\mathbb Z)$; equivalently the rational
second Betti number is zero, as in \ref{maps:src:F.2.6}.
Thus the normalized filling has rational $b_2=b_3=0$.

\item All specializations in exterior degrees two through
five are onto, by the local calculations and \ExternalPaperRef{}{6.3}{III}.
Their only point-supported kernels are
\[
\begin{array}{c|cccc}
q&2&3&4&5\\ \hline
\text{kernel}&\mathbb Z/2&\mathbb Z/2&
                 \mathbb Z\oplus\mathbb Z/2&\mathbb Z/2 .
\end{array}
\tag{F.HC9}\label{maps:eq:F.HC9}
\]
The free summand in degree four is at the cusp; the
others are at four.
Consequently $H^1(B,R^qg_*\mathbb Z)=P_q$ in the
degrees used here, while
$H^2(B,R^qg_*\mathbb Z)=G_q$.
The latter equality also holds in degree one:
changing a sheaf at isolated points changes no
degree-two cohomology.  The equality with coinvariants
follows from compactly supported top cohomology on
the punctured sphere and the sequence $j_!\to j_*$.

\item The base has real dimension two.
Thus the only possible nonzero differential in the
integral Leray spectral sequence is
$d_2:E_2^{0,q}\to E_2^{2,q-1}$.
The complete associated graded groups in total
degrees three and five are therefore
\[
\begin{split}
 \operatorname{gr}H^3&:
   G_1/\langle t_2\rangle,\quad P_2,\quad
   \ker(\mathbb Z/2\xrightarrow{\,t_3\,}G_2),\\
 \operatorname{gr}H^5&:
   G_3/\langle t_4,c_4\rangle,\quad P_4,\quad
   \ker(\mathbb Z/2\xrightarrow{\,t_5\,}G_4).
\end{split}
\tag{F.HC10}\label{maps:eq:F.HC10}
\]
These are all finite.  A finite filtered group
vanishes exactly when every associated graded group
vanishes, regardless of extension classes.
Since $\pi_1=1$ and rational $b_2=b_3=0$, the universal
coefficient theorem gives
$H^3=\operatorname{Ext}(H_2,\mathbb Z)$, so
$H_2=0$ exactly when $H^3=0$.
Integral Poincare duality gives $H^5\cong H_3$.
This proves \eqref{maps:eq:F.HC3}.

\item For reproducibility, no sheaf calculation is left
as a black box in the integer test.
Given $A=\bigwedge^qT_1$, $B=\bigwedge^qT_2$ on a
free lattice $L$, form the saturated integer kernel
\[
 K_q=\ker[\,A-I\mid A(B-I)\mid -(AB-I)\,]\subset L^3.
 \tag{F.HC11}\label{maps:eq:F.HC11}
\]
Its image under
$(x,y,z)\mapsto((A-I)x,(B-I)y)$ is precisely the
integral parabolic cocycle lattice.
Quotient by the columns of
$\left(\begin{smallmatrix}A-I\\B-I\end{smallmatrix}\right)$
to obtain $P_q$.
This uses the actual integral images of peripheral
differences, not their saturations.
Smith transformations also supply every $w_q$, the
cusp order-two representative, and the tests for
their images and generated subgroups in $G_q$.

\item The program \texttt{fourfolds\_survivor\_test.py}
implements exactly these operations.
Its executed test set is precisely $v=0,1,2$ from
\ref{maps:src:F.2.16}.  For $v=0,1$ all six tests pass;
$G_1,\ldots,G_4$ are each $\mathbb Z/2$,
$P_2=P_4=0$, all four finite-cap images are nonzero,
and the cusp image is zero.
For $v=2$ they are each $\mathbb Z/6$,
$P_2=P_4=\mathbb Z/3$, and the first, second,
fourth and fifth tests fail.  The computation thus
checks both accepted and rejected data.

\item We finish by making the convention dictionary explicit.
The imported family matrices called $A,B$ in that
program are geometric cohomology pullbacks, so the
geometric homology inputs are $(A^t,B^t)$.
Their genuine left cohomology transport in \eqref{maps:eq:F.HC2}
is $(A^{-1},B^{-1})$, which is what the code uses.
For comparison with computations written using
$(A,B)$, simultaneous reversal of the two peripheral
generators sends
\[
 g_1\mapsto g_1^{-1},\qquad
 g_2\mapsto g_2^{-1},\qquad
 g_0=(g_1g_2)^{-1}\mapsto
     g_2g_1=g_1^{-1}(g_1g_2)g_1 .
 \tag{F.HC12}\label{maps:eq:F.HC12}
\]
This is an orientation-reversing marking of the
pair of pants fixing its three ends: each peripheral
conjugacy class is inverted.
Thus the new cusp is conjugate to the inverse old
cusp, not an independently chosen third matrix.
Inversion preserves invariant and coinvariant
lattices, and the indicated conjugating generator
acts trivially on global coinvariants.
All boundary classes in \eqref{maps:eq:F.HC3} have order two, so
reversal signs do not change the tests.
The implementation uses \eqref{maps:eq:F.HC2} directly and hence
does not depend on this comparison for its validity.

\end{enumerate}
\end{proof}




\begin{corollary}[constructive classification of all survivors]\label{maps:src:F.2.30}
\StatementRecord{maps:src:F.2.30}{proved}{Historical events: Claude Fable 5.2 (Anthropic), effort high, 16 November 2026: carried as background without fault found. GPT-6 Astra (OpenAI), ChatGPT Pro mode, effort Pro, S6\_09\_PAPERS, 17 November 2026: targeted proof check. Also checked in this writing. Current writing events: GPT-6 Astra (OpenAI), ChatGPT Pro mode, effort Pro, S6\_10\_PAPERS, 18 November 2026: targeted proof check. GPT-6 Astra (OpenAI), ChatGPT Pro mode, effort Pro, S6\_10\_PAPERS continuation, 19 November 2026: used and rechecked in this run. \emph{Recorded proof scope.} Claude Fable 5.2 (Anthropic), effort high, 16 November 2026; also checked in this writing by GPT-6 Astra (OpenAI), ChatGPT Pro mode, effort Pro, 18 November 2026}{the hypotheses, notation and definitions stated here; \cref{maps:src:F.2.20}; \cref{maps:src:F.2.22}; \cref{maps:eq:F.HC1}; \cref{maps:eq:F.HC3}; \cref{maps:eq:F.T12}; \cref{maps:src:F.3.11}; \cref{maps:src:F.2.7}; \cref{maps:src:F.2.14}; \cref{maps:src:F.2.25}; \cref{maps:src:F.2.27}; \cref{maps:src:F.2.28}; \cref{maps:src:F.2.29}; \ExternalPaperRef{theorem~}{6.3}{III}; \cref{maps:src:F.2.15}; \cref{maps:src:F.2.5}; \S\ref{sec:V-known-inputs}; \cref{maps:src:F.4.1}}


For the marked Kuhn cap and affine translation clutchings
used throughout F.2, the following is a complete constructive
classification of the data with
$\pi_1=1$ and $H_2=H_3=0$, up to integral conjugacy of their
linear tuples and oriented diffeomorphism of their fillings.

There are no data with three genuine finite fillings.
With two fillings their actual affine orders are five and
four.  Order the order-five point first.
The complete, unbounded parameter set is the rational
projective normal form of \ref{maps:src:F.2.20} together with its
exceptional rational affine line.
For each parameter, run the terminating order,
conductor, stable-lattice, and primitive-extension algorithm
of \ref{maps:src:F.2.22}.  Discard precisely the integral classes failing
\eqref{maps:eq:F.HC1} or one of \eqref{maps:eq:F.HC3}.
The classes retained by this procedure are exactly all
surviving integral linear tuples.
For each retained class all simply connected affine
twists are exactly the choices satisfying \eqref{maps:eq:F.T12};
their fillings are all oriented diffeomorphic to its
normalized filling.
If ordered labels with the order-four point first are
desired, interchange the finite labels and recompute
$M_0=(M_1M_2)^{-1}$.

Every retained datum has
\[
 \pi_1=1,\qquad
 H_j=\begin{cases}
  \mathbb Z,&j=0,8,\\
  \mathbb Z^4,&j=4,\\
  0,&\text{otherwise},
 \end{cases}
 \tag{F.HC13}\label{maps:eq:F.HC13}
\]
and its oriented smooth manifold is
$3\mathbb{HP}^2\#(-\mathbb{HP}^2)$.
It has bounded-cusp holomorphic periods in the literal
sense of F.3, by \ref{maps:src:F.3.11}.
The classification here concerns the survivor subproblem.
It does not classify the nonsurviving tuples or compute
their full integral homology, as also requested in F.2.

\end{corollary}


\begin{proof}
\begin{enumerate}[label=\arabic*.]
\item The exclusions \ref{maps:src:F.2.7}, \ref{maps:src:F.2.14}, \ref{maps:src:F.2.25}, and \ref{maps:src:F.2.27}
leave exactly the stated minimal orders and rational
quotient problem.
The rational normal form is injective on conjugacy
classes, including the exceptional line.
The finite interval and extension-coset enumeration
in \ref{maps:src:F.2.22} gives every integral class over each rational
point, with all bounds proved for that input.
No distinct rational points can become integrally
conjugate, since an integral conjugacy induces a
rational quotient conjugacy.
Theorem \ref{maps:src:F.2.28} and the necessary-and-sufficient
criterion \ref{maps:src:F.2.29} prove both directions of the
retention rule and account for every affine twist.
This is also an exhaustive enumeration if desired:
enumerate primitive projective integer triples with
first nonzero coordinate positive by increasing height,
and enumerate reduced rational exceptional parameters
by increasing numerator and denominator height.
Every finite stage has finitely many inputs, each
terminates by \ref{maps:src:F.2.22}, and every rational input occurs.
No finite cutoff is imposed or needed.

\item Hurewicz makes the retained simply connected manifolds
three-connected.  The only singular contribution to
Euler number is the Kuhn cusp, whose value is six
by \ExternalPaperRef{}{6.3}{III}; torus bundles and the free finite caps have
Euler number zero.
Integral Poincare duality and the universal coefficient
theorem then give \eqref{maps:eq:F.HC13}, including torsion freeness.

\item For the smooth classification, \ref{maps:src:F.2.15} supplies a
compatible complex structure at each finite cap, and
the cusp already has its complex model.
Choose their common fibre orientation to be the one
induced by the cusp.
Commuting complex structures at the finite caps can
be chosen in that component: their two-dimensional
fixed real summands allow the orientation to be
reversed when necessary.
The fibre of the bundle of orientation-compatible
complex structures retracts to $\mathbb{CP}^3$.
Its connectedness and simple connectivity extend
the prescribed boundary structures across the
two-dimensional pair of pants, exactly as in \ref{maps:src:F.2.5}.
A compatible horizontal distribution gives an
almost complex structure on the total manifold.
Theorem \S\ref{sec:V-known-inputs} therefore gives signature two,
intersection lattice $I_{3,1}$, and $c_2^2=2$.
Since $H^2=0$, $p_1/2=-c_2$ is primitive,
characteristic, and has square two.
The complete integral reflection, Wall-form and Kramer--Stolz argument in \ref{maps:src:F.4.1}
depends only on these almost complex characteristic
numbers and three-connectivity.
It applies without alteration and gives the claimed
exact oriented diffeomorphism.
Finally \ref{maps:src:F.3.11} applies to the resulting determinant-one
monodromy and primitive cusp.
It asserts the requested punctured-base periods with
bounded correction, not an additional unproved
simultaneous good extension at all finite points.

\end{enumerate}
\end{proof}




\begin{theorem}[complete fundamental-group data for every input]\label{maps:src:F.2.31}
\StatementRecord{maps:src:F.2.31}{proved}{Historical events: Claude Fable 5.2 (Anthropic), effort high, 16 November 2026: carried as background without fault found. GPT-6 Astra (OpenAI), ChatGPT Pro mode, effort Pro, S6\_09\_PAPERS, 17 November 2026: targeted proof check. Also checked in this writing. Current writing events: GPT-6 Astra (OpenAI), ChatGPT Pro mode, effort Pro, S6\_10\_PAPERS, 18 November 2026: computation or formal certificate; targeted proof check. GPT-6 Astra (OpenAI), ChatGPT Pro mode, effort Pro, S6\_10\_PAPERS continuation, 19 November 2026: used and rechecked in this run. \emph{Recorded proof scope.} Claude Fable 5.2 (Anthropic), effort high, 16 November 2026; also checked in this writing by GPT-6 Astra (OpenAI), ChatGPT Pro mode, effort Pro, 18 November 2026}{the hypotheses, notation and definitions stated here; \ExternalPaperRef{theorem~}{6.3}{III}; \cite[Theorems 2.3--2.4, p.~425; Euler-characteristic comparison, pp.~427--428]{rootScott83}}


Let $k=2$ or $3$, and let an integral rank-six tuple satisfy the
finite-order and primitive maximal-cusp conditions of F.2.  No
irreducibility, determinant-one, or prospective-survivor hypothesis is
imposed.  Specify each actual finite affine action as
\[
 \tau_i(x)=M_ix+v_i/m_i,\qquad v_i\in\Lambda=\mathbb Z^6,
 \quad \tau_i^{m_i}=1\text{ on }\Lambda_{\mathbb R}/\Lambda,
 \tag{F.G1}\label{maps:eq:F.G1}
\]
and require its cyclic action to be free.  The numerator need not be
fixed by $M_i$.  Define its actual integral power vector and the cusp
lattice by
\[
 w_i=\frac1{m_i}\sum_{j=0}^{m_i-1}M_i^jv_i\in\Lambda^{M_i},
 \qquad U=\operatorname{im}\bigl((M_1\cdots M_k)^{-1}-I\bigr).
 \tag{F.G2}\label{maps:eq:F.G2}
\]
Use left conjugation transport for the peripheral generators $s_i$,
and record the integral translation clutching by $\kappa\in\Lambda$
in $s_1\cdots s_kx^\kappa=1$.  Then the full fundamental group is
\[
 \left\langle x^\lambda,s_1,\ldots,s_k\ \left|\
 \begin{array}{l}
 x^\lambda x^\mu=x^{\lambda+\mu},\quad x^u=1\ (u\in U),\\
 s_ix^\lambda s_i^{-1}=x^{M_i\lambda},\quad
 s_i^{m_i}=x^{w_i},\quad s_1\cdots s_kx^\kappa=1
 \end{array}\right.\right\rangle .                         \tag{F.G3}\label{maps:eq:F.G3}
\]
Only a basis of $\Lambda$ and of $U$ is needed in this finite
presentation.  The program below also determines exactly which
quotient of the apparent fibre subgroup injects.

Put
\[
 L=\sum_{g\in\langle M_1,\ldots,M_k\rangle}gU,
 \qquad A=\Lambda/L.                                      \tag{F.G4}\label{maps:eq:F.G4}
\]
This is an integral sum, without saturation.  For two points, put
$d=\gcd(m_1,m_2)$, $M=M_1$, and
$S_j(M)=I+M+\cdots+M^{j-1}$.  Choose integers $a,b$ with
$am_1+bm_2=d$, and set, in $A$,
\[
 \begin{split}
 w'_2&=-w_2-S_{m_2}(M)\kappa,\\
 e&=(m_2/d)w_1-(m_1/d)w'_2,\qquad
 z=aw_1+bw'_2,\qquad \overline A=A/\langle e\rangle .
 \end{split}                                             \tag{F.G5}\label{maps:eq:F.G5}
\]
There is an exact sequence and a complete marked extension description
\[
 1\longrightarrow\overline A\longrightarrow\pi_1
   \longrightarrow C_d\longrightarrow1,
 \qquad txt^{-1}=M(x),\quad t^d=z.                         \tag{F.G6}\label{maps:eq:F.G6}
\]
Thus its order is $d|\overline A|$ if $\overline A$ is finite;
otherwise it is infinite.  In particular $d=1$ gives precisely
the abelian group $\overline A$.

For three genuine points, so $m_i\ge2$, let
\[
 \Gamma=\langle q_1,q_2,q_3\mid q_i^{m_i}=1,
                         \ q_1q_2q_3=1\rangle,
 \qquad \epsilon=\sum_i1/m_i-1.                           \tag{F.G7}\label{maps:eq:F.G7}
\]
The action $q_i\mapsto M_i$ is well defined on $A$.
If $\epsilon\le0$, the exact sequence is
\[
                    1\longrightarrow A\longrightarrow\pi_1
                         \longrightarrow\Gamma\longrightarrow1. 
 \tag{F.G8}\label{maps:eq:F.G8}
\]
If $\epsilon>0$, $\Gamma$ is the finite spherical triangle group.
Using its action on $A$, form the exact transfer vector
\[
 e=\sum_{g\in\Gamma}g\kappa+
       \sum_{i=1}^3\ \sum_{g\in\Gamma/\langle q_i\rangle}g w_i.
 \tag{F.G9}\label{maps:eq:F.G9}
\]
The second sums use left coset representatives and are independent
of their choice.  In this case the exact sequence has kernel
$A/\langle e\rangle$, with the same lifted relations \eqref{maps:eq:F.G3}.
Consequently \eqref{maps:eq:F.G3}, \eqref{maps:eq:F.G6}, \eqref{maps:eq:F.G8}, and \eqref{maps:eq:F.G9} give every requested
fundamental group by finite operations, including the precise
injected abelian kernel and all extension data.  These statements
do not claim a conjugacy classification of the original tuples.

For every input, the full first integral homology is the Smith
cokernel of the integer matrix
\[
 \begin{pmatrix}
 M_1-I&\cdots&M_k-I&-w_1&\cdots&-w_k&\kappa\\
 0&\cdots&0&&\operatorname{diag}(m_1,\ldots,m_k)&&\boldsymbol1
 \end{pmatrix}.                                           \tag{F.G10}\label{maps:eq:F.G10}
\]
Here the bottom left zero block has $k$ rows.  The code records
the full matrix, not only its invariant factors.

\end{theorem}


\begin{proof}
\begin{enumerate}[label=\arabic*.]
\item \emph{Actual finite and cusp caps.}
The affine lift on $\mathbb R^6$ has $m_i$th power translation by
$w_i$, as the displayed sum shows.  It commutes with that power,
so $M_iw_i=w_i$.  The group generated by lattice translations
and this lift is the fundamental group of the reduced flat
six-manifold.  Its presentation consists of the lattice, its
$s_i$ conjugation relations, and $s_i^{m_i}=w_i$.
The finite cap deformation retracts onto that reduced manifold
by the equivariant radial contraction of its root disc.
Thus these are the actual cap relations, even for a numerator
which is not an invariant lattice vector.

\item For the primitive Kuhn cap, the compact toric strata in its
normalization and all their intersection components are simply
connected.  Applying van Kampen successively to their tubular
neighborhoods identifies the fundamental group of the central
fibre with that of its dual three-torus.  The nearby fibre map
kills the angular lattice $U$ and sends the long-period lattice
$\Lambda/U$ isomorphically onto this dual-torus group.  This is
also the explicit toric collapse used in \ExternalPaperRef{}{6.3}{III}.  The primitive
index-one condition gives $U=\ker(M_0-I)$ and makes $\Lambda/U$
free of rank three.  A section through a smooth central stratum
bounds the base meridian by its parameter disc.  The smooth cap
retracts onto the central fibre, so it adds exactly these
relations, with no extra meridian relation left unspecified.

\item Over the punctured sphere the bundle has a section up to homotopy,
because the base retracts to a graph.  Its group is the appropriate
semidirect product.  Choose the cap sections and retain their
relative translation clutching as $\kappa$.  Van Kampen with the
actual boundary maps just described gives exactly \eqref{maps:eq:F.G3}.
Changing the transport convention causes no group ambiguity:
reverse the ordered list and replace each affine map by its
inverse, and replace $\kappa$ by $-\kappa$.  The new linear
matrices are $M_k^{-1},\ldots,M_1^{-1}$, the power vectors are
$-w_k,\ldots,-w_1$, and the cusp image lattice is unchanged.
The substitutions $s_i\mapsto s_{k+1-i}^{-1}$ give the explicit
isomorphism between these two presentations.  This supplies the
dictionary when geometric transport is recorded with opposite
path-composition order.

\item \emph{Integral closure and two-point elimination.}
In a normal abelian fibre subgroup, the normal closure of $U$
is precisely the integer sum \eqref{maps:eq:F.G4}.  Start with its column
Hermite basis and repeatedly append its images under every
$M_i$ and $M_i^{-1}$, then take Hermite normal form again.
If the rational rank increases, it can do so at most three
times.  At the first nonincrease the rational span is invariant
under every generator and its inverse, hence never increases
again.  If the current lattice has saturation index $D$ in
that span, every subsequent strict enlargement divides this
index by an integer at least two.  There can therefore be at
most $\lfloor\log_2D\rfloor$ such enlargements.  This proves
termination and gives the computable bound used by the code;
it never replaces a finite-index lattice by its saturation.

\item On $A$ one has $M_2=M^{-1}$ in the two-point case.  The product
relation eliminates $s_2=t^{-1}x^{-\kappa}$, where $t=s_1$.
Multiplying this word $m_2$ times gives
\[
        (t^{-1}x^{-\kappa})^{m_2}
            =t^{-m_2}x^{-S_{m_2}(M)\kappa}.
\]
Thus the two power relations become $t^{m_1}=w_1$ and
$t^{m_2}=w'_2$.  Both right sides are fixed by $M$ on $A$.
Also $M^{m_1}=M^{m_2}=1$, hence $M^d=1$ there.
Bezout elimination gives exactly $e=0$ and $t^d=z$ in
\eqref{maps:eq:F.G5}.  Replacing the Bezout pair changes $z$ by a multiple
of $e$ and therefore changes no extension datum.

\item To prove the asserted injection, construct
$\overline A\rtimes_M\mathbb Z$ and quotient by its central
element $t^dz^{-1}$.  This cyclic subgroup has nonzero
projection $d\mathbb Z$ to the infinite cyclic factor, so it
intersects $\overline A$ trivially.  The quotient has precisely
the presentation obtained by elimination.  This proves \eqref{maps:eq:F.G6},
including the absence of any further hidden fibre relation.

\item \emph{The three-point kernel and the finite transfer.}
For completeness this step computes the fibre injection; it
does not infer it from the appearance of a normal subgroup in
a presentation.  Let $\mathcal B$ be the orbifold classifying
space obtained from the three-holed sphere by attaching, at
its three boundary circles, the mapping cylinders of
$S^1\to BC_{m_i}$.  Over the punctured sphere take the bundle
with fibre $BA=K(A,1)$ and the specified action.  For each
boundary, the cyclic extension
\[
 E_i=\langle A,s_i\mid s_i a s_i^{-1}=M_i(a),\ s_i^{m_i}=w_i\rangle
\]
contains $A$ injectively: the same semidirect-product construction
as above applies because $M_i^{m_i}=1$ and $M_iw_i=w_i$.
The associated fibration $BA\to BE_i\to BC_{m_i}$ restricts
on $S^1$ to the required mapping-torus fibration.  Replace the
maps by their mapping-path fibrations before homotopy gluing,
so that the fibre is preserved on every attaching cylinder.
Glue these three fibrations, using the specified relative clutching
$\kappa$ on the punctured sphere.  Van Kampen gives \eqref{maps:eq:F.G3}
for the resulting total space.  Thus the homotopy sequence
of this fibration over $\mathcal B$ calculates precisely
the image of $A$ in the group under discussion.

\item When $\epsilon\le0$, the universal orbifold cover is the plane
or hyperbolic plane.  Its contractibility identifies
$\mathcal B$ with $B\Gamma$: indeed its Borel construction
has the just-described boundary-cylinder model and contractible
universal fibre.  The Euclidean and hyperbolic alternatives
are the oriented two-orbifold classification in
\cite[Theorems 2.3--2.4, p.~425; Euler-characteristic comparison, pp.~427--428]{rootScott83}.
Since $\pi_2(B\Gamma)=0$, the homotopy sequence proves \eqref{maps:eq:F.G8}.

\item When $\epsilon>0$, the universal orbifold cover is $S^2$.
Pulling the constructed $BA$ fibration to this finite cover
trivializes its local system.  Its transgression
$\pi_2(S^2)\to A$ is its integral clutching vector.
There is one lifted ordinary base cell for each $g\in\Gamma$,
contributing $g\kappa$, and one lifted cone cap for each left
coset of $\langle q_i\rangle$, contributing $gw_i$.
To check the sign, the base relation has fibre term $+\kappa$
and each cone relation has term $-w_i$; the oriented sphere
chain subtracts the cone cells from the base cell.  The result
is therefore the plus sign in both sums of \eqref{maps:eq:F.G9}.
Changing a coset representative by $q_i^j$ does not change
its contribution because $w_i$ is fixed.  The result is
$\Gamma$ invariant, as also checked by the code.  The long
homotopy sequence now says that the only additional kernel
is exactly $\langle e\rangle$.  This proves the spherical
assertion, with the orientation fixed by this cell calculation.

\item There is no implicit infinite search for a spherical group.
After sorting its orders the possibilities are
$(2,2,n),(2,3,3),(2,3,4),(2,3,5)$, with groups of orders
$2n,12,24,60$, respectively.  These follow directly from
the positive spherical-triangle angle condition.
The code constructs the dihedral, tetrahedral, octahedral,
or icosahedral permutation group and finds the marked
peripheral generators in that finite group.  Breadth-first
words supply their actions on $A$; every generator edge is
checked modulo the actual lattice $L$.  It then computes
the two finite sums in \eqref{maps:eq:F.G9} and their integer Smith form.

\item An order-one site is not a genuine cone.  Its matrix is the
identity and its relation is $s_i=x^{w_i}$.  Eliminate it
from \eqref{maps:eq:F.G3}; moving this vector to the right of the remaining
peripheral product replaces the clutching by
\[
       \kappa+(M_{i+1}\cdots M_k)^{-1}w_i.                 \tag{F.G11}\label{maps:eq:F.G11}
\]
This reduces the three-point case with such a site to the
two-point computation.  A primitive rank-three cusp cannot
remain with only one finite linear monodromy: a finite-order
unipotent matrix is the identity.  Thus no smaller case has
been omitted within the hypotheses of this theorem.

\item \emph{Freeness, abelianization, and executed certificates.}
Freeness of \eqref{maps:eq:F.G1} is itself decidable exactly.  A nontrivial
stabilizer in a cyclic group contains a subgroup of prime
order, so it is enough to check $\tau_i^{m_i/p}$ for every
prime $p\mid m_i$.  For an affine map $Rx+b$ of a torus,
a fixed point exists precisely when every integral covector
in $\ker(R^t-I)$ takes an integral value on $b$.
Smith form proves this criterion: the image of $R-I$ on
the torus is the connected subtorus annihilated by those
primitive covectors.  The code computes the powers in the
rational affine matrix group and performs these finite tests.

\item Abelianizing \eqref{maps:eq:F.G3} yields \eqref{maps:eq:F.G10}.  No cusp column was lost:
\[
 \operatorname{im}(M_0-I)
       \subset\sum_i\operatorname{im}(M_i-I)
\]
follows by telescoping the product on the right; inversion
preserves the image of its difference from the identity.
The Smith form therefore computes every free and torsion
summand of $H_1$ exactly.

\item The program \texttt{fourfolds\_vankampen.py} saves all input
matrices, power vectors, integral closure traces, group
extension data, and homology presentation matrices.
It checks the nine explicitly listed three-point parameters
of its original test set, four two-point order/clutching
choices, a genuinely nonfixed affine numerator, and a
finite-index closure step.  It also checks a nonabelian
rank-six two-point tuple, a spherical $(2,4,2)$ example
with injected kernel $\mathbb Z^2\oplus\mathbb Z/4$, and
the infinite perfect hyperbolic example with orders
$(5,4,19)$.  The nonabelian test deliberately belongs to
the full $\operatorname{GL}_6$ scope and has orientation-
reversing finite matrices; it is not claimed to admit a
complex compactification.  All bounds above are derived
from each input, rather than imposed as a search cutoff.

\end{enumerate}
\end{proof}




\begin{theorem}[a complete three-point logarithmic family]\label{maps:src:F.2.32}
\StatementRecord{maps:src:F.2.32}{proved}{Historical events: Claude Fable 5.2 (Anthropic), effort high, 16 November 2026: carried as background without fault found. GPT-6 Astra (OpenAI), ChatGPT Pro mode, effort Pro, S6\_09\_PAPERS, 17 November 2026: targeted proof check. Also checked in this writing. Current writing events: GPT-6 Astra (OpenAI), ChatGPT Pro mode, effort Pro, S6\_10\_PAPERS, 18 November 2026: computation or formal certificate; targeted proof check. GPT-6 Astra (OpenAI), ChatGPT Pro mode, effort Pro, S6\_10\_PAPERS continuation, 19 November 2026: used and rechecked in this run. \emph{Recorded proof scope.} Claude Fable 5.2 (Anthropic), effort high, 16 November 2026; also checked in this writing by GPT-6 Astra (OpenAI), ChatGPT Pro mode, effort Pro, 18 November 2026}{the hypotheses, notation and definitions stated here; \cref{maps:src:F.3.7}; \cref{maps:src:F.2.28}; \cref{maps:eq:F.LT4}; \cref{maps:eq:F.LT5}; \cref{maps:src:F.2.5}; \cref{maps:eq:F.HC8}; \cref{maps:eq:F.HC9}; \cref{maps:eq:F.LT8}; \cref{maps:src:F.2.29}; \cref{maps:eq:F.LT6}; \cref{maps:eq:F.LT9}}


Start with the compact complex fourfold of \ref{maps:src:F.3.7} and its
normalized order-five and order-four affine fillings.
Let $\ell$ be the primitive invariant covector on fibre
homology, with the sign chosen in \ref{maps:src:F.2.28}.
At any additional ordinary fibre choose
$v\in\mathbb Z^6$ and an integer $m\ge2$ such that
\[
 r=\ell(v),\qquad \gcd(r,m)=1.
\tag{F.LT1}\label{maps:eq:F.LT1}
\]
Perform the marked logarithmic transformation with
translation $v/m$ described below.
This gives a smooth compact complex fourfold $X_{m,v}$
with Euler number six and three genuine finite affine
fillings, whose linear orders are five, four and one
and whose actual affine orders are five, four and $m$.
Its fundamental group is
\[
 \pi_1(X_{m,v})=
 \langle a,b,c,h\mid h\text{ central},\
   a^5=h,\ b^4=h^{-1},\ c^m=h^r,\ abc=1\rangle .
\tag{F.LT2}\label{maps:eq:F.LT2}
\]
This is the group of the oriented Seifert manifold
with invariants
$(0;(5,-1),(4,1),(m,-r))$, in the convention
$x_i^{\alpha_i}h^{\beta_i}=1$ and product $x_i=1$.
Its base orbifold is hyperbolic, and its central fibre
$h$ has infinite order.

Put $C_{m,r}=\mathbb Z/(20r-m)\mathbb Z$, with
$\mathbb Z/0\mathbb Z=\mathbb Z$, and let
$\varepsilon=1$ when $20r=m$ and zero otherwise.
All the integral homology groups are
\[
\begin{array}{c|c}
j&H_j(X_{m,v};\mathbb Z)\\ \hline
0,8&\mathbb Z\\
1,6&C_{m,r}\\
2&\mathbb Z^\varepsilon\oplus(\mathbb Z/m)^5\\
3&(\mathbb Z/m)^{10}\\
4&\mathbb Z^4\oplus(\mathbb Z/m)^{10}\\
5&(\mathbb Z/m)^5\\
7&\mathbb Z^\varepsilon .
\end{array}
\tag{F.LT3}\label{maps:eq:F.LT3}
\]
Under \eqref{maps:eq:F.LT1} the exceptional free-rank case is exactly
$(m,r)=(20,1)$.
In particular $(m,r)=(19,1)$ and $(21,1)$ have
$H_1=0$ but nontrivial, infinite fundamental groups.
More generally this occurs whenever $|20r-m|=1$.
The result treats every vertical component of $v$
in $\ker\ell$, not only a chosen fixed representative.

\end{theorem}


\begin{proof}
\begin{enumerate}[label=\arabic*.]
\item We first specify the complex construction and its marking.
Over a small disc at the ordinary value use a holomorphic
period marking
$T_t=\mathbb C^3/\Lambda_t$.
Let $\omega_v(t)$ be the holomorphic period vector
corresponding to $v$.
On the root disc $t=s^m$ quotient the pulled-back family by
\[
 (s,Z)\longmapsto
   (\zeta_m s,\ Z+\omega_v(s^m)/m).
\tag{F.LT4}\label{maps:eq:F.LT4}
\]
Condition \eqref{maps:eq:F.LT1} makes $v$ have exact order $m$
in $\mathbb Z^6/m\mathbb Z^6$.
Every nontrivial power therefore acts freely, including
over $s=0$.
The quotient is smooth, has a proper map to the disc,
and has central fibre of multiplicity $m$.
On the punctured disc its holomorphic identification with
the original family is
\[
 (s,Z)\longmapsto
 \left(s^m,\ Z-\frac{\log s}{2\pi i}\omega_v(s^m)\right).
\tag{F.LT5}\label{maps:eq:F.LT5}
\]
Changing the branch of $\log s$ changes the second
coordinate by a lattice period.
Changing $s$ by $\zeta_m s$ changes its logarithm by
$2\pi i/m$ and cancels the translation in \eqref{maps:eq:F.LT4}.
Thus \eqref{maps:eq:F.LT5} is a well-defined holomorphic isomorphism
of the quotient with the original torus family off zero.
Gluing on a fixed annulus supplies the asserted compact
complex manifold: its charts are the original charts
and the free finite quotient charts \eqref{maps:eq:F.LT4}, and it is
the union of compact pieces over closed base regions.
The gluing is an isomorphism on collars, so the total
space is Hausdorff and smooth.
The Euler number is unchanged, since the removed
ordinary torus cap, the new cap, and their boundary
all have Euler number zero.
This marking fixes the new relation $c^m=v$ and
zero extra cusp clutching.

\item In the original datum the invariant normal closure
of the cusp angular lattice is exactly $\ker\ell$,
as proved in \ref{maps:src:F.2.5}.
Adding the identity linear generator does not change
that lattice.
The fibre image in van Kampen therefore reduces to
the central generator $h$, while the original two
finite relations remain $a^5=h$, $b^4=h^{-1}$.
The new relation is $c^m=h^{\ell(v)}$ and the base
relation is $abc=1$.
This proves \eqref{maps:eq:F.LT2}, including the absence of any
dependence on the vertical part of $v$.

\item One may verify the Seifert identification directly:
take a product circle bundle on a three-holed sphere
and attach its three solid tori with primitive meridian
slopes $(5,-1),(4,1),(m,-r)$.
They are primitive by \eqref{maps:eq:F.LT1}, and van Kampen gives
exactly \eqref{maps:eq:F.LT2}.
Its orbifold base has
$1/5+1/4+1/m<1$ for every $m\ge2$.
Pulling back the circle fibration to its contractible
hyperbolic orbifold cover gives a circle bundle on
the plane, hence the trivial circle bundle.
The subsequent universal cover is the plane times
$\mathbb R$, and the fibre subgroup is infinite cyclic.
This also proves directly that none of these groups
is trivial.
The rational Euler number in these conventions is
$r/m-1/20$.

\item In abelianization the relations $5a=h$, $4b=-h$
give $a=4x$, $b=-5x$, $h=20x$ because five and four
are coprime.
The product relation gives $c=x$.
The remaining relation is $(m-20r)x=0$.
Thus $H_1=C_{m,r}$ exactly, including its integral
generator and the exceptional infinite-cyclic case.

\item We now compute the higher integral groups, including
all extensions.
Write $\mathcal R^q$ for the higher direct-image
sheaf of the original fourfold and
$\mathcal R_m^q$ for that of $X_{m,v}$, after identifying
the bases and their ordinary local systems.
At the new point the reduced finite quotient is a
torus with lattice
$\mathbb Z^6+\mathbb Z(v/m)$.
Its covering by the original torus is cyclic of
degree $m$.
The degree-one pullback has Smith factors
$1,1,1,1,1,m$, so its degree-$q$ cokernel is
\[
 C^q=(\mathbb Z/m)^{\binom5{q-1}},
 \qquad 1\le q\le6 .
\tag{F.LT6}\label{maps:eq:F.LT6}
\]
There is no cohomology torsion in this torus cap.
The sheaf comparison, which does not assert a map
between the total manifolds, is therefore
\[
 0\longrightarrow\mathcal R_m^q
  \longrightarrow\mathcal R^q
  \longrightarrow C^q_{\{p\}}\longrightarrow0 .
\tag{F.LT7}\label{maps:eq:F.LT7}
\]
It is the identity away from the new point and the
covering image at that point.

\item For $2\le q\le5$, all original global sections of
$\mathcal R^q$ are supported at the original finite
point or the cusp, by \eqref{maps:eq:F.HC8}--\eqref{maps:eq:F.HC9}.
Their map to $C^q_{\{p\}}$ is zero.
Consequently their $H^0$ groups are unchanged,
their $H^2$ groups are unchanged, and
\[
 0\longrightarrow C^q
  \longrightarrow H^1(\mathcal R_m^q)
  \longrightarrow H^1(\mathcal R^q)
  \longrightarrow0 .
\tag{F.LT8}\label{maps:eq:F.LT8}
\]
The original degree-one base cohomology of these
sheaves is zero for $q=2,4,5$ and is $\mathbb Z^3$
for $q=3$.
The assertion for $q=5$ also follows immediately
from the original $H^6=0$: its bidegree $(1,5)$
has no incoming or outgoing differential.
Thus \eqref{maps:eq:F.LT8} gives $C^q$ in degrees two, four and
five, and $\mathbb Z^3\oplus C^3$ in degree three.

\item The original four order-two torsion transgressions
remain isomorphisms.
Indeed their boundary suspension representatives and
their images in global coinvariants are the ones in
\ref{maps:src:F.2.29}; adding the identity monodromy changes none
of those coinvariant quotients.
These point-supported boundary maps do not depend
on the free degree-one transgression, which changes
under the new logarithmic gluing.
The free degree-four cusp column likewise remains
zero, since its boundary torsion representative has
zero image in that same unchanged global quotient.
It follows that the full filtrations in degrees
three and five have just one nonzero term:
\[
 H^3=C^2=(\mathbb Z/m)^5,\qquad
 H^5=C^4=(\mathbb Z/m)^{10}.
\tag{F.LT9}\label{maps:eq:F.LT9}
\]
In degree four there is an exact sequence
\[
 0\longrightarrow
       \mathbb Z^3\oplus(\mathbb Z/m)^{10}
  \longrightarrow H^4
  \longrightarrow\mathbb Z\longrightarrow0.
\]
It splits as an abelian-group extension because
its quotient is free, giving the stated middle group.
Thus no torsion extension has been inferred from
rational ranks.

\item For the remaining free ranks,
$H^1(X_{m,v})=\operatorname{Hom}(C_{m,r},\mathbb Z)$
has rank $\varepsilon$.
The degree-one global section group has rank one,
and the only rational differential involving it is
$d_2:H^0(\mathcal R_m^1)\to H^2(\mathcal R_m^0)$.
Its kernel has rank $\varepsilon$, hence its
cokernel also has rank $\varepsilon$.
The degree-one parabolic group is rationally zero,
and there are no rational degree-two global sections.
Therefore $b_2=\varepsilon$.
For an integral check of this degree-one row, its
finite sheaf cokernel is
$\mathbb Z/\gcd(20,m)$: the original invariant
section restricts as $20\ell$, and its value at
the new covering cokernel is $20r$ modulo $m$.
This finite group changes no rational rank.

\item Poincare duality gives
$H_3=H^5$, $H_4=H^4$, and $H_5=H^3$.
The universal coefficient theorem now identifies
the torsion of $H_2$ with $H^3$ and its free rank
with $\varepsilon$.
It also gives $H^2=C_{m,r}$: when $20r\ne m$ this
is $\operatorname{Ext}(H_1,\mathbb Z)$; when
$20r=m$ it is the free rank-one group.
Duality gives $H_6=C_{m,r}$ and
$H_7=\mathbb Z^\varepsilon$.
This proves the complete table \eqref{maps:eq:F.LT3}.

\item The program \texttt{fourfolds\_three\_point\_family.py}
checks the original four torsion images and zero
cusp image once, then computes the actual cyclic
covering Smith forms in all exterior degrees one
through six and the exact van Kampen abelianization
for the thirty specified pairs
$m\in\{2,3,4,5,6,8,10,19,20,21\}$,
$r\in\{-1,1,m+1\}$.
It includes a nonzero vertical component in every
test vector.
The unbounded assertion follows from \eqref{maps:eq:F.LT6}--\eqref{maps:eq:F.LT9},
not from extrapolation of those finite tests.

\end{enumerate}
\end{proof}




\begin{theorem}[arbitrary finite affine caps]\label{maps:src:F.2.33}
\StatementRecord{maps:src:F.2.33}{proved}{Historical events: Claude Fable 5.2 (Anthropic), effort high, 16 November 2026: carried as background without fault found. GPT-6 Astra (OpenAI), ChatGPT Pro mode, effort Pro, S6\_09\_PAPERS, 17 November 2026: targeted proof check. Also checked in this writing. Current writing events: GPT-6 Astra (OpenAI), ChatGPT Pro mode, effort Pro, S6\_10\_PAPERS, 18 November 2026: computation or formal certificate; targeted proof check. GPT-6 Astra (OpenAI), ChatGPT Pro mode, effort Pro, S6\_10\_PAPERS continuation, 19 November 2026: used and rechecked in this run. \emph{Recorded proof scope.} GPT-6 Astra (OpenAI, ChatGPT Pro mode), effort Pro, 2026-11-17 UTC; GPT-6 Astra (OpenAI, ChatGPT Pro mode, effort Pro),17November2026,/root; Claude Fable 5.2 (Anthropic), effort high, 16 November 2026; also checked in this writing by GPT-6 Astra (OpenAI), ChatGPT Pro mode, effort Pro, 18 November 2026}{\cref{maps:src:F.2.31}; the displayed affine actions and bar-chain maps}


Let $g(x)=Mx+v/m$ be a free cyclic affine action of actual order
$m\geq2$ on $T^r$, with $M\in\mathrm{GL}(r,\mathbb Z)$ and
$v\in\mathbb Z^r$.  For every specified degree bound $d$, there is
an explicit terminating integer algorithm producing chains and
boundary maps through degree $d+1$ for the pair
\[
 B=(D^2\times T^r)/\langle(s,x)\mapsto
 (e^{2\pi i/m}s,g(x))\rangle,
 \qquad Y=\partial B.                                      \tag{F.CA1}\label{maps:eq:F.CA1}
\]
The cap complex has at most $2^r$ generators in each degree and
the boundary complex at most $2^{r+1}$.  These are complexes of
finite rank in each degree; only their indicated finite prefixes
are required to compute $H_j(B)$, $H_j(Y)$ and $H_j(B,Y)$ for
$j\leq d$.  Their maps retain the affine translations and all
integral higher coherence terms.  In particular, the construction
does not replace an equivariant torus chain complex by its
homology without recording the correction terms.

\end{theorem}
\begin{proof}
\begin{enumerate}[label=\arabic*.]
\item We give the complete recursive algorithm.  All chains below are
finite integer dictionaries, so the displayed operations are
effective over $\mathbb Z$.

\item First consider the normalized bar resolution of $\mathbb Z^r$
over $R=\mathbb Z[t_1^{\pm1},\ldots,t_r^{\pm1}]$ and its Koszul
resolution $K=R\otimes\bigwedge^*\mathbb Z^r$.
The latter differential is contraction with $(t_i-1)_i$.
Write $t^a=t_1^{a_1}\cdots t_r^{a_r}$ and put
$S_n(t)=(t^n-1)/(t-1)$, interpreted as its finite integral
Laurent sum for positive, zero or negative $n$.
The following $\mathbb Z$-linear contraction of the augmented
Koszul complex fixes a completely explicit choice:
\[
 s_K(t^a e_I)=
 \sum_{j<\min I}
 t_{j+1}^{a_{j+1}}\cdots t_r^{a_r}S_{a_j}(t_j)
 e_j\wedge e_I,                                           \tag{F.CA2}\label{maps:eq:F.CA2}
\]
where $\min\varnothing=r+1$ and indices in $I$ are increasing.
The bar contraction is
$s_B(t^a[g_1|\cdots|g_q])=[a|g_1|\cdots|g_q]$; terms with a
zero bar entry are zero.  Set $p_0=1$ and define the
$R$-linear projection and homotopy on coefficient-one bar
generators recursively by
\[
 p_q=s_Kp_{q-1}d_B,
 \qquad
 h_q=s_B(1-ip_q-h_{q-1}d_B).                               \tag{F.CA3}\label{maps:eq:F.CA3}
\]
Here $i(e_I)$ is the alternating sum of the bars of the ordered
unit vectors in $I$.  The recursion decreases degree.  The
identities for the two contractions give, by induction,
$pi=1$, $dh+hd=1-ip$, $ph=hi=h^2=0$.
For example $s_K^2=s_B^2=0$, $ps_B=s_Kp$, and the alternating
unit-vector formula gives $pi=1$; substituting these into
\eqref{maps:eq:F.CA3} proves the three side identities in the same induction.
Augmenting $R\to\mathbb Z$ gives a strong chain contraction
of the ordinary torus bar complex to $\bigwedge^*\mathbb Z^r$.

\item Translations must enter before this contraction.  Let
$b_j=\sum_{a=0}^{j-1}M^av/m$ and let $p_j$ be its coordinatewise
fractional part.  The $m$ points $p_j$ are distinct because the
action is free.  Use the full fundamental groupoid of $T^r$
on these points: an arrow from object $a$ to object $b$ is an
integer vector $\lambda$, represented by the affine path from
$p_a$ to $p_b+\lambda$.  Composition adds the integer vectors.
The generator acts strictly by shifting objects and, for its
$s$th power, by
\[
 \lambda\longmapsto M^s\lambda+z_s(b)-z_s(a),
 \qquad z_s(a)=M^sp_a+b_s-p_{a+s}\in\mathbb Z^r.             \tag{F.CA4}\label{maps:eq:F.CA4}
\]
The subscripts on objects are taken modulo $m$.
The nerve of this groupoid is equivariantly homotopy equivalent
to the torus: the simplex with successive arrow labels
$\lambda_i$ maps affinely to the torus with lifted vertices
$p_{a_j}+\sum_{i\leq j}\lambda_i$.
Forgetting the objects contracts its chains to the ordinary
torus bar complex.  An explicit homotopy is the negative
alternating prism which, at each possible position, inserts
the zero-labelled arrow to object zero and replaces all
subsequent objects by zero.  Combining this with \eqref{maps:eq:F.CA3}
gives a contraction $(i,p,h)$ of the groupoid chains to
$\bigwedge^*\mathbb Z^r$ with the same side identities.

\item Use the normalized cyclic bar construction on this strict
groupoid action.  Its differential is $d_v+\delta$, where
$\delta$ is the ordinary group-bar boundary, multiplied by
$(-1)^q$ on vertical degree $q$; its first face applies \eqref{maps:eq:F.CA4}.
The vertical contraction transfers this to cyclic bar chains
with exterior coefficients by the finite formulas
\[
 I=(1+h\delta)^{-1}i,\qquad
 P=p(1+\delta h)^{-1},\qquad D=p\delta I.                  \tag{F.CA5}\label{maps:eq:F.CA5}
\]
Inverses mean their alternating geometric series.  On a chain
of cyclic bar degree $p$ each series has at most $p+1$ terms,
since $\delta$ lowers that degree.  Multiplication and the
contraction identities directly give $(d_v+\delta)I=ID$ and
$P(d_v+\delta)=DP$,
and a chain homotopy inverse.  Equivalently these are the
finite homological perturbation identities, here derived by
expanding the finite series.  The first term of $D$ is the
cyclic bar differential for the exterior action of $M$.
Every remaining term lowers cyclic bar degree by at least two.

\item For completeness, the second contraction used in the algorithm
is also explicit.  The cyclic periodic resolution has one
$\mathbb Z[C_m]$ generator in each degree, with differential
$g-1$ in odd degree and $1+g+\cdots+g^{m-1}$ in even degree.
On $g^a$, $0\leq a<m$, its augmented contraction is
$\sum_{j=0}^{a-1}g^j$ in even degree and is $1$ precisely when
$a=m-1$ in odd degree, zero otherwise.
Apply \eqref{maps:eq:F.CA3} with this contraction to compare it with the
cyclic bar resolution.  Tensor its coefficients with the
exterior action of $M$.  Transferring the higher terms of $D$
once more by \eqref{maps:eq:F.CA5} terminates: each higher term followed by
the horizontal homotopy decreases cyclic degree by at least
one.  The resulting basis in total degree $n$ is
$\{(p,I):p+|I|=n,\ I\subseteq\{1,\ldots,r\}\}$.
This proves the claimed $2^r$ bound and termination in every
specified degree.

\item \par\medskip\noindent\textbf{Explicit comparison with the periodic cyclic resolution.}\quad 
\label{maps:mapsF:repair:cyclic-inclusion}
Write $C_m=\langle g\rangle$, $R_m=\mathbb Z[C_m]$, and let
$e_n$ generate the periodic resolution in degree $n$.
Let $B_*$ be the normalized $R_m$-bar resolution, with the
insertion contraction $s_B$.  Besides the periodic contraction
$s_P$ just specified, fix the inclusion recursively by
\[
 i_0(e_0)=[],\qquad
 i_n(e_n)=s_B\bigl(i_{n-1}(d_Pe_n)\bigr),
 \quad n>0,
\]
and extend it $R_m$-linearly.  This is a finite formula.  More
explicitly, with every summation index between $1$ and $m-1$,
\[
\begin{split}
 i_{2j}(e_{2j})
   &=\sum_{a_1,\ldots,a_j}
       [g^{a_1}|g|\cdots|g^{a_j}|g],\\
 i_{2j+1}(e_{2j+1})
   &=\sum_{a_1,\ldots,a_j}
       [g|g^{a_1}|g|\cdots|g^{a_j}|g].
\end{split}
\]
The empty sum of indices gives $i_0(e_0)=[]$ and
$i_1(e_1)=[g]$.  The contraction identity for $s_B$ proves
$d_Bi=id_P$ by induction: the chain to which $s_B$ is applied
is a cycle, since the preceding comparison commutes with
differentials and $d_P^2=0$.

\item Define $p$ and $h$ on coefficient-one bar generators by the
recursion already used for the torus comparison, substituting
$s_P$ for $s_K$, and extend them $R_m$-linearly.  Their values
can be checked without an unspecified comparison theorem.
For $1\leq a_i<m$, put
$\epsilon(a,b)=\lfloor(a+b)/m\rfloor$.  The projection is
\[
\begin{split}
 p_{2j}[g^{a_1}|\cdots|g^{a_{2j}}]
   &=\prod_{\nu=1}^{j}\epsilon(a_{2\nu-1},a_{2\nu})\,e_{2j},\\
 p_{2j+1}[g^{a_1}|\cdots|g^{a_{2j+1}}]
   &=\left(\sum_{b=0}^{a_1-1}g^b\right)
       \prod_{\nu=1}^{j}\epsilon(a_{2\nu},a_{2\nu+1})\,e_{2j+1}.
\end{split}
\]
These expressions follow by induction from
$p_n=s_Pp_{n-1}d_B$: the terms other than the first bar face
are annihilated after the preceding contraction identity is
used.  In $pi(e_{2j})$ only the summand with every
$a_\nu=m-1$ survives.  The same observation, with initial
factor $1$, proves $pi(e_{2j+1})=e_{2j+1}$.
Thus $pi=1$ in every degree.

\item The recursion
$h=s_B(1-ip-hd_B)$ gives $d_Bh+hd_B=1-ip$ by induction.
Moreover $p s_B=s_Pp$ follows from the displayed projection
formulas, while $s_P^2=s_B^2=0$ and $s_Bi=0$ follow directly
from their definitions.  Substitution in the same recursion
gives $ph=hi=h^2=0$.  We therefore have the strong comparison
required for the second finite perturbation series, including
its inclusion, projection, homotopy and side identities.
After tensoring with the exterior coefficient module, the
same identities hold.  Every higher perturbation followed by
the horizontal homotopy lowers cyclic degree, so this added
description changes neither the termination proof nor the
$2^r$ generator bound.

\item The free action makes the Borel construction equivalent to
$T^r/C_m$, and radial contraction identifies this with $B$.
For $Y$ use exactly the same algorithm on $T^{r+1}$ with
matrix $\operatorname{diag}(1,M)$ and numerator $(1,v)$.
The projection forgetting its first circle is equivariant and
represents $Y\to B$ under the radial contraction.
More generally an equivariant integral matrix $L$ acts on
groupoid arrows as
$\lambda\mapsto L\lambda+k_b-k_a$, where
$k_a=Lp_a-p'_a\in\mathbb Z^{r'}$.
Transferring this strict bar map by $P'LI$ gives the required
integer chain map.  Its mapping cone computes the pair.
Every series is finite with the proved degree bound, and
integer kernels and Smith normal forms finish the calculation.
\end{enumerate}
\end{proof}


\emph{Executed certificates.}
The programs \path{code/fourfolds_bar_koszul.py} and
\path{code/fourfolds_affine_cap_chains.py} implement these
formulas, including the prime-subgroup freeness check of \ref{maps:src:F.2.31}.
The first checks all 281 abelian and 498 cyclic bar simplices
in its explicitly stated alphabets and degree ranges.
The second checks translated circles of orders two and three,
two affine presentations of the Klein bottle (one with a
nonfixed numerator), and the three-dimensional half-turn
manifold.  The separate rank-six program checks every degree
zero through seven for the order-five and order-four caps,
against the independent integral Wang sequence of their
five-torus mapping-torus descriptions.
The boundary-map program checks both translated circles and
the Klein bottle, including the full chain-map identities and
the integral Thom-shift homology of each cone.  These are
finite verification sets, not an asserted exhaustive test over
all input matrices; the terminating formulas prove generality.







\begin{theorem}[the marked Kuhn cusp pair]\label{maps:src:F.2.34}
\StatementRecord{maps:src:F.2.34}{proved}{Historical events: Claude Fable 5.2 (Anthropic), effort high, 16 November 2026: carried as background without fault found. GPT-6 Astra (OpenAI), ChatGPT Pro mode, effort Pro, S6\_09\_PAPERS, 17 November 2026: targeted proof check. Also checked in this writing. Current writing events: GPT-6 Astra (OpenAI), ChatGPT Pro mode, effort Pro, S6\_10\_PAPERS, 18 November 2026: computation or formal certificate; targeted proof check. GPT-6 Astra (OpenAI), ChatGPT Pro mode, effort Pro, S6\_10\_PAPERS continuation, 19 November 2026: used and rechecked in this run. \emph{Recorded proof scope.} GPT-6 Astra (OpenAI, ChatGPT Pro mode), effort Pro, 2026-11-17 UTC; GPT-6 Astra (OpenAI, ChatGPT Pro mode, effort Pro),17November2026,/root; Claude Fable 5.2 (Anthropic), effort high, 16 November 2026; also checked in this writing by GPT-6 Astra (OpenAI), ChatGPT Pro mode, effort Pro, 18 November 2026}{\cref{maps:src:F.2.33}; \cite[Theorem 4.1.1 and Lemmas 4.1.2--4.1.3, pp.~480--482]{ffAchingerOgus20}; \cite[Theorem 8.7, pp.~57--59, and Corollary 8.8, p.~59, attached source build]{manYang26}}


The fixed primitive rank-three Kuhn cusp admits an explicit
finite integer chain complex for its cap $W$, its boundary
$Y$, and the inclusion $Y\to W$.  It is obtained from a
category with 26 objects, 170 nonidentity morphisms, and
respectively $26,170,288,144$ strict flags of lengths
$0,1,2,3$.  The construction records the full integral
collapse map.  In particular it gives
\[
\begin{array}{c|rrrrrrrrr}
j&0&1&2&3&4&5&6&7&8\\\hline
\operatorname{rank}H_j(W)&1&3&9&9&10&3&1&0&0\\
\operatorname{rank}H_j(Y)&1&4&12&18&18&12&4&1&0\\
\operatorname{rank}H_j(W,Y)&0&0&1&3&10&9&9&3&1
\end{array}                                                \tag{F.CK1}\label{maps:eq:F.CK1}
\]
The only torsion in these three rows is
$\operatorname{Tor}H_3(Y)=\mathbb Z/2$.

\end{theorem}
\begin{proof}
\begin{enumerate}[label=\arabic*.]
\item Let $\mathcal K$ be the periodic Kuhn triangulation of
$\mathbb R^3$.  Normalize every simplex by translating its
lexicographically first vertex to zero.  The objects are its
simplex orbits; a morphism $\tau\to\sigma$ is every strict
translated face inclusion $\tau+w\subset\sigma$.
Keep distinct translated inclusions as distinct morphisms.
Composition adds the translation vectors.  Enumerating all
nonempty faces of the six chains in the unit cube gives
$(1,7,12,6)$ simplex orbits in dimensions zero through three,
and the displayed morphism and flag counts.  The category
nerve is the barycentric cell model of the dual decomposition
of $\mathbb R^3/\mathbb Z^3$; the direction of the face order
is reversed when interpreted as dual cells.

\item For an object $\tau$ let $L_\tau\subset\mathbb Z^3$ be the
primitive span of its vertex differences.  Its primitivity
follows from unimodularity of every Kuhn simplex.  Choose by
Smith reduction a quotient map
$Q_\tau:\mathbb Z^3\to\mathbb Z^{3-\dim\tau}$ with kernel
$L_\tau$, and an integral section $R_\tau$.
Define two strict diagrams of tori on this category:
\[
\begin{array}{c|c|c}
 &\text{object at }\tau&\text{map for }\tau+w\subset\sigma\\\hline
\mathcal W&T^{3-\dim\tau}&Q_\sigma R_\tau\\
\mathcal Y&T^4&S_w=\begin{pmatrix}1&0\\ w&I_3\end{pmatrix}.
\end{array}                                                \tag{F.CK2}\label{maps:eq:F.CK2}
\]
The first maps compose because $Q_\sigma$ vanishes on
$L_\tau$, and the second compose because translations add.
The maps $[\,0\ Q_\tau\,]$ form a natural transformation
$\mathcal Y\to\mathcal W$: the only additional term is
$Q_\sigma w=0$, since $w$ is a vertex of $\sigma$.

\item Here these diagrams describe the actual local topology.
The toric moment quotient on the dual cell to $\tau$ has
angular torus $T^3/(L_\tau\otimes\mathbb R/L_\tau)$;
passing to a face collapses exactly the next angular
subtorus.  Barycentrically subdividing the dual cells gives
the homotopy colimit of $\mathcal W$ for the central fibre.
For its real oriented logarithmic boundary, retain the
argument $\theta_0$ of the base parameter and the three
angular coordinates $\theta$ of the toric chart.  Translating
the chart by $w$ changes them to
$(\theta_0,\theta+w\theta_0)$, which is precisely $S_w$.
Thus the logarithmic boundary is the homotopy colimit of
$\mathcal Y$, and forgetting the logarithmic phases on each
stratum is exactly the natural transformation just given.
This can also be read directly in the unimodular charts
$x_0\cdots x_k=t$: the moment coordinates form the dual
simplex, the four logarithmic phases remain, and the ordinary
stratum retains their quotient by the vertex differences.
The finite face identifications give \eqref{maps:eq:F.CK2}, so there are
no further gluing constants in this fixed marked cap.

\item The Kuhn map with its divisor log structures is proper,
log smooth and vertical after restricting to a small disc.
Achinger--Ogus, Theorem~4.1.1, identifies its topological germ
with the mapping cylinder of this logarithmic collapse
\cite[Theorem 4.1.1 and Lemmas 4.1.2--4.1.3, pp.~480--482]{ffAchingerOgus20}.  Consequently the preceding
homotopy-colimit diagrams and natural transformation model
the pair $(W,Y)$ itself.  Fix the marking by letting $b_i$
be the category loop with translation label $e_i$.  A
homotopy-colimit edge with map $S_{e_i}$ has the relation
$b_i^{-1}\theta_0b_i=\theta_0+a_i$; equivalently
$\theta_0b_i\theta_0^{-1}=b_i+a_i$.  This specifies the
required primitive cusp
$U=\begin{pmatrix}I_3&I_3\\0&I_3\end{pmatrix}$.

\item \par\medskip\noindent\textbf{Hypotheses and marking of the logarithmic comparison.}\quad 
\label{maps:mapsF:repair:log-cylinder-scope}
Give the small disc its standard divisorial log structure and
the Kuhn total space the divisorial structure of the reduced
central fibre.  In a unimodular chart the characteristic
monoid map is
\[
 \mathbb N\longrightarrow\mathbb N^{k+1},\qquad
 1\longmapsto(1,\ldots,1).
\]
Both monoids are fine and saturated.  The group cokernel is
free, and the underlying chart is the monomial semistable
map $x_0\cdots x_k=t$ with additional smooth coordinates;
this is log smooth.  The quotient monoid is a group, since
the negative of each coordinate generator equals the sum of
the other generators after the diagonal generator is killed.
The morphism is therefore vertical.  The periodic toric
construction is proper over the sufficiently small disc.

\item These are precisely the hypotheses of
\cite[Theorem~4.1.1, pp.~480--482]{ffAchingerOgus20}.
Its proof also notes that exactness is automatic over the
standard log disc.  It identifies the germ with the mapping
cylinder of the logarithmic collapse.  This homeomorphism
depends on a trivialization; Lemma~4.1.3 fixes that
trivialization on the logarithmic central fibre.  We keep
the resulting radial boundary identification and its stated
angular and tropical generators.  This suffices to model
the pair and its inclusion up to marked homotopy.  It does
not replace the supplied cusp marking by another integral
basis.

\item We specify finite integer chains, rather than only the
homology spectral sequence of these diagrams.  Replace every
torus by its ordinary normalized bar complex.  The affine
simplex map from the bar construction to the torus is natural
for all the integer homomorphisms in \eqref{maps:eq:F.CK2}, so this replacement
does not change either homotopy colimit or its collapse map.
Take the simplicial replacement over all strict flags of the
finite category.  Its vertical differential is the torus bar
boundary; its horizontal differential composes adjacent
arrows, with the first face applying the diagram map, and
has the sign $(-1)^q$ in vertical degree $q$.
Apply the first contraction and the finite series \eqref{maps:eq:F.CA5}
object by object.  There are no flags longer than three,
so every series has at most four terms.  The resulting
finite basis consists of a strict flag beginning at $\tau$
and an exterior basis vector of its torus.  Transfer the
natural transformation by $P_{\mathcal W}FI_{\mathcal Y}$.
This proves both the chain-map identity and completeness of
the collapse map, including all integral coherence corrections.

\item The raw cap chain ranks in degrees zero through six are
\[
 (26,199,688,1231,1226,648,144),
\]
and the raw boundary ranks in degrees zero through seven are
\[
 (26,274,1124,2420,3010,2186,864,144).
\]
The program enumerates every object, every translated
incidence and every flag, verifies composition, checks
$d^2=0$ on every generator, and checks the collapse-chain
identity on every generator.  Integral cancellation of a
unit pivot is a chain homotopy equivalence: clear its row by
integer column operations, then remove the resulting
contractible two-generator summand.  The adjacent differential
updates are forced by $d^2=0$.  After respectively 2063,
4988 and 7087 such cancellations for the cap, boundary and
collapse cone, Smith reduction gives \eqref{maps:eq:F.CK1}.  All matrices
and the remaining differentials are saved by
\texttt{fourfolds\_kuhn\_cap\_chains.py}; the reduced cap and
relative differentials vanish, and the boundary has one
nonzero Smith factor, equal to two in degree four to three.
This computes the integral groups without a filtration
extension problem.
\end{enumerate}
\end{proof}


The finite-pair and cusp-pair chain models enter the global
gluing cone of \cref{maps:src:F.2.35}. Its maps to the common
torus-bundle complement retain the specified peripheral word
and clutching translation.

The periodic toric quotient used for the cusp is the construction in
Yang's theorem \emph{Higher-dimensional periodic quotient}, and its
Kuhn-grid Euler number is the formula in his corollary
\emph{The standard grid family}
\cite[Theorem 8.7, pp.~57--59, and Corollary 8.8, p.~59, attached source build]{manYang26}.
For the present primitive rank-three grid, this gives Euler number
$3!=6$.  The full cap and boundary complexes, their integral inclusion,
and the torsion groups are calculated in the proof above.







\begin{theorem}[complete integral homology of a marked gluing]\label{maps:src:F.2.35}
\StatementRecord{maps:src:F.2.35}{proved}{Historical events: Claude Fable 5.2 (Anthropic), effort high, 16 November 2026: carried as background without fault found. GPT-6 Astra (OpenAI), ChatGPT Pro mode, effort Pro, S6\_09\_PAPERS, 17 November 2026: targeted proof check. Also checked in this writing. Current writing events: GPT-6 Astra (OpenAI), ChatGPT Pro mode, effort Pro, S6\_10\_PAPERS, 18 November 2026: computation or formal certificate; targeted proof check. GPT-6 Astra (OpenAI), ChatGPT Pro mode, effort Pro, S6\_10\_PAPERS continuation, 19 November 2026: used and rechecked in this run. \emph{Recorded proof scope.} GPT-6 Astra (OpenAI, ChatGPT Pro mode), effort Pro, 2026-11-17 UTC; GPT-6 Astra (OpenAI, ChatGPT Pro mode, effort Pro),17November2026,/root; Claude Fable 5.2 (Anthropic), effort high, 16 November 2026; also checked in this writing by GPT-6 Astra (OpenAI), ChatGPT Pro mode, effort Pro, 18 November 2026}{\cref{maps:src:F.2.31}, \cref{maps:src:F.2.33}, \cref{maps:src:F.2.34}; the displayed complement and peripheral maps}


For every fully marked F.2 datum there is a terminating exact
integer algorithm computing every integral homology group of
the resulting eight-manifold, including every torsion invariant
factor.  The input consists of the finite homology matrices
$M_1,\ldots,M_k$, free affine numerators $v_i/m_i$ of their
actual orders, the integral clutch vector $\kappa$, and the
chosen primitive cusp marking
\[
 M_0=(M_1\cdots M_k)^{-1},\qquad
 M_0P=P\begin{pmatrix}I_3&I_3\\0&I_3\end{pmatrix},
 \qquad P\in\mathrm{GL}(6,\mathbb Z).                       \tag{F.GC1}\label{maps:eq:F.GC1}
\]
Here $k=2$ or $3$ and the genuine orders are at least two;
order-one sites are first eliminated as in \ref{maps:src:F.2.31}.
If a marking is not supplied, the algorithm makes and records
one by Smith reduction.  It does not identify that choice
with an unspecified prescribed marking.
The complete chain prefix through degree nine has at most
4003 generators in any degree.  No spectral-sequence
extension, saturation convention, orientation assumption or
unbounded search enters the homology calculation.

\end{theorem}
\begin{proof}
\begin{enumerate}[label=\arabic*.]
\item We supply both global boundary maps which were left outstanding
after \ref{maps:src:F.2.34}.  Let $F_k$ be the free group with ordered
generators $s_1,\ldots,s_k$, and let $M(w)$ be the product
specified by a reduced word $w$.  The torus-bundle complement
over the punctured sphere retracts to the corresponding bundle
over a graph.  Its fundamental group and multiplication are
\[
 \Gamma=\mathbb Z^6\rtimes F_k,
 \qquad (x,w)(y,u)=(x+M(w)y,wu).                            \tag{F.GC2}\label{maps:eq:F.GC2}
\]
The fibre and base are aspherical and the base has no
two-cells.  Hence a chosen section identifies the complement
up to homotopy with $B\Gamma$.  These section and peripheral
choices are fixed by the input clutch convention below.

\item \emph{A finite complement complex with its full bar projection.}
Use reduced words for $F_k$.  The tree resolution of the free
group has $K_0=\mathbb Z[F_k]$ and
$K_1=\bigoplus_i\mathbb Z[F_k]e_i$, with
$de_i=s_i-1$, and no higher terms.  Its explicit augmented
contraction sends a word to its signed path in the Cayley
tree: for each positive letter $s_i$ add its preceding prefix
times $e_i$, and for a negative letter $s_i^{-1}$ subtract
the prefix including that letter times $e_i$.
This is the Fox derivative, and its boundary is $w-1$.
Together with the bar insertion contraction it gives, by
the recursion \eqref{maps:eq:F.CA3}, a strong contraction of the free-group
bar resolution to this tree resolution.  Every recursion
decreases bar degree and every word is finite.

\item For a bar simplex $[(x_1,w_1)|\cdots|(x_n,w_n)]$ of $\Gamma$,
put $y_j=M(w_1\cdots w_{j-1})x_j$.
The exact Alexander--Whitney chain map to the torus/free-group
double bar complex is
\[
 \sum_{p+q=n}(-1)^{pq}
 [w_1|\cdots|w_p]\otimes[y_{p+1}|\cdots|y_n].              \tag{F.GC3}\label{maps:eq:F.GC3}
\]
Terms with an identity bar entry are zero.  The horizontal
first face acts on torus vectors by
$\rho(w)=M(w)^{-1}$, and the horizontal boundary has sign
$(-1)^q$.  Notice that $\rho$ is a right action:
$\rho(wu)=\rho(u)\rho(w)$.  This order is essential when
the matrices do not commute.  Formula \eqref{maps:eq:F.GC3} follows from
the diagonal bisimplicial bar construction; alternatively,
its chain-map identity follows by pairing adjacent interior
faces, with the first face using
$M(w_1)^{-1}y_j=M(w_2\cdots w_{j-1})x_j$.
The usual shuffle inverse is finite in each degree, so the
map is a chain equivalence.

\item Apply the torus contraction and \eqref{maps:eq:F.CA5}, then the free-group
tree contraction and the same finite perturbation formula.
Higher terms decrease group bar degree at least two, whereas
the tree complex only has degrees zero and one.  Thus the
resulting complement complex is
\[
 C_n(E)=\bigwedge^n\mathbb Z^6\ \oplus\
 \bigoplus_{i=1}^k\bigwedge^{n-1}\mathbb Z^6,
 \qquad
 d(e_i\otimes\xi)=(-1)^{n-1}
   (\bigwedge^{n-1}M_i^{-1}-1)\xi.                         \tag{F.GC4}\label{maps:eq:F.GC4}
\]
Its differential alone has this short form; its projection
from group bars retains the finite higher perturbation terms.
Those projection terms are used for every boundary map.

\item \emph{Finite affine boundary maps.}
For the $i$th finite cap use the orbit groupoid of \ref{maps:src:F.2.33}.
Its one-edge homotopy colimit for the affine generator is
the boundary mapping torus.  Vertically contracting it gives
a bounded complex with
$C_n(Y_i)=\bigwedge^n\mathbb Z^6\oplus
\bigwedge^{n-1}\mathbb Z^6$ and differential
$(-1)^{n-1}(\bigwedge^{n-1}M_i-1)$ on the second summand.
Its inclusion into the cyclic bar construction, followed by
the cap projection in \eqref{maps:eq:F.CA5}, is the full map $Y_i\to B_i$.
In particular this one-edge model avoids an unnecessary
periodic resolution of the seven-dimensional boundary.

\item Here is an explicit group-bar map for its other inclusion.
Send a torus-groupoid arrow labelled $\lambda$ to the
translation $x^\lambda\in\Gamma$.  With the integral offsets
$z_1(a)$ in \eqref{maps:eq:F.CA4}, put
$\eta_a=s_i^{-1}x^{z_1(a)}$.
For an arrow $f:a\to b$, the exact identity
\[
 \eta_a\,\psi(gf)=\psi(f)\,\eta_b                        \tag{F.GC5}\label{maps:eq:F.GC5}
\]
is the natural-transformation identity: it is just
$s_i^{-1}x^{M_i\lambda+z_1(b)}
 =x^\lambda s_i^{-1}x^{z_1(b)}$.
For a $q$-simplex, the edge map is $(-1)^q$ times the
alternating prism which inserts $\eta$ at every position,
uses the original arrows before the insertion and the
transformed arrows afterwards.  Its boundary is precisely
the signed difference of the two end maps, by \eqref{maps:eq:F.GC5}.
Compose this prism with the finite vertical lift of the
one-edge complex, then with \eqref{maps:eq:F.GC3} and the complement
projection.  This proves the full chain-map identity.
The offsets record the actual integer affine lift: replacing
$v_i$ by $v_i+m_i\lambda$ changes the offsets and the map,
even if the affine map on the torus has not changed.
An optional unimodular local basis simply applies its matrix
to every translation in \eqref{maps:eq:F.GC5}.

\item We make the distinction between a bar cell and the geometric
positive root explicit.  Put
$a_i=z_1(0)=\lfloor v_i/m_i\rfloor$, coordinatewise, and
let $\tau_i$ be the positive one-edge bar cell closed using
the zero-labelled return arrow to object zero.  With the
first-face convention above its geometric loop is
\[
 \tau_i=g_i^{-1}x^{a_i},\qquad
 F_E(\tau_i)=-e_i+M_i^{-1}a_i\in C_1(E),\qquad
 g_i=x^{a_i}\tau_i^{-1}.                                  \tag{F.GC5a}\label{maps:eq:F.GC5a}
\]
Thus the bar cell has the opposite peripheral orientation,
with the displayed integer correction when the numerator
has a nonzero integer part.  Formula \eqref{maps:eq:F.GC5} sends its
group element to $s_i^{-1}x^{a_i}$; consequently the
geometric positive root in the last expression of \eqref{maps:eq:F.GC5a}
maps to $s_i$ exactly.  This is the stated conversion between
the boundary and complement chain conventions.  In a local basis
$Q$, replace $a_i$ in the group and complement formulas by
$Q\lfloor Q^{-1}v_i/m_i\rfloor$.
Writing $\mathcal N_i=\sum_{j=0}^{m_i-1}M_i^j$, the cyclic
degree-two cell has boundary
$m_i\tau_i+w_i-\mathcal N_i a_i$.
Indeed $\tau_i^{m_i}=x^{\mathcal N_i a_i-w_i}$ follows
directly from $g_i^{m_i}=x^{w_i}$.
This independently checks the affine power relation and
the degree-one prism sign, including negative numerators.

\item There is an optional exact geometric shortcut.  If a verified
local integral basis gives $M_i=\operatorname{diag}(1,D)$ and
$v_i=\epsilon e_0$, $\epsilon\in\{1,-1\}$, eliminate the
first fibre generator from the cap group using
$g^{m_i}=u^\epsilon$.  The cap is the mapping torus with
group $\mathbb Z^5\rtimes_D\mathbb Z$.
The fibre vector $(u,a)$ maps to $x^a g^{\epsilon m_i u}$,
and the root generator maps to $g$.  Substituting this
homomorphism for $\psi$ in the same prism gives the boundary
map to this bounded geometric cap model.  The shortcut is
used only after checking the displayed integral normal form;
the general periodic construction remains available for
every other input.

\item To check the peripheral convention directly, the universal
cover of $Y_i$ has the root generator $g_i$ with linear part
$M_i$ and affine displacement $v_i/m_i$.  Its integral
translations and $g_i$ map to the same fibre translations
and $s_i$ in \eqref{maps:eq:F.GC2}.  The full base-circle translation obeys
$g_i^{m_i}=x^{w_i}t$, where $w_i$ is the power vector of
\ref{maps:src:F.2.31}.  Filling the disc kills $t$, giving exactly
$s_i^{m_i}=x^{w_i}$.  The straight paths between the chosen
orbit points give \eqref{maps:eq:F.GC5}, so the prism is this marked map,
including numerators with nonzero integer parts.

\item \emph{The marked cusp boundary map.}
Write $P=[A\ B]$ in \eqref{maps:eq:F.GC1}, and set
\[
 w_0=(s_1\cdots s_k)^{-1},\qquad
 c=w_0x^{-\kappa}\in\Gamma.
                                                               \tag{F.GC6}\label{maps:eq:F.GC6}
\]
Send the three angular generators to the columns of $A$,
the three tropical generators to the columns of $B$, and
the base-phase generator to $c$.  Relation \eqref{maps:eq:F.GC1} gives
$c b_i c^{-1}=b_i+a_i$ and all the other nilpotent boundary
relations.  Filling kills the angular generators and $c$,
so \eqref{maps:eq:F.GC6} gives exactly
$s_1\cdots s_kx^\kappa=1$, in agreement with \ref{maps:src:F.2.31}.

\item This homomorphism has a direct chain implementation on the
full Kuhn diagram.  For an incidence with translation label
$w$, assign the group element $x^{Bw}$.  A vertical phase
$(n,a)\in\mathbb Z\oplus\mathbb Z^3$ maps to $c^nx^{Aa}$.
The incidence shear is $S_w(n,a)=(n,a+n w)$, and
\[
 (c^nx^{Aa})x^{Bw}=x^{Bw}(c^nx^{A(a+n w)}).                \tag{F.GC7}\label{maps:eq:F.GC7}
\]
For a flag of $p$ incidence arrows and a vertical bar
simplex of length $q$, sum over all $(p,q)$ shuffles.
Each incidence step contributes its assigned group element;
each vertical step contributes its phase transformed by
all preceding incidence shears.  Give the shuffle the sign
$(-1)^{\#\{\text{incidence steps before vertical steps}\}}$.
Interior face cancellations use exactly \eqref{maps:eq:F.GC7}, and adjacent
incidence faces use addition of their translation labels.
This is therefore a chain map from the double bar complex
to $C_*(B\Gamma)$.  Precompose with the finite Kuhn lift of
\ref{maps:src:F.2.34} and postcompose with \eqref{maps:eq:F.GC3} and the complement
projection.  The other cusp map is the full logarithmic
collapse already computed in \ref{maps:src:F.2.34}.
Both the cusp boundary and the complement are aspherical,
so the displayed marked homomorphism also verifies the
homotopy class of this geometric inclusion.

\item \emph{The final finite integer calculation.}
Take the mapping cone of the just constructed maps
\[
 \bigoplus_{i=0}^k C_*(Y_i)\longrightarrow
 C_*(E)\oplus\bigoplus_{i=0}^k C_*(B_i),
 \qquad y\longmapsto(f_E(y),-f_B(y)),                     \tag{F.GC8}\label{maps:eq:F.GC8}
\]
where $B_0=W$ is the Kuhn cap.  Gluing collars makes this
the homotopy pushout model of the actual closed manifold.
Consequently its cone computes integral homology directly;
there is no remaining filtration extension to choose.

\item For clarity the exact finite size bound can be read off
before running any matrix calculation.  Put
$t_n=\sum_{q=0}^{\min(n,6)}\binom6q$, with binomial
coefficients outside their range interpreted as zero.
Let $w_n$ and $y_n$ be the raw cap and boundary ranks of
\ref{maps:src:F.2.34}, zero outside their listed ranges.  Then the cone rank
in degree $n$ is
\[
 \binom6n+k\binom6{n-1}+k t_n+w_n+
 k\binom7{n-1}+y_{n-1}.                                 \tag{F.GC9}\label{maps:eq:F.GC9}
\]
For $0\leq n\leq9$ its maximum is 3890 when $k=2$ and
4003 when $k=3$.  Every bar comparison and perturbation
series used to produce these matrices terminates by its
degree filtration.  Word reductions and affine orbit
enumerations are finite, with no imposed size cutoff.
Computing integer kernels followed by Smith normal forms
now gives $H_0,\ldots,H_8$; higher groups vanish by dimension.
Unit cancellations are optional exact integral chain
equivalences which substantially reduce the matrices before
Smith reduction.  This proves the algorithm and the bound.
\end{enumerate}
\end{proof}


\emph{Verification and implementation.}
The implementation is in \path{code/fourfolds_global_chains.py}
and \path{code/fourfolds_total_homology.py}.
The generic JSON driver \path{code/fourfolds_marked_gluing.py}
retains the complete supplied marking and affine data, computes
the full group structure of \ref{maps:src:F.2.31}, and independently compares
its abelianization with the chain calculation of $H_1$.
The Fox contraction is checked on all 156 explicitly listed
free-group bar simplices through degree three; the exact
Alexander--Whitney and finite-projection identities are
checked on 156 specified semidirect-product simplices with
two noncommuting rank-two matrices.
Four finite boundary comparisons using the independent
seven-dimensional quotient model have acyclic integral
comparison cones.  The more efficient one-edge model is
checked in six stated low-rank cases, including negative,
nonfixed, and nontrivial integer-part numerators; both its
complement comparison and its full Thom cone pass.
The program \path{code/fourfolds_peripheral_signs.py} also
checks \eqref{maps:eq:F.GC5a} and the full cyclic degree-two relation in
seven cases, including both actual rank-six caps.  It
reconstructs the positive geometric root as a group element
and checks that its complement bar projection is $+e_i$.
For the Kuhn boundary, every one of its 10048 raw generators
passes the chain-map identity in all degrees zero through
seven.  Its comparison to the corresponding one-generator
complement has an integrally acyclic cone, eliminated by
5088 unit cancellations.  Thus the comparison also checks
the nonzero boundary order-two torsion.
The full eight-dimensional integration for the original
order-$(5,4)$ survivor has passed.  Every marked boundary
column and the complete Mayer--Vietoris differential are
checked.  After 7372 integral unit cancellations, its Smith
homology is $\mathbb Z$ in degrees zero and eight,
$\mathbb Z^4$ in degree four, and zero otherwise, exactly
as in the independent calculation \ref{maps:src:F.2.6}.
The corresponding three-point run with logarithmic order
$m=2$ also passes, after 7470 integral unit cancellations.
It gives $H_1=H_6=\mathbb Z/18$,
$H_2=H_5=(\mathbb Z/2)^5$, $H_3=(\mathbb Z/2)^{10}$,
$H_4=\mathbb Z^4\oplus(\mathbb Z/2)^{10}$,
$H_0=H_8=\mathbb Z$ and $H_7=0$, exactly as in the
independent theorem \ref{maps:src:F.2.32}.  It reuses only the identical
first-two finite boundary maps: their group bars involve only
the fibre and their own peripheral generator, so adjoining
a third generator changes neither their columns nor their
target row indices.  The third cap, the new cusp word,
the full cone and every Smith group are recomputed.

A third full integration uses
$A=I_3$, $B=\operatorname{diag}(0,1,0)$ in \eqref{maps:eq:F.IV2},
actual affine orders $4,4$, numerators $a_1,a_2$, and zero
clutching.  These two linear actions reverse orientation.
Its 7345 integral unit cancellations give free ranks
\[
                  (1,0,4,11,10,1,6,3,0).
\]
The nonzero torsion groups are
\[
\begin{aligned}
 \operatorname{Tor}H_1=\operatorname{Tor}H_5
       &=(\mathbb Z/2)^3\oplus\mathbb Z/4,&
 \operatorname{Tor}H_2&=(\mathbb Z/2)^2,\\
 \operatorname{Tor}H_3=\operatorname{Tor}H_7
       &=\mathbb Z/2,&
 \operatorname{Tor}H_4&=(\mathbb Z/2)^4\oplus(\mathbb Z/4)^5.
\end{aligned}
\]
The independent van Kampen $H_1$ agrees.  Its Euler number
is six, its mod-two Betti vector is
$(1,4,10,14,20,14,10,4,1)$, and its integral top groups
have the required nonorientable form.  Thus this integration
also tests the orientation-reversing part of the stated scope.

The \texttt{--fresh} option bypasses all saved map checkpoints.
The free-bar and semidirect tests, all four periodic-boundary
comparisons, all six one-edge comparison/Thom tests, all seven
peripheral-sign checks, the complete 10048-generator cusp
comparison, and the cached generic survivor computation have
also been independently rerun successfully.  The
reduced-chain, van Kampen $H_1$, Euler and mod-two duality
checks have additionally been independently replayed on
all three complete integration outputs.  The general
algorithm is proved by the explicit chain equivalences and
finite operations above; these finite execution sets are
reported with their exact scope.







\begin{theorem}[simultaneous integral conjugacy for every input]\label{maps:src:F.2.36}
\StatementRecord{maps:src:F.2.36}{proved}{Historical events: Claude Fable 5.2 (Anthropic), effort high, 16 November 2026: carried as background without fault found. GPT-6 Astra (OpenAI), ChatGPT Pro mode, effort Pro, S6\_09\_PAPERS, 17 November 2026: targeted proof check. Also checked in this writing. Current writing events: GPT-6 Astra (OpenAI), ChatGPT Pro mode, effort Pro, S6\_10\_PAPERS, 18 November 2026: targeted proof check. GPT-6 Astra (OpenAI), ChatGPT Pro mode, effort Pro, S6\_10\_PAPERS continuation, 19 November 2026: used and rechecked in this run. \emph{Recorded proof scope.} Claude Fable 5.2 (Anthropic), effort high, 16 November 2026; also checked in this writing by GPT-6 Astra (OpenAI), ChatGPT Pro mode, effort Pro, 18 November 2026}{the hypotheses, notation and definitions stated here; \cite[Algorithm 8.3, p.~21]{rootBHJ22}; \cite[Algorithm 8.3 and its correctness proof, p.~21]{rootBHJ22}; \cite[Theorem 3.4, p.~6; Propositions 8.1--8.2, p.~21]{rootBHJ22}}


There is an unconditional exact algorithm which decides simultaneous
$\mathrm{GL}_6(\mathbb Z)$ conjugacy of any two marked matrix tuples of
F.2 and returns a conjugating matrix when one exists. The algorithm
does not impose a bound on their entries, on the affine multiplicities,
or on the integral conjugator. More generally it applies to any integral
tuple containing a matrix $1+N$ with $N^2=0$ and $\operatorname{rank}N=3$;
the primitive-image assumption is unnecessary for this statement.

\end{theorem}


\begin{proof}
\begin{enumerate}[label=\arabic*.]
\item We use column vectors and seek $P$ satisfying $PX_i=Y_iP$ for every
marked index, including the cusp index. Put
\[
\begin{aligned}
 H&=\{P\in\operatorname{Mat}_6(\mathbb Z):PX_i=Y_iP\ (\forall i)\},\\
 \mathcal O&=\{C\in\operatorname{Mat}_6(\mathbb Z):CX_i=X_iC\ (\forall i)\},
 \qquad E=\mathbb Q\mathcal O.
\end{aligned}                                                    \tag{F.IC1}\label{maps:eq:F.IC1}
\]
Both lattices are saturated integer kernels of explicitly displayed
linear equations in thirty-six unknowns. Thus Hermite or Smith
reduction computes full bases, rather than sublattices of finite index.

\item We first decide rational isomorphism. For a basis $h_1,\ldots,h_s$ of
$H$, form the polynomial $f(c)=\det(\sum c_jh_j)$. If it is zero,
there is no rational isomorphism. Otherwise successively substitute
one of $0,1,\ldots,6$ for each variable while keeping the remaining
polynomial nonzero. At each stage one of these seven values works,
because the degree in that variable is at most six. The final matrix
$D\in H$ is invertible over $\mathbb Q$. This is a terminating exact
test, including the negative case, and not a search for a lucky
invertible intertwiner.

\item The full lattice $L=D^{-1}H\subset E$ is a right $\mathcal O$-module.
Let $J$ be the Jacobson radical of $E$, let $B=E/J$, and use bars
for images under the quotient map. We next verify the splitting
hypothesis needed for the principal-ideal algorithm. In a rational
$N$-adapted basis,
\[
 N=\begin{pmatrix}0&I_3\\0&0\end{pmatrix},\qquad
 E\subseteq
 \left\{\begin{pmatrix}A&T\\0&A\end{pmatrix}:A,T\in
                                      \operatorname{Mat}_3(\mathbb Q)\right\}.
                                                               \tag{F.IC2}\label{maps:eq:F.IC2}
\]
Projection to $A$ has square-zero kernel and image an algebra
$S\subseteq\operatorname{Mat}_3(\mathbb Q)$. A nilpotent ideal is
contained in the radical, so $E/J\cong S/J(S)$.

\item Choose a composition series of the faithful $S$-module $\mathbb Q^3$.
The kernel of its action on the associated graded module lowers this
filtration and is therefore nilpotent. Conversely the radical kills
every simple factor. Thus this kernel is exactly $J(S)$ and every
simple factor of $S/J(S)$ acts nontrivially on a composition factor.
For a simple algebra $\operatorname{Mat}_n(D)$, where $D$ has centre
$K$ and degree $d$ over that centre, every nonzero simple module has
$\mathbb Q$-dimension $n d^2[K:\mathbb Q]$. If $D$ were
noncommutative this would be at least four, which is impossible in
$\mathbb Q^3$. Consequently
\[
 B\cong\prod_j\operatorname{Mat}_{n_j}(K_j),\qquad
 n_j[K_j:\mathbb Q]\le3.                                    \tag{F.IC3}\label{maps:eq:F.IC3}
\]
Every noncommutative factor is therefore
$\operatorname{Mat}_2(\mathbb Q)$ or
$\operatorname{Mat}_3(\mathbb Q)$; the remaining factors are number
fields of degree at most three. This proves hypothesis (H) of
\cite[Algorithm 8.3, p.~21]{rootBHJ22} for every tuple under consideration.

\item For clarity this decomposition is effective without a split-algebra
oracle. Rational linear algebra computes $J$ as the radical of the
regular trace pairing and computes the quotient multiplication table.
Its simple factors are obtained from its centre. A commutative factor
is a number field. In a noncommutative factor of dimension $n^2$, we
enumerate integral combinations $e$ of a rational basis until the
left ideal $Be$ has dimension $n$. Such an $e$ exists because the
factor is $\operatorname{Mat}_n(\mathbb Q)$, so its rank-one matrices
are rational points; clearing denominators expresses one among the
enumerated combinations. Left multiplication on $Be$ then supplies
an explicit isomorphism to $\operatorname{Mat}_n(\mathbb Q)$.
The program verifies multiplication on every pair of basis elements.

\item The images $\overline{\mathcal O}$ and $\overline L$ are a full order
and a full right lattice in $B$. If the right order of
$\overline L$ differs from $\overline{\mathcal O}$, this lattice
cannot be principal over that order. Otherwise Algorithm 8.3 of
\cite[Algorithm 8.3 and its correctness proof, p.~21]{rootBHJ22}, in its right-ideal form, decides principality and
returns a generator $\overline\alpha$ when one exists. The cited
algorithm passes through a maximal order and its central conductor,
tests local freeness and principality over the maximal order, and
solves the remaining reduced-norm congruence in a finite quotient.
Its correctness proof uses its Theorem 4.5(b) and Theorem 7.9; its
Theorem 8.4 specifies the number-field, finite-field and integer
factorization subproblems. These have unconditional exact algorithms.
In the recorded implementation the class and unit groups of every
number-field component are certified with \texttt{GRH=false} before
the principal-ideal decision is used. The negative decision thus
does not presume the generalized Riemann hypothesis.

\item If $\overline L$ is not principal, $L$ cannot be principal. If it is
principal, solve the integer linear equations expressing
$\overline\alpha$ as a combination of the images of the chosen basis
of $L$. This produces a lift $\alpha\in L$ and an integral
intertwiner $P=D\alpha\in H$. We claim that a simultaneous integral
conjugacy exists if and only if $|\det P|=1$.

\item The sufficient direction is immediate. For the necessary direction,
suppose $P_0\in H$ is unimodular. Every $h\in H$ has
$P_0^{-1}h\in\mathcal O$, so
$L=\alpha_0\mathcal O$, where $\alpha_0=D^{-1}P_0$.
Write $\alpha=\alpha_0u$ with $u\in\mathcal O$.
Because both image elements generate $\overline L$,
$\overline u$ is a unit of $\overline{\mathcal O}$.
Lift its inverse to $v\in\mathcal O$. Then $uv=1+j$ with
$j\in J\cap\mathcal O$, and the finite geometric series for
$(1+j)^{-1}$ shows that $u^{-1}=v(1+j)^{-1}\in\mathcal O$.
Thus $u$ is an integral unimodular matrix and
$P=P_0u$ is unimodular. This also proves that any chosen lift of the
image generator gives the same yes-or-no answer. It is the
right-module form of \cite[Theorem 3.4, p.~6; Propositions 8.1--8.2, p.~21]{rootBHJ22},
with the needed determinant condition written out.

\item Every branch either proves nonexistence or returns a matrix whose
integrality, determinant and all intertwining equations are verified
exactly. This completes the decision procedure.

\end{enumerate}
\end{proof}




\begin{theorem}[a complete canonical-representative algorithm]\label{maps:src:F.2.37}
\StatementRecord{maps:src:F.2.37}{proved}{Historical events: Claude Fable 5.2 (Anthropic), effort high, 16 November 2026: carried as background without fault found. GPT-6 Astra (OpenAI), ChatGPT Pro mode, effort Pro, S6\_09\_PAPERS, 17 November 2026: targeted proof check. Also checked in this writing. Current writing events: GPT-6 Astra (OpenAI), ChatGPT Pro mode, effort Pro, S6\_10\_PAPERS, 18 November 2026: targeted proof check. GPT-6 Astra (OpenAI), ChatGPT Pro mode, effort Pro, S6\_10\_PAPERS continuation, 19 November 2026: used and rechecked in this run. \emph{Recorded proof scope.} GPT-6 Astra (OpenAI, ChatGPT Pro mode), effort Pro, 2026-11-17 UTC; GPT-6 Astra (OpenAI, ChatGPT Pro mode, effort Pro),17November2026,/root; Claude Fable 5.2 (Anthropic), effort high, 16 November 2026; also checked in this writing by GPT-6 Astra (OpenAI), ChatGPT Pro mode, effort Pro, 18 November 2026}{the hypotheses, notation and definitions stated here; \cref{maps:src:F.2.38}; \cref{maps:eq:F.CN1}; \cref{maps:src:F.2.36}; \cref{maps:eq:F.CN2}; \cref{maps:src:F.2.16}}


Every marked linear tuple in F.2, for either $k=2$ or $k=3$,
has a computable canonical integral representative and a
returned simultaneous unimodular conjugator. Two inputs have
the same representative if and only if they are simultaneously
$\mathrm{GL}_6(\mathbb Z)$ conjugate. The algorithm uses a proved
finite bound depending on the input. No fixed global entry bound
can represent every class.

This is an effective classification, not a claim that its
infinite output has been printed in a finite table. The
ordered involution sector has the separate complete finite
normal form \ref{maps:src:F.2.38}, and all of its $56$ classes have been
computed. The general bounded fallback is potentially very
expensive; its finiteness is a theorem, not a runtime estimate.

\end{theorem}


\begin{proof}
\begin{enumerate}[label=\arabic*.]
\item We first remove the existential condition about affine lifts
from the admissibility test for a linear tuple. A finite-order
matrix $M$ admits a free finite affine lift on $T^6$ if and
only if $\ker(M-1)\ne0$. If this kernel is zero, $1-M$ is a
surjective endomorphism of the torus, so every affine map with
linear part $M$ has a fixed point. Conversely let the order of
$M$ be $d>1$ and choose a primitive $u\in\ker(M-1)\cap\mathbb Z^6$.
The map $x\mapsto Mx+u/d^2$ has actual order $d^2$: its
$d$th power is translation by $u/d$, of exact order $d$.
For each prime $p\mid d$, the generator of its order-$p$
subgroup is translation by $u/p$, which has no fixed point.
Every nontrivial stabilizer contains a prime-order subgroup,
so the whole cyclic action is free. If $M=1$, translation
by a primitive half-vector is a free lift of order two.
The invariant integral kernel is primitive, so such a
primitive vector is obtained by Smith reduction.

\item Finite order is also effectively bounded at this linear
stage. Its minimal polynomial is a product of distinct
cyclotomic polynomials, and eigenvalue one uses at least one
of the six dimensions. Every other cyclotomic factor therefore
has degree at most five and their total degree is at most
five. For a cyclotomic index $d$, each prime divisor $p$ has
$p-1\le5$, and its prime-power factors satisfy
$p^{a-1}(p-1)\le5$. Enumerating these factors and the possible
degree sums gives exactly
\[
                 d\in\{1,2,3,4,5,6,8,10,12\}.             \tag{F.CN1}\label{maps:eq:F.CN1}
\]
Here $d$ is the order of the matrix, the least common multiple
of its cyclotomic indices. Multiplicities of a factor can only
increase the dimension and introduce no new orders. Each
listed order occurs with a fixed vector after adding trivial
summands. Thus taking the powers in \eqref{maps:eq:F.CN1}, checking the
fixed kernel, the product relation, and the cusp Smith
factors gives an exact admissibility test. The unbounded
actual affine multiplicities are not bounded by \eqref{maps:eq:F.CN1}.

\item For an admissible full tuple $T=(M_1,\ldots,M_k,M_0)$ let
$H(T)$ be its maximum absolute entry. If a three-point linear
tuple has an identity at a finite position, temporarily delete
that identity, classify the remaining two-point tuple, and
reinsert the identity at the same marked position. There can
be only one such identity, since otherwise a finite matrix
would equal the inverse of the nontrivial unipotent cusp.
Deleting a fixed identity changes neither height nor the
lexicographic comparison of tuples in that conjugacy-invariant
sector. This is only a linear normalization: a nontrivial
affine order at the identity-linear site is retained in the
gluing datum. In the two-involution sector use the canonical
representative defined in \ref{maps:src:F.2.38}. This sector is invariant under simultaneous
conjugacy. In every other sector, order full integer tuples
first by $H$, then lexicographically by their row-major
entries in the displayed marked order. Define the canonical
representative to be the first admissible tuple in the
simultaneous conjugacy class.

\item To compute it, enumerate heights $1,\ldots,H(T)$ and the
integer entries at that height. There are at most
\[
                     (2H(T)+1)^{36(k+1)}                  \tag{F.CN2}\label{maps:eq:F.CN2}
\]
tuples in this complete box. The product relation allows
$M_0$ to be calculated from the first $k$ matrices, so the
implementation enumerates fewer entries. It also rejects
matrices with different characteristic polynomials or orders
before calling the conjugacy decision \ref{maps:src:F.2.36}. These filters
cannot reject an equivalent tuple. At least one comparison
succeeds: the supplied tuple itself belongs to its own box.
The first successful comparison returns the canonical
representative and its exact conjugating matrix.

\item Every simultaneous orbit has a least element in this order.
Its height is at most the height of any supplied member,
which proves that \eqref{maps:eq:F.CN2} loses no class representative for
that input. If two inputs are equivalent they therefore
return the same least element. Conversely the returned
conjugators show that equal representatives imply equivalent
inputs. The casewise involution normal form has the same
property by \ref{maps:src:F.2.38}, and an involution tuple cannot be
equivalent to a tuple in the other branch. This proves the
classification assertion, including its finite stopping rule.

\item Finally, \ref{maps:src:F.2.16} already gives an unbounded admissible integral
family. In its homology convention direct symbolic
multiplication gives
\[
             \operatorname{tr}(M_1M_2^3)
                       =-v^4+v^3-v^2+2.                 \tag{F.CN3}\label{maps:eq:F.CN3}
\]
This is unbounded even on the previously constructed free
affine subfamily $v\in20\mathbb Z$. It is invariant under
simultaneous conjugacy. A fixed finite entry box contains
only finitely many tuples and hence only finitely many such
traces; it cannot contain a representative of every class.
The input-dependent nature of \eqref{maps:eq:F.CN2} is therefore necessary
for this particular exhaustive procedure.

\end{enumerate}
\end{proof}




\begin{theorem}[the complete ordered involution stratum]\label{maps:src:F.2.38}
\StatementRecord{maps:src:F.2.38}{proved}{Historical events: Claude Fable 5.2 (Anthropic), effort high, 16 November 2026: carried as background without fault found. GPT-6 Astra (OpenAI), ChatGPT Pro mode, effort Pro, S6\_09\_PAPERS, 17 November 2026: targeted proof check. Also checked in this writing. Current writing events: GPT-6 Astra (OpenAI), ChatGPT Pro mode, effort Pro, S6\_10\_PAPERS, 18 November 2026: computation or formal certificate; targeted proof check. GPT-6 Astra (OpenAI), ChatGPT Pro mode, effort Pro, S6\_10\_PAPERS continuation, 19 November 2026: used and rechecked in this run. \emph{Recorded proof scope.} GPT-6 Astra (OpenAI, ChatGPT Pro mode), effort Pro, 2026-11-17 UTC; GPT-6 Astra (OpenAI, ChatGPT Pro mode, effort Pro),17November2026,/root; Claude Fable 5.2 (Anthropic), effort high, 16 November 2026; also checked in this writing by GPT-6 Astra (OpenAI), ChatGPT Pro mode, effort Pro, 18 November 2026}{the hypotheses, notation and definitions stated here; \cref{maps:eq:F.IV1}; \cref{maps:eq:F.IV3}; \cref{maps:eq:F.IV4}; \cref{maps:eq:F.IV2}; \cref{maps:src:F.2.31}}


There are exactly 56 simultaneous $\mathrm{GL}_6(\mathbb Z)$
conjugacy classes of ordered tuples $(M_1,M_2,M_0)$ satisfying
the primitive maximal-cusp conditions of F.2 with both finite
linear orders equal to two.  An explicit integral normalizer
returns a canonical representative and a checked unimodular
conjugating matrix for every such input, without any bound on
its entries.  Every representative has free affine lifts of
actual order four at both finite points.  None can give a
simply connected filling, for any choices of free affine
orders or twists.

\end{theorem}


\begin{proof}
\begin{enumerate}[label=\arabic*.]
\item First choose an integral cusp basis, so that
\[
 M_0=I_6+N,\qquad N=\begin{pmatrix}0&I_3\\0&0\end{pmatrix}.
                                                               \tag{F.IV1}\label{maps:eq:F.IV1}
\]
This is effective by Smith reduction.  In fact choose vectors
$\beta_i$ whose images $\alpha_i=N\beta_i$ form a primitive
basis of $\operatorname{im}N=\ker N$.
The index-one hypothesis says that their classes form a basis
of the quotient, so $[\alpha_1,\alpha_2,\alpha_3,
\beta_1,\beta_2,\beta_3]$ is unimodular.

\item Since $M_2=M_1(I-N)$ and $M_1^2=M_2^2=I$, we have
$M_1NM_1=-N$.  Solving this block equation gives precisely
\[
 M_1=\begin{pmatrix}A&B\\0&-A\end{pmatrix},\qquad
 M_2=\begin{pmatrix}A&B-A\\0&-A\end{pmatrix},\qquad
 A^2=I_3,\quad AB=BA.                                     \tag{F.IV2}\label{maps:eq:F.IV2}
\]
Every integral conjugator preserving \eqref{maps:eq:F.IV1} is
$\begin{pmatrix}C&D\\0&C\end{pmatrix}$, with
$C\in\mathrm{GL}_3(\mathbb Z)$ and $D$ arbitrary integral.
Thus one first classifies $A$ integrally; after fixing it,
$B$ is considered modulo $AD+DA$ and modulo conjugation by
the integral centralizer of $A$.

\item We include the integral involution classification needed here.
For the lattice $L=\mathbb Z^3$, let $L_+$ and $L_-$ be the
primitive kernels of $A-I$ and $A+I$.  They satisfy
$2L\subseteq L_+\oplus L_-\subseteq L$.
The quotient $L/(L_+\oplus L_-)$ is a vector space over
$\mathbb F_2$.  Its projections into $L_+/2L_+$ and
$L_-/2L_-$ are injective: a glue vector with only one
nonzero projection would give a new integral vector in the
corresponding eigensublattice, contradicting its definition.
It is consequently the graph of an isomorphism between two
subspaces of equal dimension $r$.
Unimodular changes of bases in $L_\pm$ put this graph in
coordinate form.  For each glued pair $e_+,e_-$ the vectors
$(e_++e_-)/2$ and $(e_+-e_-)/2$ form a swap block.
The unglued vectors give trivial and sign blocks.  Hence
there is a unique type
\[
 A\cong I_p\oplus(-I_q)\oplus
 \begin{pmatrix}0&1\\1&0\end{pmatrix}^{\oplus r},
 \qquad p+q+2r=3.                                         \tag{F.IV3}\label{maps:eq:F.IV3}
\]
The dimensions of the two rational eigenspaces and the
index $[L:L_+\oplus L_-]=2^r$ prove uniqueness.
Here $r$ is zero or one, so this normalization only requires
integer kernels, one Smith reduction and, in the glued case,
one explicit half-sum pair whose integrality is checked.

\item For fixed \eqref{maps:eq:F.IV3}, compute the quotient of the commuting
matrices by the subgroup $\{AD+DA\}$.
On the trivial and sign diagonal blocks the norm is
respectively $2D$ and $-2D$.  The commuting off-diagonal
blocks between the trivial and sign summands are zero.
Every commuting block involving a swap summand is a norm.
For example a trivial-to-swap block is a row $(a,a)$ and
is obtained as $(d,e)+(e,d)$; a sign-to-swap block has the
same calculation with a difference.  For two swap summands,
the commuting block is
$\begin{pmatrix}a&b\\b&a\end{pmatrix}$, and applying the
norm to $\begin{pmatrix}d&e\\f&g\end{pmatrix}$ gives
$\begin{pmatrix}e+f&d+g\\d+g&e+f\end{pmatrix}$.
The reverse mixed blocks have the transposed calculations.
Therefore the full integral quotient is exactly
\[
 \operatorname{End}_{\mathbb Z[A]}(L)/\{AD+DA\}
 \cong \operatorname{Mat}_p(\mathbb F_2)
       \oplus\operatorname{Mat}_q(\mathbb F_2).             \tag{F.IV4}\label{maps:eq:F.IV4}
\]
No finite-index sublattice has been substituted for a norm
image; the displayed block calculations are surjective
integer formulas.

\item The centralizer action on \eqref{maps:eq:F.IV4} is precisely independent
similarity in the two matrix factors.  Indeed, in a product
of commuting block matrices, every contribution factoring
through a swap block is a norm.  The remaining diagonal
blocks are just the reductions of the trivial and sign
blocks of the centralizer matrix.  If that matrix is an
integral unit, its inverse shows that both reductions are
invertible.  Conversely block-diagonal unimodular changes
on the trivial and sign summands induce every element of
$\mathrm{GL}_p(\mathbb F_2)\times
\mathrm{GL}_q(\mathbb F_2)$; elementary matrices give
unimodular lifts.  There is no additional identification
between the two factors, since an intertwiner of a trivial
and a sign lattice is zero over $\mathbb Z$.

\item It remains a genuinely finite classification.  Enumerate
all $2^{n^2}$ matrices of size $n=0,1,2,3$ over $\mathbb F_2$
and conjugate by every invertible matrix of that size.
Their respective numbers of similarity classes are
$1,2,6,14$.  Taking the smallest bit-coded representative
in each orbit gives a fixed canonical choice.
The six integral involution types and their counts are
\[
\begin{array}{c|rrrrrr}
 (p,q,r)&(0,3,0)&(1,2,0)&(2,1,0)&(3,0,0)&(0,1,1)&(1,0,1)\\\hline
 \text{classes}&14&12&12&14&2&2
\end{array}
\]
and their sum is 56.  A representative is \eqref{maps:eq:F.IV2} with
$B$ equal to the chosen $0$--$1$ representatives in its
trivial and sign blocks, and zero in every other block.

\item This also gives an effective normalizer for an arbitrary
input, not merely a list of representatives.  Normalize its
cusp and then its involution by the preceding Smith and
glue constructions.  Use the finite-field orbit lookup to
normalize its two reductions.  In dimensions at most three
the binary lift of an invertible binary matrix is already
unimodular: its determinant is odd and has absolute value
at most two.  The latter bound follows by adjoining a row
and column of ones and applying the $4\times4$ Hadamard
bound.  Finally solve the full integral linear equation
$AD+DA=B-B_{\rm can}$ by Smith reduction.  The preceding
norm calculation proves solvability.  Conjugating by
$\begin{pmatrix}I&D\\0&I\end{pmatrix}$ subtracts this norm
from the upper-right block.  The product of these explicitly
recorded conjugators is integral unimodular and carries all
three marked matrices to their canonical representative.

\item Both matrices in \eqref{maps:eq:F.IV2} have a three-dimensional fixed
space.  Choose a primitive integral fixed vector $v_i$.
Then $g_i(x)=M_ix+v_i/4$ has actual order four:
$g_i^2$ is the nontrivial translation $v_i/2$.
This order-two subgroup is free, and every nontrivial
subgroup of a cyclic group of order four contains it, so
the full action is free.  For any other finite affine lift
the actual order is still even, since its linear part has
order two.  The two-point orbifold quotient of \ref{maps:src:F.2.31} is
therefore the nontrivial cyclic group of order
$\gcd(m_1,m_2)$.  Simple connectivity is impossible for
every choice of twists and actual orders.

\end{enumerate}
\end{proof}




\begin{theorem}[the full effective classification and all marked gluings]\label{maps:src:F.2.39}
\StatementRecord{maps:src:F.2.39}{proved}{Historical events: Claude Fable 5.2 (Anthropic), effort high, 16 November 2026: carried as background without fault found. GPT-6 Astra (OpenAI), ChatGPT Pro mode, effort Pro, S6\_09\_PAPERS, 17 November 2026: targeted proof check. Also checked in this writing. Current writing events: GPT-6 Astra (OpenAI), ChatGPT Pro mode, effort Pro, S6\_10\_PAPERS, 18 November 2026: targeted proof check. GPT-6 Astra (OpenAI), ChatGPT Pro mode, effort Pro, S6\_10\_PAPERS continuation, 19 November 2026: used and rechecked in this run. \emph{Recorded proof scope.} GPT-6 Astra (OpenAI, ChatGPT Pro mode), effort Pro, 2026-11-17 UTC; GPT-6 Astra (OpenAI, ChatGPT Pro mode, effort Pro),17November2026,/root; Claude Fable 5.2 (Anthropic), effort high, 16 November 2026; also checked in this writing by GPT-6 Astra (OpenAI), ChatGPT Pro mode, effort Pro, 18 November 2026}{the hypotheses, notation and definitions stated here; \cref{maps:src:F.2.37}; \cref{maps:src:F.2.38}; \cref{maps:src:F.2.31}; \cref{maps:src:F.2.35}; \cref{maps:src:F.2.3}; \cref{maps:src:F.2.30}; \cref{maps:src:F.2.1}; \cref{maps:src:F.2.2}; \cref{maps:src:F.2.36}; \cref{maps:eq:F.CN3}; \cref{maps:src:F.2.33}; \cref{maps:src:F.2.7}}


The linear classification and all-twist topological computation
asked for in F.2 have the following complete effective answer.
For $k=2$ and $k=3$, every admissible ordered integral tuple
has the canonical representative of \ref{maps:src:F.2.37}, with an explicit
simultaneous unimodular conjugator.  Equality of these
representatives is necessary and sufficient for simultaneous
integral conjugacy.  There is no fixed global height cutoff;
the finite stopping bound for an input of height $H$ is
$(2H+1)^{36(k+1)}$ candidates before the proved rejection
filters.  The separate involution normal form has exactly
56 classes and is computed completely in \ref{maps:src:F.2.38}.

For every choice of actual free affine orders and integer
numerators, integral clutch vector, and specified primitive
Kuhn cusp marking, \ref{maps:src:F.2.31} and \ref{maps:src:F.2.35} compute the full
fundamental group and every integral homology group of the
gluing.  This is the fully marked correction to \ref{maps:src:F.2.3};
peripheral matrices are not substituted for a choice of
filling.  The algorithms have no cutoff on affine orders,
twist entries or lattice indices.  The complete list of data
with $\pi_1=1$ and $H_2=H_3=0$ is the constructive
parametrization of \ref{maps:src:F.2.30}, with its exact lattice enumeration
and necessary-and-sufficient finite tests.  Each retained
filling has the homology and oriented smooth type stated
there.

\end{theorem}


\begin{proof}
\begin{enumerate}[label=\arabic*.]
\item The input is finite, so the prime-subgroup test in \ref{maps:src:F.2.31}
decides whether every claimed affine cyclic action is free
and has the stated actual order.  The integral cusp test and
its adapted basis are \ref{maps:src:F.2.1}.  The complete linear admissibility
filter and the canonical representative, including the proved
finite stopping bound, are \ref{maps:src:F.2.2} and \ref{maps:src:F.2.36}--38.  The
characteristic-polynomial filters used by that algorithm are
conjugacy invariants, so none can remove a possible
representative.  The unbounded trace family \eqref{maps:eq:F.CN3} proves
that no fixed finite box could replace the input-dependent
bound.

\item We spell out how the linear classification is used with a
specified gluing.  If $R$ is its returned conjugator and
$P$ is the supplied cusp marking, change the data by
\[
 M_i'=RM_iR^{-1},\qquad v_i'=Rv_i,\qquad
 \kappa'=R\kappa,\qquad P'=RP.                           \tag{F.SY1}\label{maps:eq:F.SY1}
\]
The actual orders are unchanged.  The coordinate map
$x\mapsto Rx$ on the torus intertwines the finite quotient
charts, since it sends $M_ix+v_i/m_i$ to
$M_i'x'+v_i'/m_i$.  At the cusp, its coordinates in the
marked Kuhn model are unchanged:
$(P')^{-1}x'=P^{-1}x$.  Hence it intertwines every toric
chart, translated face identification and collapse map of
that fixed filling.  It also carries the clutch path to
the path with vector $\kappa'$.  These maps agree on the
collars and give a diffeomorphism of the two fully marked
gluings.  Thus homology can be computed in the original
coordinates, while the canonical data and their group
presentation are also returned.  Independently replacing
$P'$ by another Smith marking is not this comparison.

\item An actual order-one site is ordinary, even if its integer
numerator has been recorded.  Delete it by the explicit
comparison of F.2.31: if it is the $i$th site, replace
$\kappa$ by
$\kappa+(M_{i+1}\cdots M_k)^{-1}v_i$ before deletion.
The cusp matrix and its supplied marking are unchanged.
There cannot be an order-one site in an admissible two-point
tuple: its other finite matrix would be the inverse of a
nontrivial unipotent matrix.  Nor can a three-point tuple
have two such sites.  Thus this reduction leaves precisely
the genuine two- or three-point inputs of \ref{maps:src:F.2.35} and loses
no case from the specified input class.

\item For each remaining input, the exact unsaturated normal
closure, cyclic or triangle-group extension and exceptional
spherical Euler relation of \ref{maps:src:F.2.31} give the full fundamental
group.  The finite cap, complement and marked Kuhn chain
maps of \ref{maps:src:F.2.33}--35 give a finite integer Mayer--Vietoris cone
through degree nine.  Integer kernels and Smith reduction
therefore give all $H_0,\ldots,H_8$ with their complete
torsion factors.  This is a calculation for arbitrary
supplied twists, rather than a restriction to the affine
orders of one displayed example.  Its independent
abelianization check compares the resulting $H_1$ with the
van Kampen answer.

\item Finally \ref{maps:src:F.2.7}--30 prove both directions of the survivor
classification.  Genuine three-point orbifold quotients
exclude simple connectivity.  The two-point filters reduce
every survivor to the minimal order-$(5,4)$ spectrum; the
rational normal form, finite stable-lattice and extension
algorithms include every integral representation above it.
The universal affine normalization includes every twist,
and the six exact integral tests decide retention.
Conversely every retained datum satisfies all three
topological conditions by those tests.  \ref{maps:src:F.2.30} computes
its common integral homology and applies the almost-complex
and smooth-classification arguments to identify its oriented
type.  There is no residual untested rational-height bound
in this parametrization.  The canonical-representative
algorithm can remove duplicates from it if desired, but
no finite global enumeration is asserted or needed for
these complete per-input decisions.

\end{enumerate}
\end{proof}




\begin{theorem}\label{night:P1B.1}
\StatementRecord{night:P1B.1}{proved}{Historical events: GPT-6 Astra (OpenAI), ChatGPT Pro mode, effort Pro, S6\_06\_BLIND, 16--18 November 2026: targeted proof check. Also checked in this writing. Current writing events: GPT-6 Astra (OpenAI), ChatGPT Pro mode, effort Pro, S6\_10\_PAPERS, 18 November 2026: targeted proof check. GPT-6 Astra (OpenAI), ChatGPT Pro mode, effort Pro, S6\_10\_PAPERS continuation, 19 November 2026: used and rechecked in this run. }{Definitions and the displayed constructions}
Let $N$ be an endomorphism of $\Z^6$ with $N^2=0$, rank three,
and primitive image. There is a unimodular basis in which
\[
 I+N=\begin{pmatrix}I_3&I_3\\0&I_3\end{pmatrix}.
\]
Its centralizer in $\operatorname{GL}_6(\Z)$ consists exactly of
$\left(\begin{smallmatrix}P&Q\\0&P\end{smallmatrix}\right)$,
with $P\in\operatorname{GL}_3(\Z)$ and $Q\in\operatorname{Mat}_3(\Z)$.
\end{theorem}
\begin{proof}
\begin{enumerate}[label=\arabic*.]
\item The kernel of an integral homomorphism to a free group is
primitive. Both $\ker N$ and $\operatorname{im}N$ have rank three,
and the latter is contained in the former. Since the image is
primitive in $\Z^6$, their finite quotient is zero. Choose a basis
$e_1,e_2,e_3$ of this common group and choose $f_i$ with $Nf_i=e_i$.
The induced map $\Z^6/\ker N\to\operatorname{im}N$ is an
isomorphism, so the images of the $f_i$ form a quotient basis.
Thus $(e_1,e_2,e_3,f_1,f_2,f_3)$ is a unimodular basis and gives
the displayed matrix. This is also an algorithm: a Smith
decomposition gives the primitive kernel and solves the three
integral equations $Nf_i=e_i$.
Writing a potential centralizer element as
$\left(\begin{smallmatrix}P&Q\\R&S\end{smallmatrix}\right)$,
the equality $CN=NC$ gives $R=0$ and $S=P$.
Its determinant is $(\det P)^2$, so integral invertibility
is equivalent to $P\in\operatorname{GL}_3(\Z)$.

\end{enumerate}
\end{proof}




\begin{theorem}\label{night:P1B.2}
\StatementRecord{night:P1B.2}{proved}{Historical events: GPT-6 Astra (OpenAI), ChatGPT Pro mode, effort Pro, S6\_06\_BLIND, 16--18 November 2026: targeted proof check. Also checked in this writing. Current writing events: GPT-6 Astra (OpenAI), ChatGPT Pro mode, effort Pro, S6\_10\_PAPERS, 18 November 2026: targeted proof check. GPT-6 Astra (OpenAI), ChatGPT Pro mode, effort Pro, S6\_10\_PAPERS continuation, 19 November 2026: used and rechecked in this run. }{Definitions and the displayed constructions}
A finite-order integral linear action on $\R^6$ admitting a
free finite affine torus lift has characteristic polynomial
among exactly thirty-eight cyclotomic products. Their linear
orders are
\[
 1,2,3,4,5,6,8,10,12.
\]
Exactly nineteen of these polynomials admit a commuting complex
structure on $\R^6$. These are a preliminary spectral list;
no assertion about global gluings is inferred from the list.
\end{theorem}
\begin{proof}
\begin{enumerate}[label=\arabic*.]
\item The linear action has eigenvalue one. Otherwise $I-A$ is
invertible over $\R$, and the affine lift $x\mapsto Ax+b$
has a fixed point, obtained by solving $(I-A)x=b$ in $\R^6$
and projecting to the torus. A finite-order integral matrix is
semisimple and its characteristic polynomial is a product of
cyclotomic polynomials with nonnegative integer exponents.
The one-eigenvalue leaves at most five dimensions for the other
factors. A cyclotomic degree is even for every index greater
than two. The possible degrees are therefore one, two and four,
with indices respectively $1,2$; $3,4,6$; and $5,8,10,12$.
To verify that no index has been omitted, the formula
$\varphi(n)=\prod_{p^a\parallel n}p^{a-1}(p-1)$ with
$\varphi(n)\le4$ restricts $p$ to $2,3,5$, the powers to
$2^a$ with $a\le3$, $3^a$ with $a\le1$, and $5^a$ with
$a\le1$, and gives exactly those indices after multiplication.

\item The executed program \texttt{code/part1/spectral\_filter.py}
enumerates every exponent vector satisfying total degree six
and a positive exponent at $\Phi_1$, with bounds supplied by
these degrees. It writes all thirty-eight exponent vectors,
polynomials and least common multiples of their active indices.
Companion matrices in block diagonal form realize every listed
polynomial and order. A translation by a nonzero rational
fraction in a trivial eigenline makes a finite affine action
free: choose its denominator to be a multiple of the linear
order, and every nonidentity power has nonintegral translation
in that eigenline.

\item A commuting real complex structure requires even multiplicity
of each real eigenvalue $1$ and $-1$. Conversely, on these real
eigenspaces pair basis vectors into oriented complex lines;
on every conjugate nonreal eigenpair use its rotation complex
structure. This constructs a commuting complex structure.
The program checks these two parity conditions and obtains
nineteen products. Since the remaining cyclotomic factors have
even degrees and the total degree is six, the two real
multiplicities have the same parity. Thus the nineteen
complex-compatible products are exactly the orientation-preserving
ones in this preliminary list.

\end{enumerate}
\end{proof}


\section{Source obstructions and the punctured-base period decision}



\begin{proposition}[the previously open period decision]\label{maps:src:F.3.3}
\StatementRecord{maps:src:F.3.3}{proved}{Historical events: Claude Fable 5.2 (Anthropic), effort high, 16 November 2026: carried as background without fault found. GPT-6 Astra (OpenAI), ChatGPT Pro mode, effort Pro, S6\_09\_PAPERS, 17 November 2026: targeted proof check. Also checked in this writing. Current writing events: GPT-6 Astra (OpenAI), ChatGPT Pro mode, effort Pro, S6\_10\_PAPERS, 18 November 2026: targeted proof check. GPT-6 Astra (OpenAI), ChatGPT Pro mode, effort Pro, S6\_10\_PAPERS continuation, 19 November 2026: used and rechecked in this run. \emph{Recorded proof scope.} Claude Fable 5.2 (Anthropic), effort high, 16 November 2026; also checked in this writing by GPT-6 Astra (OpenAI), ChatGPT Pro mode, effort Pro, 18 November 2026}{\cref{maps:src:F.3.4}, \cref{maps:src:F.3.7}, \cref{maps:src:F.3.8}, \cref{maps:src:F.3.9}, \cref{maps:src:F.3.10}, \cref{maps:src:F.3.11}}


Theorem \ref{maps:src:F.3.11} now gives a positive answer for every topological
survivor of F.2: there is a holomorphic period family on the
punctured sphere, with nondegenerate imaginary part in every
local normalized period chart and a bounded holomorphic
correction at its primitive maximal cusp.  The proof is universal
and does not require that the remaining integral classification
be finite.  The proposed polarized abelian-surface source is
excluded by Theorem \ref{maps:src:F.3.4}; the universal construction instead
uses the explicit unpolarized LSAP domain and the relative
section principle in \ref{maps:src:F.3.8}--\ref{maps:src:F.3.11}.  Simultaneous good reduction
at every finite puncture is a stronger compactification condition
than the literal F.3 statement and is not inferred from this
argument.  It is already proved for the explicit compact
fourfold of F.4.

\end{proposition}
\begin{proof}
\begin{enumerate}[label=\arabic*.]
\item Every retained topological datum has determinant-one transport
and the marked primitive maximal cusp. Apply the explicit cusp graph
and relative holomorphic section construction of
\cref{maps:src:F.3.10}, \cref{maps:src:F.3.11}. Its target plane is transverse
to its conjugate, and its difference from the specified cusp model is
bounded in the displayed cusp coordinate. Quotienting the resulting
holomorphic rank-three vector bundle by its integral period lattice
gives the required punctured-base family, as in \cref{maps:src:F.3.9}.
\item The displayed map from a polarized source would supply the
semisimple exterior-square local system excluded in
\cref{maps:src:F.3.4}. The actual section construction instead uses
the nine spray directions and the relative section maps proved in
\cref{maps:src:F.3.8}, \cref{maps:src:F.3.10}. Its finite punctures remain
punctures; the theorem asserts precisely the bounded cusp conclusion
just obtained. The separately constructed explicit family has the
finite root-cap extensions proved in \cref{maps:src:F.3.7}.
\end{enumerate}
\end{proof}





\begin{theorem}[exclusion of the abelian-surface source]\label{maps:src:F.3.4}
\StatementRecord{maps:src:F.3.4}{proved}{Historical events: Claude Fable 5.2 (Anthropic), effort high, 16 November 2026: targeted proof check. GPT-6 Astra (OpenAI), ChatGPT Pro mode, effort Pro, S6\_09\_PAPERS, 17 November 2026: targeted proof check. Also checked in this writing. Current writing events: GPT-6 Astra (OpenAI), ChatGPT Pro mode, effort Pro, S6\_10\_PAPERS, 18 November 2026: targeted proof check. GPT-6 Astra (OpenAI), ChatGPT Pro mode, effort Pro, S6\_10\_PAPERS continuation, 19 November 2026: used and rechecked in this run. \emph{Recorded proof scope.} GPT-6 Astra (OpenAI, ChatGPT Pro mode), effort Pro, 2026-11-17; Claude Fable 5.2 (Anthropic), effort high, 2026-11-16; also checked in this writing by GPT-6 Astra (OpenAI), ChatGPT Pro mode, effort Pro, 18 November 2026}{the hypotheses, notation and definitions stated here; \cite[Definition 4.2.4 and Proposition 4.2.5, p.~44; Theorem 4.2.6, p.~45]{ffDeligne}; \cref{maps:src:F.3.1}}


The Engel construction proposed in F.3, starting with a polarized
algebraic family of abelian surfaces with a toric-rank-two cusp,
taking a fibrewise degree-zero line bundle, and quotienting its
punctured bundle by a lifted translation, cannot have $b_1=b_2=0$
with the stipulated rank-three square-zero Kuhn cusp.
The two sections need not be proportional.

\end{theorem}


\begin{proof}
\begin{enumerate}[label=\arabic*.]
\item Let $W$ be the rank-four rational homology local system of the
abelian surfaces.  It is semisimple by
\cite[Definition 4.2.4 and Proposition 4.2.5, p.~44; Theorem 4.2.6, p.~45]{ffDeligne}.
Its toric-rank-two cusp has
$N_W^2=0$ and $\ker N_W=\operatorname{im}N_W$ of dimension two.
Semisimplicity splits $W=W^\Gamma\oplus W'$.
The image of $N_W$ lies in $W'$ whereas $W^\Gamma$ lies in its
kernel.  The equality of image and kernel forces $W^\Gamma=0$.
The coinvariants vanish as well.

\item The rank-six rational homology system $T$ after the two Engel
operations has a natural filtration
\[
 0\subset L\subset E\subset T,\qquad
 L\simeq\mathbb Q,\quad E/L\simeq W,\quad T/E\simeq\mathbb Q.
\]
The lower line is the circle in the punctured line bundle and
the upper quotient is the lifted-translation generator.
No polarization on $T$ is assumed.
The triviality of $T/E$ gives $\operatorname{im}N_T\subset E$.
Maximal square-zero cusp monodromy gives
$\ker N_T=\operatorname{im}N_T$; thus every globally fixed vector
lies in $E$.  Since $W^\Gamma=0$, the fixed vectors of $E$ are
exactly $L$.  Hence $\dim T^\Gamma=1$.
The dual filtration, whose middle term is $W^*$, gives
$\dim(T^*)^\Gamma=1$ by the identical argument.

\item Consequently the two ranks $a,b$ in Theorem \ref{maps:src:F.3.1} are both one.
Its Euler-parity calculation applies without requiring self-duality
of $T$, and contradicts $b_1=b_2=0$.
This excludes the proposed algebraic abelian-surface source and,
more generally, every such filtration with semisimple middle term.
It does not exclude an arbitrary unpolarized torus period family.

\end{enumerate}
\end{proof}




\begin{theorem}[the unpolarized period domain has LSAP]\label{maps:src:F.3.8}
\StatementRecord{maps:src:F.3.8}{proved}{Historical events: Claude Fable 5.2 (Anthropic), effort high, 16 November 2026: carried as background without fault found. GPT-6 Astra (OpenAI), ChatGPT Pro mode, effort Pro, S6\_09\_PAPERS, 17 November 2026: targeted proof check. Also checked in this writing. Current writing events: GPT-6 Astra (OpenAI), ChatGPT Pro mode, effort Pro, S6\_10\_PAPERS, 18 November 2026: targeted proof check. GPT-6 Astra (OpenAI), ChatGPT Pro mode, effort Pro, S6\_10\_PAPERS continuation, 19 November 2026: used and rechecked in this run. \emph{Recorded proof scope.} GPT-6 Astra (OpenAI, ChatGPT Pro mode), effort Pro, 2026-11-17; Claude Fable 5.2 (Anthropic), effort high, 16 November 2026; also checked in this writing by GPT-6 Astra (OpenAI), ChatGPT Pro mode, effort Pro, 18 November 2026}{the hypotheses, notation and definitions stated here; \cite[Definition 2.1, p.~4]{IVVOkaOneV5}; \cite[Proposition 7.13 and preceding paragraph, pp.~23--24]{IVVOkaOneV5}; \cref{maps:eq:F.O3}; \cref{maps:eq:F.O4}}


Let $D_3$ be the component of
\[
 \{F\in\operatorname{Gr}(3,\mathbb C^6):F\cap\overline F=0\}
\]
containing the standard complex structure, so that
$D_3\cong GL^+(6,\mathbb R)/GL(3,\mathbb C)$.
At every point, $D_3$
is dominated by a spanning tube of nine complex lines.  In particular,
it has the local spray approximation property (LSAP) and is Oka--1.
These assertions concern the unpolarized domain.

\end{theorem}


\begin{proof}
\begin{enumerate}[label=\arabic*.]
\item Use the graph chart $F_Z=\operatorname{im}\binom Z{I_3}$.
The open set
\[
 \Omega=\{Z\in M_3(\mathbb C):\det(\operatorname{Im}Z)>0\}
 \tag{F.O1}\label{maps:eq:F.O1}
\]
lies in $D_3$ and contains $iI_3$.  For a real column $w$, write
$K_w x=w\times x$, put $q_w=w^tw$, and define
\[
 A_w=q_wI_3-ww^t+K_w.
 \tag{F.O2}\label{maps:eq:F.O2}
\]
Take the six vectors
\[
 e_1,\ e_2,\ e_3,\ e_1+e_2,\ e_1+e_3,\ e_2+e_3
\]
and the corresponding six matrices $A_w$, followed by the three
matrices $K_{e_1},K_{e_2},K_{e_3}$.  Call these nine matrices
$A_1,\ldots,A_9$.  Their symmetric parts are a basis of the
six-dimensional space of real symmetric matrices: the first three
give all diagonal matrices and the last three add the three distinct
off-diagonal entries.  The final three skew matrices form a basis
of the skew matrices.  Thus the nine directions are linearly
independent over $\mathbb R$ and also over $\mathbb C$.
For reference, the determinant of their nine row-major coordinate
columns is $-16$.

\item Each $A_w$ is real normal, is zero on $\mathbb Rw$, and has
eigenvalues $q_w\pm i\sqrt{q_w}$ on $w^\perp\otimes\mathbb C$.
For every real $s$ its two transverse squared singular values are
\[
 (1+q_ws)^2+q_ws^2
 =q_w(q_w+1)\left(s+\frac1{q_w+1}\right)^2
       +\frac1{q_w+1}\ge\frac13,
 \tag{F.O3}\label{maps:eq:F.O3}
\]
since $q_w\in\{1,2\}$; the remaining squared singular value is one.
For a pure skew direction $K_{e_j}$ all singular values of
$I_3+sK_{e_j}$ are at least one.  Their determinants are positive
at $s=0$ and never vanish, so they are positive for all real $s$.

\item Let $\Lambda=\bigcup_{j=1}^9(iI_3+\mathbb C A_j)$ and let
$B_{1/4}$ be the open Frobenius ball of radius $1/4$ in
$M_3(\mathbb C)$.  Then
\[
 T=\Lambda+B_{1/4}\subset\Omega.
 \tag{F.O4}\label{maps:eq:F.O4}
\]
Indeed, the imaginary part of $iI_3+tA_j+E$ is
$I_3+(\operatorname{Im}t)A_j+\operatorname{Im}E$.
Its unperturbed least singular value is at least $1/\sqrt3$,
whereas $\|\operatorname{Im}E\|_{\rm op}<1/4$.
It therefore stays invertible throughout the straight perturbation
from $E=0$, preserving the positive determinant.
The set $T$ is open and connected and is a union of affine
translates of the spanning tree $\Lambda$.  Thus it is a spanning
tube in the precise sense of \cite[Definition 2.1, p.~4]{IVVOkaOneV5},
not merely a neighborhood of one entire line.  Its inclusion in
$D_3$ is a local biholomorphism at $iI_3$.

\item The real group $GL^+(6,\mathbb R)$ acts holomorphically and
transitively on $D_3$.  Transporting this tube gives the same
domination at every point.  The observation preceding
\cite[Proposition 7.13 and preceding paragraph, pp.~23--24]{IVVOkaOneV5} gives LSAP from
this domination, and its part (b), or Theorem 2.2 of that paper,
gives Oka--1.  The exact program
\texttt{code/fourfolds\_oka\_tubes.py} checks all nine integer
matrices, their linear independence and the symbolic identity
\eqref{maps:eq:F.O3}; the displayed uniform bound proves \eqref{maps:eq:F.O4} for all complex
parameters and is not a finite positivity test.

\end{enumerate}
\end{proof}




\begin{theorem}[holomorphic periods on every open base]\label{maps:src:F.3.9}
\StatementRecord{maps:src:F.3.9}{proved}{Historical events: Claude Fable 5.2 (Anthropic), effort high, 16 November 2026: carried as background without fault found. GPT-6 Astra (OpenAI), ChatGPT Pro mode, effort Pro, S6\_09\_PAPERS, 17 November 2026: targeted proof check. Also checked in this writing. Current writing events: GPT-6 Astra (OpenAI), ChatGPT Pro mode, effort Pro, S6\_10\_PAPERS, 18 November 2026: targeted proof check. GPT-6 Astra (OpenAI), ChatGPT Pro mode, effort Pro, S6\_10\_PAPERS continuation, 19 November 2026: used and rechecked in this run. \emph{Recorded proof scope.} GPT-6 Astra (OpenAI, ChatGPT Pro mode), effort Pro, 2026-11-17; Claude Fable 5.2 (Anthropic), effort high, 16 November 2026; also checked in this writing by GPT-6 Astra (OpenAI), ChatGPT Pro mode, effort Pro, 18 November 2026}{the hypotheses, notation and definitions stated here; \cref{maps:src:F.3.8}; \cite[Proposition 7.13(d), p.~24]{IVVOkaOneV5}}


Let $B$ be an open connected Riemann surface and
$\rho:\pi_1(B)\to GL^+(6,\mathbb R)$ any representation.
Its associated flat $D_3$-bundle admits a holomorphic section.
More generally, every continuous section which is holomorphic
near a compact Runge set can be deformed through sections to a
holomorphic section, with arbitrary uniform approximation there
and any prescribed finite-order interpolation at finitely many
points of that set.  For an integral representation, this gives
a holomorphic family of complex three-tori with the prescribed
lattice local system and a zero section over $B$.

\end{theorem}


\begin{proof}
\begin{enumerate}[label=\arabic*.]
\item The flat bundle
$\mathscr D=\widetilde B\times_\rho D_3\to B$ is a holomorphic
fibre bundle: its transition maps are constant real linear
transformations, acting holomorphically on the Grassmannian.
Its fibre has LSAP by Theorem \ref{maps:src:F.3.8}, so the bundle map is an
Oka--1 map by \cite[Proposition 7.13(d), p.~24]{IVVOkaOneV5}.
The lifting property in Definition 7.7 of that paper, applied
to the identity map of $B$, gives precisely the stated section
deformation and approximation, including the finite jets.
This uses the section theorem, rather than an inference from
the Oka--1 property for maps to the fibre alone.

\item There is a continuous section to start with.  An open Riemann
surface has the homotopy type of a one-dimensional CW complex;
a bundle with connected fibre admits a section over its vertices
and edges, and the homotopy lifting property transports this
section to $B$.  Connectedness holds for the chosen component
$D_3$.  Taking the Runge set empty therefore proves existence.

\item For an integral representation let $\mathbb V_\mathbb Z$ be its
lattice local system and let $\mathcal F\subset
\mathbb V_\mathbb Z\otimes\mathcal O_B$ be the resulting
holomorphic rank-three subbundle.  Real transversality says that
projection identifies the real vector bundle with the underlying
real bundle of
$\mathcal Q=(\mathbb V_\mathbb Z\otimes\mathcal O_B)/\mathcal F$.
Hence the image of $\mathbb V_\mathbb Z$ is a full lattice in
every fibre of $\mathcal Q$.  Locally it is uniformly discrete
and acts by holomorphic translations.  The quotient
$\mathcal Q/\mathbb V_\mathbb Z\to B$ is consequently a smooth
proper holomorphic family of compact complex three-tori,
with its evident zero section and the required marking.

\end{enumerate}
\end{proof}




\begin{lemma}[retaining isolated extension points]\label{maps:src:F.3.10}
\StatementRecord{maps:src:F.3.10}{proved}{Historical events: Claude Fable 5.2 (Anthropic), effort high, 16 November 2026: carried as background without fault found. GPT-6 Astra (OpenAI), ChatGPT Pro mode, effort Pro, S6\_09\_PAPERS, 17 November 2026: targeted proof check. Also checked in this writing. Current writing events: GPT-6 Astra (OpenAI), ChatGPT Pro mode, effort Pro, S6\_10\_PAPERS, 18 November 2026: targeted proof check. GPT-6 Astra (OpenAI), ChatGPT Pro mode, effort Pro, S6\_10\_PAPERS continuation, 19 November 2026: used and rechecked in this run. \emph{Recorded proof scope.} GPT-6 Astra (OpenAI, ChatGPT Pro mode), effort Pro, 2026-11-17; Claude Fable 5.2 (Anthropic), effort high, 16 November 2026; also checked in this writing by GPT-6 Astra (OpenAI), ChatGPT Pro mode, effort Pro, 18 November 2026}{the hypotheses, notation and definitions stated here; \cref{maps:eq:F.O5}; \cite[Theorem 1.2, preprint p.~3]{ffForstnericPrezelj02}; \cite[Theorem 2.13, p.~9]{IVVOkaOneV5}}


Let $R$ be an open Riemann surface, $V\to R$ a holomorphic vector
bundle, and $\mathscr U\subset\operatorname{Gr}_r(V)$ an open
subset whose projection onto $R$ is surjective.  Let $S\subset R$
be finite.  Suppose that over $R\setminus S$ this projection is
a holomorphic fibre bundle with an LSAP fibre.  Let $K\subset R$
be a compact Runge set containing $S$ in its interior.
Every continuous section of $\mathscr U$ which is holomorphic
near $K$ can be deformed through sections to a holomorphic
section, with arbitrary uniform approximation on $K$ and
interpolation of any given finite jets at finitely many points
of its interior.  No fibre-bundle or Serre-fibration assertion
is required at the points of $S$.

\end{lemma}


\begin{proof}
\begin{enumerate}[label=\arabic*.]
\item Here are the details of the relative version of the LSAP argument
in \cite[Lemma 5.1, Theorem 6.1 and Proposition 7.13(d)]
{IVVOkaOneV5}.  They also specify why isolated extension
points cause no new analytic step.  Enlarge $K$ slightly inside
the domain of holomorphicity, retaining the Runge property.
Choose a strongly subharmonic Morse exhaustion of $R$ starting
with this compact set.  All subsequent boundary arcs, attached
discs and critical one-handles can and will be chosen disjoint
from $S$.  On them $\mathscr U$ is an ordinary holomorphic
bundle with LSAP fibre.

\item The global spray on the old compact set can be taken to be
vertical even though the fibres at $S$ differ.  Indeed, let $f$
be the current holomorphic section on a Stein neighborhood $W$
of that set.  Put $A$ equal to the finite interpolation set,
including $S$, and choose an ideal sheaf $\mathcal I$ vanishing
to one more than the prescribed jet orders at its points.
Cartan's theorem A supplies finitely many holomorphic sections
$T_1,\ldots,T_N$ of
$\operatorname{End}(V)|_W\otimes\mathcal I$ which generate
$\operatorname{End}(V)$ on the compact attaching arcs, after
shrinking $W$ if needed.  The formula
\[
 F(z,a)=\exp\left(\sum_{j=1}^N a_jT_j(z)\right)f(z)
 \tag{F.O5}\label{maps:eq:F.O5}
\]
defines a holomorphic spray of sections of the Grassmannian
bundle.  Its differential in $a$ surjects onto the vertical
tangent space at every attaching point, since
$\operatorname{End}(V_z)\to
\operatorname{Hom}(f(z),V_z/f(z))$ is onto.
For $a$ sufficiently small it has range in $\mathscr U$ on
a neighborhood of the old compact set, by openness and
compactness.  Every parameter value preserves all the chosen
jets.  Domination at $A$ is neither asserted nor used.

\item For a noncritical extension, subdivide the new annular regions
into small discs lying in bundle trivializations over
$R\setminus S$.  Restrict \eqref{maps:eq:F.O5} to a thin attaching disc.
LSAP of the fibre extends this spray across the attached disc,
approximating the old spray as closely as desired on the
intersection.  The spray gluing used in the proof of
\cite[Lemma 5.1, equations (5.6)--(5.8)]
{IVVOkaOneV5} now applies verbatim to vertical sprays:
the inverse-function step is needed only on the attaching
intersection, where \eqref{maps:eq:F.O5} is dominating.  It produces
reparametrizations of the two sprays and a glued section.
The old part remains in $\mathscr U$ because it is evaluated
in \eqref{maps:eq:F.O5}; the new part remains in $\mathscr U$ because the
LSAP spray takes values in the bundle fibre.  The fixed jets
survive the reparametrizations, being independent of the
spray parameters on the old side.

\item For a critical extension, first extend over the attached real
arc using the given continuous section and the homotopy already
constructed on the old compact set.  Apply Mergelyan
approximation with the same finite jets to a neighborhood of
the union of the old set and this arc, and then apply the
noncritical step.  To see explicitly that this approximation
is by sections, trivialize $V$ holomorphically on the open
Riemann surface: its underlying complex bundle is trivial over
a graph spine, and the frame bundle therefore has a continuous
section.  Its fibre $GL(\operatorname{rank}V,\mathbb C)$ has
the dominating spray $(g,A)\mapsto g\exp A$, so
\cite[Theorem 1.2, preprint p.~3]{ffForstnericPrezelj02} gives a holomorphic
frame.  The section is then a map to the
fixed Grassmannian, to which
\cite[Theorem 2.13, p.~9]{IVVOkaOneV5} applies.  The base
coordinate is kept unchanged.  Close approximation on the
compact arc and old set, followed by shrinking its
neighborhood, preserves membership in the open subset
$\mathscr U$.  This is the critical case of Theorem 6.1 of
that paper, with no approximation performed at a singular
point of the projection: here the projection is a submersion
everywhere, as $\mathscr U$ is open in a Grassmannian bundle.

\item Choose successive errors summably small on the preceding
compact sets and smaller than their distances to the complement
of $\mathscr U$.  The sections converge uniformly on every
compact set to a holomorphic section with range in
$\mathscr U$.  The finite jets persist.  The close homotopies
on successive stages concatenate in the usual exhaustion
construction.  All new attachment regions lie outside $S$;
the old section remains holomorphic on a full neighborhood of
each point of $S$ at every stage.  This proves the lemma.

\end{enumerate}
\end{proof}




\begin{theorem}[the bounded-cusp period decision]\label{maps:src:F.3.11}
\StatementRecord{maps:src:F.3.11}{proved}{Historical events: Claude Fable 5.2 (Anthropic), effort high, 16 November 2026: carried as background without fault found. GPT-6 Astra (OpenAI), ChatGPT Pro mode, effort Pro, S6\_09\_PAPERS, 17 November 2026: targeted proof check. Also checked in this writing. Current writing events: GPT-6 Astra (OpenAI), ChatGPT Pro mode, effort Pro, S6\_10\_PAPERS, 18 November 2026: targeted proof check. GPT-6 Astra (OpenAI), ChatGPT Pro mode, effort Pro, S6\_10\_PAPERS continuation, 19 November 2026: used and rechecked in this run. \emph{Recorded proof scope.} GPT-6 Astra (OpenAI, ChatGPT Pro mode), effort Pro, 2026-11-17; Claude Fable 5.2 (Anthropic), effort high, 16 November 2026; also checked in this writing by GPT-6 Astra (OpenAI), ChatGPT Pro mode, effort Pro, 18 November 2026}{the hypotheses, notation and definitions stated here; \cref{maps:src:F.2.1}; \cref{maps:eq:F.O7}; \cref{maps:src:F.3.8}; \cref{maps:src:F.3.10}; \cref{maps:src:F.2.7}; \cref{maps:src:F.2.15}}


Let
$B^*=\mathbb P^1\setminus\{p_0,p_1,\ldots,p_k\}$,
where $k\ge1$, and let
$\rho:\pi_1(B^*)\to GL^+(6,\mathbb R)$ be a representation.
Assume that the monodromy $M_0$ at $p_0$ satisfies
$(M_0-I)^2=0$ and $\operatorname{rank}(M_0-I)=3$.
Then there is a holomorphic period family with this monodromy
and nondegenerate real transversality everywhere on $B^*$,
equivalently a nondegenerate imaginary part in each local
normalized period chart.
In a marking with $M_0=\left(\begin{smallmatrix}I&I\\0&I\end{smallmatrix}\right)$
its cusp period is
\[
 Z(t)=\frac{\log t}{2\pi i}I_3+C(t),
 \qquad C(t)\text{ holomorphic at }t=0.
 \tag{F.O6}\label{maps:eq:F.O6}
\]
Any prescribed finite jet of $C$ at zero may be imposed.
For an integral primitive cusp the marking can be integral.
Consequently every topological survivor of F.2 admits the
period family requested literally in F.3.  No good-reduction
condition at the other, still punctured, points is included
in this assertion.

\end{theorem}


\begin{proof}
\begin{enumerate}[label=\arabic*.]
\item Write $N=M_0-I$ and choose a real basis with
$N=\left(\begin{smallmatrix}0&I\\0&0\end{smallmatrix}\right)$.
Its integral version is Theorem \ref{maps:src:F.2.1}.  Choose the component
of the period domain selected by the positive logarithmic
model in \eqref{maps:eq:F.O6}.  This choice matters only when an orientation
has been fixed independently of the marking; every element
of $GL^+(6,\mathbb R)$ preserves either component.

\item Retain the cusp as a point of the \emph{open} surface
\[
 B=\mathbb P^1\setminus\{p_1,\ldots,p_k\}.
\]
The flat complex vector bundle on $B^*$ has a holomorphic
extension $V_D$ across $p_0$: in a local coordinate $t$ use
the single-valued frame obtained by the change
$\exp(-sN)$, where $s=(\log t)/(2\pi i)$.
Changing the logarithm by $2\pi i$ is cancelled by $M_0$.
This specifies the extension directly, without a growth
assumption on a period section.

\item Over $B^*$ let $\mathscr U$ be the flat bundle of real-transverse
planes in the chosen component.  At $p_0$ adjoin only the
planes transverse to $\ker N_\mathbb C$, an affine space of
matrices $C\in M_3(\mathbb C)$.  This gives an \emph{open}
subset $\mathscr U\subset\operatorname{Gr}_3(V_D)$.
Indeed, in the chart $F_C=\operatorname{im}\binom C{I_3}$
the actual flat period is
$\exp(sN)F_C=F_{sI_3+C}$.  For $t\ne0$ membership is exactly
\[
 \det\bigl(y(t)I_3+\operatorname{Im}C\bigr)>0,
 \qquad y(t)=-\frac{\log|t|}{2\pi}.
 \tag{F.O7}\label{maps:eq:F.O7}
\]
For every fixed $C_0$, and uniformly for $C$ in a sufficiently
small bounded neighborhood of $C_0$, the matrix in \eqref{maps:eq:F.O7} is
invertible with positive determinant when $|t|$ is sufficiently
small.  This follows, for example, from
$y(t)>\|\operatorname{Im}C\|_{\rm op}$ and a straight-line
homotopy to $y(t)I_3$.  Thus a full neighborhood of
$(0,C_0)$ lies in the adjoined set, proving openness.
The projection is consequently a holomorphic submersion,
although its fibre at $p_0$ is $\mathbb C^9$ rather than $D_3$.

\item Start with a polynomial $C_0(t)$ realizing the desired finite
jet.  Formula \eqref{maps:eq:F.O7} makes it a holomorphic section on a small
closed cusp disc $K$, which is Runge in $B$.
This section extends continuously over $B$.  To justify the
relative extension, $D_3$ is connected and simply connected:
it retracts onto $SO(6)/U(3)$, and
$\pi_1(U(3))\to\pi_1(SO(6))$ is onto, since a unitary rotation
in one complex coordinate is a generator of the real
two-plane rotation loop.  A cellular decomposition of the
surface outside the cusp disc now extends its boundary
section over edges by connectedness and over two-cells by
simple connectivity.  The flat bundle is trivial over each
cell.  This argument uses no fibration property at $p_0$.

\item Over $B\setminus\{p_0\}$ the fibre is the LSAP manifold
$D_3$ from Theorem \ref{maps:src:F.3.8}.  Lemma \ref{maps:src:F.3.10}, with $S=\{p_0\}$,
produces a holomorphic section of $\mathscr U$ over all of
$B$, with the chosen jet at $p_0$.  Since its value at $p_0$
lies in the affine chart just adjoined, it has a holomorphic
matrix $C(t)$ there.  Undoing the canonical frame gives
exactly \eqref{maps:eq:F.O6}, with a bounded holomorphic correction.
Everywhere else membership in $\mathscr U$ is precisely
nondegeneracy of the complex-torus periods.
More explicitly, choose three members of a real lattice basis
whose projected periods are complex-linearly independent, and
normalize them to $I_3$.  The remaining three periods form a
matrix $Z$, and real transversality is exactly
$\det\operatorname{Im}Z\ne0$.  This normalization persists
locally.  A single Schubert chart is not imposed globally,
and none is needed to define a marked torus period family.

\item For F.2 survivors, Theorem \ref{maps:src:F.2.7} excludes $k=3$ and
Proposition \ref{maps:src:F.2.15} proves orientation preservation of the
finite matrices for $k=2$.  The cusp determinant is also one.
Thus all survivors satisfy the hypotheses.  This is a
universal existence decision and does not require completion
of their simultaneous integral classification.

\end{enumerate}
\end{proof}



\clearpage
\begin{thebibliography}{16}
\providecommand{\natexlab}[1]{#1}
\providecommand{\url}[1]{\texttt{#1}}
\expandafter\ifx\csname urlstyle\endcsname\relax
  \providecommand{\doi}[1]{doi: #1}\else
  \providecommand{\doi}{doi: \begingroup \urlstyle{rm}\Url}\fi

\bibitem[Achinger and Ogus(2020)]{ffAchingerOgus20}
Piotr Achinger and Arthur Ogus.
\newblock Monodromy and log geometry.
\newblock \emph{Tunisian Journal of Mathematics}, 2\penalty0 (3):\penalty0
  455--534, 2020.
\newblock \doi{10.2140/tunis.2020.2.455}.
\newblock URL \url{https://msp.org/tunis/2020/2-3/tunis-v2-n3-p01-p.pdf}.
\newblock Verified in this writing on 18 November 2026 at these loci only.
  Theorem 4.1.1, pp. 480–482;\allowbreak{} Lemmas 4.1.2–4.1.3, pp. 481–482, Tunisian
  J. Math. 2(3) (2020). Scope: Downloaded the original journal PDF and read the
  full theorem, hypotheses and proof pp.480–482. Exactness is automatic over
  the standard log disc. Lemma4.1.3 fixes the trivialization on the log central
  fibre;\allowbreak{} the mapping-cylinder homeomorphism is not canonical. Applied to the
  explicit diagonal fine saturated monoid chart and retained angular/tropical
  marking. Local PDF SHA256
  6546ad0b27fa27386b283b172858a4d95d6f4b317f22749a7caf186b3034c0ba Author/title
  and edition matched to the opened item. Verification applies only to its
  recorded locator and edition;\allowbreak{} aliasing does not transfer preprint pagination
  to a published edition.

\bibitem[Alarc{\'o}n and Forstneri{\v c}(2025)]{IVVOkaOneV5}
Antonio Alarc{\'o}n and Franc Forstneri{\v c}.
\newblock {Oka}-1 manifolds.
\newblock Version 5, 10 April 2025, 2025.
\newblock URL \url{https://arxiv.org/pdf/2303.15855v5}.
\newblock Verified in this writing on 18 November 2026 at these loci only.
  arXiv:2303.15855v5 (10 April 2025): Definition 2.1 p. 4;\allowbreak{} Theorem 2.13 p. 9;\allowbreak{}
  Lemma 5.1 pp. 14–17;\allowbreak{} equations (5.6) p. 16 and (5.7)–(5.8) p. 17;\allowbreak{} Theorem
  6.1 pp. 17–19;\allowbreak{} Proposition 7.13 pp. 23–24;\allowbreak{} Definition 7.7, p.22
  (additionally opened for F3.9). Scope: Downloaded the displayed v5 PDF and
  read the relevant full statements and spray construction, including its
  noncritical/critical stages. Locators refer explicitly to this preprint, not
  journal pagination. Author/title and edition matched to the opened item.
  Verification applies only to its recorded locator and edition;\allowbreak{} aliasing does
  not transfer preprint pagination to a published edition.

\bibitem[{Anthropic}(2026)]{AnthropicEffort2026}
{Anthropic}.
\newblock Effort.
\newblock Claude Platform official documentation, accessed 17 November 2026,
  2026.
\newblock URL
  \url{https://platform.claude.com/docs/en/build-with-claude/effort}.
\newblock Verified in this writing on 18 November 2026 at these loci only. How
  effort works and Effort levels;\allowbreak{} lines138–155.. Scope: Response-token effort
  control, including thinking.

\bibitem[Bley et~al.(2022)Bley, Hofmann, and Johnston]{rootBHJ22}
Werner Bley, Tommy Hofmann, and Henri Johnston.
\newblock Computation of lattice isomorphisms and the integral matrix
  similarity problem.
\newblock \emph{Forum of Mathematics, Sigma}, 10:\penalty0 e87, 2022.
\newblock \doi{10.1017/fms.2022.74}.
\newblock URL \url{https://doi.org/10.1017/fms.2022.74}.
\newblock Verified in this writing on 18 November 2026 at these loci only.
  Hypothesis (H) p. 2;\allowbreak{} Theorem 3.4 p. 6;\allowbreak{} Propositions 8.1–8.2 and Algorithm
  8.3 with correctness p. 21;\allowbreak{} Theorem 8.4 pp. 22–23. Scope: Opened full
  radical reduction and integral matrix-similarity algorithm. In F.2.36 the
  semisimple quotient is explicitly a product of matrix algebras over number
  fields, so (H) is proved for the actual input. GRH appears only in optional
  complexity estimates, not the existence/termination statement used here.
  Author/title and edition matched to the opened item. Verification applies
  only to its recorded locator and edition;\allowbreak{} aliasing does not transfer preprint
  pagination to a published edition.

\bibitem[Conway et~al.()Conway, Crowley, Powell, and Sixt]{IVVWallForm}
Anthony Conway, Diarmuid Crowley, Mark Powell, and J{\"o}rg Sixt.
\newblock Simply-connected manifolds with large homotopy stable classes.
\newblock Author version, file large-homotopy-stable-class-v3.pdf.
\newblock URL
  \url{https://www.maths.gla.ac.uk/~mpowell/large-homotopy-stable-class-v3.pdf}.
\newblock Verified in this writing on 18 November 2026 at these loci only.
  Definition 3.6 and Theorem 3.7(1), pp. 12–14. Scope: Opened Definition 3.6,
  full Theorem 3.7(1), and normal-invariant setup pp. 11–12. Supplies
  almost-diffeomorphism from the Wall form;\allowbreak{} current proof spells out the
  Euler/p1 map on pi3(SO4). Author/title and edition matched to the opened
  item. Verification applies only to its recorded locator and edition;\allowbreak{} aliasing
  does not transfer preprint pagination to a published edition.

\bibitem[{Creative Commons}(2026)]{CC01Universal}
{Creative Commons}.
\newblock Cc0 1.0 universal.
\newblock Official documentation;\allowbreak{} accessed 18 November 2026, 2026.
\newblock URL \url{https://creativecommons.org/publicdomain/zero/1.0/}.
\newblock Verified in this writing on 18 November 2026 at these loci only.
  Official deed, No Copyright, items1–2.. Scope: Public-domain dedication
  identified by the official deed.

\bibitem[Deligne(1971)]{ffDeligne}
Pierre Deligne.
\newblock Th{\'e}orie de hodge, {II}.
\newblock \emph{Publications Math{\'e}matiques de l'IH{\'E}S}, 40:\penalty0
  5--57, 1971.
\newblock \doi{10.1007/BF02684692}.
\newblock URL
  \url{https://publications.ias.edu/sites/default/files/Number15.pdf}.
\newblock Verified in this writing on 18 November 2026 at these loci only.
  Definition 4.2.4 and Proposition 4.2.5 p. 44;\allowbreak{} Theorem 4.2.6 p. 45. Scope:
  Opened definitions and full semisimplicity theorem. Used only after the
  proposed period object is shown polarizable in the source proof. Author/title
  and edition matched to the opened item. Verification applies only to its
  recorded locator and edition;\allowbreak{} aliasing does not transfer preprint pagination
  to a published edition.

\bibitem[Esole and Kang(2018)]{ffEsoleKang}
Mboyo Esole and Monica~Jinwoo Kang.
\newblock Characteristic numbers of crepant resolutions of weierstrass models,
  2018.
\newblock URL \url{https://arxiv.org/pdf/1807.08755}.
\newblock Verified in this writing on 18 November 2026 at these loci only.
  Equations(1.2)–(1.4), printed pp3–4.. Scope: Universal Todd genus,
  Pontryagin classes and signature polynomials only;\allowbreak{} no extension of an
  algebraic crepant invariance theorem to the analytic fourfold is inferred.
  Equations (1.2)–(1.4), pp. 3–4, arXiv:1807.08755v1. Scope: Opened
  characteristic index polynomials;\allowbreak{} current manuscript explicitly substitutes
  c1=c3=0 and distinguishes the almost complex symbol-index use from
  holomorphic cohomology. Reopened the exact v1 PDF in the later locator pass:
  equation(1.2), printed p.3, gives the Todd polynomial;\allowbreak{}
  equations(1.3)–(1.4), printed p.4, give Pontryagin and signature
  polynomials. Current local PDF SHA256
  69bb3819230a1ed2c7f44d1289c5cd86b5661ce90066be2430e939ed63f0aa42.
  Author/title and edition matched to the opened item. Verification applies
  only to its recorded locator and edition;\allowbreak{} aliasing does not transfer preprint
  pagination to a published edition.

\bibitem[Forstneri{\v c} and Prezelj(2002)]{ffForstnericPrezelj02}
Franc Forstneri{\v c} and Jasna Prezelj.
\newblock Oka's principle for holomorphic submersions with sprays.
\newblock \emph{Mathematische Annalen}, 322\penalty0 (4):\penalty0 633--666,
  2002.
\newblock \doi{10.1007/s002080100249}.
\newblock URL \url{https://arxiv.org/pdf/math/0101040}.
\newblock Verified in this writing on 18 November 2026 at these loci only.
  Theorem 1.2, preprint p. 3, arXiv:math/0101040v3. Scope: Opened the full
  section theorem. The manuscript obtains a holomorphic frame from a continuous
  frame and the explicit GL(n,C) spray. Theorem1.2 is printed on preprint p.3;\allowbreak{}
  interpolation in the receiving lemma is constructed separately using the
  opened Oka-1 spray proof. Author/title and edition matched to the opened
  item. Verification applies only to its recorded locator and edition;\allowbreak{} aliasing
  does not transfer preprint pagination to a published edition.

\bibitem[Kramer and Stolz(2007)]{IVVKramerStolz}
Linus Kramer and Stephan Stolz.
\newblock A diffeomorphism classification of manifolds which are like
  projective planes.
\newblock \emph{Journal of Differential Geometry}, 77:\penalty0 177--188, 2007.
\newblock URL
  \url{https://archive.ymsc.tsinghua.edu.cn/pacm_download/21/694-JDG-77-2-A1-kramer.pdf}.
\newblock Verified in this writing on 18 November 2026 at these loci only.
  Theorem A, p. 178. Scope: Opened full statement and surrounding proof
  discussion;\allowbreak{} absorption of homotopy eight-spheres by the quaternionic
  projective plane;\allowbreak{} applied after the explicit Wall normal data comparison.
  Author/title and edition matched to the opened item. Verification applies
  only to its recorded locator and edition;\allowbreak{} aliasing does not transfer preprint
  pagination to a published edition.

\bibitem[Mandal(2026)]{ffMandal}
Souvik Mandal.
\newblock Complex and stable complex structures on real vector bundles over
  connected sums of quaternionic projective planes, 2026.
\newblock URL \url{https://arxiv.org/abs/2609.27929}.
\newblock Verified in this writing on 18 November 2026 at these loci only.
  Corollary 1.4 p. 4;\allowbreak{} Lemma 2.1 p. 5;\allowbreak{} Theorem 2.4 p. 7 in the current
  attached-source build. Scope: Read the original full source, compiled its
  body unchanged using only font-map additions, and opened rendered pages
  4–7. The bibliography is absent in this build;\allowbreak{} body statements and
  displayed normalizations are intact. Build:
  analysis/V/references\_build/Mandal\_opened.pdf. Author/title and edition
  matched to the opened item. Verification applies only to its recorded locator
  and edition;\allowbreak{} aliasing does not transfer preprint pagination to a published
  edition.

\bibitem[{OpenAI}(2026{\natexlab{a}})]{OpenAIChatGPTPro2026}
{OpenAI}.
\newblock Gpt-6 and other models in chatgpt.
\newblock Official documentation;\allowbreak{} accessed 18 November 2026,
  2026{\natexlab{a}}.
\newblock URL
  \url{https://help.openai.com/en/articles/20001354-gpt-6-and-other-models-in-chatgpt}.
\newblock Verified in this writing on 18 November 2026 at these loci only.
  Choose a thinking level;\allowbreak{} lines25–36 in the opened page.. Scope: The
  documented Pro selection and GPT-6 Astra model mapping.

\bibitem[{OpenAI}(2026{\natexlab{b}})]{OpenAIReasoning2026}
{OpenAI}.
\newblock Reasoning models.
\newblock Official API documentation, accessed 17 November 2026,
  2026{\natexlab{b}}.
\newblock URL \url{https://developers.openai.com/api/docs/guides/reasoning}.
\newblock Verified in this writing on 18 November 2026 at these loci only.
  Reasoning effort and Reasoning mode;\allowbreak{} lines1274–1303 in the opened page..
  Scope: Reasoning allocation controls and the named configuration fields.

\bibitem[Peters and Sterk(2024)]{ffPetersSterk}
Chris Peters and Hans Sterk.
\newblock Symmetric and quadratic forms, with applications to coding theory,
  algebraic geometry and topology.
\newblock Authors\textquotesingle{} June 2024 manuscript, 2024.
\newblock URL
  \url{https://www-fourier.univ-grenoble-alpes.fr/~peters/Books/QuadraticForms/QuadForms.pdf}.
\newblock Verified in this writing on 18 November 2026 at these loci only.
  Theorem2.2.1, printed p61.. Scope: A rank-two positive definite unimodular
  integral lattice is diagonal with entries1,1. Theorem 2.2.1, printed p. 61
  (PDF p. 68), June 2024 author version. Scope: Opened full integral lattice
  classification theorem and used it on the unimodular indefinite rank-four
  form. Author/title and edition matched to the opened item. Verification
  applies only to its recorded locator and edition;\allowbreak{} aliasing does not transfer
  preprint pagination to a published edition.

\bibitem[Scott(1983)]{rootScott83}
Peter Scott.
\newblock The geometries of 3-manifolds.
\newblock \emph{Bulletin of the London Mathematical Society}, 15\penalty0
  (5):\penalty0 401--487, 1983.
\newblock \doi{10.1112/blms/15.5.401}.
\newblock URL \url{https://courses.mai.liu.se/GU/TA1022/Scott.pdf}.
\newblock Verified in this writing on 18 November 2026 at these loci only.
  Theorems2.3–2.5 pp425–426, orbifold Euler discussion pp426–428,
  Theorem3.6 p436.. Scope: The four-cone orbifold is good and has hyperbolic
  universal cover by its negative orbifold Euler characteristic. The
  finite-cover Euler formula for Seifert bundles is
  e(upstairs)=base-degree/fibre-degree times e(downstairs). The exact
  positive-circle sign is computed separately from the actual weighted Hopf
  connection. Theorems 2.3–2.4, p. 425;\allowbreak{} orbifold presentation pp. 423–424;\allowbreak{}
  Euler/geometry sign pp. 427–428. Scope: Opened the full statements and
  their immediate discussion in the original 1983 article. Three-cone
  orientable sphere orbifolds are good;\allowbreak{} the sign of the orbifold Euler
  characteristic selects S2, E2 or H2. Applied to the actual triangle-group
  presentation and its classifying-space fibration, with the integral transfer
  vector retained. Author/title and edition matched to the opened item.
  Verification applies only to its recorded locator and edition;\allowbreak{} aliasing does
  not transfer preprint pagination to a published edition.

\bibitem[Yang(2026)]{manYang26}
Xiufan Yang.
\newblock Compact complex threefolds fibred by two-dimensional complex tori,
  2026.
\newblock URL \url{https://arxiv.org/html/2609.29666v1}.
\newblock Verified in this writing on 18 November 2026 at these loci only.
  Propositions3.1,3.4,3.6;\allowbreak{} Lemma4.3 and Proposition4.4.. Scope: Actual modular
  and affine period maps, logarithmic root action and its exact meridian
  boundary relation;\allowbreak{} receiving all-order proof re-derived separately. Theorem
  8.7, pp.57–59, and Corollary 8.8, p.59, in the locally rendered attached
  REPAIRED\_v4\_5 source;\allowbreak{} section 8.4. Scope: Read the full theorem and proof
  in the attached source and the rendered PDF pages57–59;\allowbreak{} visually opened
  p.57. Verified the exact period action exp(2pi i
  D(t,s)nu)t\textasciicircum{}nu, uniform logarithmic properness bounds, the
  finite fundamental polydiscs, the canonical form descent, the lattice index
  and maximal-simplex orbit Euler count. For rank3 primitive standard grid this
  gives6. Only font maps were added to render the attached body;\allowbreak{} the pagination
  is explicitly this attached build, not an unverified arXiv edition. Source
  SHA256 5061ce4e6cda8ed7ba3bb2a7a597d50e018bfd53fb9e49e7b38317f628c332a5;\allowbreak{} PDF
  SHA256 7601378a2250cd9e7a92f36afbd97c542b251a57711538735f9517db443e0533
  Author/title and edition matched to the opened item. Verification applies
  only to its recorded locator and edition;\allowbreak{} aliasing does not transfer preprint
  pagination to a published edition.

\end{thebibliography}

\end{document}
